% !TeX program = xelatex
% Standalone English volume. Use XeLaTeX until the references stabilize.
\documentclass[a4paper,11pt]{article}
\usepackage[T1]{fontenc}
\usepackage{lmodern}
\usepackage[english]{babel}
\usepackage[left=26mm,right=26mm,top=25mm,bottom=25mm,headheight=15pt,headsep=7mm,footskip=11mm]{geometry}
\usepackage{amsmath,amssymb,mathtools,amsthm}
\newtheorem{proposition}{Proposition}[section]
\newtheorem{lemma}[proposition]{Lemma}
\newtheorem{theorem}[proposition]{Theorem}
\newtheorem{corollary}[proposition]{Corollary}
\newtheorem{maintheorem}[proposition]{Main Theorem}
\theoremstyle{definition}
\newtheorem{definition}[proposition]{Definition}
\theoremstyle{remark}
\newtheorem{remark}[proposition]{Remark}
\makeatletter
\renewcommand*\l@section{\@dottedtocline{1}{0em}{2.4em}}
\renewcommand*\l@subsection{\@dottedtocline{2}{1.45em}{3.45em}}
\renewcommand*\l@subsubsection{\@dottedtocline{3}{3.8em}{4.55em}}
\makeatother
\usepackage{xcolor,fancyhdr,titlesec,needspace,xurl}
\usepackage{longtable,multicol}
\usepackage{hyperref}
\definecolor{ink}{HTML}{1A3542}
\definecolor{muted}{HTML}{4F606A}
\hypersetup{unicode=true,colorlinks=true,linkcolor=ink,filecolor=ink,urlcolor=ink,citecolor=ink,bookmarksnumbered=true,
 pdfdisplaydoctitle=true,pdflang={en},
 pdftitle={The generalized Riemann hypothesis for Dirichlet L-functions. Edition 930, Volume III},
 pdfsubject={A natural-language reconstruction of the proposed proof},
 pdfauthor={A mathematical reconstruction}}
\titleformat{\section}{\Large\bfseries\color{ink}}{\thesection}{.65em}{}
\titlespacing*{\section}{0pt}{.8em}{.65em}
\setlength{\parindent}{1.2em}
\linespread{1.045}
\setlength{\parskip}{.22em plus .05em minus .03em}
\setlength{\emergencystretch}{2em}
\setcounter{topnumber}{2}
\setcounter{tocdepth}{3}\setcounter{secnumdepth}{3}
\widowpenalty=4500
\clubpenalty=4500
\pagestyle{fancy}\fancyhf{}
\fancyhead[L]{\small\color{muted}Dirichlet GRH \textperiodcentered\ Volume I}
\fancyhead[R]{\small\color{muted}Dirichlet zeros and power intervals}
\fancyfoot[C]{\small\thepage}
\renewcommand{\headrulewidth}{.3pt}
\newcommand{\lead}[1]{\par\smallskip\noindent{\bfseries\color{ink}#1.}\enspace}
\newcommand{\sour}[2]{\unskip\penalty150\hspace{.35em}{\small\hyperlink{ref#1}{[#2]}}}
\titleformat{\subsection}{\large\bfseries}{\thesubsection}{.6em}{}
\titlespacing*{\subsection}{0pt}{1.15em plus .15em}{.4em}
\setlength{\abovedisplayskip}{9pt plus 2pt minus 3pt}
\setlength{\belowdisplayskip}{9pt plus 2pt minus 3pt}
\newcommand{\RR}{\mathbb R}\newcommand{\CC}{\mathbb C}\newcommand{\QQ}{\mathbb Q}
\newcommand{\ZZ}{\mathbb Z}\newcommand{\ord}{\operatorname{ord}}
\newcommand{\Frob}{\operatorname{Frob}}\newcommand{\Gal}{\operatorname{Gal}}
\newcommand{\len}{\operatorname{length}}\newcommand{\Res}{\operatorname*{Res}}
\newcommand{\wmean}{\operatorname{av}}\newcommand{\Ker}{\operatorname{ker}}
\newcommand{\Lh}{\mathcal L}\newcommand{\Gg}{\mathcal G}
\newcommand{\real}{\operatorname{Re}}\newcommand{\imag}{\operatorname{Im}}
\newcommand{\abs}[1]{\left|#1\right|}\newcommand{\eps}{\varepsilon}
\newcommand{\anchoritem}[2]{\par\Needspace{2.8\baselineskip}\hypertarget{ref#1}{}\noindent{\bfseries\color{ink}#2.}\enspace}
\newcommand{\srcloc}[2]{{\small\nolinkurl{#1.lean}}, line~#2}
\numberwithin{equation}{section}
\raggedbottom

\newcommand{\Sel}{\mathcal S}
\newcommand{\Cl}{\operatorname{Cl}}
\newcommand{\Aut}{\operatorname{Aut}}
\newcommand{\GL}{\operatorname{GL}}
\newcommand{\OO}{\mathcal O}
\hypersetup{bookmarksdepth=3}


\addto\captionsenglish{\renewcommand{\partname}{Volume}}


\makeatletter
\expandafter\def\csname r@cor:normalizedtransfer722\endcsname{{45.17}{55}{}{proposition.45.17}{Dirichlet-GRH-v930-Volume-IX.pdf}}
\expandafter\def\csname r@cor:oddwindow677\endcsname{{11.30}{27}{}{proposition.11.30}{Dirichlet-GRH-v930-Volume-II.pdf}}
\expandafter\def\csname r@cor:quotientchainheight734\endcsname{{13.14}{73}{}{proposition.13.14}{Dirichlet-GRH-v930-Volume-II.pdf}}
\expandafter\def\csname r@cor:remainingbudget784\endcsname{{52.3}{60}{}{proposition.52.3}{Dirichlet-GRH-v930-Volume-X.pdf}}
\expandafter\def\csname r@cor:scalarfl724\endcsname{{12.16}{42}{}{proposition.12.16}{Dirichlet-GRH-v930-Volume-II.pdf}}
\expandafter\def\csname r@eq:mainloss732\endcsname{{{4}}{66}{}{AMS.5}{Dirichlet-GRH-v930-Volume-II.pdf}}
\expandafter\def\csname r@eq:projectedcups742\endcsname{{{10}}{12}{}{AMS.1}{Dirichlet-GRH-v930-Volume-X.pdf}}
\expandafter\def\csname r@lem:abelianunits744\endcsname{{50.12}{15}{}{proposition.50.12}{Dirichlet-GRH-v930-Volume-X.pdf}}
\expandafter\def\csname r@lem:archstring787\endcsname{{52.11}{69}{}{proposition.52.11}{Dirichlet-GRH-v930-Volume-X.pdf}}
\expandafter\def\csname r@lem:conductorsign783\endcsname{{51.21}{56}{}{proposition.51.21}{Dirichlet-GRH-v930-Volume-X.pdf}}
\expandafter\def\csname r@lem:graphrealization724\endcsname{{12.14}{41}{}{proposition.12.14}{Dirichlet-GRH-v930-Volume-II.pdf}}
\expandafter\def\csname r@lem:localevaluation728\endcsname{{12.27}{51}{}{proposition.12.27}{Dirichlet-GRH-v930-Volume-II.pdf}}
\expandafter\def\csname r@lem:nearweightprimes766\endcsname{{51.4}{42}{}{proposition.51.4}{Dirichlet-GRH-v930-Volume-X.pdf}}
\expandafter\def\csname r@lem:onesignmain771\endcsname{{51.9}{46}{}{proposition.51.9}{Dirichlet-GRH-v930-Volume-X.pdf}}
\expandafter\def\csname r@lem:spacing781\endcsname{{51.19}{54}{}{proposition.51.19}{Dirichlet-GRH-v930-Volume-X.pdf}}
\expandafter\def\csname r@lem:sphere48\endcsname{{6.25}{58}{}{proposition.6.25}{Dirichlet-GRH-v930-Volume-I.pdf}}
\expandafter\def\csname r@lem:sumcupheight763\endcsname{{51.1}{39}{}{proposition.51.1}{Dirichlet-GRH-v930-Volume-X.pdf}}
\expandafter\def\csname r@lem:tameorbit715\endcsname{{47.35}{101}{}{proposition.47.35}{Dirichlet-GRH-v930-Volume-IX.pdf}}
\expandafter\def\csname r@lem:wildgenerators745\endcsname{{50.14}{19}{}{proposition.50.14}{Dirichlet-GRH-v930-Volume-X.pdf}}
\expandafter\def\csname r@prop:conductorcups783\endcsname{{51.22}{56}{}{proposition.51.22}{Dirichlet-GRH-v930-Volume-X.pdf}}
\expandafter\def\csname r@prop:testedkernel785\endcsname{{52.4}{62}{}{proposition.52.4}{Dirichlet-GRH-v930-Volume-X.pdf}}
\expandafter\def\csname r@sec:levelmass678\endcsname{{130.1}{3}{}{subsection.130.1}{Dirichlet-GRH-v930-Volume-XXV.pdf}}
\expandafter\def\csname r@sec:missing\endcsname{{6.6}{32}{}{subsection.6.6}{Dirichlet-GRH-v930-Volume-I.pdf}}
\expandafter\def\csname r@seriesguide\endcsname{{1.1}{4}{}{subsection.1.1}{Dirichlet-GRH-v930-Volume-I.pdf}}
\expandafter\def\csname r@thm:fixedconductor781\endcsname{{51.20}{55}{}{proposition.51.20}{Dirichlet-GRH-v930-Volume-X.pdf}}
\expandafter\def\csname r@thm:fllift724\endcsname{{12.15}{41}{}{proposition.12.15}{Dirichlet-GRH-v930-Volume-II.pdf}}
\expandafter\def\csname r@thm:labelledlift735\endcsname{{13.15}{74}{}{proposition.13.15}{Dirichlet-GRH-v930-Volume-II.pdf}}
\expandafter\def\csname r@thm:longcharacter787\endcsname{{52.10}{68}{}{proposition.52.10}{Dirichlet-GRH-v930-Volume-X.pdf}}
\expandafter\def\csname r@thm:oldcolumn745\endcsname{{50.16}{20}{}{proposition.50.16}{Dirichlet-GRH-v930-Volume-X.pdf}}
\expandafter\def\csname r@thm:projectedcups742\endcsname{{50.11}{12}{}{proposition.50.11}{Dirichlet-GRH-v930-Volume-X.pdf}}
\expandafter\def\csname r@thm:rankoneselmer737\endcsname{{50.6}{6}{}{proposition.50.6}{Dirichlet-GRH-v930-Volume-X.pdf}}
\expandafter\def\csname r@thm:ribettrace720\endcsname{{45.11}{48}{}{proposition.45.11}{Dirichlet-GRH-v930-Volume-IX.pdf}}
\expandafter\def\csname r@thm:splitattachment722\endcsname{{45.16}{54}{}{proposition.45.16}{Dirichlet-GRH-v930-Volume-IX.pdf}}
\expandafter\def\csname r@thm:tametransfer721\endcsname{{45.14}{51}{}{proposition.45.14}{Dirichlet-GRH-v930-Volume-IX.pdf}}
\expandafter\def\csname r@thm:typedindex799\endcsname{{52.13}{71}{}{proposition.52.13}{Dirichlet-GRH-v930-Volume-X.pdf}}
\makeatother
\begin{document}
\setcounter{part}{2}
\fancyhead[L]{\small\color{muted}Dirichlet GRH \textperiodcentered\ Volume III}
\fancyhead[R]{\small\color{muted}Frobenius lifts and local cup products}
\thispagestyle{plain}\begin{center}
{\LARGE\bfseries The generalized Riemann hypothesis\\[2pt]for Dirichlet \(L\)-functions}\\[6pt]
{\small 27 November 2026 \quad\textperiodcentered\quad Edition 930}\end{center}
\part{Trace extensions, Frobenius lifts and local cup products}\label{part:III}
\begin{abstract}\noindent We analyze the integral extensions detected by the marked traces. Local Frobenius lattices, cup products and Taylor conditions give bounds for the scalar and higher-rank terms at spherical and ramified primes.\end{abstract}
\setcounter{section}{14}
\noindent The \hyperref[seriesguide]{guide in Volume~I} lists all twenty-seven volumes.
\begingroup\small\setlength{\parskip}{0pt}\tableofcontents\endgroup
\clearpage

\section{Scalar trace extensions at the original level}\label{nav:trace-extensions}\label{nav:volume-III}

\subsection{Rank-one trace classes at places with negative filtration jumps}
\label{sec:negtrace740}

The local terms in \eqref{eq:mainloss732} depend on the integral
lattice and on a range of weights of the representations giving
the trace relations. At some completions the target alone gives
an upper bound. We first compute it at a single place. We then
sum the bounds on the same fields and for the same twists as in
Theorem~\ref{thm:rankoneselmer737}.

Let \(K/\QQ_\ell\) be finite. Let \(\mathcal O\) be the integer
ring of a finite extension \(L/\QQ_\ell\) containing the images
of all embeddings of \(K\), and fix a uniformizer \(\pi\). Put
\(e_L=v_\pi(\ell)\) and normalize \(v_\ell(\ell)=1\).
For a \(G_K\)-stable free lattice \(T\) put \(V=T[1/\ell]\),
\(T'=T^\vee(1)\), and \(A'=T'[1/\ell]/T'\). The cyclotomic
character has de Rham filtration jump \(-1\). Write \(Z_f(K,T)\)
for integral cocycles whose rational class lies in \(H^1_f(K,V)\).
At precision \(a\ge1\), put
\[
 D_a(T)=
 Z^1(K,T/\pi^a T)\big/\operatorname{red}_{a} Z_f(K,T).
\]
We use continuous cocycles throughout.

\begin{proposition}[The full local lifting obstruction]
\leavevmode\label{prop:negobstruction740}
Suppose \(V\) is de Rham, every filtration jump of \(V\) at every
embedding of \(K\) is negative, and \(H^0(K,T'[1/\ell])=0\).
Then \(D_a(T)\) is canonically isomorphic to \(H^2(K,T)[\pi^a]\).
Its normalized length is that of
\(H^0(K,A')/\pi^a H^0(K,A')\), divided by \(e_L\), where length
on the latter module is taken over \(\mathcal O\). In rank one
let \(\beta:G_K\to\mathcal O^\times\) be the character of \(T'\)
and \(I(\beta)=(\beta(g)-1:g\in G_K)\). Then
\[
     \frac{\operatorname{length}_{\mathcal O}D_a(T)}{e_L}
        =\min\bigl(a/e_L,v_\ell(I(\beta))\bigr).
 \tag{6}\label{eq:negobstruction740}
\]
The same right hand side bounds the excess for any submodule of
raw trace cocycles, measured modulo its intersection with the
integral finite cocycles reduced to this precision.
\end{proposition}
\begin{proof}
The de Rham tangent space \(D_{\rm dR}(V)/{\rm Fil}^0\) has
\(L\)-dimension \([K:\QQ_\ell]\dim_L V\), since the filtration
has no graded part in nonnegative degree. The Bloch--Kato
dimension formula adds \(\dim_L H^0(K,V)\) to this tangent
dimension to give \(\dim_L H^1_f(K,V)\). Local duality and the
Euler characteristic formula give
\[
 \dim_L H^1(K,V)=[K:\QQ_\ell]\dim_L V+\dim_L H^0(K,V)
                              +\dim_L H^0(K,V^\vee(1)).
\]
The last summand is zero by hypothesis. Thus every rational
class is finite, and \(Z_f(K,T)=Z^1(K,T)\).
\cite[Corollary~3.8.4 and the proof of (3.8.5)]{III-II-II-bktate717}

Every coboundary modulo \(\pi^a\) lifts by lifting the vector
defining it. The quotient of raw cocycles is therefore the
cokernel of reduction on \(H^1\). The exact sequence for
multiplication by \(\pi^a\) on \(T\) identifies this cokernel
with \(H^2(K,T)[\pi^a]\). Local Tate duality identifies the
Pontryagin dual of \(H^2(K,T)\) with \(H^0(K,A')\) after choosing
a generator of the inverse different of \(L/\QQ_\ell\).
Indeed that generator and the field trace identify the
\(\mathcal O\)-linear and \(\ZZ_\ell\)-linear dual lattices.
Apply finite local duality to the reductions of the lattices and
pass to the inverse limit on degree two and the direct limit on
degree zero. The finite groups in the inverse systems satisfy
the Mittag--Leffler condition. The dual of the kernel of
multiplication by \(\pi^a\) is the quotient by \(\pi^a\) in
degree zero. A finite module and its Pontryagin dual have the
same length. \cite[Chapter~I, \S2]{III-II-II-duality13}

For completeness \(H^0(K,A')\) is finite here. Over \(L\) the
common kernel of \(g-1\) on \(T'[1/\ell]\) is zero. Finitely
many \(g-1\) have zero common kernel by finite dimensionality.
The map to the direct sum given by these operators has an
\(L\)-linear left inverse. Clearing its denominators bounds the
denominators modulo \(T'\) of every invariant. In rank one the
invariants are precisely \(I(\beta)^{-1}\mathcal O/\mathcal O\)
in a basis of \(T'\). The ideal is nonzero by the hypothesis on
invariants. Its single elementary divisor gives
\eqref{eq:negobstruction740}. Restricting the quotient map to
a submodule proves the last assertion.
\end{proof}

\begin{lemma}[A test on rational local units]
\leavevmode\label{lem:rationalunits740}
Suppose \(T\) has rank one and that on \(\mathcal O_K^\times\)
local reciprocity expresses its dual character in the form
\[
     \beta({\rm rec}_K u)=\eta(u)
        \prod_{\tau:K\hookrightarrow L}\tau(u)^{\epsilon b_\tau},
        \qquad \epsilon\in\{1,-1\}.
\]
Here \(\eta\) is trivial on \(1+\mathfrak m_K\), all \(b_\tau\)
are nonnegative integers, and \(B=\sum_\tau b_\tau>0\).
Put \(c_2=2\) and \(c_\ell=1\) for odd \(\ell\). Then
\[
       v_\ell(I(\beta))\le c_\ell+v_\ell(B).
\]
If \(\eta\) is trivial on \(\mathcal O_K^\times\) and
\(\ell-1\nmid B\), one has \(I(\beta)=\mathcal O\).
More generally \(I(\beta)=\mathcal O\) whenever \(\ell\) is odd,
\(\eta^Q=1\) on the units, and \(\ell-1\nmid QB\).
\end{lemma}
\begin{proof}
Embed \(\ZZ_\ell^\times\) in the local units. Every embedding
sends \(u=1+\ell^{c_\ell}\) to itself and \(\eta(u)=1\).
The valuation of \(u^{\epsilon B}-1\) is \(c_\ell+v_\ell(B)\).
To see this, raising a principal unit to a power prime to \(\ell\)
preserves its depth by the binomial expansion. Raising it to
the power \(\ell\) increases the depth by exactly one for odd
\(\ell\) at positive integral depth and at \(\ell=2\) at depth
at least two. Inversion preserves the valuation. This proves
the upper bound and also proves that \(\beta\ne1\).
For the residue test take a Teichmüller generator
\(\zeta\in\ZZ_\ell^\times\) of order \(\ell-1\). The \(Q\)-th
power of its image is \(\zeta^{\epsilon QB}\). When the latter
has nontrivial reduction, the original image has nontrivial
reduction as well. At two the divisibility \(1\mid B\) always
holds. The argument uses rational units inside \(K\) for every
ramification index of \(K/\QQ_\ell\).
\end{proof}

\begin{theorem}[A uniform bound on the negative-jump part of the main trace]
\leavevmode\label{thm:negativeinterval740}
Fix \(q_0,d\) in the shorter family of
Corollary~\ref{cor:oddwindow677}. On one of its period intervals
choose the designated data at each integer node \(x\). Take all
the main rank-one lattices of Theorem~\ref{thm:rankoneselmer737},
their critical integers \(m_x<z\le r_x\), coefficient embeddings
\(\iota\), and signed powers of the prescribed twists. The estimate
also applies with \(1\le z\le t_x\) to rank-one characters with
the same weight windows, orientations, and tame unit factors. At each node write \(b_x\) for the
large weight scale. For each such lattice \(T\) write \(W^-(T,\ell)\)
for the completions \(w\mid\ell\) at which the filtration jumps
of \(T\) are all negative and at least one jump is at most \(-2\).
Let \(d_w(T)\) be the supremum over precision of the normalized
length in \eqref{eq:negobstruction740} over \(K=E_{x,w}\).
At each of the two level primes \(p_{x,i}\),
\(d_w(T)=0\) for every \(w\in W^-(T,p_{x,i})\). If the interval
starts at \(V_0\) sufficiently large, then
\[
   \sum_{\ell}(\log\ell)
        \sup_{x,\iota,T}
        \frac{\displaystyle\sum_{w\in W^-(T,\ell)}d_w(T)}
             {\sqrt{b_x}}
       \ \ll_{q_0,d}\ V_0^{-10}.
 \tag{7}\label{eq:negativeinterval740}
\]
The bound applies to the raw trace classes at every auxiliary
level and precision. It includes places above \(2\) and the
discriminant. At each rational prime splitting in
\(\QQ(\sqrt{-d})\), \(W^-\) contains exactly one member of
every conjugate pair of places above that prime for the designated
characters.
\end{theorem}
\begin{proof}
We recall explicitly the data entering the estimate. Put
\(n_x=[E_x:\QQ]\), where
\(E_x=\QQ(\mu_{q_{1,x}},\sqrt{-d})\). The two base weights
are \(b_x\) and \(a_x/x_{1,x}\). For \(1\le j\le m_x\) the
unshifted lift weight is \(w_{j,\tau}=k^{(0)}_\tau+(2j+1)r_x^8\)
with \(k^{(0)}_\tau\in\{b_x,a_x/x_{1,x}\}\). Indeed the
construction adds half the gap \(2r_x^8\) to the base weight
and then adds \(j\) full gaps. The defining window of a main
pair is \(w_{j,\tau}+r_x-t_x\le k_\tau\le w_{j,\tau}\).
Equivalently \(k_\tau=w_{j,\tau}-h_\tau\), where each \(h_\tau\)
is an integer in \([0,4r_x]\) and is shared by the two conjugate
labels of that infinite place. This is the weight condition in
\texttt{CriticalPair.IsMainPair} in
\path{New/Determinant/Twist/FitTwists.lean} of the shorter archive.
The proof below allows the choices at distinct infinite places
to be independent. The positivity of the base weights and
\(r_x^8-4r_x>5r_x\) for \(r_x\ge2\) give \(k_\tau>t_x\). The bounds \(k_\tau\le 4b_x\) also follow from
\(k^{(0)}_\tau\le b_x\), \(m_x\le2\sqrt{x}\),
\(r_x\le3x^5\), and \(b_x\ge x^{500}\) for sufficiently large \(x\).
The last inequality is weaker than the scale conditions
\(x^{10^6}\le x_{1,x}\) and \(x_{1,x}^{10^6}\le b_x\).

For the pulled-back character \(u(a)=\nu(a/\bar a)^{-1}\) the
two exponents at conjugate embeddings have opposite signs and
magnitude \(k_\tau\). This follows by applying the defining
torus character to the quotient of the two coordinates.
Its sign is fixed on the chosen CM type, induced from the
specified embedding of \(\QQ(\sqrt{-d})\), and reversed on the
conjugate type. The main lattice has character
\(u_\ell\kappa^{1-z}\), with \(\kappa\) cyclotomic. Hence its
two jumps are \(z-1-k_\tau\) and \(z-1+k_\tau\). On a place
in \(W^-\) the dual jumps are \(b_\tau=k_\tau-z>0\). A prime
splitting in the quadratic field separates its two embeddings
at every completion, so exactly one member of each pair has
only negative jumps. The finite-order auxiliary characters
are trivial at the coefficient completions in the chosen twists.
All these statements apply to each signed power.

On the local units the avatar of an algebraic Hecke character
has its finite character multiplied by the monomial in the
labelled embeddings. Indeed its definition transports the
finite character through the coefficient embedding and replaces
each algebraic infinity factor by the corresponding local
embedding in the component above \(\ell\). Evaluation on the
idele supported at \(w\) retains precisely those labels inducing
\(w\). Dualizing inverts both factors, and the Tate twist shifts
each jump by \(-1\). Thus on \(W^-\) the unit formula of
Lemma~\ref{lem:rationalunits740} applies with \(b_\tau=k_\tau-z\).
The one overall sign depends on the reciprocity convention.
The smooth unit factor is trivial off the marked places. At the
marks its order divides \(Q\), the order of the prescribed
residue character. This follows from the tame parts
\(\omega_w^{\pm j}\) in Corollary~\ref{cor:quotientchainheight734}.
In particular the dual character is nontrivial by the rational
principal-unit test, as required in Proposition~\ref{prop:negobstruction740}.

For \(\ell=p_{x,i}\) every completion is \(\QQ_\ell\), since
\(p_{x,i}\equiv1\bmod Q^2\) and the conductor of \(E_x\)
divides \(Q\). There is just one label, so \(0<B\le 4b_x\).
We have \(Q\le\sqrt{p_{x,i}-1}\) and \(p_{x,i}\ge b_x^{10}\)
as in Section~\ref{sec:levelmass678}. Thus \(0<QB<p_{x,i}-1\)
for sufficiently large nodes. The residue test gives
\(I(\beta)=\mathcal O\) and proves the asserted vanishing.

At the remaining primes a nonzero loss requires \(\ell-1\mid B\).
At a fixed node the possible positive \(B\) lie in a finite set
\(\mathcal B_x\) formed as follows. Fix \(j,z\), and take the
sum of \(n_1\) copies of the first unshifted weight minus \(z\)
and \(n_2\) copies of the second unshifted weight minus \(z\),
with \(1\le n_1+n_2\le n_x\). Subtract an integer between zero
and \(4r_x(n_1+n_2)\), the sum of all the shifts \(h_\tau\)
at this completion. Retain the positive sums. Every completion
and every coefficient embedding occurs in this list. Thus
\[
   |\mathcal B_x|\le m_xt_x(n_x+1)^2(1+4r_xn_x).
\]
Every member is at most \(4n_xb_x\). For any one positive integer
\(B\), write \(\tau(B)\) for its number of positive divisors.
The map \(\ell\mapsto\ell-1\) sends the possible primes to
distinct divisors of \(B\). Therefore
\[
 \sum_{\ell-1\mid B}(c_\ell+v_\ell(B))\log\ell
       \le \tau(B)\log(B+1)+\log2+\log B .
\]
The extra \(\log2\) accounts for \(c_2=2\). For any fixed
\(\delta>0\) this upper bound is \(O_\delta((1+B)^\delta)\).
Indeed, for \(\alpha=\delta/2\), every prime \(p\) large
enough satisfies \(j+1\le2^j\le p^{\alpha j}\) for \(j\ge1\).
For each of the finitely many smaller primes the supremum of
\((j+1)p^{-\alpha j}\) is finite. Multiplication gives
\(\tau(B)\ll_\alpha B^\alpha\), and the remaining logarithm
is \(O_\alpha((1+B)^\alpha)\).

There are at most \(n_x\) completions over a prime. Summing over
the explicit list, which contains all the possible \(B\) even
when the maximizing lattice varies with \(\ell\), gives
\[
 \begin{split}
 &\sum_\ell(\log\ell)\sup_{\iota,T}\sum_{w\in W^-(T,\ell)}d_w(T)\\
 &\quad\ll_\delta n_x(n_x+1)^2m_xt_x(1+4r_xn_x)(1+4n_xb_x)^\delta .
 \end{split}
\]
Use \(\delta=1/4\) and \(n_x\ll_{q_0,d}x^{11}\log(2x)
\ll_{q_0,d}x^{12}\), as proved in
Corollary~\ref{cor:oddwindow677}. Since \(m_xt_x\ll x^{11/2}\),
the extra factor \(1+4r_xn_x\) is \(O_{q_0,d}(x^{17})\).
Division by \(\sqrt{b_x}\) gives \(O_{q_0,d}(x^{62}b_x^{-1/4})\).
This is \(O_{q_0,d}(x^{-63})\) by the lower bound on \(b_x\).
For each prime the supremum across nodes is at most the sum of
their nonnegative bounds before the prime sum is taken. The
tail \(\sum_{x\ge V_0}x^{-63}\) over integer nodes is
\(O(V_0^{-62})\), proving \eqref{eq:negativeinterval740}.
The final estimate for trace classes follows by restricting the
quotient map of Proposition~\ref{prop:negobstruction740} at
each place, with the very same coefficient prime and lattice.
\end{proof}

In \eqref{eq:mainloss732} one may use \(d_w(T)\) for the
contribution at each \(w\in W^-\). Indeed the maps to the full
local lifting quotients bound the images of the restricted trace
classes at every precision. The other completions still require
their integral finite-condition estimates. The projected period
in Section~\ref{sec:missing} also involves higher-rank lattices
and the comparison of its integral and archimedean realizations.
\clearpage

\subsection{Integral cup-product tests at all coefficient places}
\label{sec:allcup741}

The obstruction to lifting a torsion cocycle to an integral
cocycle and the obstruction to imposing the finite local condition
on such a lift enter the trace estimate together. We bound the
first at every completion of the shorter period fields, including
ones where a rank-one character has jumps of both signs. We also
determine how many integral cup products suffice to test the
second obstruction and how their joint image measures the loss.

\begin{proposition}[Integral tests for a finite local condition]
\leavevmode\label{prop:freetests741}
Keep the local notation \(K,L,\mathcal O,\pi,T,V,T'\) of
Section~\ref{sec:negtrace740}, with \(V\) de Rham. Write
\(F(T)\) for the inverse image of \(H_f^1(K,V)\) in \(H^1(K,T)\),
and \(F_a(T)\) for its image in \(H^1(K,T/\pi^a T)\).
The finite local Tate pairing makes \(F_a(T)\) and \(F_a(T')\)
exact annihilators. Choose lifts \(y_1,\ldots,y_u\) of a basis
of the free quotient of \(F(T')\). For every submodule
\(C\subset H^1(K,T/\pi^a T)\) form the joint map
\[
 P_a:C\longrightarrow(\mathcal O/\pi^a)^u,\qquad
 h\longmapsto
       \bigl({\rm inv}_K(h\mathbin{\cup}\operatorname{red}_a y_i)\bigr)_i .
\]
Here \({\rm inv}_K\) uses the evaluation pairing and the local
invariant in degree two. If \(b_a\) is the length of the image
of \(F(T')_{\rm tors}\) under reduction at precision \(a\), then
\[
  0\le
     \operatorname{length}_{\mathcal O}\bigl(C/(C\cap F_a(T))\bigr)
          -\operatorname{length}_{\mathcal O}P_a(C)
      \le b_a .
 \tag{8}\label{eq:freetests741}
\]
The cokernel of integral reduction
\(H^1(K,T)\to H^1(K,T/\pi^a T)\) also has length \(b_a\).
For simultaneous restrictions to several places one may take
all the displayed coordinates in a single map and sum the upper
errors. These assertions apply to submodules of raw cocycles by
passing to cohomology.
\end{proposition}
\begin{proof}
Write \(q=|\mathcal O/\pi|\) and \(h_i(V)=\dim_L H^i(K,V)\).
Continuous cohomology of a local lattice is finitely generated over
\(\mathcal O\), and inverting \(\pi\) gives rational cohomology.
These facts follow in degrees zero and one from the inverse limit
of finite local cohomology, finiteness modulo \(\pi\), and compact
Nakayama. Rational cocycles have bounded image by compactness of
\(G_K\), so clearing a denominator computes the rationalization
in degree one. Write \(h_f(V)=\dim_L H_f^1(K,V)\).
The quotient \(H^1(K,T)/F(T)\) is torsion free, since it maps
injectively into the \(L\)-vector space \(H^1(K,V)/H_f^1(K,V)\).
Thus the kernel of reduction on \(F(T)\) is \(\pi^a F(T)\).
If \(u_a(T)\) is the length of \(H^1(K,T)[\pi^a]\), the structure
theorem gives
\[
  |F_a(T)|=q^{a h_f(V)+u_a(T)},\qquad b_a=u_a(T').
\]
The second equality holds since all torsion classes belong to
\(F(T')\) and the torsion submodule is a direct summand of
this finitely generated module.

The product of the orders of \(F_a(T)\) and \(F_a(T')\) is
\(|H^1(K,T/\pi^a T)|\). To verify this explicitly, the finite
Euler characteristic and local duality give
\[
 |H^1(K,T/\pi^a T)|=
 q^{a[K:\QQ_\ell]\dim_L V}
       |H^0(K,T/\pi^a T)|\,|H^0(K,T'/\pi^a T')|.
\]
The first invariant group has order \(q^{a h_0(V)+u_a(T)}\)
by the multiplication sequence in degrees zero and one. Use the
same formula for \(T'\). The Bloch--Kato dimension formula for
the two de Rham representations gives
\[
 h_f(V)+h_f(V^\vee(1))=[K:\QQ_\ell]\dim_L V
                          +h_0(V)+h_0(V^\vee(1)).
\]
Indeed at each embedding the paired jumps are \(j\) and \(-1-j\),
so the two tangent dimensions add to \([K:\QQ_\ell]\dim_L V\).
This proves the order identity. Integral lifts in the two finite
conditions pair to zero. Their rational classes pair to zero by
Bloch--Kato duality, and \(H^2(K,\mathcal O(1))\simeq\mathcal O\)
is torsion free. Finite local duality is perfect, hence the
order identity proves the assertion about annihilators.
The \(\mathcal O/\pi^a\)-valued pairing may equivalently be
tested through its trace pairing with a generator of the inverse
different as in Proposition~\ref{prop:negobstruction740}.
\cite[Theorem~3.8 and Corollary~3.8.4]{III-II-II-bktate717}
\cite[Chapter~I, \S2]{III-II-II-duality13}

Every element of \(\ker P_a\) pairs trivially with the free
generators. Its remaining pairings define a homomorphism from the
reduction of \(F(T')_{\rm tors}\) to \(\mathcal O/\pi^a\). The
kernel of this test on \(\ker P_a\) is \(C\cap F_a(T)\) by the
annihilator equality. Homomorphisms from a finite module killed by
\(\pi^a\) to \(\mathcal O/\pi^a\) form a module of the same length,
as is seen on cyclic summands. This proves \eqref{eq:freetests741}.
The cokernel of integral reduction is \(H^2(K,T)[\pi^a]\).
Its dual is \(H^0(K,A')/\pi^a H^0(K,A')\). The divisible image of
\(H^0(K,T'[1/\ell])\) in \(H^0(K,A')\) has quotient
\(H^1(K,T')_{\rm tors}\) by the coefficient sequence. Reducing
modulo \(\pi^a\) therefore gives length \(u_a(T')\). For a family
of places apply the same kernel argument to the direct sum of
the torsion tests on the common subgroup \(C\). Finally every
finite coboundary lifts to an integral coboundary, which belongs
to \(Z_f(K,T)\). This justifies passage from raw cocycles.
\end{proof}

\begin{theorem}[Torsion bounds in both CM orientations]
\leavevmode\label{thm:alltorsion741}
Take the actual rank-one main lattices and the period intervals
of Theorem~\ref{thm:negativeinterval740}. Write \(n_x=[E_x:\QQ]\).
For a lattice character \(M\) put
\[
 h_w(M)=e_L^{-1}\operatorname{length}_{\mathcal O}
              H^0(E_{x,w},M[1/\ell]/M)
\]
whenever the group is finite. For both \(M=T\) and
\(M=T^\vee(1)\) these lengths are finite at every coefficient
place. Both vanish at every place over either level prime.
Moreover on any of the indicated intervals starting at \(V_0\)
one has
\[
 \sum_\ell(\log\ell)\sup_{x,\iota,T}
 \frac{\displaystyle\sum_{w\mid\ell}\bigl(h_w(T)+h_w(T^\vee(1))\bigr)}
      {\sqrt{b_x}}\ll_{q_0,d}V_0^{-10}.
 \tag{9}\label{eq:alltorsion741}
\]
The twists and their signed powers are included in the
supremum. Every completion in each of these sums is taken on
the field of the same period node.
\end{theorem}
\begin{proof}
Put \(K_0=\QQ(\sqrt{-d})\). Up to a common choice of reciprocity
sign, the exponents on the labelled local units of each of the
two characters are \(c+\epsilon_\tau k_\tau\). Here
\(c=z-1\) or \(-z\), \(\epsilon_\tau\) is constant on each
orientation of \(K_0\) with opposite values on the two orientations,
and \(t_x<k_\tau\le4b_x\). The two base weights, the allowed
labelled shifts, and the tame unit components are computed in the proof of
Theorem~\ref{thm:negativeinterval740}. In particular all
principal units kill the smooth factor. For any nontrivial
lattice character \(\beta\), taking invariants in \(L/\mathcal O\)
gives \(I(\beta)^{-1}\mathcal O/\mathcal O\). Thus one local
unit giving a value different from one bounds \(h_w\) by the
valuation of that difference.

At a prime splitting in \(K_0\) test on the rational units as
in Lemma~\ref{lem:rationalunits740}, allowing the common
exponent to have either sign. For the set of labels at \(w\)
use the positive integer
\[
       B=\left|\sum_{\tau\mid w}(c+\epsilon_\tau k_\tau)\right|.
\]
It is positive since the signs agree and \(k_\tau>|c|\).
At each old level prime the field splits completely. There
\(0<B\le5b_x\) and \(QB<p_{x,i}-1\) for sufficiently large
nodes, so raising a residue-unit value to \(Q\) proves that
the full invariant group vanishes for both lattices. At the
other split primes a positive contribution requires \(\ell-1\mid B\)
and is at most \(c_\ell+v_\ell(B)\), in the notation of
Lemma~\ref{lem:rationalunits740}.

At a nonsplit prime let \(k\) be the quadratic completion of
\(K_0\) inside \(E_{x,w}\). For a norm-one unit \(u\in k\), its
conjugate is \(u^{-1}\). The two embeddings of \(k\) have
equally many extensions to \(E_{x,w}\), by separability.
Consequently the factors with exponent \(c\) cancel and the
character on these units is, up to inversion, \(u^B\) with
\[
                          B=\sum_{\tau\mid w}k_\tau>0 .
\]
The tame factor is trivial because both old level primes split.
If \(\ell\) is inert, take in \(\mathcal O_k\) an element
\(\alpha\) whose residue is outside \(\mathbb F_\ell\).
Then \(u=(1+\ell^{c_\ell}\alpha)/(1+\ell^{c_\ell}\bar\alpha)\)
has norm one and \(v_\ell(u-1)=c_\ell\). Its \(B\)-th power
gives the same upper bound \(c_\ell+v_\ell(B)\) by the binomial
argument of Lemma~\ref{lem:rationalunits740}. The Teichmüller
lift of a generator of the norm-one residue group has order
\(\ell+1\). Thus a positive contribution requires \(\ell+1\mid B\).
This argument also applies when \(E_{x,w}/k\) is ramified.

The only prime ramifying in \(K_0\) is \(d\). Here \(d\) is odd,
and \(v_d(\sqrt{-d})=1/2\). The norm-one unit
\[
     u=\frac{1+d\sqrt{-d}}{1-d\sqrt{-d}}
\]
satisfies \(v_d(u-1)=3/2\). For an element of depth greater
than \(1/(d-1)\), the linear binomial term in its \(d\)-th
power has strictly smaller valuation than all later terms.
Powers prime to \(d\) preserve depth. We obtain
\(h_w\le 3/2+v_d(B)\) at this last prime. The principal-unit
tests in all three cases also prove the asserted finiteness.

For each \(j\) write \(w_1,w_2\) for the two values of the
unshifted lift weight at its node. If \(n_i\) labels at the
completion use \(w_i\), their actual magnitudes sum to
\(S=n_1w_1+n_2w_2-H\), with the integer \(H\) between zero
and \(4r_x(n_1+n_2)\). Thus all the positive integers \(B\)
above occur among
\[
  S,\qquad \bigl|(n_1+n_2)c\ \pm S\bigr|,
            \qquad 1\le n_1+n_2\le n_x,\quad n_1,n_2\ge0.
\]
The common list of actual sums for all \(w,\ell,\iota,T\) at
this node has size \(O(m_xt_x(n_x+1)^2(1+4r_xn_x))\) and
maximum at most \(5n_xb_x\).
For a fixed \(B\) each of \(\ell-1\mid B\) and \(\ell+1\mid B\)
gives at most \(\tau(B)\) primes, all at most \(B+1\).
Counting the ramified prime separately for its constant part,
the weighted sum of the bounds for that \(B\) is at most
\[
 2\tau(B)\log(B+1)+\log2+\tfrac32\log d+\log B .
\]
The valuations of \(B\) sum to at most \(\log B\) even when
the union of these prime sets contains \(d\). This bound is
\(O_d((1+B)^{1/4})\) by the divisor estimate in the proof of
Theorem~\ref{thm:negativeinterval740}. For each prime there
are at most \(n_x\) places and two characters to count. Thus
after division by \(\sqrt{b_x}\) the per-node prime sum is
\(O_{q_0,d}(x^{62}b_x^{-1/4})\), including the independent
weight shifts exactly as in that proof. The scale inequality
\(b_x\ge x^{500}\) bounds this by \(O_{q_0,d}(x^{-63})\).
Summing over integer nodes \(x\ge V_0\) proves
\eqref{eq:alltorsion741} with the supremum inside the prime sum.
\end{proof}

\begin{corollary}[The remaining integral tests in a torus lift]
\leavevmode\label{cor:torustests741}
At every precision the local finite conditions on the rank-one
main trace in Theorem~\ref{thm:ribettrace720} can be tested by
\(n_x/2\) free cup-product coordinates across all \(w\mid\ell\),
with an additional normalized length at most
\(\sum_{w\mid\ell}h_w(T^\vee(1))\) as in \eqref{eq:freetests741}.
This additional length vanishes at both level primes. For each
scalar trace with lattice \(\mathcal O(-e)\), \(0\le e<t_x\),
there are \(n_x\) free coordinates. Summed over these scalar
weights and places, the weighted prime sum of the upper errors
at a fixed node is
\[
       O\bigl((n_xt_x+1)^2\log(2+n_xt_x)\bigr).
\]
After division by \(\sqrt{b_x}\), the sum of the scalar and
main upper errors, with the supremum over a whole interval
taken inside the prime sum, is \(O_{q_0,d}(V_0^{-10})\).
The same assertions hold for any common submodule of the
restricted raw trace classes.
\end{corollary}
\begin{proof}
For the torus character the dual has no rational invariants by
the local-unit computation. Its tangent dimension is the number
of embeddings at which the original character has nonnegative
jump. Exactly one of the two jumps \(z-1\pm k_\tau\) in each
conjugate pair is positive. Counting over all coefficient
completions therefore gives \(n_x/2\). The dimension formula
used in Proposition~\ref{prop:freetests741} makes this the total
rank of the free local test modules. In the coefficient sequence
the torsion in \(H^1(K,T')\) is \(H^0(K,A')\), because
\(H^0(K,T'[1/\ell])=0\). The reduced torsion has normalized
length at most \(h_w(T')\), proving the claims for the main trace.

For the scalar trace the dual is \(\mathcal O(e+1)\). Its
jumps are negative at all embeddings and it has no rational
invariants. The same dimension formula gives \(n_x\) free
generators over all \(w\). On rational units in \(E_{x,w}\)
the cyclotomic power acts by the power of the local norm with
absolute exponent \(B=(e+1)[E_{x,w}:\QQ_\ell]\). Thus rational
principal and residue units bound the normalized torsion length
by \(\mathbf 1_{\ell-1\mid B}(c_\ell+v_\ell(B))\). Put
\(Y=n_xt_x\). All these \(B\)'s lie between \(1\) and \(Y\).
A positive contribution requires \(\ell\le Y+1\), and the
largest possible \(v_\ell(B)\) is at most \(v_\ell(Y!)\).
Consequently, before counting the choices of place and scalar
weight, the weighted sum of simultaneous upper bounds is at most
\(\log((Y+1)!)+\log2+\log(Y!)\).
There are at most \(Y\) place and weight pairs per prime,
giving the asserted per-node estimate. Using \(n_x\ll_{q_0,d}x^{12}\)
and \(t_x\ll x^5\) bounds it by \(O_{q_0,d}(x^{35})\).
Divide by \(\sqrt{b_x}\), use \(b_x\ge x^{400}\), and sum over
the nodes together with \eqref{eq:alltorsion741}. The joint-map
assertion in Proposition~\ref{prop:freetests741} applies even
when the restricted classes at different places are dependent.
\end{proof}

The surviving quantity in \eqref{eq:freetests741} is the length
of the image of the trace classes under the free cup products.
Its rank bound alone allows length proportional to the precision.
The period estimate requires a bound for this image independent
of precision, with the growth in \(x\) prescribed in
Section~\ref{sec:missing}.
\clearpage

\subsection{Trace extensions at the two original level primes}
\label{sec:tameends743}

At a level prime the torus character has a single filtration
jump at each completion. Its sign changes at the conjugate
completion. We use the ends of the interval of jumps of the
auxiliary representations to determine the local classes at
every congruence precision. We use the filtration convention
in which \(\kappa\), the cyclotomic character, has jump \(-1\).

\begin{lemma}[Extensions at an end of a small-height interval]
\leavevmode\label{lem:tameends743}
Let \(K=\QQ_p\), \(p>2\), and let \(K'/K\) be a totally ramified
Galois extension of degree \(q\) prime to \(p\). Fix nonnegative
integers \(a,B,h\) with \(0\le h\le B\) and \(q(a+B)<p-1\).
Let \(T\) be a free line over a finite coefficient ring of integers
\(\mathcal O\), with uniformizer \(\pi\). Suppose its character
on \(G_{K'}\) is \(\xi\kappa^{-h}\) for an unramified character
\(\xi\). For \(h=0\) take \(T=\mathcal O\) with trivial action.
Let \(z:G_K\to T/\pi^jT\) be a continuous cocycle, \(j\ge1\).
Suppose its extension of \(\mathcal O/\pi^j\) by \(T/\pi^jT\)
is, on restriction to \(G_{K'}\), a subquotient of a lattice in
a crystalline representation with filtration jumps in \([-a,B]\).
Here and below this condition concerns the underlying \(p\)-adic
representation at all embeddings. If \(a=0\) or \(h=B\), then
the class of \(z\) is zero for \(h>0\). For \(h=0\), the cocycle
is unramified and lifts to an integral unramified cocycle.
\end{lemma}
\begin{proof}
First take \(j=1\) and write \(k=\mathcal O/\pi\). Over \(K'\)
take the \(\mathbb F_p\)-linear dual of the extension and twist by
\(\bar\kappa^a\), where the bar denotes reduction modulo \(p\).
The dual of a finite subquotient of a lattice is a subquotient of
its dual lattice: for two nested full lattices take their dual
lattices in reverse order using the pairing into
\(\QQ_p/\ZZ_p\). Subquotients inside smaller finite groups follow
by taking preimages. The twist is the reduction of the corresponding
integral Tate twist. Thus the resulting module \(U\) is torsion
crystalline of weights in \([0,H]\), \(H=a+B\), in the
Breuil--Kisin convention in which \(\kappa\) has weight \(1\).
It is an extension with subobject \(k^*(a)\) and quotient
\((k(\bar\xi))^*(a+h)\), where the star means the
\(\mathbb F_p\)-linear dual and \(k^*\) has trivial action.

Choose a uniformizer of \(K'\) and a compatible sequence of its
\(p\)-power roots. Write \(G_\infty\) for the absolute Galois group
of the resulting union of fields. We recall the integral comparison
facts we use. The restriction to \(G_\infty\) of a torsion crystalline
module killed by \(p\) of weights in \([0,H]\) has a contravariant
Kisin model of height at most \(H\). After inverting the parameter
\(u\), Fontaine's equivalence with \'etale Frobenius modules is
exact. Here the residue field of \(K'\) is \(\mathbb F_p\), so the
model is a free \(\mathbb F_p[[u]]\)-lattice with semilinear
Frobenius \(\varphi(u)=u^p\). Its linearized Frobenius matrix
\(F\) is integral and \(u^{qH}F^{-1}\) is integral, since the
Eisenstein polynomial reduces to \(u^q\). In this convention an
unramified module has a constant invertible Frobenius matrix.
Tensoring by \(\bar\kappa^s\), \(s\ge0\), multiplies the matrix by
a scalar unit times \(u^{qs}\). This is the reduction of the
rank-one Tate model with multiplier \(E(u)/c_0\), where
\(c_0=E(0)/p\). The formula for an unramified module follows by
extending its constant Frobenius module from \(\mathbb F_p\) to
\(\mathbb F_p[[u]]\).
~\cite[\S\S2.1, 2.4, 3.3]{III-II-cltame743}\cite[\S2]{III-II-ozekitame743}

In this range the sequence of the three models is exact already
over \(\mathbb F_p[[u]]\). We give the lattice argument. A matrix
\(A\) for a map after inverting \(u\) satisfies
\(AF_1=F_2\varphi(A)\). If its largest pole has order \(s>0\),
then \(\varphi(A)=F_2^{-1}AF_1\) has largest pole at most \(s+qH\),
but also exactly \(ps\). This contradicts \(qH<p-1\). In particular
two height-\(H\) lattices in the same generic module agree, by applying
the assertion to the identity in both directions. The image of a
map of models has height at most \(H\), by applying the map to the
linearized Frobenius inclusions for its domain. Its kernel also
has height at most \(H\): the inverse linearized Frobenius of a
kernel vector multiplied by \(u^{qH}\) is integral and belongs to
the Frobenius pullback of the kernel. Indeed its image in the
pullback of the target vanishes by injectivity of linearized
Frobenius and flatness of \(u\mapsto u^p\). Images and kernels
are finite free over the power series discrete valuation ring.
Exactness of the generic sequence and uniqueness of the lattices
now give the asserted integral exactness.

Choose bases of the end models and lift the basis of the quotient
of the middle model. Contravariance puts the positive additional
weight \(h\) in the submodel. In these bases Frobenius has matrix
\[
    \begin{pmatrix}
       u^{q(a+h)}D & X(u)\\
          0        & u^{qa}C
    \end{pmatrix},
    \qquad C,D\in\mathrm{GL}_{[k:\mathbb F_p]}(\mathbb F_p).
 \tag{11}\label{eq:tameends743}
\]
All scalar unit factors of the Tate model are included in \(C,D\).
The upper right block of the inverse matrix is
\(-u^{-q(2a+h)}D^{-1}X(u)C^{-1}\). Its product with \(u^{qH}\)
is integral. If \(h=B\), this says that \(u^{qa}\) divides every
entry of \(X\). If \(a=0\), this divisibility just says \(X\)
is integral. Consequently the matrix after removing the twist
\(\bar\kappa^a\) still has integral entries. Its two diagonal
blocks are \(u^{qh}D_1\) and \(C_1\), with \(D_1,C_1\) constant
and invertible. Denote its upper right block by \(X_1\).

For \(h>0\), change the lifted quotient basis by a matrix \(Y\)
in the submodel. The equation for zero upper right block is
\[
       Y=X_1 C_1^{-1}+u^{qh}D_1\varphi(Y)C_1^{-1}.
\]
The operator in the second term sends matrices divisible by \(u^m\)
to matrices divisible by \(u^{qh+pm}\). Iterating from zero therefore
converges in \(\mathbb F_p[[u]]\) to a solution. Retwisting gives
a splitting over \(G_\infty\) of the sequence defining \(U\).
Ozeki's full faithfulness theorem for the restriction of torsion
crystalline representations gives the same splitting over \(G_{K'}\):
the source and target of the section are torsion crystalline in
\([0,H]\) and \(qH<p-1\). Dualizing back gives a splitting of the
original sequence over \(\mathbb F_p\). The invariant lift of
\(1\in k\) then yields a \(k\)-linear section by scalar
multiplication. Restriction followed by corestriction multiplies
\(H^1(K,T/\pi T)\) by \(q\), which is invertible. Thus the
original class is zero over \(K\).~\cite[Theorem~1]{III-II-ozekitame743}

For \(h=0\), \(T\) is trivial. After removing the twist we may
use identity matrices at both ends. Changing basis replaces the
off-diagonal block by \(X_1+\varphi(Y)-Y\). Take
\(Y=\sum_{i\ge0}\varphi^i(X_1-X_1(0))\); this converges and
leaves the constant block \(X_1(0)\). A constant invertible
Frobenius matrix gives an unramified representation of \(G_\infty\)
obtained by restriction from \(G_{K'}\), since the uniformizer tower
is totally ramified. Its finite unramified representation lifts to
a weight-zero crystalline lattice by lifting its Frobenius matrix
to \(\mathrm{GL}_{2[k:\mathbb F_p]}(\ZZ_p)\). Retwist by
\(\bar\kappa^a\). The resulting lift has weight \(a\in[0,H]\),
so Ozeki's theorem lifts the isomorphism over \(G_\infty\) to
\(G_{K'}\). Undoing the twist and the dual, \(z\) kills inertia
over \(K'\). Its values on inertia over \(K\) factor through a
group of order \(q\), which has no nonzero homomorphism to the
additive group of \(k\). An integral lift is obtained by lifting
its Frobenius value and extending continuously from
\(\widehat{\ZZ}\) through \(\ZZ_p\).

For \(j>1\), reduce modulo \(\pi\). For \(h>0\), subtract a
coboundary given by any lift of a splitting element. For \(h=0\),
subtract the reduction of the integral unramified cocycle just
constructed. In both cases the resulting extension still has the
required crystalline presentation: changing a section gives an
isomorphic extension, and taking a difference of extensions uses
the diagonal pullback and the subtraction pushout of their direct
sum. The subquotient condition is stable under these operations.
The new cocycle takes values in \(\pi(T/\pi^jT)\). In its extension
restrict the first summand to this submodule, then divide out the
fixed submodule \(0\oplus\pi^{j-1}(\mathcal O/\pi^j)\).
The quotient is the extension for the cocycle divided by \(\pi\)
at precision \(j-1\). Induction gives a coboundary, or an integral
unramified lift, respectively. Multiplying back by \(\pi\) and
restoring the subtracted term proves both assertions.
\end{proof}

\begin{theorem}[Local trace classes at both original levels]
\leavevmode\label{thm:oldtorustraces743}
Fix one of the shorter intervals and its nonzero designated torus
lifts as in Theorem~\ref{thm:projectedcups742}. Use the separating
systems in the image of the rank-one marked idempotent, at any of
the permitted auxiliary levels. At each \(p=p_{x,0},p_{x,1}\)
all the main trace classes belong to the integral finite local
condition at every \(w\mid p\) and every precision. Put
\(t'=2z-1\). At a completion where the main jump is positive the
same assertion holds for all the scalar weights \(0\le h<t'\).
At its conjugate it holds for \(h=t'-1\) as well.
\end{theorem}
\begin{proof}
We verify the common field, tame type and jumps. Both primes split
completely in \(E=E_x\), so \(E_w=\QQ_p\). Put \(Q=(q_0d)^{h_x+6}\),
where \(h_x\) is the level exponent of the node. We have
\(Q^2\mid p-1\) and \(p\ge b_x^{10}\). Use the degree \(Q\)
subextension \(K'\) of \(\QQ_p(\mu_p)/\QQ_p\). It is totally
ramified and cyclic. Its inertia subgroup kills every character
of \(\mathbb F_p^\times\) of order dividing \(Q\), by local
reciprocity and the cyclotomic action on \(\mu_p\).

At a chosen old place the target level fixes a specified residue
vector. Its line stabilizer is generated by this vector stabilizer
and the diagonal matrices scaling that vector. In the marked
projector the normalized action of the latter group is trivial:
the source label \(X-1\) is the trivial character of
\(\mathrm{GL}_1(\mathbb F_p)\). In the translation action this
leaves the inverse of the prescribed residue character on that
line, and trivial action on the complementary block. Its tame
character is a nontrivial power \(\omega_w^{\pm j_0}\) of order
dividing \(Q\), as in Theorem~\ref{thm:rankoneselmer737}.
The auxiliary twist splits above \(p\) and the later trace places
are away from it. The characteristic-zero classification of the
nontrivial line-character space used in
Theorem~\ref{thm:tametransfer721} gives one tame line and an
unramified spherical complement, with zero monodromy. The local
Weil parameter thus has inertia of exponent dividing \(Q\).
The opposite place has the dual type with the unramified norm
twist. The normalizing splitting character has unramified smooth
part here, so it preserves this exponent and zero monodromy.

We need the compatibility at the coefficient characteristic itself.
In the base change of Theorem~\ref{thm:splitattachment722}, the
normalized cuspidal factors are regular algebraic essentially
conjugate self-dual, as verified in
Corollary~\ref{cor:scalarfl724}. The parameter of the isobaric
sum is the direct sum of their parameters with their indicated
norm twists. Every summand therefore has zero monodromy.
Theorem~1.1 of Barnet-Lamb, Gee, Geraghty and Taylor identifies
the semisimple Weil parameter at \(w\mid p\) and bounds the
monodromy Jordan blocks by those of the automorphic parameter.
The upper bound has only blocks of length one, so the Galois
monodromy vanishes. The restriction of the finite Weil inertia
to \(I_{K'}\) is trivial since its characters have order dividing
\(Q\). A de Rham representation with these properties is crystalline
over \(K'\): in a larger semistable period module the vanishing of
monodromy gives crystalline periods, trivial inertia descends
them through the ramified part, and unramified descent preserves
their full dimension by semilinear Galois descent. This reasoning
applies to all the trace matrices in the projected space.
~\cite[Theorem~1.1]{III-II-blggttame743}

All the systems at one such level have the same algebraic
coefficient representation at infinity and the same normalizing
character. Their labelled Hodge multisets are identical by the
weight formula in Corollary~\ref{cor:scalarfl724}. To evaluate
this common multiset use the nonzero base lift itself: its
normalized representation is the sum of \(T\) and the lines
\(\mathcal O(-h)\), \(0\le h<t'\), by the parameter calculation
in Theorem~\ref{thm:rankoneselmer737}. Hence the jumps at \(w\) are
\[
              0,1,\ldots,t'-1,\qquad J_w=z-1\pm k_\tau .
\]
The choice of coefficient conjugate fixing \(\QQ_p\) does not
change this multiset at a completely split place, since the image
of the local embedding of \(E\) is in \(\QQ_p\). Restriction to
\(K'\) repeats the same jumps. Set
\(a=\max(0,-J_w)\), \(B=\max(t'-1,J_w)\). From
\(t_0<k_\tau\le4b_x\) and \(t_0\le b_x\) we have
\(a+B\le 6b_x\). Therefore \(Q(a+B)<p-1\) for the indicated
intervals with sufficiently large initial node.

By Lemma~\ref{lem:graphrealization724} each local trace extension
is a subquotient of an integral lattice with rational constituents
among these systems, the target and the trivial line. All are
crystalline over this single \(K'\) in \([-a,B]\). If \(J_w>0\)
then \(a=0\). On \(G_{K'}\) the main line is an unramified
character times \(\kappa^{-J_w}\), since on \(\QQ_p^\times\) its
unit character is the product of its smooth factor and the
algebraic character of jump \(J_w\). Lemma~\ref{lem:tameends743}
applies to this line and to every scalar line. For \(h=0\) the
unramified lifts are finite classes because unramified extensions
of the trivial representation are crystalline. When \(J_w<0\),
the main finite-condition quotient vanishes by
Theorem~\ref{thm:negativeinterval740}; its dual has no residual
invariants at this level prime. Now \(B=t'-1\), so the upper
endpoint in Lemma~\ref{lem:tameends743} gives the assertion for
the last scalar line at the same completion.
\end{proof}

Thus at either old coefficient prime the main free cup products
in \eqref{eq:projectedcups742} vanish for the marked separating
families. The only possible scalar coordinates there occur on the
negative side at \(0\le h<t'-1\). There are at most
\([E_x:\QQ](t'-1)/2\) of them, because the prime splits completely
and one completion in each conjugate pair has negative main jump.
The joint images are still taken at the original precision on
the original field. The argument gives no bound for their lengths
from the number of coordinates alone.
\clearpage

\subsection{The marked character as a subrepresentation or quotient}
\label{sec:markedline747}
At an original level prime \(p\), let \(K=E_{x,w}=\QQ_p\).
The normalized representations used for the torus pairing are
crystalline over the totally ramified extension \(F/K\) of degree
\(Q\) in Theorem~\ref{thm:oldtorustraces743}. Their smooth
inertial type, namely the inertia action on potentially crystalline
periods, is the sum of a trivial representation and one copy
of a character \(\xi\). This character has order greater than two.
Indeed its order is \(Q/\gcd(Q,j)\) for \(1\le j\le m_x\), as in
Theorem~\ref{thm:rankoneselmer737}. We prove a full-precision
local condition when the character of this type occurs on a
rational line at either end of the representation. The remaining
representation can have noncommutative image on every open subgroup.

\begin{lemma}[Removing outer constituents from a finite subquotient]
\leavevmode\label{lem:discardouter747}
Let \(V\) be a characteristic-zero representation of a group on
a finite-dimensional space over a local coefficient field \(L\),
with a stable lattice and a stable filtration
\(0\subset V_1\subset V_2\subset V\). Fix a collection \(S\)
of simple residual representations over its residue field \(k\).
Suppose no reduction of a stable lattice in \(V_1\) or \(V/V_2\)
has a composition factor in \(S\).
A finite subquotient \(U\) of a lattice in \(V\), killed by
\(\pi^a\) and with all composition factors in \(S\), is a
subquotient of the reduction modulo \(\pi^a\) of a stable
lattice in \(V_2/V_1\).
\end{lemma}
\begin{proof}
Write \(U=N/N'\) with \(N,N'\) stable submodules of a lattice.
By adding a common sufficiently deep multiple of that lattice
to \(N\) and \(N'\) one may take them to be full lattices because \(N\cap\pi^b\Lambda\subset\pi^aN\subset N'\) for large
\(b\), by elementary divisors. The new quotient is the original
one. There is now a surjection \(N/\pi^aN\to U\).
Intersect \(N\) with \(V_1,V_2\). The resulting filtration of
\(N\) has free successive quotients, so reduces exactly modulo
\(\pi^a\). The outer factors have no simple in \(S\), by filtering
their reductions by powers of \(\pi\). The bottom factor therefore
has zero image in \(U\). The image of the middle step is all of
\(U\), since its cokernel would be a quotient of the top factor.
Thus the map from the middle step induces a surjection from the
middle graded lattice modulo \(\pi^a\) to \(U\).
\end{proof}

\begin{proposition}[Full scalar local conditions with a marked line]
\leavevmode\label{prop:markedFL747}
Fix \(K\) as above. Suppose that each member \(V_y\) of a finite
local trace family of the specified type and Hodge multiset has
a \(G_K\)-stable line or a one-dimensional quotient of smooth
inertial type \(\xi\), after extending coefficients if necessary.
For every \(0\le h<t'\) and every precision, all the local raw
trace cocycles with target \(\mathcal O(-h)/\pi^a\) satisfy
the integral finite local condition. For \(h>0\) their classes
vanish and for \(h=0\) they are unramified.
\end{proposition}
\begin{proof}
A line of this type has inertia character \(\xi\kappa^{-d}\),
where \(d\) is one of the jumps of \(V_y\), with the common
identification of smooth characters by reciprocity. This is
the local formula for a potentially crystalline character used
in Lemma~\ref{lem:abelianunits744}. Its reduction differs from
the trivial character and from \(\bar\kappa^{-h}\). To see this,
evaluate at a prime-to-\(p\) lift \(s\) of a tame generator.
An equality with \(\bar\kappa^{-e}\), \(e=0\) or \(h\), would give
\(p-1\mid Q(d-e)\) after raising to \(Q\). The common jump
interval \([A,B]=[\min(0,J_w),\max(t'-1,J_w)]\) has
\(Q(B-A)<p-1\). Thus \(d=e\), which would make the nontrivial
prime-to-\(p\) value \(\xi(s)\) trivial modulo \(\pi\).

After taking the quotient by a marked subline, or the kernel of
a marked quotient, the remaining representation is crystalline
over \(K\). Indeed it is crystalline over \(F\), and exactness
of crystalline periods there removes the only nontrivial smooth
inertia character. The extension \(F/K\) is totally ramified,
so its maximal unramified subfield equals that of \(K\).
The full period space over \(F\) is now fixed by
\(\operatorname{Gal}(F/K)\), and supplies all required periods
over \(K\). The jumps of the remaining representation lie in
\([A,B]\) by filtered exactness.

Use the graph construction of Lemma~\ref{lem:graphrealization724}
for a raw cocycle \(c\) of weight \(h\). The extension of
\(\mathcal O/\pi^a\) by \(\mathcal O(-h)/\pi^a\) is a finite
subquotient of a lattice whose rational representation is a
direct sum of copies of the members, the target and the trivial
line. This construction works with any \(\pi^a\) by using that
ideal in its graph and pushout. Give this direct sum a filtration
of length three. Put the distinguished sublines at the bottom
in members with sublines, and put the distinguished quotients
at the top in members with quotients. In each case put the
remaining representation in the middle, together with the
target and the trivial line. Taking the direct sum of these
filtrations gives one crystalline middle representation over
\(K\), with the same bound on its jumps. Both outer terms are
direct sums of lines with the residual exclusion just proved.
All stable lattices in those direct sums have the same residual
composition factors. For example one may filter each lattice
by its intersections with a flag of these lines and reduce the
saturated filtration. The residual composition factors of the
extension for \(c\) all belong to
\(\{k,k(-h)\}\), by its defining exact sequence. Apply
Lemma~\ref{lem:discardouter747} to see that this extension is
torsion Fontaine--Laffaille over the absolutely unramified field
\(K\). Here \(B-A\le p-2\), and torsion objects of the common
Fontaine--Laffaille interval are closed under subquotients.

For precision, the lifting argument in
Theorem~\ref{thm:fllift724} applies once this torsion presentation
and the two ends are given. Twist the filtration interval to
\([0,B-A]\), with the quotient a rank-one object of pure degree
\(-A\). In the filtered module choose a lift of its generator
in that filtration degree, by strictness. Its divided Frobenius
differs from itself by a vector in the reduced kernel. Lifting
that vector to the integral kernel with the same split filtration
gives an integral Fontaine--Laffaille extension reducing to the
original one with both ends identified. Undoing the twist puts
the class of \(c\) in the integral finite condition. This argument
uses the same equations for \(\pi^a\) since Frobenius fixes the
coefficient ring. At \(K=\QQ_p\) this condition is zero for
\(h>0\) and the full unramified subgroup for \(h=0\), as in
Theorem~\ref{thm:oldcolumn745}. All assertions descend through
a coefficient extension by an \(\mathcal O\)-linear projection
from the enlarged ring onto \(\mathcal O\) sending one to one.
\end{proof}

\begin{lemma}[A single nontrivial tame eigenvalue]
\leavevmode\label{lem:singleinertia747}
Let \(V\) have the local Hodge multiset and the smooth type above.
If \(V\) is potentially abelian then its marked line is a direct
summand over \(G_K\), after extending coefficients. In general,
an absolutely irreducible local constituent of dimension greater
than one whose smooth inertia contains \(\xi\) stays absolutely
irreducible on every open subgroup of \(G_K\).
\end{lemma}
\begin{proof}
We first record the restriction on permutations of summands.
Suppose \(V=\bigoplus_i V_i\) and \(G_K\) permutes these subspaces.
Choose a finite normal extension \(L/K\) containing \(F\) for
which every summand is \(G_L\)-stable. All are crystalline over
\(L\). Their period spaces form a direct sum on which inertia
permutes the summands in the same way, by functoriality of
periods. In passing to a coefficient vector space one can
use any embedding of \(L_0\); inertia fixes \(L_0\). On a cycle
of length \(j\) for an inertia element the eigenvalues on these
blocks include all \(j\)-th roots of any eigenvalue of its
\(j\)-th power on the first block. This follows, for instance,
by following an eigenvector round the cycle. Wild inertia acts
as the identity on these periods. On a generator of tame inertia
the only eigenvalues available are \(1\) and \(\xi(s)\), and
\(\xi(s)\ne-1\). Thus no cycle of length \(j\ge2\) is possible.
All of inertia stabilizes the blocks. In each Frobenius orbit
of blocks the dimensions of their nontrivial smooth inertia
parts are equal. Conjugation and semilinear transport preserve
the dimension of inertia invariants. Moreover the multiplicity
can be computed in any coefficient embedding of \(L_0\), since
the bijective crystalline Frobenius identifies the factors
and commutes with inertia. As the total nontrivial dimension
is one, the orbit which carries it consists of a single block.

If \(V\) is potentially abelian, choose \(L\) normal such
that the matrices of \(G_L\) commute. Their joint generalized
character spaces are permuted by \(G_K\). Each space has all
composition factors equal to its one character. At every label
of \(L\) the jumps in \(V\) are pairwise distinct, by the scalar
string and the separated main jump of
Theorem~\ref{thm:oldtorustraces743}. Filtered exactness therefore
forces each generalized character space to have dimension one.
The preceding orbit argument singles out one line fixed by
all of \(G_K\). The sum of the other orbits is a stable complement.

Now take an absolutely irreducible constituent \(W\) containing
\(\xi\) and restrict to an open normal subgroup \(H\). Work
over an algebraic closure of the coefficient field. The algebra
spanned by \(H\) in \(\operatorname{End}W\) is finite-dimensional
and normalized by \(G_K\). Its radical is nilpotent and its image
on \(W\) is \(G_K\)-stable. Irreducibility and nilpotence force
that image to be zero. Thus \(W|_H\) is semisimple, and \(G_K\)
permutes its isotypic spaces transitively. The finitely many
summands and a decomposition into simples can be defined over
a finite coefficient extension by choosing bases. Regularity
of the jumps on restriction to \(H\) forces multiplicity one
of each simple, since every repetition would repeat a jump after
a further extension containing \(F\). The permutation argument
applies to this transitive collection in \(W\), whose nontrivial
smooth inertia part still has dimension one. There is therefore
one block and \(W|_H\) is absolutely irreducible. Every open
subgroup contains an open normal subgroup, which proves the claim.
\end{proof}

\begin{theorem}[The old-prime trace contribution with a marked line]
\leavevmode\label{thm:markedloss747}
At either original level prime of Theorem~\ref{thm:projectedcups742},
the quantity \(B_{\rm old}(D)\) of Theorem~\ref{thm:oldcolumn745}
is zero whenever every member at each negative completion has
its marked character on a subline or on a quotient line. This
includes mixtures of such members with potentially abelian members.
It also suffices for a marked character to occur as a rank-one
composition factor with the smallest or the largest labelled
jump of the member. On this subfamily the full old-prime sum
of \(p(D)/\sqrt{b_x}\), with the same prime weights and suprema
as in Theorem~\ref{thm:oldcolumn745}, is
\(O_{q_0,d}(V_0^{-10})\).

For an arbitrary family at an old prime, all surviving joint
cup coordinates are determined by local trace evaluations on
integral tuples whose components vanish on every member meeting
these line conditions at the place of evaluation.
\end{theorem}
\begin{proof}
If a marked line \(\chi\) is a composition factor, write
\(V_-\subset V_1\subset V_y\) with \(V_1/V_-=\chi\), and
write \(V_+=V_y/V_1\). If the jump of \(\chi\) is smallest,
all jumps of \(V_-\chi^{-1}\) are strictly positive. Over \(F\)
this crystalline representation has zero de Rham tangent space
and zero invariants, hence zero \(H^1_f\). The extension of
\(\chi\) by \(V_-\) is crystalline there, so splits there.
Restriction on rational \(H^1\) is injective by corestriction
and \([F:K]\ne0\). It splits over \(K\), providing a subline of
\(V_y\). If the jump is largest, apply the same argument to
\(\operatorname{Hom}(V_+,\chi)\) and the extension
\(0\to\chi\to V_y/V_-\to V_+\to0\). A splitting provides a
quotient line of \(V_y\). A zero-dimensional outer term in
either argument already gives the required line.

Apply Proposition~\ref{prop:markedFL747} at each negative
completion. It applies simultaneously to all the indicated
members, and Lemma~\ref{lem:singleinertia747} includes the
potentially abelian ones. All remaining coordinates in
Theorem~\ref{thm:oldcolumn745} are zero. Restricting its prime
sum proves the estimate for this subfamily.

Finally at one completion let \(\mathcal N_w\) be the span of the
local matrix differences for the full family, and project it
to the matrix components of the members with the stated line
property. The image is precisely their local difference lattice,
since the same group elements generate both. Two local trace
functionals of the same weight which agree on the kernel of
this projection have difference factoring integrally through
its image. Their cocycle difference obeys all the relations of
that projected family, so Proposition~\ref{prop:markedFL747}
puts its class in the finite condition. For an empty projected
family the assertion follows directly from the definition.
Thus each surviving cup coordinate is determined on the kernel,
including when the original global classes have dependent
restrictions. This proves the last assertion.
\end{proof}

A local representation may have its marked smooth type in an
irreducible constituent of rank greater than one. The existence
of a one-dimensional eigenspace on its period module alone
does not supply a rational Galois subline. The integral estimates
for such constituents, and for interior marked factors outside
the line conditions above, remain necessary for the unrestricted
prime sum.
\clearpage

\subsection{The weight-zero trace at the original level primes}
\label{sec:logold748}
We bound the weight-zero scalar trace for the entire marked
family at either original level prime. The bound is independent
of the distance to the negative main jump and of the congruence
precision. The proof compares two descriptions of the ramified
additive character on the uniformizer tower.

Let \(p\ge5\), \(K=\QQ_p\), and let \(F/K\) be a totally ramified
cyclic tame extension of degree \(q\mid p-1\). Take a uniformizer
of \(F\) with Eisenstein polynomial \(E(u)=u^q-cp\), where
\(c\in\mathbb Z_p^\times\) is a Teichm\"uller unit. Here is one
construction. Start with a uniformizer, and project it to its
tame character on the first graded piece of the maximal ideal
by averaging against the inverse Teichm\"uller lift of that
character. The projection is unchanged modulo the square of
the maximal ideal. Its \(q\)-th power is fixed by Galois, with
\(K\)-valuation one. The principal-unit part of its quotient by
\(p\) has a \(q\)-th root in \(K\), which can be divided out of
the uniformizer. Write \(F_\infty\) for the field obtained by
adjoining compatible \(p\)-power roots of this uniformizer.

For \(n\ge1\) put \(S_n=(\mathbb Z/p^n)[[u]]\),
\(D_n=S_n[1/u]\) and \(c_0=E(0)/p=-c\).
Frobenius \(\varphi\) fixes coefficients and sends \(u\) to
\(u^p\). We use the integral comparison for \(G_{F_\infty}\)
in the convention of Lemma~\ref{lem:tameends743}. Fontaine's
equivalence identifies finite representations with torsion
\'etale Frobenius modules over the \(p\)-adic completion of
\(\mathbb Z_p[[u]][1/u]\), whose reduction is \(D_n\). In the
contravariant convention the Tate character \(\kappa\) has the
rank-one model with multiplier \(E(u)/c_0\).
A finite subquotient of a lattice in a crystalline representation
of weights in \([0,H]\), in the convention giving \(\kappa\)
weight one, has a finite-height model for its restriction to
\(G_{F_\infty}\). Such a model \(M\) is a finitely generated
\(S_n\)-submodule spanning the \'etale module after inverting
\(u\), and the image \(\Phi(\varphi^*M)\) under the linearized
Frobenius satisfies
\[
                 E^H M\ \subset\ \Phi(\varphi^*M)\ \subset\ M.
\]
These are the torsion Kisin comparison and rank-one Tate model
used here.~\cite[\S\S2--3, especially Example~2.3.5]{III-II-liutorsion748}

\begin{lemma}[Removing the common Tate twist at finite precision]
\leavevmode\label{lem:shortpole748}
Let \(U\) be an extension of the trivial \(\mathbb Z/p^n\) by
itself over \(G_F\). Suppose it is a finite subquotient of a
crystalline lattice in the filtration interval \([-a,B]\),
in which \(\kappa\) has filtration jump \(-1\). Put \(H=a+B\),
\(D=\min(a,B)\) and assume \(q(H+n-1)<p-1\).
The \'etale Frobenius module on \(F_\infty\) of its dual in
the contravariant convention has an extension matrix with
diagonal entries one and upper right entry in \(E^{-D}S_n\).
The entries refer to the fixed trivializations of the two ends.
\end{lemma}
\begin{proof}
Two height-\(H\) models for the same \'etale module coincide
under this inequality. Indeed, if \(M,N\) are such models,
let \(s\) be the smallest nonnegative integer with \(u^sM\subset N\).
If \(s>0\), the smallest exponent for the corresponding inclusion
after Frobenius pullback is \(ps\). To see that \(ps-1\) fails,
choose \(m\in M\) with \(u^{s-1}m\notin N\). The coefficient
of \(u^{p-1}\) in \(u^{ps-1}(1\otimes m)\), using the basis
\(1,u,\ldots,u^{p-1}\) over \(\varphi(S_n)\), detects
\(u^{s-1}m\). Since \(u^{q(H+n-1)}=(E+cp)^{H+n-1}\) is
divisible by \(E^H\) modulo \(p^n\), the height inclusions give
\[
 u^{s+q(H+n-1)}\Phi(\varphi^*M)
       \subset u^{q(H+n-1)}N\subset\Phi(\varphi^*N).
\]
Thus \((p-1)s\le q(H+n-1)<p-1\), a contradiction. Reversing
the inclusion proves uniqueness.

Dualize \(U\) and twist by \(\kappa^a\). This is a torsion
crystalline subquotient of weights in \([0,H]\) with both end
characters \(\kappa^a\bmod p^n\), just as in the proof of
Lemma~\ref{lem:tameends743}. Take a model \(M\) of height \(H\).
Intersect it with the first end module after inverting \(u\)
and take its image in the other end module. Both models have
height \(H\). For the image this follows by projecting the
two inclusions. For the intersection, the preimage of \(E^Hm\)
under linearized Frobenius in \(\varphi^*M\) maps to zero in
the generic quotient. Flatness of \(\varphi\), and injectivity
of linearized Frobenius on the generic quotient, put this
preimage in the pullback of the intersection. The intersection
and image span their generic modules because localization is
exact. By uniqueness they are the standard free rank-one
models with multiplier \(t=(E/c_0)^a\). Lifting the basis of
the quotient now gives a basis of \(M\), in which
\[
             \Phi=\begin{pmatrix}t&X\\0&t\end{pmatrix},
             \qquad X\in S_n,\qquad E^H t^{-2}X\in S_n .
\]
The last inclusion uses the upper right entry of \(E^H\Phi^{-1}\).
Undoing the Tate twist replaces \(X\) by \(X/t\), which lies in
\(E^{-a}S_n\cap E^{-B}S_n=E^{-D}S_n\). All divisions take
place in \(D_n\), where \(E\) is invertible.
\end{proof}

\begin{lemma}[The cyclotomic logarithm in the Laurent module]
\leavevmode\label{lem:logmatrix748}
Set \(\mathcal L(\sigma)=p^{-1}\log(\kappa(\sigma)^{p-1})\).
For its additive extension modulo \(p^n\), a Frobenius entry on
\(F_\infty\), in the convention of Lemma~\ref{lem:shortpole748},
is, up to multiplication by a constant unit and a change
\((\varphi-1)Y\), the reduction of
\[
 C(u)=\frac{p-1}{p}\log(1-cp\,u^{-q})
      =-(p-1)\sum_{i\ge1}\frac{c^i p^{i-1}}{i}\,u^{-qi}.
\]
\end{lemma}
\begin{proof}
For the calculation compose the contravariant equivalence with
duality. This gives a covariant exact equivalence sending
\(\kappa^{-1}\) to the line with multiplier \(E/c_0\).
Take \(j\ge n\), \(m=(p-1)p^j\), and the character
\(\theta=\kappa^{-m}\). It is trivial modulo \(p^{j+1}\).
The connecting class of the invariant generator modulo \(p^{j+1}\)
in this line has cocycle \((\theta(\sigma)-1)/p^{j+1}\).
Modulo \(p^n\) it represents \(-\mathcal L\), by applying the exponential
to \(-p^j\log(\kappa^{p-1})\) and dividing by \(p^{j+1}\).
Terms after the linear one are zero modulo \(p^n\).

The Frobenius multiplier of the line is \((E/c_0)^m\).
In its \'etale module change basis by \(u^{-q p^j}\).
The multiplier becomes \(f=(1-cp u^{-q})^m\), since the
constant unit \(c_0\) has order dividing \(p-1\).
It is one modulo \(p^{j+1}\). Lifting the fixed generator in
this basis and taking its Frobenius difference gives
\((f-1)/p^{j+1}\). The connecting extension in either category is obtained by
pulling back the reduction map of the line along the
map from the trivial object carrying its generator to the
invariant generator, identifying the kernel by division by
\(p^{j+1}\). The two choices of generator differ by a constant
unit modulo \(p^{j+1}\). Indeed Frobenius-fixed Laurent series
over \(\mathbb F_p\) are constants by comparing coefficients
under \(u\mapsto u^p\), and reduction and division by \(p\)
prove the same fact over \(\mathbb Z/p^{j+1}\). A change of
lift adds a Frobenius coboundary. Finally the logarithm and
exponential on the \(p\)-adically complete Laurent ring give
\((f-1)/p^{j+1}\equiv C(u)\bmod p^n\), by the same estimate
on the exponential terms. This proves the formula.
\end{proof}

\begin{theorem}[A precision-independent bound for weight zero]
\leavevmode\label{thm:oldweightzero748}
Let an additive cocycle \(z:G_K\to\mathbb Z/p^\nu\) have an
extension which over \(F\) is a torsion crystalline subquotient
in \([-a,B]\), with \(a,B\ge0\). Put \(D=\min(a,B)\).
If \(D+1\le p-2\), \(q(a+B+D)<p-1\), and
\(q\max(D+1,2D)<p\), the order of \(z(I_K)\) is at most
\(p^D\), independently of \(\nu\).
\end{theorem}
\begin{proof}
Local reciprocity identifies the pro-\(p\) units of \(K\) with
\(1+p\mathbb Z_p\), on which \(\mathcal L\) is a generator of
the additive character module. Subtract an unramified integral
character modulo \(p^\nu\) and multiply by a unit so that the
ramified character is \(p^{\nu-k}\mathcal L\bmod p^\nu\), where
its image has order \(p^k\). These operations preserve the
torsion subquotient condition by the direct sum, diagonal pullback
and addition pushout of extensions. The unramified extension
is crystalline of weight zero. If \(k>D\), the extension for
\(\mathcal L\bmod p^{D+1}\) is such a subquotient over \(F\).
For detail, first restrict the kernel of the extension with
the displayed cocycle to \(p^{\nu-k}(\mathbb Z/p^\nu)\), retaining
the full quotient coordinate. Then divide by the fixed subgroup
\(0\oplus p^k(\mathbb Z/p^\nu)\). Identifying the kernel with
\(\mathbb Z/p^k\) gives the extension for
\(\mathcal L\bmod p^k\), which can be reduced modulo \(p^{D+1}\).

Set \(n=D+1\). Lemmas~\ref{lem:shortpole748}
and~\ref{lem:logmatrix748} give an element \(X_0\in E^{-D}S_n\)
whose difference from a constant unit times \(C(u)\) is a
Frobenius coboundary in \(D_n\). Both representatives have pole
order strictly below \(p\). For \(C\) only the terms \(i\le n\)
survive, since \(i-1-v_p(i)\ge n\) for \(i>n\) and \(n\le p-2\).
For \(n<i<p\) this is immediate, and for \(i\ge p\) use
\(i-1-\log_p i\ge p-2\). For \(X_0\) the pole order is at most
\(q(D+n-1)=2qD\) if \(D>0\), by expanding \(E^{-D}\) as a
binomial series in \(cp\,u^{-q}\) modulo \(p^n\). When \(D=0\)
it is integral. If \(Y\in D_n\) has positive pole order \(s\),
\(\varphi(Y)-Y\) has pole order exactly \(ps\), since its
leading coefficient remains nonzero. A coboundary of pole order
strictly below \(p\) must consequently be integral. We obtain
\(E^D C(u)\in S_n\).

Its coefficient at \(u^{-q}\), before reduction, is
\[
 -(p-1)c^{D+1}p^D
         \sum_{j=0}^{D}\frac{(-1)^{D-j}}{j+1}{D\choose j}
       =-(p-1)c^{D+1}p^D\frac{(-1)^D}{D+1}.
\]
The identity follows by integrating \((x-1)^D\) from zero
to one. The coefficient has valuation exactly \(D\), hence
is nonzero modulo \(p^{D+1}\). This contradicts integrality
and proves \(k\le D\).
\end{proof}

\begin{corollary}[Removal of the trivial scalar from the old-prime sum]
\leavevmode\label{cor:oldlogperiod748}
In Theorem~\ref{thm:oldcolumn745} one may replace
\(B_{\rm old}(D)\) by
\[
 B_{\rm old}^{>0}(D)=\frac1{e_L}
       \sum_{1\le h<t'-1}\operatorname{length}_{\mathcal O}\Gamma_h(D)
\]
in the prime-sum inequality, for all the marked families,
including those with strongly irreducible local constituents.
\end{corollary}
\begin{proof}
At a negative completion use \(a=-J_w\) and \(B=t'-1\).
Theorem~\ref{thm:oldtorustraces743} and the graph presentation
give the required subquotient condition. Here \(D=\min(a,B)
\le t_0\). The hypotheses of Theorem~\ref{thm:oldweightzero748}
hold uniformly, since \(a+B\le6b_x\), \(t_0\le b_x\),
\(p\ge b_x^{10}\) and \(Q\le\sqrt{p-1}\).
At positive completions the ramified additive character already
vanishes. For a coefficient ring \(\mathcal O\) apply the theorem
to the \(\mathbb Z_p\)-coordinates of the cocycle modulo \(p^\nu\).
Pull back its quotient along the copy of \(\mathbb Z/p^\nu\)
generated by one, then push out the kernel along each coordinate
functional. These operations give exactly the scalar extensions
used in the theorem and preserve the torsion subquotient condition.
Their ramified images are all killed by \(p^D\), so the same is
true over \(\mathcal O\).

The evaluation at \(g\) of Lemma~\ref{lem:wildgenerators745}
computes the ramified coordinate. In one copy of
\(\mathcal O/p^\nu\) the subgroup killed by \(p^D\) has length
at most \(e_L D\). Thus the joint normalized length omitted
from \(B_{\rm old}\) is at most \(n_xt_0/2\) at either prime,
even when the restrictions at distinct places are dependent.
At a node there are two old primes, both at most \(b_x^{11}\).
The sum of these errors with the prime weights, divided by
\(\sqrt{b_x}\), is \(O(n_xt_0\log b_x/\sqrt{b_x})\).
Use \(n_x\ll_{q_0,d}x^{12}\), \(t_0\ll x^5\) and
\(\log b_x\ll b_x^{1/4}\) to bound the sum over nodes \(x\ge V_0\)
by \(O_{q_0,d}(V_0^{-107})\). Restricting the original
suprema and using \((u+v)_+\le u_++v\) for \(v\ge0\)
gives the asserted prime sum.
\end{proof}

The bound removes the weight-zero ramification term on the same
fields and intervals as the projected period. A uniform estimate
for the positive intermediate scalar weights at the negative
completions would remove the remaining old-prime term.
\clearpage

\subsection{Positive scalar weights on the rational uniformizer tower}
\label{sec:rationaltower749}
We pass from the varying tame extension at an original level
prime to a single tower over \(\QQ_p\). For each positive scalar
weight this gives a necessary local equation independent of the
large negative jump. We compute its first residue class.

Fix compatible \(p\)-power roots of \(-p\) and write
\(K_\infty=\QQ_p((-p)^{1/p^\infty})\). For brevity put
\(S_n(v)=(\mathbb Z/p^n)[[v]]\), \(\mathcal E_n(v)=S_n(v)[1/v]\)
and \(e(v)=v+p\). Let \(1\le h\le p-3\), and choose a generator
of \(H^1(\QQ_p,\mathbb Z_p(-h))\), which is free of rank one
by Lemma~\ref{lem:wildgenerators745}. At precision \(p^n\)
denote by \(C_{h,n}\) an upper right Frobenius entry of the
corresponding extension restricted to \(G_{K_\infty}\).
Use the covariant convention of Lemma~\ref{lem:logmatrix748},
so the diagonal entries are \(e(v)^h\) and one. Thus \(C_{h,n}\)
is well-defined up to a constant unit and addition of
\[
           L_h(Y)=e(v)^h\varphi(Y)-Y,\qquad
           \varphi(v)=v^p,\quad Y\in\mathcal E_n(v).
\]
Membership tests below are independent of these choices.

\begin{proposition}[Descent of the pole bound]
\leavevmode\label{prop:tamedescent749}
Let \(F\) be the degree \(q\) subfield of \(\QQ_p(\mu_p)\),
and let \(0\le h\le B\), \(a\ge0\). Suppose an extension of
\(\mathbb Z/p^n\) by \((\mathbb Z/p^n)(-h)\) over \(\QQ_p\)
is a crystalline lattice subquotient over \(F\) in the
filtration interval \([-a,B]\). If \(q(a+B+n-1)<p-1\), its
entry \(C\) over \(K_\infty\) satisfies
\[
     C\ \in\ e(v)^{-d} S_n(v)+L_h(\mathcal E_n(v)),
                 \qquad d=\min(a,B-h).
\]
\end{proposition}
\begin{proof}
If \(\zeta_p\) is a primitive \(p\)-th root, expansion of
\((1+(\zeta_p-1))^p=1\) gives
\((\zeta_p-1)^{p-1}/(-p)\equiv1\bmod(\zeta_p-1)\).
The map raising to \(p-1\) is an automorphism of principal
units, so \(\QQ_p(\zeta_p)\) contains a uniformizer \(\lambda\)
with \(\lambda^{p-1}=-p\) and
\(\lambda/(\zeta_p-1)\equiv1\bmod(\zeta_p-1)\).
The field \(F\) has a uniformizer \(\lambda^{(p-1)/q}\)
whose \(q\)-th power is \(-p\). Choose its roots compatibly
with the fixed tower. On its tower the Kisin parameter \(u\)
satisfies \(v=u^q\) and \(E(u)=u^q+p=e(v)\).
Indeed the parameters are Teichm\"uller lifts of the compatible
uniformizer sequences. Restriction of \'etale Frobenius modules
is scalar extension for the resulting inclusion of Laurent
rings, by the field-of-norms equivalence. The tame group acts
by \(u\mapsto\eta u\) for \(\eta\in\mu_q\). Its action commutes
with Frobenius since \(q\mid p-1\).

As in Lemma~\ref{lem:shortpole748}, dualize the extension
and twist by \(\kappa^a\). The intersection and image models
agree with the standard end models by the uniqueness proved
there. In their bases the matrix has diagonal entries
\(E^{a+h},E^a\) and upper right entry \(X\in S_n(u)\).
The inverse height condition gives
\(E^{a+B-(2a+h)}X\in S_n(u)\).
After undoing the twist the entry \(X'=E^{-a}X\) is in both
\(E^{-a}S_n(u)\) and \(E^{-(B-h)}S_n(u)\). This proves the
pole bound \(E^{-d}\) on that tower.

The end bases can be taken to be the base changes of the two
standard end bases over \(K_\infty\), since the rank-one Tate
multipliers are already \(e(v)^h\) and one. Therefore
\(C(u^q)-X'=E^h\varphi(Y)-Y\) for a Laurent series \(Y\)
over the larger tower. Average this equality over \(\mu_q\),
dividing by the unit \(q\). The average of \(X'\) belongs to
\(e(v)^{-d}S_n(v)\), and the average of \(Y\) belongs to
\(\mathcal E_n(v)\). Indeed the average kills exactly the powers
of \(u\) whose exponent is not divisible by \(q\).
This proves the assertion over the original tower.
\end{proof}

\begin{lemma}[A finite system of coefficient equations]
\leavevmode\label{lem:finitepoles749}
Let \(h<p-1\) and let \(C\in\mathcal E_n(v)\) have pole order
\(N\), counting zero for an integral series. In the test of
Proposition~\ref{prop:tamedescent749} one need only use
\[
  Y=\sum_{m=1}^{M}y_m v^{-m},\qquad
  M=\left\lfloor\frac{\max(N,d+n-1)+h}{p}\right\rfloor .
\]
The test is equivalent to the vanishing of all negative-power
coefficients of \(e(v)^d(C-L_h(Y))\) for some such \(Y\).
\end{lemma}
\begin{proof}
Integral terms of \(Y\) give only integral terms of \(L_h(Y)\).
An element of \(e(v)^{-d}S_n(v)\) has pole order at most
\(d+n-1\) by the binomial expansion, with pole order zero
if \(d=0\). If \(m>0\) is the largest pole of \(Y\), its
coefficient occurs unchanged at exponent \(-pm+h\) in
\((v+p)^h\varphi(Y)\), since the leading coefficient of the
polynomial is one. A different Frobenius monomial is separated
by at least \(p>h\), and \(-Y\) cannot reach that exponent
because \(pm-h>m\). Thus \(pm-h\le\max(N,d+n-1)\) whenever
the membership holds. Conversely, any solution to the finite
equations gives the required membership since \(e(v)\) is a
unit in \(\mathcal E_n(v)\). Only the negative part of the input
series \(C\) is needed in these equations.
\end{proof}

\begin{theorem}[The first residue of a positive scalar generator]
\leavevmode\label{thm:residuepositive749}
For \(1\le h\le p-3\) the class of \(C_{h,1}\) in
\(\mathbb F_p((v))/L_h(\mathbb F_p((v)))\) has a representative
\(A v^{-1}\) with \(A\in\mathbb F_p^\times\). Consequently the
pole test at \(n=1\) excludes a primitive trace class precisely
when \(d=0\).
\end{theorem}
\begin{proof}
Put \(L=\QQ_p(\zeta_p)\) and use \(\lambda\) from the preceding
proof. Over \(G_L\) a nonzero residual class becomes a homomorphism
to \(\mathbb F_p\). Its restriction is nonzero since \([L:\QQ_p]\)
is prime to \(p\). Local reciprocity identifies it with a homomorphism
of \(L^\times\). Conjugation on its argument by
\(\delta\in\operatorname{Gal}(L/\QQ_p)\) scales its value by
\(\bar\kappa(\delta)^{-h}\). The quotient of consecutive principal
unit groups \(U_j/U_{j+1}\), \(U_j=1+\lambda^j\mathcal O_L\),
has character \(\bar\kappa^j\), since
\(\delta(\lambda)/\lambda\) is the Teichm\"uller lift of
\(\bar\kappa(\delta)\). The unramified direction is trivial
under conjugation, and the residue unit group has order prime
to \(p\). Moreover \(U_{p+1}=U_2^p\) by the logarithm on
\(\lambda^2\mathcal O_L\). Among the layers \(1,\ldots,p\)
only \(i=p-1-h\) has the required character, with \(2\le i\le p-2\).
Semisimplicity for the group of order \(p-1\), or the successive
unit quotients, shows that the homomorphism is nonzero on \(U_i\)
and zero on \(U_{i+1}\). Its cyclic degree \(p\) extension
therefore has upper jump \(i\) by the ramification theorem of
local reciprocity. We use the upper numbering in which inertia
is the zeroth group.

Set \(L_m=L(\lambda^{1/p^m})\), choosing compatible roots.
Each step over \(L_m\) is cyclic of degree \(p\) with unique
jump \(p^{m+1}\). Indeed a nonidentity automorphism multiplies
its new uniformizer by a nontrivial \(p\)-th root of unity,
and \(v_{L_{m+1}}(\zeta_p-1)=p^{m+1}\), with the uniformizer
of \(L_{m+1}\) having valuation one. The degree \(p\) extension
cut out by our character remains cyclic with jump \(i\) over
every \(L_m\). To check the induction step, the two cyclic
extensions of \(L_m\) have different jumps, hence are linearly
disjoint. In their compositum every character with nonzero
component from the higher-jump extension has that higher jump,
since the lower character vanishes on that last group.
The upper group of the compositum is therefore the full group
through \(i\) and thereafter intersects the subgroup over
\(L_{m+1}\) trivially. The Herbrand function for the base change
to \(L_{m+1}\) is the identity up to \(i\), proving the assertion.

The same jump occurs in the field of norms with uniformizer \(z\)
for the sequence of roots of \(\lambda\), where \(v=z^{p-1}\).
One can see this directly in the field-of-norms construction.
Take a uniformizer in the field of norms of the composite tower
and its norm-compatible components in the finite layers. These
components have valuation one because the extensions are totally
ramified. For a generator \(\sigma\) of the cyclic group the
differences of each component and its \(\sigma\)-translate have
valuation \(i+1\), by the preceding calculation and the equality
of the two numberings for the single jump of a cyclic group
of order \(p\). The components of the difference in the field
of norms are limits of the norms of these differences. Those
norms have the same normalized valuation, so their limits do also.
This gives the lower, hence upper, jump \(i\) in the norm field.
~\cite{III-II-wintenberger749}

Modulo \(p\) replace the basis of the subline by \(z^{-h}\)
times its original basis. Its Frobenius multiplier becomes one,
and the upper right entry becomes \(z^h C_{h,1}(z^{p-1})\).
It is now an Artin--Schreier class, up to the harmless choice
of sign for extensions. Such a class over \(\mathbb F_p((z))\)
has a unique representative consisting of a constant and terms
\(a_j z^{-j}\) with \(j>0\) prime to \(p\). Indeed remove
positive powers by successive Frobenius, and replace \(z^{-pj}\)
by \(z^{-j}\) modulo a coboundary. Uniqueness follows by looking
at the largest pole of \(\varphi(Y)-Y\), and then at the constant.
The greatest remaining pole gives the ramification jump of a
ramified Artin--Schreier extension. Explicitly, for greatest pole
\(j\), a root \(y\) of \(y^p-y=f\) has valuation \(-j\) upstairs,
where \(z\) has valuation \(p\). Choose integers \(A,B\) with
\(-jA+pB=1\). The uniformizer \(y^Az^B\) is multiplied under
\(y\mapsto y+1\) by \((1+1/y)^A\), which differs from one by
an element of valuation \(j\), since \(p\nmid A\).

The largest pole is therefore \(i\). Under \(z\mapsto\eta z\),
\(\eta\in\mu_{p-1}\), the class has character \(\eta^h\).
Uniqueness of the reduced representative forces the constant
to be zero and every pole to satisfy \(-j\equiv h\bmod p-1\).
The only possibility is \(j=i\), with nonzero coefficient \(A\).
Take an Artin--Schreier coboundary giving the difference from
\(A z^{-i}\). Average its primitive against the character
\(\eta^h\) to keep only that eigenspace. Its primitive then
has the form \(z^hY(z^{p-1})\). Dividing its Frobenius
difference by \(z^h\) gives \(v^h\varphi(Y)-Y\). The reduced
term becomes \(A v^{-1}\), as required. This entry belongs to
\(v^{-d}\mathbb F_p[[v]]\) for \(d\ge1\). For \(d=0\) use
Lemma~\ref{lem:finitepoles749} with \(N=1\), which gives \(M=0\),
so its nonzero pole cannot be removed.
\end{proof}

\begin{corollary}[A finite-precision test for the remaining old-prime loss]
\leavevmode\label{cor:finitepositive749}
At a negative completion of either old prime, take \(B=t'-1\)
and \(1\le h<B\). For any \(n\ge1\) with \(Q(-J_w+B+n-1)<p-1\),
if the test of Lemma~\ref{lem:finitepoles749} fails for
\(C_{h,n}\) and \(d=B-h\), the image of the scalar cup coordinate
at this completion is killed by \(p^{n-1}\), for the full marked
family at every original trace precision.
\end{corollary}
\begin{proof}
First take \(\mathbb Z_p\) coefficients. If the local image of a
class at precision \(p^\nu\) has order \(p^k\), it is a unit
times \(p^{\nu-k}\) times the fixed generator in cohomology,
by Lemma~\ref{lem:wildgenerators745}. Subtract a coboundary
and multiply by a unit, so this identity holds on cocycles.
Restricting the kernel coordinate and dividing out the redundant
quotient coordinate, as in Theorem~\ref{thm:oldweightzero748},
gives the extension for the generator modulo \(p^k\). These
operations preserve the torsion subquotient condition. If
\(k\ge n\), reduce modulo \(p^n\) and apply
Proposition~\ref{prop:tamedescent749}. This contradicts the
assumed failure. Over \(\mathcal O/p^\nu\) apply the same
argument to each \(\mathbb Z_p\)-coordinate by pullback and
pushout. The evaluation isomorphism of
Lemma~\ref{lem:wildgenerators745} identifies the exponent bound
with the one for the free cup coordinate.
\end{proof}

In this test the exponent \(B-h\) and the Laurent ring on the uniformizer
tower of \(\QQ_p\) do not depend on \(-J_w\) or on
the degree of the tame extension. The first residue calculation
shows why the intermediate weights require information at higher
precision. An estimate for the precision of failure, uniform
over the primes and weights in the interval, is still needed.
\clearpage

\subsection{Finite obstruction depth for the scalar traces}
\label{sec:finiteobstruction750}
The Laurent test on \(\QQ_p\) in Section~\ref{sec:rationaltower749}
eventually fails at each positive scalar weight. We prove this
for its actual integral cohomology generator and remove the
restriction on its precision in the tame descent argument.

\begin{lemma}[Uniqueness of the model at every precision]
\leavevmode\label{lem:allmodels750}
Use \(S_n,\mathcal E_n\) and the linearized Frobenius of
Section~\ref{sec:logold748}, with Eisenstein polynomial of degree
\(q\) and reduction \(u^q\). If \(qH<p-1\), an \'etale
Frobenius module killed by \(p^n\) has at most one model of
height at most \(H\) inside it, for every \(n\ge1\).
Consequently Proposition~\ref{prop:tamedescent749} holds for
every \(n\) under the single inequality \(q(a+B)<p-1\).
\end{lemma}
\begin{proof}
At \(n=1\) this is the pole comparison in
Lemma~\ref{lem:shortpole748}. The same proof applies to any
module killed by \(p\) inside its generic module. For general
\(n\) write \(D\) for the \'etale module. Given a model \(M\),
intersect it with \(p^iD\) and take its image in
\(p^iD/p^{i+1}D\), for \(0\le i<n\). This image spans its
generic module and has height at most \(H\). In detail,
intersection with the \'etale submodule preserves the lower
Frobenius inclusion. Indeed the preimage of \(E^H x\) for an
intersection vector \(x\) maps to zero in the generic quotient
by injectivity of the linearized Frobenius there. Flatness of
the Frobenius pullback therefore puts that preimage in the
pullback of the intersection. The upper inclusion follows from
stability. Taking the image in a generic quotient preserves
both inclusions by projecting them. Localization preserves
each kernel and image used here.

For two models \(M,N\) their sum is also a model of height
at most \(H\) by adding the Frobenius inclusions. In each
graded piece just considered its image equals the image of
\(M\) by the uniqueness at \(n=1\), since the ambient graded
piece is killed by \(p\). Now descend from \(i=n\) to \(i=0\).
For \(x\in(M+N)\cap p^iD\) the equality of graded images
gives \(m\in M\cap p^iD\) such that \(x-m\) belongs to the
next intersection, already in \(M\) by induction.
This gives \(M=M+N\), and the same argument gives \(N=M+N\).
In Proposition~\ref{prop:tamedescent749} the restriction on
\(n\) was used only to identify the two end models. The
uniqueness just proved identifies them at every precision.
\end{proof}

For \(1\le h\le p-3\) and \(d\ge0\), let \(N(p,h,d)\) be the
first \(n\) at which the Laurent class \(C_{h,n}\) fails the
test with exponent \(d\) in Proposition~\ref{prop:tamedescent749}.
For \(h=0\) use the ramified generator \(\mathcal L\) of
Lemma~\ref{lem:logmatrix748} on the same \(-p\) tower.
Failure at one precision implies failure at every larger one,
by reduction of the module and its basis. The next theorem proves
existence without using any estimate for the auxiliary systems.

\begin{theorem}[Finite depth on the rational tower]
\leavevmode\label{thm:finiteobstruction750}
If \(1\le h\le p-3\) and \(h+2d<p-1\), then \(N(p,h,d)\)
exists and is finite. It is one for \(d=0\) and at least two
for \(d>0\). For \(h=0\) and \(2d<p-1\) it is exactly \(d+1\).
\end{theorem}
\begin{proof}
Let \(U\) be the integral extension of \(\mathbb Z_p\) by
\(\mathbb Z_p(-h)\) representing a generator. Suppose its
test succeeds at every \(n\). Put \(H'=h+2d\) and apply the
contravariant \'etale module functor to the restrictions of
\(U^\vee(d)\). At precision \(p^n\) the assumed test and a
change of lifted quotient basis give a Frobenius matrix
\[
 A_n=
 \begin{pmatrix}
       e(v)^{h+d} & e(v)^d X_n\\
            0     & e(v)^d
 \end{pmatrix},\qquad   X_n\in e(v)^{-d}S_n(v).
\]
Indeed before twisting the dual the diagonal is \(e(v)^h,1\)
in the covariant convention for \(U\), and the extra twist
multiplies the whole Frobenius by \(e(v)^d\). The entries of
\(A_n\) and of \(e(v)^{H'}A_n^{-1}\) are integral, including
the upper right entry \(-e(v)^dX_n\) of the latter. The span
of this basis over \(S_n(v)\) is thus a free model \(M_n\) of
height \(H'\) inside the fixed generic module.

These models are compatible, regardless of which successful
solution was used at each level. Indeed the image of \(M_{n+1}\)
at level \(n\) has height \(H'\) by the image argument above.
It is \(M_n\) by Lemma~\ref{lem:allmodels750} applied with
\(q=1\). The induced map \(M_{n+1}/p^nM_{n+1}\to M_n\)
is injective as well, since a basis of \(M_{n+1}\) is a
generic basis, remains one after reduction and is therefore
still independent over \(S_n(v)\). Lift a basis successively.
The lifts are bases at every step by reduction through a nilpotent
ideal. Their inverse limit identifies \(M=\varprojlim M_n\)
with two copies of \(\mathbb Z_p[[v]]\). The compatible
Frobenius matrices and their integral multiples
\(e(v)^{H'}A_n^{-1}\), expressed in these bases, give the two
height inclusions on \(M\). After localization and \(p\)-adic
completion its basis gives the original integral \'etale module, since it gives an isomorphism modulo every \(p^n\) by construction.

This is finite \(E(v)\)-height on the fixed uniformizer tower
for the restriction of a lattice of the full \(G_{\QQ_p}\).
The finite-height theorem of Gao implies that the associated
rational representation is potentially semistable. One may
pass to the dual convention in that theorem, since rational
potential semistability is preserved by duality and Tate twist.
We conclude that \(U[1/p]\) is semistable over some finite
extension \(L/\QQ_p\).~\cite[Theorem~1.1.16]{III-II-gaofinite750}

Over \(L\) the monodromy \(N\) of this extension is zero.
Indeed it kills the subline and is zero on the quotient,
so could only induce a map from the trivial line to the
line for \(\QQ_p(-h)\). If \(f\) is the residue degree of
\(L\), the linear operator \(\varphi^f\) is one on the first
line and \(p^{hf}\) on the second. The relation
\(N\varphi^f=p^f\varphi^f N\) makes the putative map equal
to \(p^{(h+1)f}\) times itself. Hence it is zero. Semistability
with zero monodromy gives crystallinity. But
\(H^1_f(L,\QQ_p(-h))=0\) for \(h>0\) by the crystalline
dimension formula, since all its filtration jumps are positive
and its invariants vanish. The extension splits over \(L\).
Restriction followed by corestriction multiplies rational
cohomology by \([L:\QQ_p]\), so it would split over \(\QQ_p\)
as well. This contradicts the choice of the integral generator
in its torsion-free cohomology group. We obtain finite failure.
The first-layer assertions follow from Theorem~\ref{thm:residuepositive749}.

For \(h=0\) take \(c=-1\) in Lemma~\ref{lem:logmatrix748}.
For \(n\le d\), only its first \(n\) series terms survive
modulo \(p^n\) since \(n<p-1\). Every negative coefficient of
\((v+p)^d C(v)\) is zero there. In fact its terms involve
\(v^j v^{-i}\) with \(j<i\le n\), and the accompanying power
of \(p\) is \(d-j+i-1\ge d\), with denominator \(i\) a unit.
The test succeeds. At \(n=d+1\), the nonzero coefficient at
\(v^{-1}\) of \((v+p)^d C(v)\) has valuation exactly \(d\)
by the computation of Theorem~\ref{thm:oldweightzero748}.
The original logarithm and every representative in
\((v+p)^{-d}S_n(v)\) have pole order less than \(p\) since
\(2d<p-1\). A Frobenius coboundary of pole order less than
\(p\) is integral at weight zero. Thus no change of basis
removes that coefficient. The argument also gives failure at
\(n=1\) when \(d=0\).
\end{proof}

\begin{corollary}[A bound independent of the tame field and the far jump]
\leavevmode\label{cor:rationaldepth750}
For either old coefficient prime in the full marked family,
put \(B=t'-1\) at a nonzero lift. The positive scalar term
remaining in Corollary~\ref{cor:oldlogperiod748} satisfies
\[
     B_{\rm old}^{>0}(D)
       \ \le\ \frac{[E_x:\QQ]}2
          \sum_{h=1}^{B-1}\big(N(p,h,B-h)-1\big).
\]
Every number on the right is finite and depends only on its
three displayed arguments. The estimate holds at arbitrary
trace precision and for every auxiliary family at the node.
\end{corollary}
\begin{proof}
At each negative completion \(-J_w>B\), and
\(Q(-J_w+B)<p-1\) by Theorem~\ref{thm:oldtorustraces743}.
Also \(h+2(B-h)\le2t_0<p-1\) by the size of the primes.
We may therefore apply Corollary~\ref{cor:finitepositive749}
with \(n=N(p,h,B-h)\), using Lemma~\ref{lem:allmodels750}
to remove its precision restriction. This kills each cup
coordinate by \(p^{N(p,h,B-h)-1}\). A subgroup killed by
\(p^s\) in one copy of \(\mathcal O/p^\nu\) has normalized
length at most \(s\). Add these bounds over the negative
completions and the positive intermediate weights. The joint
image is a submodule of their direct sum, so this addition
also applies when the local restrictions are dependent.
\end{proof}

The corollary replaces the finite local defect for a long
interval of jumps over the ramified field by a finite depth
over \(\QQ_p\) depending only on the scalar weight and the
short string at that prime. An upper bound on these positive-weight
depths as \(p\) and the string vary is still required to sum
the unrestricted projected period on the full interval.
\clearpage

\subsection{Cyclotomic units detecting the positive scalar classes}
\label{sec:unitcoordinates751}
For \(1\le h\le p-3\), the Laurent test in
Section~\ref{sec:rationaltower749} uses a primitive class of
\(H^1(\QQ_p,\ZZ_p(-h))\). We give its local reciprocity character
at every finite precision. The unit on which the character is
normalized can be chosen through one polynomial calculation over
\(\mathbb F_p\).

Choose compatible primitive roots \(\zeta_{p^n}\), with
\(\zeta_{p^{n+1}}^p=\zeta_{p^n}\). Put \(L_n=\QQ_p(\zeta_{p^n})\),
\(\lambda_n=\zeta_{p^n}-1\),
\(\Gamma_n=\operatorname{Gal}(L_n/\QQ_p)=\Delta\times P_n\), where
\(\Delta\) is the Teichm\"uller subgroup and \(P_n\) is the kernel
of reduction to \(\Gamma_1\). Write
\(\omega:\Delta\longrightarrow\ZZ_p^\times\) for the action on roots
of unity, and write \([b]\) for the Teichm\"uller lift of
\(b\in\mathbb F_p^\times\). On a pro-\(p\) abelian group with
\(\Delta\)-action use the projector
\[
 e_i=\frac1{p-1}\sum_{\delta\in\Delta}\omega(\delta)^{-i}\delta .
\]
In a multiplicative group the coefficients of this operator act by
continuous \(\ZZ_p\)-powers. Set \(U_n=1+\lambda_n\mathcal O_{L_n}\).

\begin{lemma}[A unit basis at every cyclotomic level]
\leavevmode\label{lem:unitbasis751}
For each \(2\le i\le p-2\) there is
\(b\in\mathbb F_p^\times\setminus\{1\}\) such that, at every
\(n\ge1\), the \(\ZZ_p[P_n]\)-module \(e_i U_n\) is free of rank
one with basis
\[
 w_n=e_i u_n(b),\qquad
 u_n(b)=\frac{1-[b]\zeta_{p^n}}{[1-b]}.
\]
Both sequences are compatible under the norm.
\end{lemma}
\begin{proof}
The numerator has residue \(1-b\), so the quotient is a principal
unit. Taking the product over the \(p\) choices of a \(p\)-th root
of \(\zeta_{p^n}\) gives \(1-[b]^p\zeta_{p^n}=1-[b]\zeta_{p^n}\).
The denominator also has \(p\)-th power equal to itself. Norms
commute with the action of \(\Delta\).

We first select \(b\) at level one. In
\(\mathcal O_{L_1}/\lambda_1^{i+1}\) the characteristic is \(p\)
because \(i+1\le p-1=v_{L_1}(p)\). The truncated exponential
and logarithm through degree \(i\) are mutually inverse.
In the coordinate \(X=\log(1+\lambda_1)\), an element
\(\delta\in\Delta\) acts by \(X\mapsto\bar\omega(\delta)X\).
This follows by expanding \((1+\lambda_1)^a\), with any integer
\(a\) representing \(\delta\) modulo \(p\). Thus in the logarithm
of a unit the projector \(e_i\) retains exactly the coefficient
of \(X^i\) through this degree. For \(u_1(b)\) that coefficient is
\[
 A_i(\beta)=-[X^i]\sum_{j=1}^{i}
               \frac{\beta^j}{j}(\exp X-1)^j,\qquad
 \beta=\frac{b}{1-b}.
\]
Indeed the constant factor \((1-[b])/[1-b]\) is \(1\) in the
truncated ring. The polynomial \(A_i\) has degree \(i\), leading
coefficient \(-1/i\) and constant term zero. The admissible
\(\beta\)'s are all elements of \(\mathbb F_p\) except \(0,-1\).
Vanishing at every admissible value would give at least \(p-1\)
roots, including \(0\). Choose a value with \(A_i(\beta)\ne0\).
It follows that \(w_1-1\) has valuation exactly \(i\).

The only torsion in \(U_n\) is \(\mu_{p^n}\), on which \(\Delta\)
acts by \(\omega\). After tensoring with \(\QQ_p\), the logarithm
identifies \(U_n\otimes_{\ZZ_p}\QQ_p\) with \(L_n\), since it
is an isomorphism on sufficiently deep units. The normal basis
theorem applied to \(L_n/\QQ_p\) now gives
\(\operatorname{rank}_{\ZZ_p} e_i U_n=|P_n|\).
This module is finite free over \(\ZZ_p\), since its projection
kills the torsion. Every \(p\)-th power in \(U_1\) is
\(1\bmod\lambda_1^p\), by the binomial formula. Hence \(w_1\)
is a basis of the rank-one module \(e_i U_1\).

For \(n>1\) let \(\sigma\) generate \(P_n\). Norm induces an
isomorphism from the coinvariants \((e_i U_n)_{P_n}\) to \(e_i U_1\).
Surjectivity follows from the norm formula for \(w_n\). For
injectivity, a unit of norm one is \(\sigma(y)/y\) for some
\(y\in L_n^\times\) by Hilbert's Theorem 90. In the pro-\(p\)
completion of \(L_n^\times\), the valuation quotient is \(\ZZ_p\)
with trivial \(\Delta\)-action and the kernel is \(U_n\). Thus
\(e_i y\) belongs to \(e_i U_n\). Applying \(e_i\) to the
identity proves injectivity. The ring \(\ZZ_p[P_n]\) is local
with maximal ideal \((p,\sigma-1)\), since modulo \(p\) its group
ring is \(\mathbb F_p[T]/T^{|P_n|}\) with \(T=\sigma-1\) and
the ring is finite over \(\ZZ_p\). Nakayama's lemma now shows
that \(w_n\) generates \(e_iU_n\) over this ring. Both are free
of the same finite rank over \(\ZZ_p\), so the resulting
surjection is an isomorphism.
\end{proof}

\begin{theorem}[An integral reciprocity coordinate for each scalar weight]
\leavevmode\label{thm:unitcoordinate751}
Fix \(1\le h\le p-3\), take \(i=p-1-h\), and choose \(b\) as in
Lemma~\ref{lem:unitbasis751}. Write \(R_n=\ZZ/p^n\ZZ\), fixing
the usual scalar basis of \(R_n(-h)\). Restrict a cohomology
class to \(G_{L_n}\) and use local reciprocity there. Evaluation
of the resulting character on \(u_n(b)\) is an isomorphism
\[
 H^1(\QQ_p,R_n(-h))\ \longrightarrow\ R_n .
\]
These isomorphisms commute with reduction. In particular the
inverse images of \(1\) are the reductions of a single generator
of \(H^1(\QQ_p,\ZZ_p(-h))\).
\end{theorem}
\begin{proof}
The restriction of \(R_n(-h)\) to \(G_{L_n}\) is trivial. On
\(\Delta\) it has nontrivial character \(\omega^i\), so its
\(\Delta\)-invariants are zero. All higher cohomology of
\(\Delta\) with these coefficients vanishes by averaging. The
extension of \(P_n\) by \(\Delta\) therefore gives
\(H^j(\Gamma_n,R_n(-h))=0\) for every \(j\). By inflation and
restriction, followed by local reciprocity, the group on the
left is precisely
\[
 \operatorname{Hom}^{\mathrm{cont}}_{\Gamma_n}
   \big((L_n^\times)^{\wedge p},R_n(-h)\big).
\]
Here the superscript on the source denotes pro-\(p\) completion.
We use one fixed convention for reciprocity at every level.
Its compatibility with norms identifies the map of characters
under field restriction with precomposition by the norm
\cite{III-II-V-cft101}.

Only \(e_i(L_n^\times)^{\wedge p}=e_iU_n\) contributes. A
homomorphism on this component is determined by any assigned
value \(t\in R_n\) on \(w_n\). In fact the value on \(\sigma w_n\)
must be \(\kappa(\sigma)^{-h}t\) for \(\sigma\in P_n\), where
\(\kappa\) in this formula is reduced modulo \(p^n\). These
assignments extend uniquely on the group-ring basis of the
lemma and give an equivariant continuous homomorphism. It takes
the same value on \(u_n(b)\) as on \(w_n\). The norm formula
for \(u_{n+1}(b)\) proves compatibility with reduction, even
though the evaluating field varies with \(n\). Finally continuous
cohomology with coefficients in the inverse limit agrees here
with the inverse limit of the \(H^1\)'s: the finite coefficient
transition maps are surjective and the groups of invariants
are zero. This proves the integral statement.
\end{proof}

\begin{corollary}[Finite unit calculation for a trace class]
\leavevmode\label{cor:finiteunits751}
For \(e_n=(p-1)p^{n-1}\), put \(m_n=p^{n-1}+1+n e_n\).
The finite reciprocity character giving the primitive class
above can be calculated on units in
\[
 \ZZ_p[T]\big/\big(\Phi_{p^n}(1+T),T^{m_n}\big).
\]
If a local scalar trace class at a split place over an old prime
has coefficient ring \(\mathcal O/p^n\mathcal O\), its normalized
length in the free scalar coordinate equals the normalized length
of the ideal generated by its value on \(u_n(b)\).
\end{corollary}
\begin{proof}
The indicated cyclotomic polynomial is Eisenstein and \(T\) maps
to \(\lambda_n\). For \(s=1+p^{n-1}>e_n/(p-1)\), logarithm and
exponential identify \(1+\lambda_n^s\mathcal O_{L_n}\) with
\(\lambda_n^s\mathcal O_{L_n}\). Raising to \(p^n\) carries it
onto \(1+\lambda_n^{m_n}\mathcal O_{L_n}\). Every \(R_n\)-valued
character kills this last subgroup. One can therefore decompose
the projected unit in the finite group modulo \(p^n\)-th powers
using the basis \(\{\sigma w_n:\sigma\in P_n\}\) and evaluate
each basis vector by the formula in the theorem.

The coefficient ring \(\mathcal O\) is finite free over \(\ZZ_p\).
Choosing a basis and applying the theorem to each component
identifies the scalar cohomology with \(\mathcal O/p^n\mathcal O\)
by evaluation on this one unit. Any two free scalar coordinates
differ by a unit, so they give the same ideal length. This
applies also to the ideal of values of a family of classes.
\end{proof}
This determines the primitive character and its normalization at
all precisions. To bound the period defect by this calculation
one still needs to bound the orders of the values attained by
the actual trace classes, uniformly over the primes and weights
of Corollary~\ref{cor:rationaldepth750}.
\clearpage

\subsection{Crystalline monodromy at the original level primes}
\label{sec:oldmonodromy752}
The Frobenius condition of Section~\ref{sec:rationaltower749}
allows a simple pole modulo \(p\) at the intermediate weights.
The crystalline lattice also gives a monodromy operator on its
residual Breuil module. We use this operator to exclude that pole
for the actual trace classes.

Let \(p\ge5\), let \(q\mid p-1\), and use the totally ramified
field \(F=\QQ_p((-p)^{1/q})\) and uniformizer from
Proposition~\ref{prop:tamedescent749}. The Eisenstein polynomial
is \(u^q+p\). We use the covariant Frobenius-module functor:
the representation on the uniformizer tower consists of
Frobenius-fixed vectors after separable scalar extension.
The filtration jump of \(\kappa\) is \(-1\).
The residual module of \(\kappa^{-j}\) for \(j\ge0\) has
Frobenius multiplier \(u^{qj}\) in this convention. Equivalently
one may first dualize the representation and then apply the
contravariant functor of Lemma~\ref{lem:tameends743}.

\begin{lemma}[The lowest coefficient allowed by crystalline monodromy]
\leavevmode\label{lem:lowestmonodromy752}
Let a crystalline Breuil module in residue characteristic \(p\)
be an extension of rank-one objects. Suppose its underlying
Breuil--Kisin module has Frobenius
\[
         \varphi(e)=u^t e,\qquad
         \varphi(f)=u^s f+y(u)e,\qquad 0\le s\le t<p-1 ,
\]
where \(e\) spans the submodule, and \(y\in k'[[u]]\) for a
finite coefficient field \(k'\) of characteristic \(p\). The
residue field of \(F\) is \(\mathbb F_p\), so Frobenius acts
trivially on the coefficient field. Then
\[
           \operatorname{ord}_u(y)\ \ge\ \max(0,s-q+1),
\]
with the order of zero taken to be infinite.
\end{lemma}
\begin{proof}
We give the one-residue-embedding calculation underlying
\cite[Lemma~2.2.7]{III-II-bcgntmonodromy752}. Put \(d=q(p-2)\) and
\(R=k'[u]/u^{qp}\). Write \(e_0,f_0\) for the two basis vectors
of the Frobenius pullback modulo \(u^{qp}\). Its filtered part
in degree \(p-2\) contains
\[
 E=u^{d-t}e_0,\qquad
 G=u^{d-s}f_0-u^{d-s-t}y\,e_0 .
\]
Indeed linearized Frobenius sends these to \(u^d e,u^d f\).
The indicated coefficient of \(y\) is integral by the height
condition, or by the upper right entry of \(u^d\) times the
inverse Frobenius matrix. The intersection of the filtered part
with \(R e_0\) is \(R E\), because a lift along this line has
linearized image divisible by \(u^d\) exactly when its
coefficient is divisible by \(u^{d-t}\). The divided Frobenius
\(\varphi_{p-2}\) takes \(E,G\) to \(e_0,f_0\), respectively.

The monodromy on scalars is \(-u\,d/du\). On the rank-one ends
it kills the displayed Frobenius-pullback generators by
\cite[Lemma~2.2.3]{III-II-bcgntmonodromy752}. Thus \(N(e_0)=0\) and
\(N(f_0)=\mu e_0\) with \(\mu\in uR\), the last inclusion being
the crystalline condition. For every filtered vector \(z\),
\(u^qN(z)\) is filtered and
\(\varphi_{p-2}(u^qN(z))=N(\varphi_{p-2}(z))\).
In particular the vector
\[
 U=u^qN(G)+(d-s)u^qG
\]
is in \(R E\). Applying divided Frobenius to it gives
\(\mu e_0\bmod u\mathcal M\), hence zero modulo \(u\), since
\(\varphi(u^q)=u^{pq}=0\) in \(R\). Here \(\mathcal M\) is
the Breuil module. If \(U=zE\), it follows that \(z\) has
zero constant coefficient.

The coefficient of \(e_0\) in \(U\), expanded with the Leibniz
rule, is
\[
     u^{q+d-s}\mu+u^{q+d-s-t}(u y'-t y).
\]
If \(y\) had least degree \(l\le s-q\), the term of degree
\(d-t+q-s+l\le d-t\) in the second summand would be nonzero:
\(0\le l<s\le t<p-1\), so \(l-t\ne0\) in \(k'\).
This degree is nonnegative since \(u^{d-s-t}y\) is integral,
and \(d-t<qp\), so the truncation preserves it.
The first summand has strictly greater degree since \(t\ge s\)
and \(\mu\in uR\). This contradicts \(U\in u R E\).
When \(s<q\) the asserted bound is just integrality.
\end{proof}

\begin{theorem}[Integral scalar trace conditions throughout the tame range]
\leavevmode\label{thm:oldscalars752}
Let \(a,B\) be nonnegative integers with \(q(a+B)<p-1\).
Suppose an extension of \(\ZZ/p^\nu\) by
\((\ZZ/p^\nu)(-h)\) over \(\QQ_p\), \(0\le h\le\min(B,p-3)\),
restricts to a finite subquotient of a lattice crystalline over
\(F\) with filtration jumps in \([-a,B]\). For \(h>0\) its
class is zero. For \(h=0\) its class is unramified.
The same conclusions hold with coefficients in
\(\mathcal O/p^\nu\mathcal O\) for every finite extension of
coefficient rings \(\mathcal O/\ZZ_p\).
\end{theorem}
\begin{proof}
Since \(q\mid p-1\) and \(qh\le q(a+B)<p-1\),
\(q(h+1)\le p-1\).
First suppose \(h>0\). A nonzero class gives as a
subquotient the primitive nonsplit scalar extension modulo \(p\).
Indeed \(H^1(\QQ_p,(\ZZ/p^\nu)(-h))\) is free of rank one by
Theorem~\ref{thm:unitcoordinate751}. After changing the section
and a unit in the kernel, the cocycle is \(p^{\nu-k}\) times a
primitive integral cocycle, where \(p^k>1\) is its order. Take
the submodule with kernel \(p^{\nu-k}(\ZZ/p^\nu)(-h)\), still
containing the chosen quotient generator. In it the subgroup
generated by \(p^k\) times that quotient generator is fixed.
Divide out this subgroup and then reduce modulo \(p\). In the
result the kernel and quotient each have dimension one and the
cocycle is the primitive reduction.

If \(h=0\) and the class is ramified, first subtract its
unramified part. This is a Baer sum with the reduction of an
integral unramified extension, crystalline with jumps zero.
The diagonal pullback and addition pushout are subquotients of
the direct sum of the two extensions. The ramified character of
\(\QQ_p\) is now a multiple of
\(\mathcal L=p^{-1}\log(\kappa^{p-1})\) by local reciprocity.
The same kernel and quotient construction yields its nonzero
primitive reduction.

Denote the resulting two-dimensional residual extension by \(V\).
The twist \(V(-a)|_{G_F}\) is a subquotient of a crystalline
lattice with jumps in \([0,H]\), \(H=a+B\le p-2\). A killed-\(p\)
subquotient can be obtained from the reduction of a lattice:
in the inverse images of its numerator and denominator take the
numerator as the new lattice, whose \(p\)-multiple lies in the
denominator. We can extend coefficients to contain all embeddings
of \(F\) and take the residue field there, since extending this
lattice and then reducing gives scalar extension of its reduction.
By the crystalline Breuil comparison and the correspondence for
subobjects and quotients, \(V(-a)|_{G_F}\) has a crystalline
Breuil module, with a crystalline subline and quotient line.
The underlying Kisin modules form an exact sequence. Their
heights are at most \(H\): the parent has this height and the
linearized Frobenius inclusions pass to saturated intersections
and images as in Lemma~\ref{lem:tameends743}.
\cite[Theorems~2.1.4 and~2.1.5 and \S2.1]{III-II-bcgntmonodromy752}

Over \(k'[[u]]\) the two end lattices are the standard ones for
the characters \(\bar\kappa^{-(a+h)}\) and
\(\bar\kappa^{-a}\). Indeed any two height-\(H\) lattices in the
same generic module agree when \(qH<p-1\), by
Lemma~\ref{lem:allmodels750}. Thus in the order subline, quotient
the exponents in Lemma~\ref{lem:lowestmonodromy752} are
\[
                         t=q(a+h),\qquad s=qa .
\]
Lift a basis of the quotient. After undoing the twist, its
upper right entry is \(u^{-s}y\). By the lemma this has pole
order at most \(q-1\).

On the other hand the primitive class on the rational
uniformizer tower has upper right entry \(A v^{-1}\), where
\(A\ne0\), by Theorem~\ref{thm:residuepositive749}. At \(h=0\)
this follows from the first coefficient of the logarithm in
Lemma~\ref{lem:logmatrix748}. On the tower over \(F\) the
substitution is \(v=u^q\), so the entry is \(A u^{-q}\).
Two entries for the same extension with the same end bases
differ by \(u^{qh}\varphi(Y)-Y\) in \(k'((u))\).
If \(Y\) has a pole of order \(m\ge1\), the pole of this
difference has order \(pm-qh\), since \(pm-qh>m\).
In fact \(pm-qh\ge p-qh>q\).
If \(Y\) is integral its difference is integral.
Neither case can account for the difference between a pole of
order \(q\) and a series with pole order at most \(q-1\).
This contradicts the existence of \(V\).

Finally for coefficients \(\mathcal O/p^\nu\mathcal O\), pull
back the quotient along the rational scalar copy and push out
the kernel along each vector coordinate of a \(\ZZ_p\)-basis
of \(\mathcal O\). Each resulting extension satisfies the same
subquotient condition, by adding copies of the endpoint
characters if needed. Its class is zero, or unramified for
\(h=0\), by what we have proved. The coordinate maps identify
the original cohomology group with their direct sum over
\(\ZZ/p^\nu\), giving the conclusion.
\end{proof}

\begin{corollary}[The projected period at the two original level primes]
\leavevmode\label{cor:oldperiod752}
On the fields \(E_x\), nodes \(x\), coefficient embeddings and
shorter intervals of Theorem~\ref{thm:oldcolumn745}, all the
local scalar trace classes at the two original coefficient
primes belong to the integral finite condition at every
precision. For every marked family there one has \(r(D)=0\).
With \(p(D)\) the normalized integral projected period valuation
of Theorem~\ref{thm:projectedcups742}, one consequently has
\[
 \sum_p(\log p)\!
       \sup_{\substack{D\ {\rm at}\ p\\
               p\in\{p_{x,0},p_{x,1}\}}}
       \frac{p(D)}{\sqrt{b_x}}
                 \ll_{q_0,d} V_0^{-10}.
\]
An empty supremum contributes zero.
\end{corollary}
\begin{proof}
At a negative completion the local field is \(\QQ_p\).
The graph lattice of Lemma~\ref{lem:graphrealization724}
is crystalline over the common degree-\(Q\) tame subfield of
\(\QQ_p(\mu_p)\). Theorem~\ref{thm:oldtorustraces743} gives
the common window with \(a=-J_w\), \(B=t'-1\), and
\(Q(a+B)<p-1\). This is the field \(F\) above with \(q=Q\).
Theorem~\ref{thm:oldscalars752} therefore applies to every
scalar restriction, including the trivial character. Positive
completions and the main summand at both signs were covered
in Theorem~\ref{thm:oldtorustraces743}. Zero lifts have zero
defect by definition. All the groups \(\Gamma_h(D)\) in
Theorem~\ref{thm:oldcolumn745} thus vanish, including the
ramified coordinate at \(h=0\). Substitute in its prime sum.
This uses the same set of families at every node, including
those with a locally irreducible constituent of dimension
greater than one.
\end{proof}
The corollary settles the contribution at the two original
level primes of each node. The other coefficient primes in
the full projected period and the remaining comparisons in
the global argument still require estimates.
\clearpage

\subsection{Tame inertia characters in a scalar trace extension}
\label{sec:tameshifts753}
At a primitive label of rank \(f\), the characters of tame
inertia can have order \(p^f-1\). The small-ramification inequality
of Theorem~\ref{thm:oldscalars752} then fails even for a short
weight interval. We give a necessary condition on the inertial
type of a lattice that can contribute a scalar trace class.
The condition involves the two rank-one ends of the extension
at every residue embedding.

Let \(U/\QQ_p\) be unramified of degree \(f\), let \(e=p^f-1\),
and put \(F=U((-p)^{1/e})\). Fix the compatible uniformizer
towers, with parameters \(v\) on \(U\) and \(u\) on \(F\), so
that \(v=u^e\). Let \(\Delta=\operatorname{Gal}(F/U)\).
On the \(i\)-th residue embedding let
\(\theta_i(\delta)=\delta(u)/u\) in that component, for
\(\delta\in\Delta\). The \(\theta_i\) have order \(e\) and
are the Frobenius conjugates of one another. Index components
cyclically so that Frobenius goes from \(i-1\) to \(i\).
We extend the coefficient field to split the residue embeddings.

\begin{theorem}[Necessary tame characters for a nonfinite scalar class]
\leavevmode\label{thm:tameorder753}
Let \(H=a+B\) with \(a,B\ge0\) and \(3H+1<p\), \(p\ge5\).
Let \(0\le h\le\min(B,p-3)\).
Suppose an extension of \(\ZZ/p^\nu\) by \((\ZZ/p^\nu)(-h)\)
over \(\QQ_p\) is a Galois-stable subquotient of a lattice over
\(\QQ_p\) which is crystalline over \(F\), with jumps in
\([-a,B]\). Write \(\mathcal T\) for the set of its tame
inertial characters together with the trivial character, closed
under inverses and residue Frobenius. If the class is nonzero
for \(h>0\), or ramified for \(h=0\), there exist cyclic lists
of integers \(b_i,c_i\) satisfying, at every \(i\),
\[
\begin{gathered}
 |b_i|,\ |c_i|\le\frac{eH}{p-1},\qquad b_i>c_i,\qquad
       \theta_i^{b_i},\theta_i^{c_i}\in\mathcal T,\\
 0\le ea+p b_{i-1}-b_i\le eH,\qquad
 0\le e(a+h)+p c_{i-1}-c_i\le eH .
\end{gathered}
\]
The same test applies at all precisions over finite extensions
of the coefficient ring, with the tame types of the parent
lattices and their local coefficient conjugates.
\end{theorem}
\begin{proof}
The extraction in Theorem~\ref{thm:oldscalars752} gives the
primitive residual extension. All the submodules and quotients
in that construction are stable already over \(\QQ_p\).
Any added unramified or Tate representation has trivial smooth
inertial character. Twist by \(\bar\kappa^{-a}\) and take
the crystalline Breuil module over \(F\) of the extracted
subquotient; its underlying Kisin module has height at most \(H\).
Here we need its tame descent as well. For the new numerator
lattice in the characteristic-zero parent, the integral Kisin
module carries the descent action, and its reduction modulo
\(u\) has in each component the tame type of the parent.
This is the integral comparison with descent data of
\cite[Definition~3.6 and Propositions~5.16--5.17]{III-II-caraianilevin753}.
It uses the canonical integral Kisin module of the lattice over
\(F\), whose reduction also underlies the crystalline Breuil
module by \cite[Theorem~2.1.5]{III-II-bcgntmonodromy752}. Either choice
of dual convention only reverses characters in the type test.

By the subobject correspondence of
\cite[Theorem~2.1.4]{III-II-bcgntmonodromy752}, the subobjects defining
our quotient and its two ends have saturated underlying Kisin
submodules. They are preserved by tame descent: they are the
intersections of the parent model with their generic subspaces,
and those subspaces descend. Thus the underlying Kisin sequences
remain exact after reduction modulo \(u\), since their quotients
are free over the component power series rings. Every tame
character of an end modulo \(u\) occurs in the parent. Taking
subobjects and quotients preserves the bound on height, by the
linearized Frobenius argument of Lemma~\ref{lem:allmodels750};
that argument for intersections and images does not require
uniqueness of the models.

In the canonical generic bases pulled up from \(U\), the
quotient and subline lattices in component \(i\) are generated
by \(u^{b_i}\) and \(u^{c_i}\) times those bases. Indeed a
lattice in a Laurent line is a power of \(u\) times the standard
lattice. These generators have tame characters
\(\theta_i^{b_i}\) and \(\theta_i^{c_i}\). Their Frobenius
exponents are, respectively,
\[
                  s_i=ea+p b_{i-1}-b_i,\qquad
                  t_i=e(a+h)+p c_{i-1}-c_i .
\]
The height condition puts both in \([0,eH]\). In particular
\(\lvert p b_{i-1}-b_i\rvert\le eH\). At an index with
\(\lvert b_{i-1}\rvert\) maximal, this implies
\((p-1)\max_i\lvert b_i\rvert\le eH\). Apply the same argument
to \(c_i\). It remains to prove the strict ordering.

Take a generic basis for the primitive extension over the
rational tower in which its untwisted off-diagonal term is
\(A v^{-1}\), \(A\ne0\), as in Theorem~\ref{thm:residuepositive749}
and Lemma~\ref{lem:logmatrix748}. After the unramified extension
this term occurs in every component. Lift the chosen quotient
generator in the middle lattice, and project the lift to its
tame eigencharacter by averaging over \(\Delta\), of order
prime to \(p\). Dividing the lift by \(u^{b_i}\) gives a generic
tame-invariant lift of the original quotient generator.
Its difference from the original lift is \(Y_i(v)\) times
the canonical subline generator, with \(Y_i(v)\in k'((v))\),
since only powers of \(u^e\) are invariant. In the resulting
lattice basis the off-diagonal term is
\[
 y_i=u^{ea+p b_{i-1}-c_i}
       \big(A v^{-1}+v^h\varphi(Y_{i-1})-Y_i\big).
\]
This follows by applying Frobenius to the lifted quotient
generator and subtracting \(u^{s_i}\) times its next lift.
Frobenius raises the parameter to its \(p\)-th power.

If some \(Y_j\) has a pole, let \(m\) be the largest pole order
and take \(i-1=j\) attaining it. In component \(i\) the
parenthesis then has pole order \(pm-h>\max(m,1)\).
Integrality of \(y_i\) bounds this order above by
\[
 \frac{s_i+b_i-c_i}{e}\le H+\frac{2H}{p-1}<p-h.
\]
For the strict inequality use \(h\le H\) and \(3H+1<p\).
This contradicts \(pm-h\ge p-h\). Consequently all
\(Y_i\) are integral. In particular the least degree of \(y_i\)
and its residue relative to \(t_i\) are
\[
 l_i=ea+p b_{i-1}-c_i-e=s_i-e+b_i-c_i,\qquad
 l_i-t_i\equiv -e(h+1)\not\equiv0\pmod p .
\]

We spell out the part of the monodromy calculation needed when
\(s_i,t_i\) can exceed \(p\). Use \(d=e(p-2)\) in the proof
of Lemma~\ref{lem:lowestmonodromy752}. In component \(i\) the
filtered vectors \(E_i,G_i\) in that proof generate the
filtered submodule: the condition for a pair of power series
is exactly that its linearized image be divisible by \(u^d\),
which gives those two generators by solving the triangular
system. Their divided Frobenius images are the next basis
vectors. For \(N(f_{0,i})=\mu_i e_{0,i}\), the commutation
relation therefore gives \(\mu_i\in k'[u^p]/(u^{ep})\) in the
sense that only degrees divisible by \(p\) occur. Indeed it
expresses this vector as divided Frobenius of a combination of
the two filtered generators in the preceding component.
Crystallinity gives \(\mu_i\in uR\) as well.

Exactly as in that lemma, the following coefficient of the
subline lies in \(u^{d-t_i+1}R\), \(R=k'[u]/(u^{ep})\):
\[
       u^{e+d-s_i}\mu_i+
              u^{e+d-s_i-t_i}(u y_i'-t_i y_i).
\]
If \(l_i\le s_i-e\), the lowest term of \(y_i\) contributes
nontrivially in degree \(d-t_i+e-s_i+l_i\le d-t_i<ep\).
Its degree is nonnegative by the inverse height condition.
Cancellation by the first summand would require a term of
degree \(l_i-t_i\) in \(\mu_i\), impossible by the congruence
above. This contradiction proves \(l_i>s_i-e\), hence
\(b_i>c_i\) at every index. Coordinate pullbacks and pushouts
as in Theorem~\ref{thm:oldscalars752} give the coefficient-ring
assertion. On forgetting the coefficient action the additional
tame characters are local conjugates, included in the test.
\end{proof}

\begin{corollary}[A count for a primitive marked orbit]
\leavevmode\label{cor:primitiveorbit753}
If, for some \(i\), the only character in \(\mathcal T\) of the
form \(\theta_i^j\), \(|j|\le eH/(p-1)\), is the trivial one,
all the scalar classes in the theorem satisfy the integral
finite condition. Suppose the nontrivial types consist of a
single primitive orbit of order \(e\), its inverses, and
characters of orders dividing \(Q\) with \(QH<p-1\).
Among all embeddings of the field of values of the primitive
character into \(\overline{\QQ}_p\), the proportion for which
this vanishing test can fail is at most
\[
             \frac{2f\lfloor eH/(p-1)\rfloor}{\varphi(e)}
                  \le\frac{2 f H(1+\log_2 e)}{p-1}.
\]
\end{corollary}
\begin{proof}
The absolute value bound for either shift is strictly less
than \(e\). A trivial character therefore forces that shift
to be zero, and both shifts cannot be zero at the same index.
An element of order dividing \(Q\) other than the identity
has least absolute exponent at least \(e/Q\); this exceeds
\(eH/(p-1)\).

Fix one component. The image of a primitive character, as the
coefficient embedding varies, has exponent running over the
units modulo \(e\), each once for its field of values. The
orbit gives the exponents \(p^j k\) for \(0\le j<f\). For each
\(j\), multiplication by \(p^j\) permutes the unit classes.
At most \(2\lfloor eH/(p-1)\rfloor\) of them have a nonzero
representative in the indicated interval. The interval is
symmetric, so inversion gives no further exceptions. Take the
union over \(j\). Finally if \(e\) has \(m\) distinct prime
divisors, its \(j\)-th smallest prime divisor is at least
\(j+1\). Thus
\(e/\varphi(e)\le\prod_{j=1}^m(1+1/j)=m+1\le1+\log_2 e\).
\end{proof}
At a primitive marked place of the shorter construction the
orbit in Lemma~\ref{lem:tameorbit715} has precisely this order
with \(f=r\), even after the prescribed scalar twist. For
families with the indicated types the bound therefore remains
effective when \(e\) is as large as \(p^r-1\). The count is
over coefficient embeddings. A bound for the supremum in the
period inequality must also control the exceptional embeddings
and the main representation of higher rank.
\clearpage

\section{Cup products along Frobenius orbits}\label{nav:frobenius-cups}

\subsection{Frobenius and the integral cup product}
\label{sec:cuppath757}
We give a cochain proof of the local pairing argument of
\path{New/Congruences/Lifting/CupProductBound.lean}.
It applies to matrices of arbitrary rank and retains the precision
of the original torsion cocycle. The geometric input to this
argument consists of lifts satisfying explicit inequalities.
Keeping track of the choices will let us apply the result to
the same dual lattice that defines the projected period.

Let \(K/\QQ_p\) be finite, \(G=G_K\), \(F=\mathbb C_p^\flat\),
and \(n\ge1\). The tilt \(F\) is an algebraically closed complete
valued field of characteristic \(p\), with continuous action of
\(G\). These are the usual tilting and finite-extension theorems
for perfectoid fields applied to \(\mathbb C_p\).
~\cite{III-II-scholzetilt757}
Write \(A_n=W_n(F)\) for the length-\(n\) Witt vectors, with
Witt Frobenius \(\phi\). Let \(S_n\subset A_n\) be the additive
subgroup of Witt vectors all of whose coordinates lie in the
maximal ideal \(\mathfrak m_F\). It is the kernel of residue
reduction on \(W_n(\mathcal O_F)\). We use \(A_n/S_n\) as an
additive group. For a Tate twist \((1)\), \(G\) acts also on
the twist and \(\phi\) acts on the Witt vector alone.

\begin{lemma}[Constants modulo the maximal ideal of Witt vectors]
\leavevmode\label{lem:wittconstants757}
The map \(\phi-1\) is onto \(A_n\) and bijective on \(S_n\).
The fixed subgroup of \(A_n/S_n\) under \(\phi\) is exactly the
image of \(\ZZ/p^n\), which embeds there. These statements also
hold after twisting by one. The quotient has the discrete
topology for the coordinate topology on \(A_n\).
\end{lemma}
\begin{proof}
On \(F\), surjectivity is the solvability of \(X^p-X=t\).
On successive quotients of the filtration of \(A_n\) by powers
of \(p\) the same assertion holds, since \(F\) is perfect.
Lifting a solution one quotient at a time proves surjectivity.
For \(z\in S_n\) a solution in \(S_n\) is
\(-\sum_{r\ge0}\phi^r z\). In Witt coordinates \(\phi\) is
the \(p\)-power map. All these coordinate valuations therefore
tend to infinity. For each positive lower bound on the
valuations the corresponding coordinate neighborhood of zero
is an additive subgroup, by the integral Witt addition and
inverse polynomials. Completeness gives the stated sum, still
in \(S_n\). The fixed vectors in \(A_n\) have coordinates
in \(\mathbb F_p\), so are \(W_n(\mathbb F_p)=\ZZ/p^n\).
None except zero belongs to \(S_n\). This proves injectivity
on \(S_n\) and on the constants. If \((\phi-1)t\in S_n\),
subtract its unique preimage in \(S_n\) from \(t\) to obtain
a fixed vector. Finally \(S_n\) is open, since \(\mathfrak m_F\)
is open and there are only finitely many coordinates.
\end{proof}

Let \(Y,Y'\) be finite free \(\ZZ_p\)-lattices with continuous
\(G\)-action, \(a,a'\) continuous integral 1-cocycles on these
lattices, and
\[
                P:Y\times Y'\longrightarrow(\ZZ/p^n)(1)
\]
a \(G\)-equivariant bilinear map. For nonnegative integers
\(i,j,s\), put \(b=p^ia,\ b'=p^ja'\) and extend \(B=p^sP\)
bilinearly over \(A_n\). Thus \(B\) takes values in \(A_n(1)\).
For \(u\in A_n\otimes_{\ZZ_p}Y\) and
\(u'\in A_n\otimes_{\ZZ_p}Y'\) use the displacements
\[
 d_u(g)=b(g)+gu-u,\qquad d'_{u'}(g)=b'(g)+gu'-u'.
\]
They are 1-cocycles for the ordinary linear action. Frobenius
fixes the integral terms in these formulas.

\begin{theorem}[A Frobenius path that annihilates a cup product]
\leavevmode\label{thm:cuppath757}
Suppose \(H\subseteq G\) is open of index prime to \(p\), and
there exist \(x,x'\) as above with
\[
     B\big(d_x(g),g d'_{x'}(h)\big)\in S_n(1)
                                 \qquad(g,h\in H).
\]
Suppose further that a finite sequence of pairs joins
\((x,x')\) to \((\phi x,\phi x')\), changing one component
at a time. Require the following condition on each step,
for every \(g\in H\).
\begin{enumerate}
\item In a step from \((u,u')\) to \((u+v,u')\) require
\(B(v,d'_{u'}(g))\in S_n(1)\).
\item In a step from \((u,u')\) to \((u,u'+v')\) require
\(B(d_u(g),gv')\in S_n(1)\).
\end{enumerate}
Then \(p^{i+j+s}[P(a\cup a')]=0\) in
\(H^2(K,(\ZZ/p^n)(1))\).
\end{theorem}
\begin{proof}
Use the left inhomogeneous cochain differential
\(\partial t(g)=gt-t\) and
\(\partial f(g,h)=g f(h)-f(gh)+f(g)\). Define the following
continuous cochain on \(H\), initially with values in \(A_n(1)\),
\[
             f_{u,u'}(g)=B(b(g),gu')-B(u,d'_{u'}(g)).
\]
Its differential is
\[
 \partial f_{u,u'}(g,h)=B(b(g),g b'(h))
                              -B(d_u(g),g d'_{u'}(h)).
\]
Indeed the first term in \(f\) has differential
\(-B(b(g),g(hu'-u'))\), by the cocycle identity for \(b\).
The differential of \(B(u,d'_{u'}(g))\) is
\(B(gu-u,g d'_{u'}(h))\), by the cocycle identity for
\(d'_{u'}\). Combining gives the formula. At \(u=x,u'=x'\)
its second term vanishes modulo \(S_n(1)\).

In a first-component step the change of \(f\) is
\(-B(v,d'_{u'}(g))\). In a second-component step it is
\[
          B(d_u(g),gv')-\partial B(u,v')(g).
\]
The second identity follows by expanding \(d_u(g)\) and
using \(B(gu,gv')=gB(u,v')\). Consequently, modulo \(S_n(1)\),
the first steps contribute zero and the second steps contribute
coboundaries. If \(t\) is the sum of \(B(u,v')\) over the
second steps, we obtain, for \(f=f_{x,x'}\),
\[
                 (\phi-1)f \equiv -\partial t\pmod{S_n(1)}.
\]
Here \(f_{\phi x,\phi x'}=\phi f_{x,x'}\) because \(\phi\)
commutes with \(H\), fixes \(b,b'\), and respects the pairing.

Choose \(y\in A_n(1)\) with \((\phi-1)y=t\) by
Lemma~\ref{lem:wittconstants757}. The image of \(f+\partial y\)
in the quotient is now fixed by \(\phi\) pointwise. Hence it is
the image of a unique cochain \(\kappa:H\to(\ZZ/p^n)(1)\).
It is continuous because \(f+\partial y\) is continuous and
the quotient is discrete. Taking its differential leaves the
differential of \(f\) unchanged in the quotient. Both that
differential and \(B(b(g),g b'(h))\) have the same reduction
there. The latter is constant in the Witt direction. Injectivity
on constants therefore gives the equality in finite coefficients
\[
       \partial\kappa(g,h)=p^{i+j+s}P(a(g),g a'(h))
                                            \qquad(g,h\in H).
\]
Thus the indicated class restricts to zero on \(H\).
Corestriction composed with restriction multiplies the class by
\([G:H]\). Multiplication by this index is invertible on the
\(p^n\)-torsion group, proving the result over \(K\).
\end{proof}

\begin{corollary}[A bound on the joint tests at the original precision]
\leavevmode\label{cor:pathperiod757}
Fix an actual datum \(D\) of Theorem~\ref{thm:projectedcups742},
a spherical coefficient place \(w\mid p\), and one target lattice
\(T\). Write \(K=E_{x,w}\). Let \(L_w\) be the free sublattice
spanned by the chosen lifts of the free finite dual tests in
Proposition~\ref{prop:freetests741}, with \(T'=T^\vee(1)\).
For a test \(y\in L_w\), use an integral cocycle representing it.
For each restricted global trace \(c\) at precision \(\pi^a\),
write \(c=t\circ a_Y\) as in Lemma~\ref{lem:graphrealization724},
with \(a_Y\) an integral cocycle on \(Y\) and
\(t:Y\to T/\pi^aT\) equivariant.

Form an \(\mathcal O\)-sublattice \(L'_w\subset L_w\) from
the elements \(p^k y\) for which every such trace admits a
path of Theorem~\ref{thm:cuppath757} with \(i+j+s\le k\),
for the pairing
\[
          (v,\lambda)\longmapsto \lambda(t(v))
                              \quad\hbox{in }(\mathcal O/\pi^a)(1)
\]
and the two cocycles on \(Y,T'\). Use one path for each cyclic
\(\ZZ_p\)-coordinate of the output if necessary. Any collection
of these generators may be used. The normalized length of the
local image in the free coordinates of this target is at most
\[
          \frac{1}{e_{\mathcal O}}\,
              \operatorname{length}_{\mathcal O}
                        \big(L_w/(L'_w+\pi^a L_w)\big),
             \qquad e_{\mathcal O}=v_\pi(p).
\]
This bound is valid for simultaneous tests and for arbitrary
local trace representations.
\end{corollary}
\begin{proof}
The graph construction realizes \(c\) at precisely the stated
modulus by a map of lattices, with no division of a local
length. The displayed pairing is the pairing computing
\(c\cup\operatorname{red} y\), by functoriality of cup products.
As a constant \(\ZZ_p\)-module \(\mathcal O/\pi^a\) is a finite
direct sum of cyclic \(p\)-groups. After twisting by one, each
coordinate projection is \(G_K\)-equivariant. Apply the theorem
to each output coordinate, choosing its order \(p^n\) as
modulus. Since their joint map is an isomorphism on coefficients,
their vanishing gives vanishing of the full cup with \(p^k y\).
The same holds on the entire \(\mathcal O\)-span of these tests.
Hence every local evaluation functional factors through
\(L_w/(L'_w+\pi^aL_w)\). For a cyclic module
\(\mathcal O/\pi^d\) with \(d\le a\), the maps into
\(\mathcal O/\pi^a\) form the ideal
\(\pi^{a-d}\mathcal O/\pi^a\) of length \(d\).
The elementary divisor theorem gives the bound after summation
and division by \(e_{\mathcal O}\).

In the period estimate the source of the local maps is the
actual joint global image. If any set of local targets is
discarded, the kernel of the projection to those retained
injects into the product of the discarded coordinate images.
Thus \(r(D)\) in Theorem~\ref{thm:projectedcups742} is bounded
by the length on the retained targets plus the sum of the
displayed bounds on the others. No field, test precision or
interval in that theorem changes. Obtaining the full period
bound this way requires quantitative control of these indices
for the unrestricted crystalline lattices.
\end{proof}
\clearpage

\subsection{Integral lifts of quotient generators at every precision}
\label{sec:kisinquotient759}
The two cocycles in Theorem~\ref{thm:cuppath757} define
extensions of a trivial lattice. To construct its Frobenius
path one needs control of lifts of the trivial quotient inside
the integral Breuil--Kisin modules. We bound the power of
the uniformizer parameter needed for each such lift. This is
the estimate underlying the lift ideals in
\path{New/Congruences/Lifting/CupProductBound.lean}. It applies
to both cocycles on the actual period fields, including when
those fields are ramified at the coefficient prime.

Let \(K/\QQ_p\) have residue field \(k\) and absolute ramification
index \(e_K\). Use the ring \(\mathfrak S=W(k)[[u]]\) with
Witt Frobenius on coefficients and \(\phi(u)=u^p\). The
Eisenstein polynomial \(E(u)\) of a chosen uniformizer is
monic of degree \(e_K\) and reduces to \(u^{e_K}\) modulo \(p\).
In particular \(u^{e_K}=E(u)-p q(u)\) for some \(q\in\mathfrak S\).

\begin{lemma}[A Frobenius bound for an ideal modulo \(p^\nu\)]
\leavevmode\label{lem:frobeniusideal759}
Suppose \(I\subseteq\mathfrak S\) is an ideal with nonzero
reduction in \(k[[u]]\) and
\[
                 E(u)^h I\ \subseteq\ \mathfrak S\,\phi(I)
                                    \qquad(h\ge0).
\]
For \(\nu\ge1\), let \(a_\nu\) be the least \(a\ge0\) for
which \(u^a\in I+p^\nu\mathfrak S\). Then \(a_\nu\) exists and
\[
                      (p-1)a_\nu\le e_K(h+\nu-1).
\]
If \(e_K h<p-1\), then \(I=\mathfrak S\) and \(a_\nu=0\)
for every \(\nu\).
\end{lemma}
\begin{proof}
The reduction of \(I\) is \((u^m)\) for some \(m\ge0\).
In \(\mathfrak S/I\) the element \(u^m\) is divisible by \(p\).
Raising it to the \(\nu\)-th power proves existence.
Work henceforth in \(\mathfrak S_\nu=\mathfrak S/p^\nu\)
and write \(I_\nu\) for the image of \(I\).
In this ring \(u^{e_K(h+\nu-1)}\) is divisible by \(E^h\).
Indeed, in the expansion of
\((E-pq)^{h+\nu-1}\) every term with fewer than \(h\) factors
of \(E\) has at least \(\nu\) factors of \(p\). Consequently
\[
       u^{a_\nu+e_K(h+\nu-1)}
                       \in E^h I_\nu\subseteq\mathfrak S_\nu\phi(I_\nu).
\]
Since \(k\) is perfect, Witt Frobenius is an automorphism on
the coefficient ring. Every series in \(\mathfrak S_\nu\)
therefore has a unique expression
\(\sum_{i=0}^{p-1}u^i\phi(f_i)\). In that expression the
extended ideal \(\mathfrak S_\nu\phi(I_\nu)\) consists exactly
of series with \(f_i\in I_\nu\) for all \(i\). Applying this
to the displayed monomial shows that
\(u^{\lfloor(a_\nu+e_K(h+\nu-1))/p\rfloor}\in I_\nu\).
Minimality of \(a_\nu\) gives the asserted inequality.
If \(e_Kh<p-1\), take \(\nu=1\) to get \(a_1=0\). Thus \(I\)
contains an element congruent to one modulo \(p\). It is a
unit in the local ring \(\mathfrak S\).
\end{proof}

\begin{theorem}[Lifts in a crystalline extension with a trivial quotient]
\leavevmode\label{thm:integrallifts759}
Let \(0\to Y\to V\to\ZZ_p\to0\) be an exact sequence of
\(G_K\)-lattices with crystalline rational middle term. Suppose
the jumps of \(V\), in the convention where the cyclotomic
character has jump \(-1\), lie in \([-r,h]\), \(r,h\ge0\).
Take the covariant Breuil--Kisin modules after twisting the
sequence by \(-r\), and choose a generator \(\mathbf t_r\)
of the rank-one module of \(\ZZ_p(-r)\). For the resulting map
\[
        \beta:\mathfrak M(V(-r))\longrightarrow\mathfrak S\mathbf t_r
\]
at precision \(p^\nu\), the quotient generator multiplied by
\(u^{a_\nu}\) has an integral lift, with
\((p-1)a_\nu\le e_K(h+\nu-1)\).
If \(e_K h<p-1\), the generator itself has a lift in
\(\mathfrak M(V(-r))\) before any reduction.
\end{theorem}
\begin{proof}
The covariant module of \(V(-r)\) has height at most \(r+h\).
This is the integral Breuil--Kisin correspondence for crystalline
lattices: in the convention giving \(\ZZ_p(1)\) Hodge--Tate
weight \(1\), the twist has weights in \([-(r+h),0]\). The
height assertion means that \(E^{r+h}\mathfrak M(V(-r))\)
is contained in the \(\mathfrak S\)-span of its Frobenius
image. The quotient line has Frobenius
\(\phi(\mathbf t_r)=(E/c_0)^r\mathbf t_r\), where
\(c_0=E(0)/p\in W(k)^\times\). We also use the integral
comparison after passage to the étale Frobenius module.
It identifies the scalar extension of \(\beta\) to
\(W(\mathbb C_p^\flat)\) with the given quotient map after
the same scalar extension. In particular it is surjective
there, also modulo \(p\). These standard properties of the
covariant functor, including this normalization of the Tate
line, are recorded in Sections~2--3 of the cited paper.
~\cite{III-II-cgmexact759}

Write the image as \(I\mathbf t_r\). Its reduction modulo
\(p\) is nonzero by the étale comparison: otherwise the
scalar extension modulo \(p\) of the map to the line would
vanish. Applying \(\beta\) to the height inclusion gives
\[
                E^{r+h}I\subseteq (E/c_0)^r\mathfrak S\phi(I).
\]
Cancel \(E^r\) in the integral domain \(\mathfrak S\) and
apply Lemma~\ref{lem:frobeniusideal759}. The resulting
membership \(u^{a_\nu}\in I+p^\nu\mathfrak S\) is exactly
the assertion of a lift at that precision. In the case \(h=0\)
the integral surjectivity is also the conclusion about the
quotient in the exactness theorem of \v{C}oupek, Gazaki and
Marmora. The argument here gives the bound at positive \(h\)
as well.
\end{proof}

\begin{corollary}[The two lifts for the period data]
\leavevmode\label{cor:periodlifts759}
Let \(K=E_{x,w}\) be any spherical completion of the marked
period data and apply Corollary~\ref{cor:pathperiod757}.
For every scalar target \(\mathcal O(-d)\), \(0\le d\le B\),
the integral crystalline extension attached to a finite test
of the dual \(\mathcal O(d+1)\) has a lift of its quotient
generator with \(a_\nu=0\) at every precision, over every \(K\).
The same holds for the main dual when \(J_\tau\ge-1\) at every
embedding of \(K\).

For the trace side use the extension from
Lemma~\ref{lem:graphrealization724}. At every spherical place it
admits the bound of Theorem~\ref{thm:integrallifts759} with
\[
                 h=\max\big(\{0,B\}\cup\{J_\tau:\tau\mid w\}\big).
\]
In particular, when all \(J_\tau<0\) and \(e_K B<p-1\), the
quotient generator on the trace side also lifts integrally for
arbitrary ranks, even with arbitrarily large negative jumps.
\end{corollary}
\begin{proof}
For a single integral cocycle on an \(\mathcal O\)-lattice,
forget coefficients to \(\ZZ_p\) and take the extension of
\(\ZZ_p\) defined by that cocycle. This is the pullback of its
\(\mathcal O\)-extension along the inclusion of \(\ZZ_p\) in
the trivial quotient. It is crystalline for a finite dual
test, by the defining rational finite condition. The trace
side extension is rationally split with crystalline
constituents. These constructions keep the original cocycles
on the two kernel lattices. Under restriction of coefficients
the jumps are taken over all labelled embeddings. The jumps
on each extension are the union of those on the kernel with
zero. For the scalar dual the kernel jump is \(-d-1\); for
the main dual its jumps are \(-1-J_\tau\). The respective
sign conditions thus give \(h=0\). For the trace side the
kernel has weights from the scalar string and the main jump
by the graph description. This proves all the claims by the
theorem, including the assertion at ramified places.

To relate these lifts precisely to the ambient modules of
Theorem~\ref{thm:cuppath757}, send \(u\) to the Teichm\"uller
lift \([\underline\varpi]\) of a compatible tower of roots
of the uniformizer in \(F=\mathbb C_p^\flat\). Choose a basis
of \(\ZZ_p(-r)\) in the comparison and write \(\xi_r\) for the
coefficient of the image of \(\mathbf t_r\) there. It is a
unit in \(W(F)\), since it generates the quotient line after
this base change. Frobenius gives
\(\phi(\xi_r)=(E([\underline\varpi])/c_0)^r\xi_r\).
With \(v_F(\underline\varpi)=1/e_K\), reduction modulo \(p\)
shows \(v_F(\bar\xi_r)=r/(p-1)\), because
\(\bar E(\underline\varpi)=\underline\varpi^{e_K}\).
Thus division by \([\underline\varpi]^{a_\nu}\xi_r\) in
\(W_\nu(F)\), and removal of the chosen basis of the twist,
turn the integral lift into an affine lift over one. The
valuation of the reduced divisor is exactly
\(a_\nu/e_K+r/(p-1)\). This constructs starting vectors for
the original precision without any triangularity requirement.
The estimates on their pairings with Galois displacements,
which occur in the path hypotheses of
Theorem~\ref{thm:cuppath757}, are additional requirements.
\end{proof}
\clearpage

\subsection{Uniform integral lifts from Galois displacements}
\label{sec:uniformlifts760}
Retain the notation of Theorem~\ref{thm:integrallifts759}.
We convert a power of \(u\) in the image ideal at finite
precision to a power of \(p\) in the integral image. The
result is uniform in the torsion precision and in the rank
of the crystalline representation. The conversion uses
the Galois action on the Breuil--Kisin module after scalar
extension.

Put \(F=\mathbb C_p^\flat\) and \(A=W(\mathcal O_F)\), and
embed \(\mathfrak S\) in \(A\) using \(u=[\underline\varpi]\).
Choose compatible primitive \(p^n\)-th roots of unity and
write \(\epsilon\in\mathcal O_F\) for the resulting element.
Set \(\mu=[\epsilon]-1\) and \(\tau=\phi^{-1}(\mu)\).
For a crystalline lattice of effective height let \(M\) be
its covariant Breuil--Kisin module and \(M_A=A\otimes_{\mathfrak S}M\).
The crystalline comparison supplies the semilinear \(G_K\)-action
on \(M_A\) and gives the divisibility
\[
                 (g-1)M\subseteq u\tau M_A\qquad(g\in G_K).
\]
This is the necessity direction of the crystalline descent
theorem of Gee and Liu, in its covariant convention, for
the chosen tower of uniformizer roots. It also applies to
the trivial module and to the rank-one quotient of a twisted
extension.~\cite[Appendix~F, (F.10) and Theorem~F.11]{III-II-geeliudesc760}

\begin{lemma}[Trading one parameter for a power of the prime]
\leavevmode\label{lem:kummerlower760}
Let \(I\mathbf t_r\) be the image of the map \(\beta\) in
Theorem~\ref{thm:integrallifts759}. If \(p^c u^m\in I+p^\nu
\mathfrak S\), where \(m\ge1\), \(c\ge0\) and \(\nu\ge1\), then
\[
             p^{c+1+v_p(m)}u^{m-1}\in I+p^\nu\mathfrak S .
\]
\end{lemma}
\begin{proof}
For \(f\in I\) lift \(f\mathbf t_r\) to \(M\). Put
\(g(\mathbf t_r)=b_g\mathbf t_r\). The factor \(b_g\) is a
unit of \(A\), congruent to one modulo \(u\tau A\) by crystalline
descent. Equivariance of \(\beta\) and the same divisibility
for \(M\) give
\[
                     b_g g(f)-f\in u\tau IA.
\]
In particular \(g(f)\in IA\), so subtraction gives
\((g(f)-f)/(u\tau)\in IA\). For \(t\in\mathfrak S\), descent
for the trivial module gives \((g(t)-t)/(u\tau)\in A\).
Applying the two statements to \(f+p^\nu t\) proves that
the divided displacement of every element of \(I+p^\nu\mathfrak S\)
belongs to \((I+p^\nu\mathfrak S)A\).

Choose \(g\in G_K\) with \(g(u)=[\epsilon]u\). Indeed
\(X^{p^n}-\varpi\) is Eisenstein for every \(n\), and the
embeddings sending the chosen root to its product with
\(\zeta_{p^n}\) are compatible. Their union extends to the
algebraic closure over \(K\). Applying the divided displacement
to the monomial in the hypothesis gives
\[
        p^c u^{m-1}\frac{[\epsilon]^m-1}{\tau}
                                  \in (I+p^\nu\mathfrak S)A .
\]
To compute the factor modulo \(u\), write \(m=p^v m_0\) with
\(p\nmid m_0\). The geometric sum for \(m_0\) evaluated at
\([\epsilon]^{p^v}\) is a unit of \(A\), since its residue in
the residue field of \(A\) is \(m_0\). The remaining factor
\(([\epsilon]^{p^v}-1)/\tau\) equals
\(\prod_{j=0}^{v}\phi^j(\mu/\tau)\). Both \(\mu/\tau\) and
\(E(u)\) are generators of the kernel of Fontaine's map
\(\theta:A\to\mathcal O_{\mathbb C_p}\). Thus each factor
is a unit times \(\phi^j(E(u))\), which is congruent modulo
\(u\) to \(p\phi^j(c_0)\) for \(c_0=E(0)/p\). It follows
that the entire factor is \(p^{v+1}\) times a unit, modulo
\(uA\). The error after multiplication by \(p^c u^{m-1}\)
is in \(p^cu^mA\), already in the extended ideal. This gives
the required membership over \(A\).

It descends modulo \(p^\nu\) to \(\mathfrak S\). Indeed
\(A/p^\nu\) is faithfully flat over \(\mathfrak S/p^\nu\).
Modulo \(p\) this is the local injection of the discrete
valuation ring \(k[[u]]\) into the valuation ring \(\mathcal
O_F\) taking \(u\) to an element of positive valuation.
It is flat by torsion freeness and faithful since the maximal
ideal stays proper. For higher \(\nu\) use the flatness
criterion for a nilpotent ideal with the ideal generated by
\(p\). Its required first Tor vanishes by the resolution with
alternating differentials \(p\) and \(p^{\nu-1}\), since
\(A\) is \(p\)-torsion free. Faithfulness follows from the
same reduction modulo \(p\).
\end{proof}

\begin{theorem}[A lift independent of the precision]
\leavevmode\label{thm:uniformlift760}
Write \(e=e_K\) and keep the upper jump bound \(h\) of
Theorem~\ref{thm:integrallifts759}. If \(eh<p-1\) take \(C=0\).
Otherwise suppose \(e\le p-2\) and take any nonnegative
integer \(C\) satisfying
\[
             C\big((p-1)^2-pe\big)\ \ge\ peh.
\]
Then \((p,u)^C\subseteq I\). In particular the constant
\(p^C\mathbf t_r\) has an integral lift, and so does
\(p^s E(u)^t\mathbf t_r\) whenever \(s+t\ge C\).
At an unramified place with \(p\ge5\) one may always take
\(C=\lceil h/(p-3)\rceil\).
\end{theorem}
\begin{proof}
The \(C=0\) case is Lemma~\ref{lem:frobeniusideal759}. In
the second case the coefficient of \(C\) is positive.
Work modulo \(p^{C+1}\) and apply that lemma to find \(u^a\)
in the reduced image with \((p-1)a\le e(h+C)\). Legendre's
formula gives \(v_p(a!)\le a/(p-1)\). Hence
\[
                a+v_p(a!)\le \frac{pe(h+C)}{(p-1)^2}\le C .
\]
For \(m<a\), apply Lemma~\ref{lem:kummerlower760} at
\(a,a-1,\ldots,m+1\). The total exponent paid is
\(a-m+v_p(a!)-v_p(m!)\le C-m\).
For \(a\le m\le C\), simply multiply \(u^a\) by \(u^{m-a}\).
Thus \(p^{C-m}u^m\in I+p^{C+1}\mathfrak S\) for every
\(0\le m\le C\).

At \(m=0\) this means \(p^C(1-pb)\in I\) for some
\(b\in\mathfrak S\). The factor in parentheses is a unit.
Thus \(p^C\in I\), which absorbs the modulus in all the other
memberships as well. These monomials generate \((p,u)^C\).
Since \(E\in(p,u)\), the stated products lie in the image.
For \(e=1\) one has \((p-1)^2-p=p(p-3)+1\), proving the
last assertion.
\end{proof}

\begin{corollary}[Uniform lifts on the marked period fields]
\leavevmode\label{cor:uniformperiodlift760}
At each spherical place of the data in
Theorem~\ref{thm:projectedcups742} to which the preceding
theorem applies, the graph extension of every restricted trace
admits an integral lift of \(p^C\mathbf t_r\) with
\(h=\max(0,B,\max_{\tau\mid w}J_\tau)\). For a scalar dual
test take \(C=0\) at every spherical place, including the
places of larger ramification. For the main dual use
\(h=\max(0,\max_{\tau\mid w}(-1-J_\tau))\) in the theorem.
These choices are independent of the test precision and the
number of trace matrices.
\end{corollary}
\begin{proof}
These are the two crystalline integral extensions of
Corollary~\ref{cor:periodlifts759}; their quotient is the
trivial lattice before twisting. It suffices to apply the
theorem to each one, using the labelled jumps computed in
that corollary. The resulting lift can be reduced to every
precision in the same period datum.
\end{proof}
\clearpage

\subsection{Uniform integral lifts over ramified fields}
\label{sec:ramifiedlifts761}
For large absolute ramification index, passing repeatedly through
Frobenius decreases the parameter exponent before we convert
it to a power of \(p\). The cost of a passage depends on the
different of the local field and the upper Hodge jump. We keep
the same period field, uniformizer tower and coefficient prime
as in Theorem~\ref{thm:integrallifts759}.

Write \(e=e_K\), \(d_K=\lceil v_p(E'(\varpi))\rceil\), and
\(I_\nu=I+p^\nu\mathfrak S\), where \(I\) is the ideal of
that theorem. In particular \(d_K=0\) if \(e=1\), and
\(d_K\le1+v_p(e)\) otherwise. Indeed the nonzero terms in
the derivative of an Eisenstein polynomial, evaluated at
\(\varpi\), have valuations with distinct fractional parts
\((i-1)/e\), \(1\le i\le e\). Their minimum is at most the
valuation \(v_p(e)+(e-1)/e\) of the leading term.

\begin{lemma}[One passage through Frobenius]
\leavevmode\label{lem:transportpower761}
Let \(h\) be the upper jump bound from
Theorem~\ref{thm:integrallifts759}. For \(\nu\ge1\) one has
\[
 p^cu^a\in I_\nu
 \quad\Longrightarrow\quad
 p^{c+d_K h+v_p(h!)}u^{\lfloor(a+h)/p\rfloor}\in I_\nu .
\]
\end{lemma}
\begin{proof}
Let \(J=\mathfrak S\phi(I)\) and \(J_\nu=J+p^\nu\mathfrak S\).
We first prove an operation in \(J_\nu\). Choose
\(g(u)=q u\) with \(q=[\epsilon]\) as in
Lemma~\ref{lem:kummerlower760}, and retain \(A,\mu,\tau\)
from that lemma. Then for \(f\in J_\nu\) one has
\[
                     \frac{g(f)-f}{u\mu}\in J_\nu A.
\]
Indeed \(J\) is the image ideal on the \(\mathfrak S\)-span
of the Frobenius transforms of the integral modules, using
\(\phi(\mathbf t_r)\) as the generator on the quotient.
For a generator \(\phi(m)\), the displacement is
\(\phi((g-1)m)\), divisible by \(u^p\mu\) in the \(A\)-span
of these Frobenius transforms. For a coefficient in
\(\mathfrak S\) the displacement is divisible by \(u\mu\),
since \(g\) fixes \(W(k)\) and \(g(u)-u=u\mu\).
For a power series this follows term by term, with its divided
series converging in the Witt coordinate topology. The induced
action on the quotient generator differs from one by a multiple
of \(u\mu\). The lift and subtraction used in
Lemma~\ref{lem:kummerlower760} now prove the inclusion, including
the summand \(p^\nu\mathfrak S\).

The element \(g(E)=\omega E\) has \(\omega\in A^\times\),
because both elements generate the kernel of \(\theta\). Write
\(E=\sum_i c_i u^i\) and \([i]_q=1+\cdots+q^{i-1}\) for
\(i\ge1\). If \(\eta=\mu/(E\tau)\in A^\times\), then
\[
     D=\frac{\omega-1}{u\tau}
         =\eta\sum_{i\ge1} c_i[i]_q u^{i-1},\qquad
     \theta(D)=\theta(\eta)E'(\varpi).
\]
The last identity uses \(\theta(q)=1\). Since
\(\theta(\tau)=\zeta_p-1\), we have
\[
   v_p(\theta(\omega)-1)
      =1/e+1/(p-1)+v_p(E'(\varpi))>1/(p-1).
\]
For any \(n\ge1\) set \(D_n=(\omega^n-1)/(u\tau)\).
At that depth the \(p\)-th power adds one to the valuation
of a principal unit minus one, by the binomial formula.
Raising to an exponent prime to \(p\) preserves the valuation.
Applying these facts to \(\theta(\omega)\) gives
\[
           v_p(\theta(D_n))=v_p(E'(\varpi))+v_p(n).
\]
Thus some multiple of \(D_n\) in \(A\) is congruent to
\(p^{d_K+v_p(n)}\) modulo \(E\): the required multiplier
modulo \(E\) is integral in \(\mathcal O_{\mathbb C_p}\),
and \(\theta:A\to\mathcal O_{\mathbb C_p}\) is onto.

Suppose \(f=p^c u^a E^n\in J_\nu\), with \(n\ge1\). Multiplying
its divided displacement by \(u\) gives the identity
\[
 \frac{g(f)-f}{\mu}
  =p^c u^{a+1}E^{n-1}\frac{E\tau}{\mu}D_n
                       + f\,\omega^n\frac{q^a-1}{\mu}.
\]
The second term belongs to \(J_\nu A\), as does the left side.
Multiply the first term by a unit and use the congruence for
\(D_n\). Its error is in \(ufA\). Hence
\(p^{c+d_K+v_p(n)}u^{a+1}E^{n-1}\in J_\nu A\).
Faithful flatness modulo \(p^\nu\), proved in
Lemma~\ref{lem:kummerlower760}, returns this membership to
\(\mathfrak S\).

Now \(p^c u^a E^h\in J_\nu\) by the height inclusion
\(E^h I\subset J\). Apply the operation \(h\) times. It
gives \(p^{c+d_Kh+v_p(h!)}u^{a+h}\in J_\nu\).
Decomposing the series modulo \(p^\nu\) according to the
exponent of \(u\) modulo \(p\), as in
Lemma~\ref{lem:frobeniusideal759}, returns the assertion
in \(I_\nu\).
\end{proof}

\begin{theorem}[Integral quotient lifts at every ramification index]
\leavevmode\label{thm:allramifiedlifts761}
Keep the crystalline extension and the ideal \(I\) of
Theorem~\ref{thm:integrallifts759}, with arbitrary \(K/\QQ_p\).
For \(J\ge0\) put
\[
 \alpha_J=\frac{pe}{(p-1)^2p^J},\qquad
 b_J=J\big(d_Kh+v_p(h!)\big).
\]
Choose any \(J\) with \(\alpha_J<1\). The ideal \(I\) contains
\(p^C\), independently of the torsion precision, for every
integer \(C\) satisfying
\[
 (1-\alpha_J)C\ \ge\
       b_J+\left(\alpha_J+
              \frac{p(1-p^{-J})}{(p-1)^2}\right)h.
\]
More generally \(p^{b_J}(p,u)^{C-b_J}\subset I\).
Denote the least nonnegative \(C\) allowed by these choices
by \(C_K(h)\), using \(C_K(h)=0\) when \(eh<p-1\).
\end{theorem}
\begin{proof}
Work at precision \(\nu=C+1\). Lemma~\ref{lem:frobeniusideal759}
starts with \(u^{a_0}\in I_\nu\) and
\((p-1)a_0\le e(h+C)\). Successive applications of
Lemma~\ref{lem:transportpower761} give \(p^{b_J}u^{a_J}\in
I_\nu\), where \(a_{j+1}=\lfloor(a_j+h)/p\rfloor\). Summing
the geometric progression yields
\[
 a_J\le p^{-J}\frac{e(h+C)}{p-1}
                     +(1-p^{-J})\frac h{p-1}.
\]
Multiplying by \(p/(p-1)\) bounds \(a_J+v_p(a_J!)\).
The assumed inequality on \(C\) is precisely sufficient for
\(b_J+a_J+v_p(a_J!)\le C\). Applying
Lemma~\ref{lem:kummerlower760} to remove each remaining
parameter factor now gives \(p^C\in I_\nu\). As in
Theorem~\ref{thm:uniformlift760} this gives \(p^C\in I\).
For \(m<a_J\) one may stop the removals at \(m\), obtaining
\(p^{C-m}u^m\in I_\nu\). For \(a_J\le m\le C-b_J\) this
membership follows by multiplying \(p^{b_J}u^{a_J}\) itself.
The modulus is absorbed by \(p^C\in I\), proving the ideal
inclusion. Such \(J\) exist for every \(e\), including
when \(p=2\).
\end{proof}

\begin{corollary}[The three integral extensions used by a cup test]
\leavevmode\label{cor:threeextensions761}
At any spherical coefficient place of the marked period data,
let \(Y,Y'\) be the graph and dual-test kernel lattices of
Corollary~\ref{cor:pathperiod757}, considered over \(\ZZ_p\).
Write \(X,X'\) for the extensions of \(\ZZ_p\) defined by
their two integral cocycles. For a cyclic coordinate of the
finite output of order \(p^n\), write
\(P:Y\otimes Y'\to(\ZZ/p^n)(1)\) for its equivariant
pairing. There is a third crystalline extension \(X''\)
of \(\ZZ_p\) with kernel \(Z=(Y\otimes Y')^\vee(1)\) and an
integral equivariant map \(X''\to Z\) whose reduction modulo
\(p^n\) is \(P\) times projection onto the trivial quotient.

For each of \(X,X',X''\), take \(h\) to be the maximum of zero
and its kernel jumps, over the coefficient embeddings at
that place. The corresponding power \(p^{C_K(h)}\) has an
integral quotient lift in its normalized Kisin module, at
every spherical place and for every value of \(n\).
\end{corollary}
\begin{proof}
Choose a \(\ZZ_p\)-linear lift \(\psi\in Z\) of \(P\). The
cocycle \(b(g)=(g\psi-\psi)/p^n\) is continuous and integral
because \(P\) is equivariant. Its class is rationally zero.
On \(Z\oplus\ZZ_p\) define
\(g(z,t)=(gz+t b(g),t)\). The map \((z,t)\mapsto p^n z+t\psi\)
has the asserted equivariance and reduction. The rational
kernel is crystalline by tensor products and Tate duality,
and the extension is rationally split. Its kernel jumps are
the sums \(-1-j-j'\) for the jumps \(j,j'\) in the tensor
product, at the same embedding. Adding zero gives all the
jumps on each extension. Choose the twist by \(-r\) in
Theorem~\ref{thm:integrallifts759} separately on the three
extensions, then apply Theorem~\ref{thm:allramifiedlifts761}.
Dividing by the common rank-one Tate generator restores the
untwisted normalization and projection to the stated power
of \(p\). All three constructions use the original \(K\)
and the original cyclic output.
\end{proof}
\clearpage

\subsection{A local cup bound from two adjacent Frobenius lattices}
\label{sec:adjacentcups762}
Let \(K/\QQ_p\) be finite. We use the rings \(F,A,\mathfrak S\)
and elements \(u,E,\mu,\tau\) of Lemma~\ref{lem:kummerlower760}.
For a crystalline \(\ZZ_p\)-lattice \(V\) write \(N(V)\) for
its covariant Breuil--Kisin module in \(W(F)\otimes_{\ZZ_p}V\).
Concretely, twist by \(-r\) until the jumps are nonnegative,
then tensor the effective module with the inverse of the
rank-one module of \(\ZZ_p(-r)\). In the comparison this
divides by the generator of the same Tate line on both
sides. Put
\[
       L(V)=\mathfrak S\,\phi N(V),\qquad
       N_A(V)=A N(V),\qquad L_A(V)=A L(V).
\]
Here \(\phi\) acts on the Witt vector and fixes \(V\). These
are two lattices in a common ambient module. Changing \(r\)
preserves the normalized module by compatibility with Tate
twists. More generally \(V\mapsto N(V)\) respects tensor
products, the unit and integral equivariant maps. One
description is to take the inverse of the tensor equivalence
between crystalline lattices and prismatic \(F\)-crystals,
then evaluate on the Breuil--Kisin prism. Both operations
respect tensor products. The covariant comparison identifies
the result with the modules just described.
~\cite[Theorems~5.6 and~7.9]{III-II-prismatic762}~\cite[Sections~2--3]{III-II-cgmexact759}

\begin{lemma}[The two lattices, their displacements and the Tate line]
\leavevmode\label{lem:adjacentlattice762}
If the jumps of \(V\) lie in \([-r,h]\), then
\[
      E^h N(V)\ \subseteq\ L(V)\ \subseteq\ E^{-r}N(V).
\]
Both \(A\)-lattices are stable under \(G_K\). For every
\(g\in G_K\) the displacements of the indicated
\(\mathfrak S\)-lattices satisfy
\[
   (g-1)N(V)\subseteq u\tau N_A(V),\qquad
   (g-1)L(V)\subseteq u\mu L_A(V).
\]
In the Witt comparison the two \(A\)-lattices of \(\ZZ_p(1)\)
are \(\tau^{-1}A(1)\) and \(\mu^{-1}A(1)\), respectively.
For an exact sequence \(0\to Y\to X\to\ZZ_p\to0\) of
crystalline lattices, the kernel of projection on each
\(A\)-lattice of \(X\) is the corresponding \(A\)-lattice of \(Y\).
\end{lemma}
\begin{proof}
Before untwisting, the span of Frobenius is contained in the
effective module and contains its multiple by \(E^{r+h}\).
Untwisting divides Frobenius on the line by \((E/c_0)^r\);
this gives both inclusions. For an effective module the first
displacement assertion is the crystalline Galois descent
formula used in Lemma~\ref{lem:kummerlower760}. It also gives
stability after extending scalars to \(A\). On the Tate
generator the ratio of its translate to itself is a unit
congruent to one modulo \(u\tau\). Its inverse has the
same property, so the assertion passes to the untwisted
module.~\cite[Appendix~F, (F.10) and Theorem~F.11]{III-II-geeliudesc760}

For an element \(\phi(m)\) the first displacement bound gives
\(\phi(u\tau)=u^p\mu\) at the second lattice. For a coefficient
\(a\in\mathfrak S\), one has \((g-1)a\in u\mu A\). Indeed \(g\)
fixes \(W(k)\) and \(g(u)=[\epsilon^{c}]u\) for some \(c\in\ZZ_p\).
The difference \([\epsilon^{c}]-1\) is divisible by \(\mu\).
To see this, expand \((1+\mu)^c\) by the binomial series,
first for nonnegative integers and then by continuity for
\(c\in\ZZ_p\). The series converges in the Witt coordinate
topology since \(\mu\) reduces to zero in the residue Witt
vectors. Subtracting the two power series \(a(g(u)),a(u)\)
now gives the claim, with convergence in the same topology.
Applying the product rule to sums of \(a\phi(m)\) gives the
second displacement bound. Stability also follows by commuting
the Galois action with \(\phi\).

For the Tate line let \(\xi\) be the coefficient of a generator
of \(N(\ZZ_p(-1))\) in the basis of the comparison lattice.
Then \(\phi(\xi)/\xi=E/c_0\) and
\(\phi(\tau)/\tau=\mu/\tau\). The ratio \(v\) of these
two elements belongs to \(A^\times\), since \(E\) and
\(\mu/\tau\) generate \(\ker\theta\). There is \(b\in A^\times\)
with \(\phi(b)/b=v\). Modulo \(p\) choose a \((p-1)\)-st root
of \(\bar v\) in \(\mathcal O_F^\times\). Given a solution
modulo \(p^m\), correcting it by \(1+p^m t\) requires solving
\(\bar t^p-\bar t=a\) with \(a\in\mathcal O_F\). A root exists
in \(F\) by algebraic closure and has nonnegative valuation
by the equation. Completeness gives the required unit.
Now \(\xi/(b\tau)\) is fixed by \(\phi\) and is a unit in
\(W(F)\), so it belongs to \(\ZZ_p^\times\). Taking the
inverse Tate line gives \(N_A(\ZZ_p(1))=\tau^{-1}A(1)\);
applying \(\phi\) gives the formula for \(L_A\).

The covariant Kisin functor preserves kernels on these
sequences. This is Proposition~3.4 of the exactness paper,
applied after a common effective twist. The map
\(\mathfrak S\to A\) is faithfully flat. Indeed modulo \(p\)
it is the faithfully flat map from \(k[[u]]\) to the valuation
ring \(\mathcal O_F\), and the \(p\)-adic flatness criterion
applies to these \(p\)-complete, \(p\)-torsion-free rings with
noetherian source. Thus the kernel identity holds for \(N_A\).
Applying the automorphism \(\phi\) of \(A\) gives the identity
for \(L_A\).~\cite[Sections~2.1 and~3, Proposition~3.4]{III-II-cgmexact759}
\end{proof}

\begin{proposition}[A quotient lift in the intersection]
\leavevmode\label{prop:commonlift762}
Let \(\beta:X\to\ZZ_p\) be a surjection of crystalline
lattices, \(h\) the largest of zero and the jumps of \(X\),
and \(d_K,C_K(h)\) as in Theorem~\ref{thm:allramifiedlifts761}.
Set
\[
 Q_K(X)=
 \begin{cases}
 C_K(h),&\text{if all jumps are nonnegative or all are nonpositive},\\
 C_K(h)+d_Kh+v_p(h!),&\text{otherwise}.
 \end{cases}
\]
The intersection \(N(X)\cap L(X)\) contains a vector over
\(p^{Q_K(X)}\).
\end{proposition}
\begin{proof}
By Theorem~\ref{thm:allramifiedlifts761} choose
\(x_0\in N(X)\) over \(p^C\), with \(C=C_K(h)\).
If all jumps are nonpositive, \(N(X)\subseteq L(X)\) by
the lemma; if all are nonnegative, \(L(X)\subseteq N(X)\)
and \(\phi(x_0)\) is in both lattices over the same constant.

For arbitrary jumps put \(M=N(X)\), \(L=L(X)\) and
\(H_n=E^n M\cap L\) for \(0\le n\le h\). Flatness identifies
\(A H_n\) with \(E^n M_A\cap L_A\), by taking the kernel
of subtraction in \(E^{-r}M\). We claim that whenever
\(n\ge1\) and \(p^cE^n\in\beta(H_n)\), one has
\[
                 p^{c+d_K+v_p(n)}E^{n-1}\in\beta(H_{n-1}).
\]
Choose \(y=E^n m\in H_n\) lifting \(p^cE^n\). Use the
element \(g\) of Lemma~\ref{lem:transportpower761}, with
\(\omega=g(E)/E\) and \(D_n=(\omega^n-1)/(u\tau)\).
The vector \((gy-y)/(u\mu)\) belongs to \(L_A\) by the
displacement assertion for \(L\). It also belongs to
\(E^{n-1}M_A\), since it equals
\[
 E^{n-1}\frac{E\tau}{\mu}
      \left(\omega^n\frac{gm-m}{u\tau}+D_n m\right).
\]
Its projection is \(p^c E^{n-1}(E\tau/\mu)D_n\).
The factor \(E\tau/\mu\) is a unit. By the calculation of
\(\theta(D_n)\) in Lemma~\ref{lem:transportpower761}, some
multiple of \(D_n\) equals \(p^{d_K+v_p(n)}\) modulo \(E\).
Consequently \(p^{c+d_K+v_p(n)}E^{n-1}\) belongs to the
extended image ideal \(\beta(H_{n-1})A\). Indeed the error
is a multiple of \(p^cE^n=\beta(y)\), already in that ideal.
Faithful flatness contracts the ideal to \(\beta(H_{n-1})\).
This proves the claim.

The height inclusion puts \(E^hx_0\) in \(H_h\) over
\(p^CE^h\). Apply the claim for \(n=h,h-1,\ldots,1\).
The resulting exponent is \(C+d_K h+v_p(h!)\), and
\(H_0=M\cap L\).
\end{proof}

\begin{theorem}[A precision-independent bound for crystalline cup tests]
\leavevmode\label{thm:unrestrictedcups762}
Let \(a,a'\) be integral cocycles on \(\ZZ_p\)-lattices
\(Y,Y'\), whose extensions \(X,X'\) of \(\ZZ_p\) are
crystalline. Let \(P:Y\otimes Y'\to(\ZZ/p^n)(1)\) be any
equivariant pairing. Put \(Z=(Y\otimes Y')^\vee(1)\), and
let \(h_Z\) be the maximum of zero and its jumps. Then
\[
   p^{\,Q_K(X)+Q_K(X')+C_K(h_Z)}[P(a\cup a')]=0
                   \quad\text{in }H^2(K,(\ZZ/p^n)(1)).
\]
The bound holds for every \(K/\QQ_p\), every \(n\ge1\),
and arbitrary ranks.
\end{theorem}
\begin{proof}
Write \(i=Q_K(X),\ j=Q_K(X'),\ s=C_K(h_Z)\). Construct
the third extension \(X''\) and its map to \(Z\) as in
Corollary~\ref{cor:threeextensions761}. That construction
uses just the lattices and \(P\), so it applies to the data
here. Evaluation gives an integral equivariant map
\[
       X''\otimes Y\otimes Y'\longrightarrow\ZZ_p(1)
\]
whose reduction modulo \(p^n\) is projection on the first
quotient followed by multiplication by \(P\). Take
\(z\in N(X'')\) over \(p^s\). Tensor compatibility sends
\(z,N_A(Y),N_A(Y')\) to \(N_A(\ZZ_p(1))\). Applying
\(\phi\), it sends \(\phi z,L_A(Y),L_A(Y')\) to
\(L_A(\ZZ_p(1))\). In both cases reduce the integral
identity in \(W_n(F)\). The scaled pairing \(B=p^sP\)
on the reductions of the two kernel lattices takes values
in the image of \(\tau^{-1}A(1)\) at the first level and
of \(\mu^{-1}A(1)\) at the second. These conclusions use
integral representatives of vectors in each level before
reducing the pairing.

Choose vectors
\[
 y\in N(X)\cap L(X),\quad y'\in N(X')\cap L(X')
\]
over \(p^i,p^j\) by Proposition~\ref{prop:commonlift762}.
Set \(x=\phi^{-1}y,\ x'=y'\) in the Witt comparison. Both
belong to their respective first \(A\)-lattices, because
\(\phi^{-1}L_A(X)=N_A(X)\). Under fixed integral
splittings of the two extensions they specify the affine
points in Theorem~\ref{thm:cuppath757} after reduction.
For example \(gx-x\) in the extension has kernel component
\(p^ia(g)+g x_{\rm ker}-x_{\rm ker}\). We denote these
displacements by \(d_x\) and \(d'_{x'}\).

Since \(y\in L(X)\), its displacement lies in \(u\mu L_A(Y)\).
Indeed it belongs to the indicated multiple of \(L_A(X)\)
by Lemma~\ref{lem:adjacentlattice762}; its quotient coordinate
vanishes, and division by \(u\mu\) still gives zero on the
quotient since \(A\) is a domain. Applying \(\phi^{-1}\)
gives \(d_x(g)\in\phi^{-1}(u)\tau N_A(Y)\).
Likewise \(d'_{x'}(h)\in u\tau N_A(Y')\).
Both assertions hold for all \(g,h\in G_K\).
Stability of \(N_A(Y')\) then implies that
\(B(d_x(g),g d'_{x'}(h))\) lies in the reduction of
\(\phi^{-1}(u)A(1)\). This is contained in \(S_n(1)\):
multiplication by the Teichm\"uller vector of any element of
positive valuation sends integral Witt vectors to zero in
the residue Witt vectors.

Use the path
\[
       (x,x')\ \longrightarrow\ (\phi x,x')
                    \ \longrightarrow\ (\phi x,\phi x').
\]
In the first move \(\phi x-x=y-x\in N_A(Y)\), since both
endpoints belong to \(N_A(X)\) over the same quotient.
Pairing this difference with \(d'_{x'}(g)\) thus gives an
element of the reduction of \(uA(1)\). In the second move
\(\phi x'-x'\in L_A(Y')\), since \(x'=y'\in L(X')\)
and \(x'\in N(X')\). Its \(g\)-translate also lies in
\(L_A(Y')\). The displacement \(d_{\phi x}(g)\) lies in
\(u\mu L_A(Y)\) as shown above. Their pairing is again in
the reduction of \(uA(1)\). These two subgroups lie in
\(S_n(1)\). The cochain argument of
Theorem~\ref{thm:cuppath757}, with \(H=G_K\), gives the
stated annihilator over the original field.
\end{proof}

\begin{corollary}[A local contribution to the original joint image]
\leavevmode\label{cor:sphericalcupbound762}
For the data \(D\) of Theorem~\ref{thm:projectedcups742},
choose any set \(\mathcal W\) of pairs consisting of a
spherical coefficient place \(w\mid p\) and a target \(T\).
Let \(r_{\rm rest}(D)\) be the normalized length on the
retained coordinates, with the original global source.
Then
\[
 r(D)\ \le\ r_{\rm rest}(D)+
      \sum_{(w,T)\in\mathcal W}
          [E_{x,w}:\QQ_p]\,
           \bigl(Q_K(X)+Q_K(X')+C_K(h_Z)\bigr)_{w,T},
       \qquad K=E_{x,w}.
\]
In each summand the exponent may be bounded using the
labelled weights of the local trace family and the dual
target. The same bound works at every trace precision.
\end{corollary}
\begin{proof}
At a fixed \(w,T\), use the graph of
Lemma~\ref{lem:graphrealization724} for each restricted trace
and the integral extension of its finite dual test.
The jumps of their extensions are the jumps of their
kernels together with zero. For a nonempty fixed trace
family the rational graph kernel is the same direct sum
for all the cocycles, and for a fixed target the rational
dual kernel is fixed as well. Thus the exponents just
proved work for every test and trace simultaneously. For
an empty kernel the cup vanishes. The conclusion applies
to every cyclic \(\ZZ_p\)-coordinate of the output with
its original modulus. Hence it applies to the full
\((\mathcal O/\pi^a)(1)\)-valued evaluation.

In a single free output coordinate the subgroup annihilated
by \(p^k\) has \(\mathcal O\)-length at most
\(k v_\pi(p)\). Each target uses at most \([K:\QQ_p]\)
such coordinates. As in Corollary~\ref{cor:pathperiod757},
the kernel of projection of the joint image injects into
the sum of the discarded coordinate images. Dividing their
lengths by \(v_\pi(p)\) proves the estimate. To sum it
over primes and nodes requires control of the sizes of
the exponents relative to \(\sqrt{b_x}\). The local
annihilator alone imposes no such restriction on the
largest Hodge jumps.
\end{proof}
\clearpage

\subsection{Integral lifts along an unramified Frobenius orbit}
\label{sec:labelledlifts764}
Suppose that \(K/\QQ_p\) is unramified of degree \(f\), and
that \(\mathcal O\) is the ring of integers in a coefficient
field containing all images of \(K\). Use the uniformizer \(p\)
of \(K\), so \(E(u)=u-p\). Put
\[
 R=\mathfrak S\otimes_{\ZZ_p}\mathcal O,\qquad
 R_0=W(k_K)\otimes_{\ZZ_p}\mathcal O.
\]
The factors of \(R_0=\prod_i\mathcal O\) are indexed by
the embeddings \(i:K\hookrightarrow\operatorname{Frac}(\mathcal O)\);
write \(\mathbf e_i\) for the idempotents. Thus
\(R_i=\mathbf e_i R=\mathcal O[[u]]\). Let \(F i\) be the index
of \(\phi(\mathbf e_i)\). This permutation is a single cycle
of length \(f\); on the power series coefficients in
\(\mathcal O\), Frobenius acts as the identity.

Let \(X\) be a crystalline \(\mathcal O\)-lattice with an
equivariant surjection \(\beta:X\to\mathcal O\) to the
trivial representation. Use the normalized module
of Lemma~\ref{lem:adjacentlattice762} with the action of
\(\mathcal O\) on it induced by its action on \(X\).
Put \(M_i=\mathbf e_i N(X)\) and
\(L_i=R_i\phi(M_{F^{-1}i})\). At \(i\), let \(h_i\) and
\(r_i\) be the least nonnegative integers such that all
labelled jumps belong to \([-r_i,h_i]\).

\begin{lemma}[The height at one embedding]
\leavevmode\label{lem:labelledheight764}
Both \(M_i,L_i\) are free over \(R_i\), and
\[
                 E^{h_i}M_i\ \subseteq\ L_i\
                        \subseteq\ E^{-r_i}M_i .
\]
\end{lemma}
\begin{proof}
Write \(\pi\) for a uniformizer of \(\mathcal O\), with
\(p\mathcal O=\pi^{e_{\mathcal O}}\mathcal O\).
The underlying \(N(X)\) is free over \(\mathfrak S\).
Multiplication by \(\pi\) on it is injective since a unit
times its \(e_{\mathcal O}\)-th power acts as \(p\). The
module \(M_i/\pi M_i\) embeds into \(N(X)/pN(X)\) by
multiplication by \(\pi^{e_{\mathcal O}-1}\). It has no
\(u\)-torsion and is finite over the discrete valuation
ring \((\mathcal O/\pi)[[u]]\). Lift a basis of this free
module to \(M_i\). Nakayama's lemma gives generators, and
any relation has all coefficients divisible by \(\pi\).
Division by \(\pi\) and iteration show that every
coefficient is zero. This proves freeness. Frobenius base
change of a basis gives a basis for \(L_i\), by linear
independence in the Witt comparison.

Choose a single twist making all jumps nonnegative. For
the resulting effective module, the filtration on the
linearized Frobenius image modulo \(E\), after inverting
\(p\), is the de Rham filtration. Its \(n\)-th piece is
the image of the intersection of that Frobenius image
with \(E^n\) times the original module. This is the
Hodge comparison for crystalline Breuil--Kisin modules
with coefficients, in the convention giving the cyclotomic
character jump \(-1\).
~\cite[Section~10.3, Theorem~10.5 and Remark~10.6]{III-II-bartletthodge764}
Over \(\operatorname{Frac}(\mathcal O)[[u-p]]\) the elementary
divisor \(E^j\) contributes a line in degree \(j\) to this
filtration. Thus on the \(i\)-component these exponents
are exactly its labelled jumps. Removing the common
twist subtracts its jump from every exponent.

The unlabelled height bounds in
Lemma~\ref{lem:adjacentlattice762} put both transition
matrices between the two bases in \(R_i[1/E]\). If
\(a\in R_i\) is divisible by \(E^m\) in
\(\operatorname{Frac}(\mathcal O)[[u-p]]\), it is already
divisible in \(R_i\). Indeed the substitution \(u=(u-p)+p\)
identifies \(R_i\) with \(\mathcal O[[u-p]]\), and the
first \(m\) coefficients vanish in \(\mathcal O\).
Apply this to the entries of the transition matrices
after clearing denominators. The two labelled inclusions
over the discrete valuation ring therefore hold over \(R_i\).
\end{proof}

\begin{theorem}[Frobenius-weighted upper jumps]
\leavevmode\label{thm:orbitlift764}
For \(p\ge3\) put
\[
       H_i=\sum_{j=1}^{\infty}\frac{h_{F^ji}}{p^j},\qquad
       c_i=
       \begin{cases}
       0,&H_i<1,\\[2pt]
       \left\lceil\dfrac{p(p-1)}{(p-1)^2-p}H_i\right\rceil,&H_i\ge1 .
       \end{cases}
\]
The image of \(M_i\) contains \(p^{c_i}\mathbf e_i\).
The image of \(M_i\cap L_i\) contains \(p^{q_i}\mathbf e_i\),
where one may take \(q_i=c_i+v_p(h_i!)\). If \(r_i=0\)
one may lower \(q_i\) to the minimum of this integer and
\(c_{F^{-1}i}\).
\end{theorem}
\begin{proof}
Omit the factor \(\mathbf e_i\) from the notation for scalars
in component \(i\). Write \(I_i=\beta(M_i)\subset R_i\).
Its image modulo \(\pi\) is nonzero. In fact \(\beta\) is
surjective after passing to the Witt comparison, and stays
so on each component after reduction modulo \(\pi\). The
corresponding component of the quotient there is nonzero
since the two copies of the residue field of \(K\) split
over the algebraically closed tilt. In the discrete valuation
ring \(R_i/\pi\), the image ideal therefore contains a
power of \(u\). Raising an expression \(u^a\in I_i+(\pi)\)
to the \(e_{\mathcal O}\nu\)-th power gives a pure power of
\(u\) in \(I_i+(p^\nu)\) for every \(\nu\ge1\).

Let \(a_i\) be the least exponent of such a power at
precision \(p^\nu\). The height inclusion on component
\(F i\), after projection, is
\[
          E^{h_{F i}}I_{F i}\subseteq R_{F i}\phi(I_i).
\]
Modulo \(p^\nu\), \(u^{h_{F i}+\nu-1}\) is divisible by
\(E^{h_{F i}}\), by expanding \(u=E+p\). Thus the ideal
on the right contains \(u^{a_{F i}+h_{F i}+\nu-1}\)
at this precision. Decompose its power series by their
exponents modulo \(p\), as in Lemma~\ref{lem:frobeniusideal759}.
The coefficient map between these two copies of \(\mathcal O\)
is the identity. It follows that
\[
                 p a_i\ \le\ a_{F i}+h_{F i}+\nu-1 .
\]
Iterating the inequality and using periodicity gives
\(a_i\le H_i+(\nu-1)/(p-1)\).

The operation of Lemma~\ref{lem:kummerlower760} works in
each \(I_i+(p^\nu)\). To verify this, \(G_K\) fixes the
idempotents and the coefficients \(\mathcal O\), so a lift
in \(M_i\) has displacement in \(u\tau A M_i\). For
the chosen \(g(u)=[\epsilon]u\) the divided displacement
of \(p^c u^m\) therefore gives the same scalar
\(([\epsilon]^m-1)/\tau\) as in that lemma. Its congruence
modulo \(u\) trades one \(u\) for \(p^{1+v_p(m)}\).
The extension \(\mathfrak S\to A\) stays faithfully flat
after tensoring with \(\mathcal O\) and on each idempotent,
so the ideal membership descends to \(R_i/(p^\nu)\).

If \(H_i<1\), take \(\nu=1\). Then \(a_i=0\) and
\(I_i=R_i\). Otherwise take \(\nu=c_i+1\).
Our choice of \(c_i\) ensures that
\[
      a_i+v_p(a_i!)\le\frac{p}{p-1}
             \left(H_i+\frac{c_i}{p-1}\right)\le c_i .
\]
Remove the \(a_i\) parameter factors. This gives
\(p^{c_i}\in I_i+(p^{c_i+1})\), hence \(p^{c_i}\in I_i\)
by division by a unit congruent to one modulo \(p\).

For the intersection, choose \(x_i\in M_i\) over
\(p^{c_i}\). Then \(E^{h_i}x_i\in L_i\). Apply the
divided-displacement argument of Proposition~\ref{prop:commonlift762}
to \(E^nM_i\cap L_i\), for \(n=h_i,\ldots,1\).
All its operators preserve the label, since the chosen
Galois element fixes \(K\). In the scalar congruence for
\(D_n\) one now has \(d_K=0\) because \(E'=1\). Thus
the successive costs are \(v_p(n)\). This proves the bound
for \(q_i\). When \(r_i=0\), \(L_i\subset M_i\), and
\(\phi(x_{F^{-1}i})\) is a common lift over
\(p^{c_{F^{-1}i}}\), proving the improvement.
\end{proof}

\begin{corollary}[Length of the missing constants]
\leavevmode\label{cor:orbitconstants764}
Let \(U=R_0\cap\beta(N(X)\cap L(X))\). Then
\[
        \frac{\operatorname{length}_{\mathcal O}(R_0/U)}
             {v_\pi(p)}
          \le \sum_i q_i
          \le \frac{8}{p-1}\sum_i h_i .
\]
For the first lattice alone, \(R_0\cap\beta(N(X))\),
one may replace 8 by 7. If the next \(m\) labels after
\(i\) have upper jumps at most \(B\), and every upper jump
is at most \(H\), the entire \(i\)-component of the
constant lattice lifts into \(M_i\) whenever
\[
               B(1-p^{-m})+H p^{-m}<p-1 .
\]
\end{corollary}
\begin{proof}
Each component of \(U\) contains \(p^{q_i}\mathcal O\);
the corresponding normalized length is at most \(q_i\).
For \(p\ge3\) one has
\(p(p-1)/((p-1)^2-p)\le6\). On \(H_i\ge1\) rounding up
thus gives \(c_i\le7H_i\); on the other components \(c_i=0\).
Since \(F\) permutes the indices,
\(\sum_i H_i=(p-1)^{-1}\sum_i h_i\). The estimate
\(v_p(h_i!)\le h_i/(p-1)\) proves the intersection bound.
Splitting the geometric series for \(H_i\) after \(m\)
terms gives the asserted sufficient condition for \(H_i<1\).
\end{proof}

\begin{corollary}[The labelled lifts of a period pairing]
\leavevmode\label{cor:periodlabels764}
At an unramified spherical completion \(K=E_{x,w}\) with
\(p\ge3\), use the original \(\mathcal O\)-linear graph map
\(t:Y\to T/\pi^aT\) and the test on \(Y'=T^\vee(1)\), with
the dual taken over \(\mathcal O\). After a coefficient
extension containing \(K\) if necessary, the preceding
theorem and corollary apply separately to the two extensions
of \(\mathcal O\) and to a third extension with kernel
\(Z=\operatorname{Hom}_{\mathcal O}(Y,T)\). At each label
\(i\), the jumps of this last kernel are \(j_T(i)-j\) for
the jumps \(j\) of \(Y\) at that same label.
\end{corollary}
\begin{proof}
Use the original graph cocycle before restricting scalars.
The integral dual-test cocycle likewise defines an extension
of \(\mathcal O\), crystalline by the finite condition.
Lift \(t\) to an \(\mathcal O\)-linear map \(\psi:Y\to T\),
which is possible since \(Y\) is free. The integral cocycle
\((g\psi-\psi)/\pi^a\) defines the required rationally
split extension of \(\mathcal O\) by \(Z\). The map from
this extension to \(Z\) sending \((z,b)\) to
\(\pi^a z+b\psi\) is equivariant. Evaluating it on
\(Y\) and then on \(Y'\) recovers the original pairing
times \(b\) modulo \(\pi^a\). The labelled jumps follow
from taking the dual and tensor product over \(\mathcal O\).
The same constructions commute with finite free extension
of the coefficient ring, using the extended ideal
\(\pi^a\mathcal O'\). The local field, cocycles and test
precision in the original ring have all been kept fixed.
The length in Corollary~\ref{cor:orbitconstants764} refers
to constants in a Kisin-module quotient. Passing from this
length to the joint image of the local cup tests requires
compatibility of the points on successive Frobenius components.
\end{proof}
\clearpage

\subsection{Comparing the cup cochains at successive embeddings}
\label{sec:orbitcups765}
Keep the field \(K\), the coefficient ring \(\mathcal O\), the
idempotents \(\mathbf e_i\) and the cyclic permutation \(F\)
of Theorem~\ref{thm:orbitlift764}. The exponents supplied at
individual embeddings may be different. To apply Frobenius
descent to one cup class at its original precision, we match
the exponents on each move of the path. We write \(i^-=F^{-1}i\)
and \(i^+=F i\).

\begin{theorem}[An annihilator using adjacent labelled exponents]
\leavevmode\label{thm:labelledcup765}
Let \(a,a'\) be \(\mathcal O\)-valued lattice cocycles with
crystalline extensions \(X,X'\) of \(\mathcal O\) by
\(Y,Y'\), respectively. Let \(P:Y\otimes_{\mathcal O}Y'
\to(\mathcal O/\pi^m)(1)\) be an equivariant pairing at
precision \(m\ge1\). Write \(X''\) for the third extension with
kernel \(\operatorname{Hom}_{\mathcal O}(Y\otimes_{\mathcal O}Y',
\mathcal O(1))\) constructed by lifting \(P\) and dividing
its displacement by \(\pi^m\). Let \(q_{X,i},q_{X',i}\)
be the common-lift exponents of Theorem~\ref{thm:orbitlift764}
for the first two extensions, and \(s_i=c_{X'',i}\) its
single-lattice exponent for the third. Put
\[
              A_i=q_{X,i^+},\qquad B_i=q_{X',i}.
\]
Then \(p^k\) annihilates the class of \(P(a\cup a')\), where
\[
 k=\max_i\max\left\{
     \max(A_i,A_{i^-})+B_i+s_i,\ \
     A_{i^-}+\max(B_i,B_{i^-})+s_{i^-}\right\}.
\]
\end{theorem}
\begin{proof}
Write \(D=(W(\mathbb C_p^\flat)\otimes_{\ZZ_p}\mathcal O)/
\pi^m\), and let \(D_{\rm sm}\) be the image of
\(W(\mathfrak m_{\mathbb C_p^\flat})\otimes_{\ZZ_p}\mathcal O\)
there. The idempotents act on the small subgroup because
they are integral. The assertions of
Lemma~\ref{lem:wittconstants757} extend to \(D\): the map
\(\phi-1\) is onto, and the fixed points in \(D/D_{\rm sm}\)
are precisely the embedded copy of \(\mathcal O/\pi^m\).
Indeed, as a \(\ZZ_p\)-module \(\mathcal O/\pi^m\) is a direct
sum of finite cyclic \(p\)-groups, and on each cyclic
coordinate the two groups in question are \(W_n(F_0)\)
and its quotient by \(W_n(\mathfrak m_{F_0})\), with
\(F_0=\mathbb C_p^\flat\). The assertion commutes with
twisting by one. Since \(G_K\) fixes \(W(k_K)\), it fixes
every \(\mathbf e_i\).

Choose \(y_i,y'_i\) in the common lattices for \(X,X'\)
over \(p^{q_{X,i}}\mathbf e_i,p^{q_{X',i}}\mathbf e_i\).
Take \(z_i\) in the first lattice for \(X''\) over
\(p^{s_i}\mathbf e_i\). Contraction of \(z_i\) with the
two first kernel lattices at \(i\) shows that \(p^{s_i}P\)
on their reductions takes values in the reduction of
\(\tau^{-1}\mathbf e_i(A\otimes\mathcal O)(1)\). Here and
below the unmarked tensor with \(A\) is over \(\ZZ_p\).
In fact the trilinear contraction is an integral
\(\ZZ_p\)-multilinear equivariant map. It respects the
three \(\mathcal O\)-actions by balancing. Hence the
tensor and Tate-line assertions of
Lemma~\ref{lem:adjacentlattice762} apply component by
component. At the second lattices at \(i\) use
\(\phi z_{i^-}\) instead. The bound there is
\(\mu^{-1}\mathbf e_i(A\otimes\mathcal O)(1)\) for
\(p^{s_{i^-}}P\).

At \(i\) start with the points \(x_i=\phi^{-1}y_{i^+}\)
and \(x'_i=y'_i\). Their constant coordinates are
\(p^{A_i}\mathbf e_i,p^{B_i}\mathbf e_i\). Fix
\(\mathcal O\)-linear splittings of the two extensions and
write \(u_i,v_i\) for the kernel coordinates. Put
\[
 \begin{aligned}
 b_i(g)&=p^{A_i}\mathbf e_i a(g),&
 b'_i(g)&=p^{B_i}\mathbf e_i a'(g),\\
 d_i(g)&=b_i(g)+gu_i-u_i,&
 d'_i(g)&=b'_i(g)+gv_i-v_i .
 \end{aligned}
\]
Use the continuous \(D(1)\)-valued cochain
\[
 f_i(g)=p^{k-A_i-B_i}
          \big(P(b_i(g),g v_i)-P(u_i,d'_i(g))\big).
\]
The two terms are bilinear in the first affine point and
the second affine point, including their constant
coordinates. Consequently we may first multiply either
point by an integral power of \(p\) and subtract the same
exponent from the prefactor whenever that prefactor
remains integral.

The cochain calculation of Theorem~\ref{thm:cuppath757}
gives, modulo \(D_{\rm sm}(1)\),
\[
 \partial f_i(g,h)\equiv
       p^k\mathbf e_i P(a(g),g a'(h)).
\]
Indeed \(d_i(g)\) lies in the first kernel lattice times
\(\phi^{-1}(u)\tau\), by applying \(\phi^{-1}\) to the
displacement of \(y_{i^+}\) in its second lattice. The
translate \(g d'_i(h)\) belongs to the first kernel
lattice by stability. Since \(k-A_i-B_i\ge s_i\),
their scaled pairing lies in the reduction of
\(\phi^{-1}(u)(A\otimes\mathcal O)(1)\), which is small.

We compare \(f_i\) with \(\phi f_{i^-}\) in two moves.
For the first move fix the second point \(x'_i\). Put
\(I=\max(A_i,A_{i^-})\). Move the first point from
\(p^{I-A_i}x_i\) to \(p^{I-A_{i^-}}\phi x_{i^-}\),
using the prefactor \(p^{k-I-B_i}\). The initial cochain
is \(f_i\) by bilinearity. Both first points are in the
first lattice at \(i\) over \(p^I\mathbf e_i\), since
\(\phi x_{i^-}=y_i\). Their difference is in the kernel.
The displacement of \(x'_i\) is divisible by \(u\tau\)
there. Our prefactor contains \(p^{s_i}\); the change in
cochain is therefore small, by the first-move identity
of Theorem~\ref{thm:cuppath757}.

For the second move use the point \(\phi x_{i^-}\) with
its original constant \(p^{A_{i^-}}\mathbf e_i\), and
put \(J=\max(B_i,B_{i^-})\). Move the second point from
\(p^{J-B_i}x'_i\) to \(p^{J-B_{i^-}}\phi x'_{i^-}\),
using \(p^{k-A_{i^-}-J}\). By bilinearity the initial
cochain here is exactly the terminal cochain of the
first move: both equal the cochain on
\((\phi x_{i^-},x'_i)\) with prefactor
\(p^{k-A_{i^-}-B_i}\). Now both second points are in the
second lattice at \(i\) with the same constant coordinate.
Their difference and its translates are in the second
kernel lattice. The displacement of \(\phi x_{i^-}=y_i\)
is divisible by \(u\mu\) there. The prefactor contains
\(p^{s_{i^-}}\). Thus the second-move identity says that
the change in cochain is a coboundary modulo the small
subgroup. The terminal cochain is \(\phi f_{i^-}\), since
\(\phi\) carries the idempotent at \(i^-\) to the one at
\(i\) and fixes the original coefficients and cocycles.

Sum over \(i\). For \(f_*=\sum_i f_i\), we have
\((\phi-1)f_*\equiv-\partial t\) modulo \(D_{\rm sm}(1)\)
for a zero-cochain \(t\in D(1)\), while
\(\partial f_*\equiv p^k P(a\cup a')\) since
\(\sum_i\mathbf e_i=1\). Choose \(z\) with
\((\phi-1)z=t\). The cochain \(f_*+\partial z\) now
reduces to a continuous cochain in the fixed subgroup
\((\mathcal O/\pi^m)(1)\); the quotient topology is
discrete by Lemma~\ref{lem:wittconstants757}. Injectivity
on these constants identifies its differential with
\(p^kP(a\cup a')\) in the original finite coefficients.
This proves the claim.
\end{proof}

\begin{corollary}[The labelled bound on a spherical period coordinate]
\leavevmode\label{cor:jointorbitcup765}
At an unramified spherical coefficient place \(w\mid p\),
\(p\ge3\), the contribution of a discarded target to the
drop in the normalized lengths of the joint images in
Corollary~\ref{cor:sphericalcupbound762} is at most
\([E_{x,w}:\QQ_p]\,k\), with \(k\) as above using the
three \(\mathcal O\)-linear lattices of
Corollary~\ref{cor:periodlabels764}. In particular
\[
        k\le \frac{8}{p-1}\sum_i
                     (h_{X,i}+h_{X',i}+h_{X'',i}).
\]
The maximum formula in the theorem retains the positions
of the large jumps on the orbit.
\end{corollary}
\begin{proof}
The graph and the dual test give the two cocycles at the
original precision. For the third extension evaluation of
\(\operatorname{Hom}_{\mathcal O}(Y,T)\) on \(Y\) followed
by the dual test is the required integral contraction.
The kernels have fixed rational jumps for the given trace
family and target. The resulting \(k\) works simultaneously
for the corresponding cocycles. If coefficients were
enlarged, \(1\) is primitive in \(\mathcal O'\) over
\(\mathcal O\), so an \(\mathcal O\)-linear retraction
\(\mathcal O'\to\mathcal O\) sending \(1\) to \(1\)
descends a coboundary at the extended modulus
\(\pi^m\mathcal O'\). Counting the free coordinates and
intersecting with the joint global image now gives the
bound just as in Corollary~\ref{cor:sphericalcupbound762}.
Finally each expression in the maximum uses at most
\(\sum_i q_{X,i}+\sum_i q_{X',i}+\sum_i s_i\).
Apply Corollary~\ref{cor:orbitconstants764} to the three
sums. A prime sum for arbitrary large jumps would require
additional savings beyond this last inequality.
\end{proof}
\clearpage

\subsection{Removing nonconstant terms at the Eisenstein prime}
\label{sec:taylordepth767}
Keep the unramified field \(K\), coefficient ring \(\mathcal O\),
and component \(i\) of Section~\ref{sec:labelledlifts764}.
Write \(C=\operatorname{Frac}(\mathcal O)\), \(E=u-p\),
\(\delta=u\tau\), \(R_i=\mathcal O[[u]]\), and
\(A_i=\mathbf e_i(A\otimes_{\ZZ_p}\mathcal O)\), with
\(A=W(\mathcal O_{\mathbb C_p^\flat})\). We identify the
unit of \(R_i\) with \(\mathbf e_i\). For a crystalline
lattice \(X\) write \(M=M_i(X)\) and \(L=L_i(X)\) for the
two lattices in Lemma~\ref{lem:labelledheight764}. A subscript
\(A\) denotes extension of scalars to \(A_i\). The construction
below raises the order of \(gx-x\) for all \(g\in G_K\)
at once. It uses Taylor coefficients at \(u=p\).

\begin{lemma}[Horizontal coordinates at the uniformizer]
\leavevmode\label{lem:horizontalcoordinates767}
In the \(E\)-adic comparison over \(C((E))\) there is a
basis \(d_1,\ldots,d_m\) fixed by \(G_K\), with labelled
jumps \(j_1,\ldots,j_m\), such that
\[
    L\otimes_{R_i}C[[E]]=\bigoplus_k C[[E]]d_k,\qquad
    M\otimes_{R_i}C[[E]]=\bigoplus_k E^{-j_k}C[[E]]d_k.
\]
Here \(R_i\to C[[E]]\) is expansion at \(u=p\).
In the de Rham period ring, \(g\) acts on a coordinate
\(f(E)d_k\) by \(f(gE)d_k\), with the coefficient field
fixed. This assertion applies to a finite Laurent principal
part and to every finite Taylor expansion.
\end{lemma}
\begin{proof}
First twist into effective height. In Kisin's rational
comparison the module \(D_{\rm cris}(X[1/p])\), tensored
with the ring of analytic functions on the open unit disc,
maps to the analytic span of the linearized Frobenius image.
This becomes an isomorphism upon inverting \(\phi(\lambda)\),
where \(\lambda=\prod_{m\ge0}\phi^m(E/E(0))\). In the period
comparison the constant vectors in this description are
the \(G_K\)-invariant crystalline periods. At \(u=p\) the
factors of \(\phi(\lambda)\) have values
\(1-p^{p^m-1}\), \(m\ge1\). Their product converges to a
nonzero element of \(K\). Taylor expansion of the analytic
comparison at this point therefore gives the first lattice
in any basis of \(D_{\rm cris}\). Because \(K\) is unramified,
its scalar field is already \(K\) and
\(D_{\rm cris}=D_{\rm dR}\) under comparison. Taking the
component selects a \(C\)-basis of the corresponding
filtered space. Untwisting is compatible with all these
comparisons by the tensor construction.
~\cite[Lemma~2.1.8, (2.1.9), and Proposition~2.1.10]{III-II-bartlettcrys767}

To identify the first Frobenius source lattice \(M\), let
\(\mathcal B_i^+=\mathbf e_i(B_{\rm dR}^+\otimes_{\QQ_p}C)\).
The integral comparison for the Frobenius image identifies
its \(A\)-span after inverting \(\mu\) with the representation
lattice over the same ring. Applying \(\phi^{-1}\) gives
the identification for \(M\) after inverting \(\tau\).
Since \(\theta(\tau)=\zeta_p-1\ne0\), \(\tau\) is a unit
in \(B_{\rm dR}^+\). Thus \(M\) extends over \(\mathcal B_i^+\)
to the representation tensored with \(B_{\rm dR}^+\), on
this component. Choose a basis \(d_k\) adapted to the
de Rham filtration, with \(d_k\) of degree \(j_k\).
The filtered de Rham comparison makes this last lattice
the degree-zero part of \(B_{\rm dR}\otimes_K D_{\rm dR}\).
In the adapted basis it is
\(\bigoplus_k E^{-j_k}\mathcal B_i^+ d_k\), since \(E\)
and Fontaine's \(t\) generate the same maximal ideal of
\(B_{\rm dR}^+\). The two rational lattices over \(C[[E]]\)
can be compared there by the height inclusions of
Lemma~\ref{lem:labelledheight764}. The extension to
\(\mathcal B_i^+\) is faithfully flat: it is a nonempty
product of extensions of discrete valuation rings sending
the parameter to a uniformizer. Hence the equality descends.
The basis vectors are horizontal, and \(G_K\) fixes the
scalar coefficient factor. Substitution in finite Taylor
expansions, using \(g(E)B_{\rm dR}^+=E B_{\rm dR}^+\),
gives the asserted action.
\end{proof}

\begin{lemma}[Retracting finite Taylor conditions]
\leavevmode\label{lem:integraljet767}
There is an \(R_i\)-linear retraction \(\rho:A_i\to R_i\).
Let \(\mathcal M\) denote \(M\) or \(L\), and apply \(\rho\)
coefficientwise to \(E^{-a}\mathcal M_A\), for \(a\ge0\).
This map fixes \(E^{-a}\mathcal M\) and commutes with
integral maps of these lattices. If a horizontal coordinate
of \(v\in E^{-a}\mathcal M_A\) agrees with a prescribed
Laurent polynomial in \(C[E,E^{-1}]\) modulo \(E^n\mathcal B_i^+\),
the same congruence holds for that coordinate of \(\rho(v)\).
It holds simultaneously for any finite collection of these
conditions, with arbitrary integer orders \(n\).
\end{lemma}
\begin{proof}
The ring \(A_i\) is faithfully flat over \(R_i\) by base
change from \(\mathfrak S\to A\) and passage to the
idempotent. Prescribe unknowns in \(R_i\) for the images of
all elements of \(A_i\), with the equations of linearity
and the requirement that the images of elements of \(R_i\)
equal themselves. Each finite collection is a system of
linear equations over \(R_i\) with a solution in \(A_i\),
given by the elements themselves. Its right-hand side is
therefore zero in the cokernel of its coefficient matrix
after faithfully flat extension, so is already zero in
the original cokernel. Each finite collection is solvable
in \(R_i\). This ring is compact in the \((\pi,u)\)-adic
topology, with finite residue field. Compactness of the
product of all the copies of \(R_i\) gives a solution to
every equation, since each equation is closed. This gives
\(\rho\). Tensoring with the indicated free lattice gives
the coefficientwise retraction independently of its basis.

In a basis of \(E^{-a}\mathcal M\) each horizontal coordinate
is a linear form with coefficients in \(C((E))\) by the
previous lemma. Multiply the required congruence by a power
of \(E\) so all the coefficients and its prescribed value
are power series and the new precision \(N\) is positive.
Truncate those series below exponent \(N\) and clear their
finitely many coefficient denominators by \(p^b\). The
result is a linear congruence over \(R_i\) in the entries
of \(v\), with known right-hand side. It holds modulo
\(E^N A_i\): indeed \(A/E^N\) is \(p\)-torsion-free, by
the successive quotients isomorphic to \(A/E=\mathcal O_{\mathbb C_p}\).
After inverting \(p\), its injection into \(B_{\rm dR}^+/E^N\)
follows from the definition of \(B_{\rm dR}^+\) as the
\(E\)-adic completion of \(A[1/p]\). Tensoring with the
finite free coefficients and taking the idempotent preserves
these injections. Apply \(\rho\) to the integral congruence.
In the period ring divide back by \(p^b\) and the initial
power of \(E\). This proves the assertion. The same
injection shows, entrywise, that a vector already in
\(E^{-a}\mathcal M_A\) and lying in \(E^D\mathcal M\)
after extension to \(\mathcal B_i^+\) lies in \(E^D\mathcal M_A\)
for \(D\ge-a\).
\end{proof}

\begin{theorem}[An integral lift with deep Galois displacements]
\leavevmode\label{thm:deepintegrallift767}
Let \(\beta:X\to\mathcal O\) be a crystalline lattice
surjection with trivial quotient. Suppose
\(x_0\in E^{-a}M\) has constant coordinate
\(\beta(x_0)=\mathbf e_i\), where \(a\ge0\). For any
integer \(D\ge0\), form the finite set
\[
       \mathcal I=\bigcup_k
         \{n\in\ZZ\setminus\{0\}:-a-j_k\le n<D-j_k\},
       \qquad s=\sum_{n\in\mathcal I}v_p(n).
\]
Here \(v_p(n)=v_p(|n|)\). There exists \(x\in E^{-a}M\)
such that
\[
       \beta(x)=p^s\mathbf e_i,\qquad
       (g-1)x\in\delta E^D M_A\quad(g\in G_K).
\]
If \(x_0\in M\), one has \(x\in M\). The same conclusion
holds for an initial point in \(E^{-a}M\) if \(X\) has
no labelled jump in \([-a,0)\). Each integer in the union
is counted just once, irrespective of the multiplicities.
\end{theorem}
\begin{proof}
Every \(g\in G_K\) sends \(E\) to \(\omega_g E\) for a
unit \(\omega_g\in A\) with \(\omega_g-1\in\delta A\).
In fact \(g(u)/u=[\epsilon^c]\) for some \(c\in\ZZ_p\),
\([\epsilon^c]-1\in\mu A\), and \(\mu/(E\tau)\) is a
unit. Apply this identity and
Lemma~\ref{lem:adjacentlattice762} to an element of
\(E^{-a}M\). Its divided displacement by \(\delta\)
belongs to \(E^{-a}M_A\), including for \(a=0\).

Choose \(g_0\) with \(g_0(u)=[\epsilon]u\) and put
\(\omega=\omega_{g_0}\). For \(n\ne0\) write
\(D_n=(\omega^n-1)/\delta\). The calculation in
Lemma~\ref{lem:transportpower761} gives
\(v_p(\theta(D_n))=v_p(n)\) since \(E'=1\).
For negative \(n\) use \(D_{-n}=-\omega^{-n}D_n\).
Surjectivity of \(\theta\) supplies \(\alpha_n\in A\)
with \(\alpha_nD_n\equiv p^{v_p(n)}\pmod E\).
On \(E^{-a}M\) use the additive operator
\[
     T_n v=p^{v_p(n)}v-
             \rho\left(\alpha_n\frac{g_0v-v}{\delta}\right).
\]
If \(\beta(v)\) is constant, the correction has zero
projection, by the preceding lemma applied to the
projection map. Thus \(T_n\) multiplies this coordinate
by \(p^{v_p(n)}\). It preserves \(E^{-a}M\); on inputs
in \(M\) it also preserves \(M\).

Consider a horizontal coordinate \(f(E)\in C((E))\) of
\(v\). Suppose \(f-c\in E^b C[[E]]\) for some \(c\in C\).
For every \(g\), \(f(gE)-f(E)\in E^b\mathcal B_i^+\).
Division by \(\delta\) preserves this order in
\(\mathcal B_i^+\), since \(\theta(\delta)=p(\zeta_p-1)\ne0\).
The correction after applying \(\rho\) still has order
at least \(b\) by Lemma~\ref{lem:integraljet767}. Hence
the new coordinate equals \(p^{v_p(n)}c\) modulo \(E^b\).
If \(b=n\ne0\), write \(f-c=a_n E^n+O(E^{n+1})\),
where \(a_n\in C\) is fixed, possibly zero. Then
\[
    \alpha_n\frac{f(g_0E)-f(E)}{\delta}
          \equiv p^{v_p(n)}a_n E^n\pmod {E^{n+1}\mathcal B_i^+}.
\]
The prescribed term on the right is rational in these
coordinates, so the same congruence holds after the
retraction. Subtraction raises the order of the difference
from \(p^{v_p(n)}c\) to at least \(n+1\). At \(n=0\) use
the identity operator: a power series of order at least
zero is its constant coefficient modulo \(E\).

Initially each coordinate has order at least \(-a-j_k\).
Process the union of all integer intervals
\([-a-j_k,D-j_k)\) in increasing order, inserting the
identity at zero. Whenever a coordinate reaches an integer
in its interval, its agreement with a constant gains one
order by the preceding calculation. All other coordinates
retain their already attained orders of agreement.
In the resulting point each coordinate agrees with a
constant modulo \(E^{D-j_k}\). Thus for all \(g\) the
horizontal coordinates of \((gx-x)/\delta\) have orders
at least \(D-j_k\). The description of \(M\) in
Lemma~\ref{lem:horizontalcoordinates767} puts this vector
in \(E^D M\) over \(\mathcal B_i^+\). It already belongs
to \(E^{-a}M_A\), so Lemma~\ref{lem:integraljet767} gives
membership in \(E^D M_A\).

It remains to verify the assertion about a pole when no
jump belongs to \([-a,0)\). For \(j_k<-a\), the initial
order is strictly positive. Its constant stays zero
throughout: before or after any raising step a correction
respecting a positive order also has zero constant term.
The final order on this coordinate is at least
\(D-j_k\ge-j_k\). For \(j_k\ge0\), agreement with a
constant modulo \(E^{D-j_k}\) gives order at least
\(\min(0,D-j_k)\ge-j_k\). These are all the coordinates
under the gap hypothesis, so \(x\) belongs to \(M\)
over \(C[[E]]\). Its known membership in \(E^{-a}M\)
then gives \(x\in M\) over \(R_i\), by the divisibility
test in Lemma~\ref{lem:labelledheight764}.
\end{proof}

\begin{corollary}[Deepening at one prime while retaining the other]
\leavevmode\label{cor:secondlatticedepth767}
Put \(P=\phi(E)=u^p-p\). Let \(y\in L\) have a constant
quotient coordinate and suppose
\((g-1)y\in\delta P^a L_A\) for all \(g\in G_K\).
For \(B\ge0\) let \(s_0=v_p(B!)\). There exists \(q\in L\)
with
\[
   q-p^{s_0}y\in P^a L\cap\ker\beta,\qquad
   (g-1)q\in u\mu\, E^B P^a L_A\quad(g\in G_K).
\]
\end{corollary}
\begin{proof}
Use the same operators on \(L\), whose horizontal coordinates
are power series with no shifts by Lemma~\ref{lem:horizontalcoordinates767}.
The initial Taylor step at zero is free. Process degrees
\(1,\ldots,B\). At every step the correction lies in
\(P^a L\), because \(\rho\) is \(R_i\)-linear and the
divided displacement belongs to \(P^a L_A\); the quotient
of the correction is zero. This last divisibility persists
after the correction. In fact \((g-1)L\subseteq u\mu L_A
\subseteq\delta L_A\), and
\(g(P)/P=\phi(\omega_g)\equiv1\pmod{\phi(\delta)}\),
where \(\phi(\delta)=u^p\mu\in\delta A\). The product
rule therefore gives \((g-1)(P^a L)\subseteq\delta P^a L_A\).
The final divided displacement is also in \(E^{B+1}L_A\),
by the Taylor calculation and integral divisibility.
These two divisibilities give their product: modulo \(E\),
\(P\) has nonzero scalar value \(p^p-p\), a nonzero
divisor on \(A_i/EA_i\). Successively cancelling \(E\)
shows \(P^aL_A\cap E^{B+1}L_A=P^aE^{B+1}L_A\).
Finally \(\delta E\) is a unit times \(u\mu\).
This gives both asserted conclusions.
\end{proof}
\clearpage

\subsection{The local cup product near the five weight intervals}
\label{sec:localweightmultiples768}
We use the two lattices of Section~\ref{sec:taylordepth767}
to compare the cup cochains for the original graph and
test. The order in which we change their two arguments
depends on the signs of the labelled jumps at the embedding
where the comparison takes place.

\begin{theorem}[A cup bound from two adjacent lattices]
\leavevmode\label{thm:twolatticecup768}
Let \(K/\QQ_p\) be unramified, \(p\ge5\), and take coefficients
containing \(K\). Use the crystalline extensions \(X,X'\) of
\(\mathcal O\) by \(Y,Y'\), their cocycles \(c,c'\), and an
equivariant pairing
\[
   Q:Y\otimes_{\mathcal O}Y'\longrightarrow(\mathcal O/\pi^m)(1).
\]
Assume a third crystalline extension \(X''\) of \(\mathcal O\)
admits an integral equivariant contraction with \(Y,Y'\)
to \(\mathcal O(1)\) which reduces to its quotient coordinate
times \(Q\). Let \(a,B\ge0\) be integers with \(B<p\), put
\(\eta=B+1\) and \(D=a+\eta\). Assume at every label \(i\)
that \(E^a\mathbf e_i\) lifts in each of the three lattices
\(M_i\). For each of the first two extensions at each label,
assume that either \(\mathbf e_i\) itself lifts in \(M_i\)
or no jump of that extension belongs to \([-a,0)\).
At each label assume at least one of the following conditions.
\begin{enumerate}
\item All jumps of \(X\) are nonnegative and all jumps
of \(X'\) are nonpositive.
\item All jumps of \(X\) are at most \(B\) and all jumps
of \(X'\) are at least \(-\eta\).
\end{enumerate}
For \(U=X,X'\), let \(s_U\) be the maximum over \(i\) of
\[
 \sum_{\substack{n\ne0:\ \exists\, j\ {\rm a\ jump\ of}\ U\
                      {\rm at}\ i\\-a-j\le n<D-j}}v_p(|n|).
\]
Then \(p^{s_X+s_{X'}}[Q(c\cup c')]=0\) in
\(H^2(K,(\mathcal O/\pi^m)(1))\).
\end{theorem}
\begin{proof}
Set \(P=\phi(E)\), \(\delta=u\tau\) as before. At each
label divide a lift of \(E^a\mathbf e_i\) by \(E^a\).
When a unit lift is available, start with that lift instead.
Apply Theorem~\ref{thm:deepintegrallift767} to the first
two extensions and multiply by the additional power of \(p\)
needed to attain the respective maximum. This gives
\(x_i\in M_i(X)\), \(x'_i\in M_i(X')\) over
\(p^{s_X}\mathbf e_i,p^{s_{X'}}\mathbf e_i\).
Write \(d_v(g)=gv-v\). Their displacements satisfy
\[
 d_{x_i}(g)\in\delta E^D M_{i,A}(Y),\qquad
 d_{x'_i}(g)\in\delta E^D M_{i,A}(Y').
\]
Indeed the quotient coordinate vanishes, and the factors
\(\delta,E\) are nonzero divisors on the quotient ring,
also after the finite free coefficient extension.
Use the kernel assertion of Lemma~\ref{lem:adjacentlattice762}.
For the third extension choose \(z_i\in E^{-a}M_i(X'')\)
over \(\mathbf e_i\). The contraction shows that pairing
the two first kernel lattices over \(A_i\) by \(Q\) takes
values in the image of \(E^{-a}\tau^{-1}A_i(1)\); for
the two second lattices it takes values in the image of
\(P^{-a}\mu^{-1}A_i(1)\), using \(\phi z_{i^-}\).
These are statements in the finite Witt comparison of
Theorem~\ref{thm:labelledcup765}. For example evaluate
the integral contraction on \(E^a z_i\) before reducing,
then divide by \(E^a\) there. The elements \(E,P,\tau,\mu\)
are units in \(W(\mathbb C_p^\flat)\), since their
reductions are nonzero in \(\mathbb C_p^\flat\).
Thus both pairing bounds concern the given modulus.

We verify the cochain comparisons of
Theorem~\ref{thm:cuppath757}. Throughout, any value in
the reduction of \(u A_i(1)\) is small in the sense of
Lemma~\ref{lem:wittconstants757}. At the initial pair
\((x_i,x'_i)\) the value \(Q(d_{x_i}(g),g d_{x'_i}(h))\)
is small. The first input has factor \(\delta E^D\);
the second belongs to the stable first kernel lattice.
The first pairing bound leaves the factor \(u E^{D-a}\).
The quotient coordinates match those of
\(\phi x_{i^-},\phi x'_{i^-}\), respectively.
The latter points belong to the second lattices, with
displacements in \(\phi(\delta)P^D\) times their second
kernel lattices; here \(\phi(\delta)=u^p\mu\).

Under condition (1), move the first argument from \(x_i\)
to \(\phi x_{i^-}\) and then move the second from \(x'_i\)
to \(\phi x'_{i^-}\). The first difference lies in the
first kernel lattice of \(Y\), since \(L_i(X)\subseteq
M_i(X)\) by Lemma~\ref{lem:labelledheight764}.
Pair it with \(d_{x'_i}(g)\) and use \(\delta E^D\) and
the first pairing bound. The second difference and all
its translates lie in the second kernel lattice of \(Y'\),
since \(M_i(X')\subseteq L_i(X')\). Pair a translate
with \(d_{\phi x_{i^-}}(g)\) and use \(u^p\mu P^D\)
and the second pairing bound. Both results are small.

Under condition (2), put \(y=\phi x'_{i^-}\). Its
displacement belongs to \(\delta P^a L_{i,A}(Y')\)
because \(\phi(\delta)/\delta\in A\) and \(D\ge a\).
Apply Corollary~\ref{cor:secondlatticedepth767} to
\(y\), with \(v_p(B!)=0\). We obtain \(q'_i\in L_i(X')\)
over the same constant as \(x'_i\) and with displacement
in \(u\mu E^B P^a L_{i,A}(Y')\). Take the path
\[
 \begin{aligned}
 (x_i,x'_i)&\longrightarrow(x_i,q'_i)
            \longrightarrow(\phi x_{i^-},q'_i)\\
            &\longrightarrow(\phi x_{i^-},\phi x'_{i^-}).
 \end{aligned}
\]
For the first move, the difference of the second points and
its translates lie in \(E^{-\eta}M_{i,A}(Y')\) by the
lower jump bound. Every translate of \(E\) is \(E\)
times a unit. The fixed first displacement supplies
\(\delta E^{a+\eta}\), so the first pairing bound leaves
\(u A_i(1)\). For the middle move, the difference of the
first points lies in \(E^{-B}L_{i,A}(Y)\) by the upper
jump bound. The fixed second displacement supplies
\(u\mu E^B P^a\), so the second pairing bound again
leaves \(u A_i(1)\). For the last move, both second
points are in \(L_i(X')\); pair a translate of the
kernel difference with the displacement of the fixed
first point, using \(u^p\mu P^D\). This proves smallness
for every move also under condition (2).

For completeness, the cochain of
Theorem~\ref{thm:cuppath757} at \((x_i,x'_i)\) consequently
has differential \(p^{s_X+s_{X'}}\mathbf e_i Q(c\cup c')\)
modulo the small subgroup. A first-argument move changes
it by a small cochain; a second-argument move changes it
by a coboundary modulo that subgroup. The terminal
cochain is the Frobenius transform of the initial cochain
at \(i^-\). Summing over the cycle, \(\sum_i\mathbf e_i=1\),
and the Frobenius difference of the sum is a coboundary
modulo the small subgroup. Surjectivity of \(\phi-1\)
on the finite Witt comparison corrects this difference
by a coboundary, exactly as in
Theorem~\ref{thm:labelledcup765}. The corrected cochain
takes values modulo the small subgroup in its Frobenius
constants \((\mathcal O/\pi^m)(1)\), which inject into
that quotient. Its differential is the required multiple
of the original cup.
\end{proof}

\begin{corollary}[The five intervals for the marked torus data]
\leavevmode\label{cor:actualperiodmultiples768}
For the data of Theorem~\ref{thm:projectedcups742} at a
sufficiently large node \(x\), write \(b=b_x\) and
\(b_2=a_x/x_{1,x}=xb/x_{1,x}\). There is a constant
\(C_0\ge1\), depending only on \(q_0,d\), with the following
property. For any rational prime \(p\ge x_{1,x}\) put
\[
 \begin{aligned}
 \mathcal N_{x,p}&=\{n\in\ZZ\setminus\{0\}:
       \min_{c\in\{0,b,-b,b_2,-b_2\}}|n-c|
                            \le C_0(x^{42}+b/p)\},\\
 S_{x,p}&=\sum_{n\in\mathcal N_{x,p}}v_p(|n|).
 \end{aligned}
\]
At every spherical place \(w\mid p\), each original local
cup between a trace and a free dual test is annihilated by
\(p^{2S_{x,p}}\). The exponent is simultaneous in the trace
ranks, targets, powers, precision and coefficient embeddings.
If the five intervals contain no nonzero multiple of \(p\),
each of these cups vanishes at its original precision.
\end{corollary}
\begin{proof}
The bound \(x_{1,x}\ge x^{10^6}\) puts \(p\) above the
ramified primes of \(E_x\) for large \(x\), by
Lemma~\ref{lem:sumcupheight763}. Work over \(K=E_{x,w}\).
Enlarge coefficients to contain \(K\) as in
Corollary~\ref{cor:periodlabels764}. The graph, the dual
test and its integral contraction there provide the three
extensions of the theorem. All their upper jumps are at
most \(12b\), as calculated in
Lemma~\ref{lem:nearweightprimes766}. Put
\[
       a=\left\lceil \frac{12b}{p-3}\right\rceil,\qquad
       B=t'-1 .
\]
Theorem~\ref{thm:uniformlift760} supplies the lift of
\(E^a\mathbf e_i\) on all three extensions. To see this
with the coefficient action, pull the quotient back along
\(\ZZ_p\to\mathcal O\), \(1\mapsto1\). The rational
pullback is a subrepresentation of the sum of the
crystalline extension and a trivial line, and has the
same upper bound. Apply the theorem to its quotient
\(\ZZ_p\), map the resulting lift to the original module,
and multiply by \(\mathbf e_i\).

At label \(i\), the graph jumps belong to
\(\{0,\ldots,B,J_i\}\), with \(J_i=z-1+\epsilon_i k_i\).
The dual-test extension adds zero to the jump \(-1-j_T\),
where \(j_T=J_i\) for the main target and \(0\le j_T\le B\)
for a scalar target. The weight formulas in
Theorem~\ref{thm:negativeinterval740} and the shifts of
Lemma~\ref{lem:nearweightprimes766} put \(J_i\) within
\(O_{q_0,d}(x^{41})\) of one of \(\pm b,\pm b_2\).
Indeed \(r_x\ll x^5,\ m_x\le2\sqrt{x}\), the shift from
a base angular magnitude is \((2j+1)r_x^8-h_i\) with
\(0\le h_i\le4r_x\), and \(|z-1|\le t_0\ll x^5\).
Here \(0<b_2\le b\), while
\(b_2=xb/x_{1,x}\ge x b/p\) and
\(b_2\ge b^{1-10^{-6}}\) by the scale conditions.
Consequently \(|J_i|>a+B+2\) for large \(x\), simultaneously
for all \(p\ge x_{1,x}\). Both the graph and the main
dual therefore have the required gap \([-a,0)\). For a
scalar dual every jump is nonpositive at every embedding,
so a unit already lifts by Theorem~\ref{thm:integrallifts759},
using the same pullback for coefficients.

If \(J_i>0\), the graph has only nonnegative jumps and the
dual test only nonpositive jumps. If \(J_i<0\), the graph
has upper bound \(B\), and the dual has lower bound
\(-B-1\) for either target. These give conditions (1) and
(2), respectively, with \(B<p\). Every nonzero integer in
either sum of the theorem lies within \(a+B+2\) of the
negative of one of the graph or dual jumps. Using
\(a\ll 1+b/p\) and the displayed bounds on the jumps,
all such integers belong to \(\mathcal N_{x,p}\) for one
constant \(C_0\). Thus each exponent is at most \(S_{x,p}\).
The theorem applies over \(K\). An \(\mathcal O\)-linear
retraction of the extended coefficients fixing \(1\),
applied to the cobounding cochain, descends the assertion
to the original coefficient ring and modulus.
\end{proof}
\clearpage

\subsection{Scalar tests when the graph lift has a pole}
\label{sec:scalarcupall770}
At primes below \(x_{1,x}\), a negative main jump on the
graph can lie inside the interval \([-a,0)\) of
Theorem~\ref{thm:deepintegrallift767}. The graph point
then supplied can have an \(E\)-pole. A scalar dual has
bounded jumps at both ends and admits an integral point
at each label. We use this point to absorb both the pole
of the graph and the pole at \(\phi(E)\) of its translate.

\begin{theorem}[A cup annihilator for a bounded nonpositive test]
\leavevmode\label{thm:scalarmeromorphic770}
Keep the local field, coefficients, contraction and cocycles
of Theorem~\ref{thm:twolatticecup768}, with \(p\ge5\).
Let \(a,B\ge0\), \(\eta=B+1\), and \(d=2a+2\eta+1\).
Assume \(E^a\mathbf e_i\) lifts in the first lattices for
\(X,X''\) and \(\mathbf e_i\) lifts in the first lattice
for \(X'\), at every label. All jumps of \(X'\) are to
belong to \([-\eta,0]\). At each label the jumps of \(X\)
are either all nonnegative or all at most \(B\).
For \(U=X,X'\), put \(a_X=a,\ a_{X'}=0\), and set
\[
 s_U=\max_i\sum_{\substack{n\ne0:\ \exists\,j\ {\rm a\ jump\
              of}\ U\ {\rm at}\ i\\-a_U-j\le n<d-j}}v_p(|n|).
\]
Then \(p^{s_X+s_{X'}+v_p(d!)}\) annihilates
\([Q(c\cup c')]\) at its original modulus.
\end{theorem}
\begin{proof}
Write \(P=\phi(E)\), \(e_0=v_p(d!)\). Take the graph
point over \(1\) on component \(i\) by dividing its
given lift by \(E^a\); take the integral unit lift for
the test. Theorem~\ref{thm:deepintegrallift767} gives
\[
     x_i\in E^{-a}M_i(X),\qquad x'_i\in M_i(X')
\]
over \(p^{s_X}\mathbf e_i,p^{s_{X'}}\mathbf e_i\), after
rescaling to the maximum for each extension. For \(x_i\)
and \(x'_i\) their respective displacements lie in
\(\delta E^d M_{i,A}(Y)\) and \(\delta E^d M_{i,A}(Y')\).
In particular the former displacement has no pole. Write
\(v_i=p^{e_0}x'_i\). The first and second lattice pairing
bounds of Theorem~\ref{thm:twolatticecup768} hold with
factors \(E^{-a}\tau^{-1}\) and \(P^{-a}\mu^{-1}\),
using the given lift in \(X''\). We use those bounds also
on scalar multiples of their inputs inside the Witt
comparison, where \(E\) and \(P\) are invertible.

The point \(y_i=\phi x'_{i^-}\) belongs to \(L_i(X')\),
and \(d_{y_i}(g)\in u^p\mu P^d L_{i,A}(Y')\).
Apply Corollary~\ref{cor:secondlatticedepth767} with both
of its exponents equal to \(d\): indeed
\(u^p\mu/\delta\in A\). We obtain \(q_i\in L_i(X')\)
with the same constant as \(v_i\) and with
\[
       d_{q_i}(g)\in u\mu E^d P^d L_{i,A}(Y')
                      \quad(g\in G_K).
\]
The power \(p^{e_0}\) used here is the same at all labels.
The graph translate \(\phi x_{i^-}\) is in
\(P^{-a}L_i(X)\), with displacement in
\(u^p\mu P^d L_{i,A}(Y)\). Thus at every label the two
pairs \((x_i,v_i)\) and \((\phi x_{i^-},\phi v_{i^-})\)
have the same two constants.

At the first pair the two displacements give a small cup:
the graph supplies \(\delta E^d\), and the translate of
the other displacement belongs to the stable first test
lattice. Move first from \(v_i\) to \(q_i\) with the
graph fixed. The difference and its translates belong to
\(E^{-\eta}M_{i,A}(Y')\) by the lower jump bound on \(X'\).
The first pairing bound and the graph displacement leave
\(uE^{d-a-\eta}A_i(1)\), which is small. Next move from
\(x_i\) to \(\phi x_{i^-}\) with \(q_i\) fixed. We check
both possible bounds for the graph at this label.

If its jumps are nonnegative, \(L_i(X)\subseteq M_i(X)\).
The difference of the two graph points belongs to
\(E^{-a}P^{-a}M_{i,A}(Y)\). Indeed after multiplication
by \(E^aP^a\) both endpoints are in the first lattice
and have equal quotient. The displacement of \(q_i\)
belongs to \(u\mu E^{d-\eta}P^d M_{i,A}(Y')\).
Under the first pairing bound the product is in
\[
        u(\mu/\tau) E^{d-\eta-2a}P^{d-a} A_i(1).
\]
Both exponents are nonnegative and \(\mu/\tau\in A\),
so this is small.

If the graph jumps are at most \(B\), use
\(M_i(X)\subseteq E^{-B} L_i(X)\).
The difference of the graph points now belongs to
\(E^{-a-B}P^{-a} L_{i,A}(Y)\) by the same kernel
argument. Use the displacement of \(q_i\) in its second
lattice and the second pairing bound. The product is in
\[
                    u E^{d-a-B}P^{d-2a} A_i(1),
\]
again small. Finally move the second point from \(q_i\)
to \(\phi v_{i^-}\). Their difference and its translates
belong to \(L_{i,A}(Y')\). Pair a translate with the graph
displacement \(u^p\mu P^d L_{i,A}(Y)\), using the second
pairing bound. This is small as well. Applying the two
cochain identities and summing over Frobenius as in
Theorem~\ref{thm:twolatticecup768} now gives the
coboundary at the product of the two constants, which
is the asserted exponent.
\end{proof}

\begin{corollary}[All scalar targets at unramified spherical places]
\leavevmode\label{cor:allscalartests770}
In the marked torus data, for any sufficiently large node,
let \(w\mid p\) be spherical with \(E_{x,w}\) unramified
and \(p\ge5\). Use the same five centers as in
Corollary~\ref{cor:actualperiodmultiples768}. With a fixed
enlargement of \(C_0\), define \(S_{x,p}\) by its formula
at this prime, also when \(p<x_{1,x}\). Then \(p^{3S_{x,p}}\)
annihilates every free cup evaluation of every scalar
target, in arbitrary graph rank and at the given precision.
\end{corollary}
\begin{proof}
Use \(B=t'-1\), \(a=\lceil12b_x/(p-3)\rceil\). The
height and pullback argument of
Corollary~\ref{cor:actualperiodmultiples768} gives the
\(E^a\)-lifts over this unramified field for every \(p\ge5\).
The scalar target of jump \(j_T\in[0,B]\) has test
extension of jumps \(0,-1-j_T\). Its unit lift follows
from its nonpositive jumps by
Theorem~\ref{thm:integrallifts759}. The graph satisfies
the first alternative of the theorem where \(J_i>0\)
and the second where \(J_i<0\). Since
\(d\ll_{q_0,d_{\rm quad}}1+x^5+b_x/p\), with \(d_{\rm quad}\)
the fixed discriminant parameter, every integer in the
two valuation sums and every positive integer up to \(d\)
belongs, when nonzero, to the five intervals with the
enlarged constant. This follows from the shift
\(O_{q_0,d_{\rm quad}}(x^{41})\) about each base weight
as in Corollary~\ref{cor:actualperiodmultiples768}.
The three costs are each at most \(S_{x,p}\). Descent
of the cobounding cochain from the enlarged coefficients
preserves the original modulus.
\end{proof}
\clearpage

\subsection{Taylor expansion at the next Frobenius prime}
\label{sec:nexttaylormixed772}
Keep the unramified field \(K\), the coefficients
\(\mathcal O\), and \(p\ge5\) of
Theorem~\ref{thm:twolatticecup768}. Set \(P=\phi(E)=u^p-p\).
At \(E=0\) the parameter \(u\) specializes to \(p\); at
\(P=0\) it specializes instead to a root \(\pi_1\) of
\(T^p-p\). Fix the root in the compatible system defining
\(u\), and write \(H=G_{K(\pi_1)}\). We first construct
points on which the displacements for \(H\) have prescribed
orders at both primes.

\begin{lemma}[Increasing the order at \(P\) with the first order fixed]
\leavevmode\label{lem:nextjet772}
Let \(M=M_i(U)\) for a crystalline extension with trivial
quotient map \(\beta\). Suppose \(v\in E^{-e}M\), \(e\ge0\),
has constant quotient, and \(d_v(g)=(g-1)v\) belongs to
\(\delta E^D M_A\) for all \(g\in H\), with \(D\ge0\).
For \(t\ge0\), let \(\mathcal J\) be the union of the
integer intervals \([-j,t-j)\), over the jumps \(j\) of
\(U\) at \(i^-\), with zero removed. Put
\[
              c=\sum_{n\in\mathcal J}(2+v_p(|n|)).
\]
There is a point \(q\in E^{-e}M\) such that
\[
 q-p^c v\in E^D M\cap\ker\beta,\qquad
       (g-1)q\in\delta E^D P^t M_A\quad(g\in H).
\]
\end{lemma}
\begin{proof}
Use the period ring obtained by applying \(\phi^{-1}\)
before the de Rham comparison, with residue map
\(\theta_1=\theta\phi^{-1}\) on \(A\). Its uniformizer
is \(P\). We describe the rational Taylor expansions.
For \(C=\operatorname{Frac}(\mathcal O)\) put
\(C_1=C[T]/(T^p-p)\). This is a product of fields.
At every finite order \(N\), the expansion
\[
             u=T(1+P/p)^{1/p}\pmod {P^N}
\]
identifies \(R_i[1/p]/P^N\) with \(C_1[P]/P^N\).
Indeed division by the distinguished monic polynomial
\((u^p-p)^N\) gives the left side a basis of \(pN\)
monomials over \(C\). In its quotient \(u^p=p+P\),
so \(u(1+P/p)^{-1/p}\) gives the inverse image of \(T\)
by the truncated binomial expansion. In the period ring,
this element at infinite order is the unique root of
\(T^p-p\) reducing to \(\pi_1\) in the internal period
factor. Uniqueness follows by Hensel's lemma after inverting
\(p\). Hence \(H\) fixes these Taylor constants as well
as the external coefficients.

The element \(E\) is a unit in the \(P\)-adic rational
completion: no root of \(T^p-p\) equals \(p\).
Thus \(M\) and \(L_i(U)\) agree there by the height
inclusions. The latter is spanned by \(\phi M_{i^-}(U)\).
Transporting Lemma~\ref{lem:horizontalcoordinates767}
from \(i^-\) by \(\phi\) shows that its horizontal
coordinates are \(P^{-j}C_1[[P]]\) in directions indexed
by precisely the jumps \(j\) at \(i^-\).
This description also applies to \(E^{-e}M\). The transported
horizontal basis vectors are \(H\)-invariant since \(G_K\)
commutes with Frobenius. In these coordinates \(g\in H\)
acts on \(f(P)\) by substitution of \(gP\) for \(P\).

The retraction \(\rho\) in Lemma~\ref{lem:integraljet767}
preserves finite rational Taylor congruences in these
coordinates. To check this, multiply by a power of \(P\)
to remove poles and truncate the coordinate matrices and
the prescribed terms modulo \(P^N\). The isomorphism above
represents the resulting coefficients by polynomials in
\(u\) over \(C\); clear their denominators by a power of
\(p\). The integral congruence holds over \(A_i/P^N\)
because \(\phi^{-1}\) takes \(P\) to \(E\), and the
injection into the rational period ring modulo \(E^N\)
was proved in Lemma~\ref{lem:integraljet767}. Apply
\(\rho\) in the free lattice \(E^{-e}M_A\) and divide back
in the rational comparison. This also proves that
divisibility by \(P^t\) for an element of \(M_A\) can
be tested in this rational comparison.

There is an element \(g_0\in H\) fixing all \(p\)-power
roots of unity and taking \(u\) to \([\epsilon]^p u\).
In fact \(K(\pi_1)/K\) is Eisenstein of degree \(p\)
and is not Galois: a normal extension containing \(\pi_1\)
would contain the \(p\)-th roots of unity, whose degree
over the unramified \(K\) is \(p-1\). The abelian extension
\(K(\mu_{p^\infty})\) cannot contain \(\pi_1\).
The Kummer image over this extension acting on the
compatible roots of \(p\) is therefore a closed subgroup
of \(\ZZ_p\) nonzero modulo \(p\), hence all of \(\ZZ_p\).
The element corresponding to \(p\) is as required.
Write \(\omega=g_0E/E\). Using
\((g_0u-u)=u\mu\sum_{h=0}^{p-1}[\epsilon]^h\)
and the unit \(\mu/(E\tau)\), we get
\[
        v_p\big(\theta(\omega-1)\big)
             =2+\frac1{p-1}.
\]
Also \(g_0P/P=\phi(\omega)\), and
\(\theta_1(\delta)=\pi_1(\zeta_{p^2}-1)\) has valuation
\(1/(p-1)\). For \(n\ne0\) it follows that
\[
 Z_n=\frac{(g_0P/P)^n-1}{\delta}\in A,\qquad
                v_p(\theta_1 Z_n)=2+v_p(|n|).
\]
Integrality follows from \(\omega-1\in\delta A\) and
\(\phi(\delta)\in\delta A\). For the valuation, a unit
\(1+y\) with \(v_p(y)>1/(p-1)\) gains exactly \(v_p(|n|)\)
in \(v_p((1+y)^n-1)\) by the binomial formula, first
for \(n\) prime to \(p\), then for powers of \(p\);
inversion handles negative \(n\).

Lift \(p^{2+v_p(|n|)}/\theta_1 Z_n\) through the
surjection \(\theta_1\) to \(\alpha_n\in A\). At degree
\(n\ne0\) apply
\[
       v_1\longmapsto p^{2+v_p(|n|)}v_1-
                   \rho\big(\alpha_n(g_0v_1-v_1)/\delta\big).
\]
At every step the correction belongs to \(E^D M\)
and has zero quotient. This remains true at the next step,
since \((g-1)(E^D M)\subseteq\delta E^D M_A\).
In a horizontal coordinate agreeing with a constant through
degree \(n-1\), the divided difference multiplies the term
of degree \(n\) by \(Z_n\) modulo the next degree.
Multiplication by \(\alpha_n\) and the Taylor retraction
therefore remove this term after the indicated scaling.
Previously attained orders persist, since \(\delta\) is a
unit in the rational \(P\)-adic comparison. At degree zero
take the actual Taylor constant; it is fixed by \(H\).
Processing the union of the intervals in increasing order
gives agreement with a constant modulo \(P^{t-j}\) in
each direction, just as in Theorem~\ref{thm:deepintegrallift767}.
Thus \((g-1)q/\delta\in P^t M_A\) for \(g\in H\), by
the integral divisibility test. This vector is also in
\(E^D M_A\) by construction. The image of \(E\) modulo
\(P\) is the nonzero divisor \(\pi_1-p\) on the internal
integral period ring, and remains a nonzero divisor after
finite free scalar extension and projection to \(i\).
Cancelling powers of \(P\) gives
\(E^D M_A\cap P^t M_A=E^D P^t M_A\), as required.
\end{proof}

\begin{lemma}[The pole remaining after the first expansion]
\leavevmode\label{lem:remainingpole772}
For the point produced in Theorem~\ref{thm:deepintegrallift767}
let \(e\) be the largest \(-j\) for jumps \(j\in[-a,0)\)
at its label, or zero for the empty set. Then the point
belongs to \(E^{-e}M\). If every jump at that label is at
most \(D\), it also belongs to \(L\).
\end{lemma}
\begin{proof}
For a jump \(j<-a\) the horizontal coordinate has zero
constant and final order at least \(D-j\), as proved in
that theorem. For any other jump its order is at least
\(\min(0,D-j)\). In every case this gives order at least
\(-j-e\), which is the bound for \(E^{-e}M\). When \(j\le D\)
all these orders are nonnegative, the condition for \(L\).
After clearing an \(E\)-power these divisibilities descend
from \(C[[E]]\) to \(R_i\) by Lemma~\ref{lem:labelledheight764}.
\end{proof}

\begin{theorem}[The main cup with both signs and both base weights]
\leavevmode\label{thm:mixedcup772}
At any sufficiently large node let \(p\ge x^{50}\) and
let \(w\mid p\) be spherical. Put \(b=b_x\) and
\(b_2=xb/x_{1,x}\). For some constant \(C\) depending
only on \(q_0,d\), every original cup for the main target
is killed by \(p^h\) with an integer \(h\) satisfying
\[
 h\le 2S_{x,p}
     +C\left(b_2+x^{42}+1+\frac{\log(2b)}{\log p}\right).
\]
Here the radius constant in \(S_{x,p}\) is enlarged once
if necessary. Both signs and base weights may vary at a
single place; the graph rank is arbitrary.
\end{theorem}
\begin{proof}
Work first over \(K=E_{x,w}\) with coefficients containing
\(K\). This field is unramified by our bound on \(p\).
Set \(a=\lceil12b/(p-3)\rceil,\ B=t'-1,\ D=2a+B+3\).
The three extensions of Corollary~\ref{cor:periodlabels764}
have \(E^a\)-lifts as in Corollary~\ref{cor:actualperiodmultiples768}.
Deepen the first two at \(E\) by
Theorem~\ref{thm:deepintegrallift767}. The valuation sums
for depth \(D\) are at most \(S_{x,p}\) since their radius
is \(O_{q_0,d}(x^{42}+b/p)\). By scaling, take every
initial point on both sides over \(p^{S_{x,p}}\mathbf e_i\).
All their displacements are in \(\delta E^D\) times the
respective first kernel lattices.

If the base weight at a label is \(b\), the far jump and
its dual lie at distance greater than \(a+B+2\) from zero,
since \(p\ge x^{50}\) and the shifts are \(O_{q_0,d}(x^{41})\).
At a label with the other base weight, their absolute
values are at most \(b_2+O_{q_0,d}(x^{41})\). By
Lemma~\ref{lem:remainingpole772} all of the initial points
belong to \(E^{-e}M_i\) for one integer
\[
       0\le e\le a,\qquad e\ll_{q_0,d} b_2+x^{42}.
\]
One may take the largest remaining pole bound from that
lemma over both sides and all labels.

At label \(i\) name the two sides \(U,V\) temporarily.
Choose \(U\) to have nonnegative jumps: take the graph
when \(J_i>0\), the main dual when \(J_i<0\). The other
side \(V\) has upper jump at most \(B\). Write \(u_i,v_i\)
for their initial points, and use the same side names for
their antecedents at \(i^-\) in this comparison.
We have \(L_i(U)\subseteq M_i(U)\), whereas \(v_i\in L_i(V)\)
by Lemma~\ref{lem:remainingpole772} since \(D\ge B\).
Apply Lemma~\ref{lem:nextjet772} to \(v_i\) with \(t=e\).
Let \(k\) be the maximum of its costs over all labels.
Multiplying its output at each label by the remaining
power of \(p\) gives \(q_i\) with
\[
 \begin{gathered}
 q_i-p^k v_i\in E^D M_i(V)\cap\ker\beta,\qquad
 q_i\in L_i(V),\\
 d_{q_i}(g)\in\delta E^D P^e M_{i,A}(\ker\beta)
                                      \quad(g\in H).
 \end{gathered}
\]
The membership in \(L_i(V)\) follows from
\(E^D M_i(V)\subseteq L_i(V)\). The kernel equality in
the displacement follows by dividing out the nonzero
factors in the trivial quotient.

Use the cochains on \(H\), scaling \emph{both} original
sides by \(p^k\) at every label. The initial cup of
displacements is small by the first pairing bound
\(E^{-a}\tau^{-1}A_i(1)\) from
Theorem~\ref{thm:twolatticecup768}. First change
\(p^kv_i\) to \(q_i\); their difference is in the first
kernel lattice, while the other displacement supplies
\(\delta E^D\). Next move \(p^k u_i\) to
\(p^k\phi u_{i^-}\). Their difference belongs to
\(E^{-e}P^{-e}M_{i,A}(\ker\beta_U)\), since the translate
lies in \(P^{-e}L_i(U)\) and the two constants agree.
The displacement of \(q_i\) just obtained clears both
poles in the first pairing bound and leaves
\(u E^{D-e-a} A_i(1)\). Finally move \(q_i\) to
\(p^k\phi v_{i^-}\). Both belong to \(P^{-e}L_{i,A}(V)\).
The fixed point \(p^k\phi u_{i^-}\) has displacement in
\(u^p\mu P^D\) times its second kernel lattice. Under the
second pairing bound \(P^{-a}\mu^{-1}A_i(1)\) the result
lies in \(u^p P^{D-e-a}A_i(1)\).

Every displayed bound is small and persists for the
translates of the differences required when the second
argument moves: \(H\) preserves the lattices over \(A_i\)
and the principal ideals of \(E,P\). Thus the same
computations apply with either assignment of the graph
to \(U\) or \(V\). The terminal pair is Frobenius of the
initial pair at \(i^-\) with the original argument order
and the same common scaling. The cochain descent in
Theorem~\ref{thm:twolatticecup768}, now for \(H\), shows
that \(p^{2S_{x,p}+2k}\) kills the restriction of the cup.
Corestriction composed with restriction multiplies group
cohomology by \([G_K:H]=p\), so \(h=2S_{x,p}+2k+1\)
works over the original local field. Retracting the
coefficient extension descends the coboundary at the
original precision.

It remains to estimate \(k\). At the preceding label the
graph jumps are contained in \([0,B]\) and one far jump.
For the dual they consist of zero and one far jump.
Thus each set of integers in Lemma~\ref{lem:nextjet772}
is covered by two intervals with a total of
\(O(e+B+1)\) integers of absolute value at most \(16b\).
In an interval of \(m\) successive integers there are
at most \(m/p^l+1\) nonzero multiples of \(p^l\).
Sum over \(1\le l\le\log(16b)/\log p\) to bound their
total valuation by \(m/(p-1)+\log(16b)/\log p\).
It follows that \(k\ll e+B+1+\log(2b)/\log p\).
The bound on \(e\) proves the theorem.
\end{proof}
\clearpage

\subsection{Horizontal Taylor coefficients after Frobenius}
\label{sec:horizontalradius774}
Keep the unramified field \(K\), \(p\ge5\), \(E=u-p\),
\(A_i\), \(R_i=\mathcal O[[u]]\), and \(C=\operatorname{Frac}(\mathcal O)\)
of Section~\ref{sec:taylordepth767}. For a crystalline
lattice \(U\) put \(M_{0,i}=M_i(U)\) and, for \(l\ge1\),
\[
      M_{l,i}=R_i\,\phi^l M_{0,F^{-l}i}.
\]
In particular \(M_{1,i}=L_i(U)\). To construct a point in
\(M_{2,i}\) deeply fixed at \(E\), we need a bound for
many Taylor coefficients at once. The bound comes from the
series at \(u=0\) before substituting \(u^{p^l}\) for \(u\).

\begin{lemma}[The horizontal matrix at the origin]
\leavevmode\label{lem:originradius774}
Choose an \(R_i\)-basis of \(M_{0,i}\). There is a horizontal
basis of the rational crystalline periods in which the
coordinate matrix of the chosen basis is
\[
           H(u)=I+\sum_{n\ge1}H_n u^n .
\]
This is a matrix of meromorphic functions regular at the
origin. For every \(n\ge1\), both \(p^{n+v_p(n!)}H_n\) and
the coefficient of \(u^n\) in \(p^{n+v_p(n!)}H(u)^{-1}\)
are integral over \(\mathcal O\), entry by entry (that is,
their entries lie in \(\mathcal O\)).
\end{lemma}
\begin{proof}
The analytic crystalline comparison used in
Lemma~\ref{lem:horizontalcoordinates767} identifies the
first Frobenius image rationally with the span of the
constant crystalline periods on the open unit disc after
inverting \(\phi(\lambda)\), and is an isomorphism at \(u=0\).
The change between \(M_{0,i}\) and \(M_{1,i}\) is over
\(R_i[1/E]\), with \(E(0)=-p\ne0\). Hence its horizontal
matrix and inverse are also regular at the origin. Choose
the constant basis to make its value there equal to \(I\).
These assertions commute with a common effective twist and
its inverse. The comparison identifies the constant vectors
with \(G_K\)-invariants inside the crystalline period ring.
~\cite[(2.1.7)--(2.1.9) and Proposition~2.1.10]{III-II-bartlettcrys767}

Take \(g(u)=q u\), \(q=[\epsilon]\), as in
Lemma~\ref{lem:kummerlower760}. By crystalline descent on
the chosen basis its matrix is \(G=I+u\tau D\) with
\(D\in\operatorname{Mat}(A_i)\). In horizontal coordinates
this means
\[
                         H(qu)=H(u)G.
\]
We explain how to compare finite powers of \(u\) in this
identity despite the period coefficients in \(G\).
Put \(\mathcal R_i=\varprojlim_m(A_i/u^m A_i)[1/p]\).
The map from \(A\) into the corresponding rational
completion extends to \(B_{\max}^+\): the latter is the
\(p\)-adic completion of \(A[u/p]\) with \(p\) inverted,
for our unramified parameter. Modulo \(u^m\), terms of
degree at least \(m\) in \(u/p\) vanish and the remaining
denominators are bounded by \(p^{m-1}\). The ring \(A/u^m\)
is \(p\)-adically complete and has no \(p\)-torsion. Indeed
multiplication by the Teichm\"uller element \(u^m\) multiplies
each coefficient in a \(p\)-adic Teichm\"uller expansion
by \(\underline p^{\,m}\), where \(u=[\underline p]\);
the resulting ideal is closed and \(p\)-saturated.
Finite free coefficients and the idempotent preserve these
facts. In particular if an element of \(A_i\) vanishes at
order \(m\) in this rational completion, it lies in
\(u^m A_i\). Here vanishing at order \(m\) means that its
image modulo \(u^m\) is zero.

An analytic function on the open disc evaluates in
\(B_{\max}^+\) and its image in this completion is its
formal series in \(u\). The same holds with \(qu\) in
place of \(u\). Denominators with nonzero constant term
are units in the rational completion. Thus identities of
such meromorphic expressions in the period ring may be
transferred by first clearing the analytic denominators in
\(B_{\max}^+\). This also justifies the divided differences
used below. For any one of these functions \(W\),
\((W(X)-W(Y))/(X-Y)\) is a meromorphic function regular
at the origin of the bidisc: multiply by the two
denominators and use the divided difference of each
analytic numerator or denominator. After substitution
\(X=qu,\ Y=u\), it follows that
\((W(qu)-W(u))/\tau\) is \(u(\mu/\tau)\) times a regular
formal series in the rational completion. The cancellation
by \(\tau\) here takes place in the original period ring
before mapping to this completion.

Induct on \(m\ge1\). Set \(c=m-1+v_p((m-1)!)\).
By induction write
\[
                  p^c H(u)=T(u)+u^m W(u),
\]
where \(T\) is the polynomial of degree less than \(m\)
with the indicated coefficients and is integral. The matrix
\(J=(T(u)G-T(qu))/\tau\) belongs to \(\operatorname{Mat}(A_i)\),
since \(G-I\) is divisible by \(\tau\), and
\((q^h-1)/\tau\in A\) for all \(h\ge0\).
The horizontal identity gives
\[
  J=u^m\left(
       \frac{q^m-1}{\tau}W(u)
       +q^m\frac{W(qu)-W(u)}{\tau}-uW(u)D
     \right).
\]
The expression in parentheses is in \(\operatorname{Mat}(\mathcal R_i)\).
Therefore \(J/u^m\) is integral over \(A_i\) in the sense
that its entries lie in \(A_i\). Its residue modulo \(u\),
after inverting \(p\), is \(((q^m-1)/\tau)W(0)\).
Cancellation of \(u^m\) in the completion follows by
reducing modulo \(u^{m+n}\) for every \(n\), since \(u\)
is a nonzero divisor before completion.

By the factorization in Lemma~\ref{lem:kummerlower760},
\((q^m-1)/\tau\) modulo \(u A\) is \(p^{1+v_p(m)}\)
times a unit. In fact its factors up to units are
\(\phi^j(\mu/\tau)\) for \(0\le j\le v_p(m)\),
each associated to \(\phi^j(E)\equiv-p\bmod u\).
Multiplying by the inverse unit shows
\(p^{1+v_p(m)}W(0)\) is in the image of \(A_i/u A_i\).
These entries are rational constants in \(C\). The
\(R_i\)-linear retraction of Lemma~\ref{lem:integraljet767},
reduced modulo \(u\), maps integral entries to \(\mathcal O\)
and fixes the constants after inverting \(p\). Hence
\(p^{1+v_p(m)}W(0)\) has entries in \(\mathcal O\).
Since \(W(0)=p^c H_m\), this completes the induction.
The series with these coefficient bounds are closed under
multiplication: \(j!(n-j)!\) divides \(n!\). The recursion
for the inverse of a matrix series with constant \(I\)
then gives the same assertion for \(H^{-1}\).
\end{proof}

\begin{theorem}[Integral approximation of a horizontal coefficient]
\leavevmode\label{thm:frobeniusjet774}
Let \(Q=p^l\ge3,\ A_0\ge1,\ d_0=2\lceil A_0/Q\rceil\).
In the horizontal basis transported by \(\phi^l\) from the
origin-normalized basis at \(F^{-l}i\), the following hold
for \(M_{l,i}\) at \(u=p\).
\begin{enumerate}
\item For \(y\in M_{l,i}\), every coefficient of \(E^j\),
\(0\le j<A_0\), in its horizontal expansion belongs to
\(p^{-d_0}\mathcal O^{\operatorname{rank}U}\).
\item For an integral constant vector \(v\) in that basis,
there is \(z\in M_{l,i}\) with horizontal expansion
\(p^{d_0}v+O(E^{A_0})\). In particular, for \(y\in M_{l,i}\)
and \(N<A_0,\ N\ge0\), one can realize
\(p^{2d_0}[E^N]y\) up to \(O(E^{A_0})\) in \(M_{l,i}\),
where brackets denote the horizontal Taylor coefficient.
\end{enumerate}
\end{theorem}
\begin{proof}
The basis obtained by applying \(\phi^l\) to the original
lattice basis has horizontal matrix \(H_l(u)=H(u^Q)\),
with the corresponding Frobenius transport of the constants.
It is a basis of \(M_{l,i}\) by its definition. Its inverse
has the same substitution. The lemma and \(v_p(n!)\le n\)
bound the coefficients of \(u^{Qn}\) in both matrices
below by valuation \(-2n\). In particular these series
converge near \(u=p\). The coefficient of \(E^j\) from
the term of index \(n\), for \(Qn\ge j\), has valuation at
least
\[
       Qn-j-2n=(1-2/Q)(Qn-j)-2j/Q\ \ge\ -2j/Q.
\]
The terms tend to zero \(p\)-adically as \(n\) grows.
Thus the Taylor polynomials of both \(p^{d_0}H_l\) and
\(p^{d_0}H_l^{-1}\) below degree \(A_0\) have integral
coefficients. Multiplying an integral lattice coordinate
vector by \(H_l\) proves (1), since the expansion of
\(R_i=\mathcal O[[E]]\) is integral. For (2), use as lattice
coordinates the Taylor polynomial of \(p^{d_0}H_l^{-1}v\)
below that degree. Multiplying back by \(H_l\) gives the
asserted congruence. Apply (1) to \(y\) and then (2) to
\(v=p^{d_0}[E^N]y\) for the last assertion.
\end{proof}
\clearpage

\subsection{A common point in two Frobenius lattices}
\label{sec:secondpoint775}
Let \(p\ge5\) and \(K/\QQ_p\) be unramified. Use the
lattices \(M_{l,i}(U)\) of Section~\ref{sec:horizontalradius774}
for a crystalline extension \(\beta:U\to\mathcal O\) with
trivial quotient. We suppress \(\mathbf e_i\) on its scalar
coordinate. Put \(i^-=F^{-1}i,\ i^{--}=F^{-2}i\),
\(P=u^p-p\) and \(H=G_{K(p^{1/p})}\), with the choice of
root in Section~\ref{sec:nexttaylormixed772}. We will use
an integral point in \(M_{2,i}\) in a comparison that
passes over the intervening label \(i^-\).

\begin{lemma}[Two divisibilities for the second Frobenius lattice]
\leavevmode\label{lem:secondorders775}
Let \(z\in M_{2,i}\) have constant quotient, and suppose
\((g-1)z\in\delta E^A M_{2,i,A}\) for \(g\in H\), \(A\ge0\).
For \(t\ge1\) set \(c_t=2(t-1)+v_p((t-1)!)\).
There exists \(q\in M_{2,i}\) such that
\[
 q-p^{c_t}z\in E^A M_{2,i}\cap\ker\beta,\qquad
 (g-1)q\in\delta E^A P^t M_{2,i,A}\quad(g\in H).
\]
\end{lemma}
\begin{proof}
The \(A_i\)-span of \(M_{2,i}\) is stable. Its generators
\(\phi^2 m\) have displacement divisible by
\(\phi^2(\delta)\in\delta A_i\), since \(g\) commutes with
\(\phi\). For coefficients \(r\in R_i\), \(gr-r\in\delta A_i\)
by crystalline descent on the trivial line. As \(gE/E\)
belongs to \(1+\delta A_i\), the same argument gives
\((g-1)(E^A M_{2,i})\subseteq\delta E^A M_{2,i,A}\).

At \(P=0\) all the horizontal coordinates of \(M_{2,i}\)
are regular and form the full power series lattice.
Indeed the span of \(\phi M_{0,i^{--}}\) at \(i^-\)
is \(M_{1,i^-}\), whose \(E\)-adic rational lattice is
the span of a constant basis by
Lemma~\ref{lem:horizontalcoordinates767}. Apply \(\phi\)
and extend scalars to the \(P\)-adic rational completion
at \(i\). A basis and its inverse remain regular, since
\(\phi(E)=P\). The Taylor constants are in
\(C[T]/(T^p-p)\) in these coordinates, and are fixed by
\(H\) as in Lemma~\ref{lem:nextjet772}.

Apply the coefficientwise \(R_i\)-linear retraction in
a basis of \(M_{2,i}\). It respects congruences in
horizontal coordinates at \(P\): truncate the regular
change of basis modulo the required \(P^n\), represent
the coefficients by polynomials over \(C\) using the
finite Taylor isomorphism of Lemma~\ref{lem:nextjet772},
then clear a common power of \(p\) and retract. The
injection of \(A_i/P^n\) in the rational period quotient
justifying this step is the one in that lemma.
Starting with the constant term in each coordinate,
process \(n=1,\ldots,t-1\). With its \(g_0,\alpha_n\)
replace the current point \(v\) by
\[
 p^{2+v_p(n)}v-
      \rho\big(\alpha_n(g_0-1)v/\delta\big).
\]
The correction lies in \(E^A M_{2,i}\) and has zero quotient.
The established stability preserves the hypothesis about
the divided displacement for the next step. Since
\(\alpha_n((g_0P/P)^n-1)/\delta\equiv
p^{2+v_p(n)}\bmod P\), the term of degree \(n\) beyond
the constant is removed, and all lower orders are preserved.
The point obtained agrees with an \(H\)-fixed constant
modulo \(P^t\). For every \(g\in H\), \(\delta\) is a unit
in the rational \(P\)-adic comparison. Thus its divided
displacement lies in \(P^t M_{2,i,A}\) by the same
integral divisibility test. It also belongs to
\(E^A M_{2,i,A}\). The intersection is the product,
since \(E\) is a nonzero divisor modulo \(P\), as in
Lemma~\ref{lem:nextjet772}. The sum of the costs is \(c_t\).
\end{proof}

\begin{theorem}[A lift fixed to high order across an intervening label]
\leavevmode\label{thm:secondpoint775}
Let \(s,\kappa,h_0,h_1\) be nonnegative integers.
Assume \(E^s\) lifts to \(M_{0,i}\) and \(p^\kappa\) lifts
to \(M_{0,i^{--}}\). Suppose all jumps at \(i\) are at
most \(h_0\), and those at \(i^-\) are at most \(h_1\).
Put \(N=s+h_0+p h_1\). For integers \(A>N,\ t\ge1\)
there is \(q\in M_{2,i}\) with
\[
 \begin{split}
 \beta(q)&=p^{\kappa+4\lceil A/p^2\rceil+c_t},\\
 (g-1)q&\in\delta E^A P^t M_{2,i,A}\qquad(g\in H).
 \end{split}
\]
In a horizontal basis adapted to the jumps at \(i\), the
coordinate of \(q\) at every jump \(j\) with \(h_0-j>N\)
has order at least \(A\) at \(E=0\). In particular
\(q\in M_{0,i}\cap M_{2,i}\) if all jumps at \(i^-\)
are nonnegative and every negative jump \(j\) at \(i\)
satisfies \(h_0-j>N\) and \(-j\le A\).
\end{theorem}
\begin{proof}
Take \(v\in E^{-s}M_{0,i}\) over \(1\). With \(E=u-p\),
the polynomial
\[
          f(E)=E^{s+h_0}\big((p+E)^p-p\big)^{h_1}
\]
is monic of degree \(N\). The successive upper-height
inclusions put \(y=f(E)v\) in \(M_{2,i}\): first
\(E^{s+h_0}v\in M_{1,i}\), then
\(P^{h_1}M_{1,i}\subseteq M_{2,i}\). At \(E=0\), \(M_{2,i}\)
has the same rational lattice as \(M_{1,i}\), since \(P\)
is a unit there and their two transition matrices have
only powers of \(P\) as denominators. These are the regular
horizontal coordinates of Lemma~\ref{lem:horizontalcoordinates767}.
In the direction of jump \(j\), the coordinate of \(y\)
has order at least \(h_0-j\), using the description of
\(M_{0,i}\) in the same lemma. Thus \([E^N]y\) is zero
in all the directions specified in the theorem.

With \(d=4\lceil A/p^2\rceil\),
Theorem~\ref{thm:frobeniusjet774} supplies \(\xi\in M_{2,i}\)
equal to the constant \(p^d[E^N]y\) modulo order \(A\).
This congruence is unchanged by any constant change of
horizontal basis, including the change to one adapted to
the filtration. The quotient map in horizontal bases is
constant by functoriality of crystalline periods. Since
\(\beta(y)=f(E)\) and its coefficient of \(E^N\) equals
one, \(\beta(\xi)=p^d+E^A w\) with \(w\in R_i\);
integral divisibility follows from \(R_i=\mathcal O[[E]]\).
Furthermore \((g-1)\xi/\delta\in M_{2,i,A}\) before any
Taylor comparison. In that comparison \(g\) fixes the
indicated constant, \(gE/E\) is a unit, and \(\delta\)
specializes to \(p(\zeta_p-1)\ne0\). Thus this divided
displacement lies in \(E^A M_{2,i,A}\) by the integral
test from Lemma~\ref{lem:integraljet767}, applied entrywise
in a basis of \(M_{2,i}\).

Apply \(\phi^2\) to a lift at \(i^{--}\) to get
\(y_0\in M_{2,i}\) over \(p^\kappa\). Set
\[
                   z=p^\kappa\xi-E^A w y_0 .
\]
Its quotient is \(p^{\kappa+d}\). The second term has
displacement in \(\delta E^A M_{2,i,A}\) by the stability
proved above. All its horizontal coordinates have order
at least \(A\) at \(E=0\). Lemma~\ref{lem:secondorders775}
applied to \(z\) gives the stated \(q\), preserving the
orders in the specified directions.

For the last assertion, nonnegative previous jumps give
\(M_{2,i}\subseteq M_{1,i}\) by Frobenius transport.
Consequently \(q\) is in \(E^{-r_i}M_{0,i}\), where
\(r_i=\max(0,\max_j(-j))\). In every nonnegative direction
its coordinate is regular, hence has order at least \(-j\).
In every negative direction the hypothesis gives order
at least \(A\ge-j\). The horizontal description now puts
\(q\) in \(M_{0,i}\) rationally at \(E\). Divisibility after
clearing \(E^{r_i}\) descends in \(\mathcal O[[E]]\).
\end{proof}
\clearpage

\subsection{The main cup on an orbit with both base weights}
\label{sec:mainorbit776}
Write \(b=b_x,\ G=xb/x_{1,x}\) for the two integral base
weights of the shorter torus family. At each labelled
embedding there is one of these two weights. For a fixed
datum let \(B=t'-1\). Choose an integer
\(W=\lceil C_0 x^{42}\rceil\), with \(C_0\) depending only on
\(q_0,d\), large enough for the following bounds at every
label. If its base weight is \(\beta\in\{b,G\}\), the far
jump of the graph and the far jump of the main dual, when
present, have absolute values in \([\beta-W,\beta+W]\).
The other jumps on these two extensions are in \([0,B]\).
Every upper jump on all three extensions, including the
extension defining the integral contraction, is at most
\(\beta+W\). Indeed the graph jumps belong to
\(\{0,\ldots,B,J_i\}\), and the main dual jumps belong to
\(\{0,-1-J_i\}\). The third kernel has jumps \(J_i-j\)
at the same label by Corollary~\ref{cor:periodlabels764};
these are zero or have absolute value at most \(|J_i|+B\).
The shifts of \(J_i\) are \(O_{q_0,d}(x^{41})\) and
\(B\ll x^5\). We take \(W\ge B+10\).

\begin{theorem}[A uniform bound for the main cup at large primes]
\leavevmode\label{thm:mainorbit776}
At sufficiently large nodes, for every prime \(p\ge x^{50}\)
and every spherical place \(w\mid p\), each original
cup for the main target is killed by \(p^h\), where one
can take a nonnegative integer \(h\) satisfying
\[
 h\ll_{q_0,d}
       W+1+\frac{b}{p^2}+\frac{G}{p}
                    +\frac{\log(2b)}{\log p}.
\]
The estimate is uniform in all twists in the designated
family, trace ranks, tests and precisions.
\end{theorem}
\begin{proof}
Let \(K=E_{x,w}\), unramified in this range. Extend
coefficients to contain \(K\), using the original ideal
\(\pi^m\mathcal O'\). Denote the graph and dual extensions
by \(X,X'\), their finite pairing by \(Q\), and the
extension defining the contraction by \(X''\). Put
\[
 a=\left\lceil\frac{b+W}{p-3}\right\rceil,\quad
 H_p=W+1+\left\lceil b/p^2\right\rceil+\left\lceil G/p\right\rceil,
 \qquad D=2a+B+4 .
\]
Since \(p\ge7\), \(a\le b/p+8H_p\) and \(a/p\le4H_p\);
these follow by subtracting \(b/p\) before estimating the
ceiling. At every label the image of each of the three
lattices \(M_0\) contains \(p^k E^s\) for nonnegative
\(k+s\ge a\). This is Theorem~\ref{thm:uniformlift760}
with upper jump bound \(b+W\), pulling back along the
trivial \(\ZZ_p\)-line, then projecting by the labelled
idempotent as in Corollary~\ref{cor:actualperiodmultiples768}.

If the sign or the base weight is constant on the orbit,
Lemma~\ref{lem:onesignmain771} gives \(2S_{x,p}\).
Each of its five intervals has \(O(W+1+b/p)\) integers
with absolute value at most \(16b\). In an interval of
\(m\) consecutive integers, the number of nonzero
multiples of \(p^r\) is at most \(m/p^r+1\).
Summing this inequality for \(1\le r\le\log(16b)/\log p\)
bounds the sum of valuations by \(m/(p-1)+
\log(16b)/\log p\). This proves the assertion in this case.

It remains to consider a place with both signs and both
weights. The rational prime is inert in \(\QQ(\sqrt{-d})\):
it is unramified, and a split prime has constant sign.
Consecutive labels restrict to conjugate embeddings on
the unramified quadratic completion, since \(F\) acts by
one Frobenius step. The sign is determined by this
restriction by Theorem~\ref{thm:negativeinterval740}.
In particular the signs alternate. At a positive label
the graph is nonnegative and the dual has upper jump
zero. At a negative label the dual is nonnegative and
the graph has upper jump at most \(B\).

We choose a set of labels at which to start integral
points. Write \(s_i\) for a permitted pole order at a
chosen label and \(\kappa\) for a common exponent. The
following two choices will be used.
\begin{enumerate}
\item If \(b/p-G\le40H_p\), choose every label, put
\(s_i=s=\min(a,G-W-2)\) and \(\kappa=a-s\). We have
\(G\ge b^{1-10^{-6}}\), so \(s\ge0\) for large \(x\).
The inequality for \(a\) gives \(\kappa\le50H_p\).
The images contain \(p^\kappa E^s\). A negative far
jump has size at least \(G-W>s\), so neither of the
first two extensions has a jump in \([-s,0)\).
\item Suppose \(b/p-G>40H_p\). Choose every label of
weight \(b\) and every label whose successor has weight
\(G\). At a chosen label with successor of weight \(G\)
put \(s_i=0\). At each such label a pure power \(p^\kappa\)
lifts on all three extensions for one \(\kappa\ll H_p\).
Indeed the sum of the successive upper jumps divided by
the respective powers of \(p\) is bounded by
\[
                 (G+W)/p+(b+W)/(p(p-1)).
\]
The estimate for \(c_i\) in Theorem~\ref{thm:orbitlift764}
is at most a constant times \(1\) plus this sum. Take
the largest of those bounds, rounded up, for \(\kappa\).
At every remaining chosen label put \(s_i=a\).
Its weight is \(b\), so the negative far jumps have
size at least \(b-W>a\). Also \(E^a\) lifts there on all
three extensions. Use the same \(\kappa\) by scaling.
A label not chosen has weight \(G\) and successor of
weight \(b\); its predecessor has a successor of weight
\(G\), hence was chosen. Thus successive chosen labels
are one or two steps apart.
\end{enumerate}

At a chosen label take, on \(X\) and \(X'\), lifts of
\(p^\kappa E^{s_i}\) and divide by \(E^{s_i}\).
Apply Theorem~\ref{thm:deepintegrallift767} at depth \(D\).
Its operators multiply a constant quotient by the same
powers for any initial constant, hence apply also to
\(p^\kappa\). The resulting points are in \(M_{0,i}\):
for \(s_i=0\) they began there, and otherwise the required
gap was just checked. Their displacements lie in
\(\delta E^D\) times their first kernel lattices. The cost
in powers of \(p\) beyond \(\kappa\) is
\[
                       O\big(H_p+\log(2b)/\log p\big).
\]
Indeed the union of the intervals \([-s_i-j,D-j)\) is
covered by the short cluster and one far interval with
a total of \(O(a+B+1)\) consecutive indices, counted
within each interval. Their nonzero absolute values are
at most \(16b\). Apply the preceding valuation count
and use \(a/p\le4H_p\). On \(X''\) choose
\(z_i\in E^{-s_i}M_{0,i}(X'')\) over \(p^\kappa\).
Contraction with \(z_i\) gives the bound
\(E^{-s_i}\tau^{-1}A_i(1)\) for \(p^\kappa Q\) on the
first kernel lattices, taken at the original precision.
Transport by \(\phi^l\) gives
\[
        \phi^l(E)^{-s_i}\phi^l(\tau)^{-1}A_{F^li}(1)
\]
on their \(l\)-th Frobenius lattices at \(F^li\).
These statements take place after reduction in the
finite Witt comparison, where the displayed denominators
are units before imposing the indicated lattice bounds.

For each chosen endpoint \(i\) two steps from its
preceding chosen label, name the nonnegative side at
\(i\) by \(U\) and the other by \(V\). Here \(i\) has
weight \(b\), its predecessor \(i^-\) has weight \(G\),
and \(s_{i^{--}}=0\). The extension \(V\) has upper
jump \(B\) as a bound at \(i\), is nonnegative at
\(i^-\), and has upper bound \(G+W\) there. Each of
its negative jumps \(j\) at \(i\) satisfies
\(b-W\le-j\le b+W\). Apply
Theorem~\ref{thm:secondpoint775} to \(V\) with
\[
 s=a,\quad h_0=B,\quad h_1=G+W,\quad
 A=2(b+W)+B+4,\quad t=B+1 .
\]
The pure lifts required are \(E^a\) at \(i\) and
\(p^\kappa\) at \(i^{--}\). To verify the degree condition,
\(p(G+W)<b-39pH_p\), and \(a+W< 4pH_p\) by the bound
for \(a\) and \(b/p\le pH_p\). Hence
\[
             a+p(G+W)< b-W,\qquad
 N=a+B+p(G+W)<A.
\]
In each negative direction \(B-j>N\) and \(A\ge-j\).
The theorem supplies \(q_i\in M_{0,i}(V)\cap M_{2,i}(V)\)
with displacement in \(\delta E^A P^{B+1}M_{2,i,A}\)
on \(H=G_{K(p^{1/p})}\). The exponent of its quotient
is \(\kappa+4\lceil A/p^2\rceil+2B+v_p(B!)\ll H_p\);
here \(B<p\).

Choose \(e_*\) to be the maximum of the quotient
exponents of the points on the first two extensions at
all chosen labels and of all the \(q_i\) just constructed.
Multiply each point by the additional power of \(p\)
needed for quotient \(p^{e_*}\). This preserves all
membership and displacement bounds. Thus
\(e_*\ll H_p+\log(2b)/\log p\). Use the pairing
\(\widetilde Q=p^\kappa Q\) throughout. At each unchosen
label use, for both initial points, the Frobenius of the
initial points at its predecessor. For the following
comparisons a value in the image of \(uA_i(1)\) at the
given precision will be called small, as in
Theorem~\ref{thm:twolatticecup768}.

At a chosen label the initial cup of displacements is
small since one supplies \(\delta E^D\) and the first
pairing bound costs \(E^{s_i}\tau\), with \(D\ge s_i\).
At an unchosen label the initial cup is the Frobenius
of the previous one and is still small since
\(\phi(u)=u^p\); its comparison with the preceding
Frobenius is the identity.

First suppose that the preceding label \(i^-\) and \(i\)
are both chosen. Name their two sides by reference to
\(i\), taking \(U\) nonnegative there and \(V\) with
upper jumps bounded by \(B\). Denote the points at
\(i\) by \(u_i,v_i\). Move \(u_i\) to \(\phi u_{i^-}\)
at fixed \(v_i\). The difference is in the first kernel
lattice since \(M_{1,i}(U)\subseteq M_{0,i}(U)\).
Pairing with \(d_{v_i}\) leaves \(uE^{D-s_i}A_i(1)\).
Next move \(v_i\) to \(\phi v_{i^-}\). Both are in
\(M_{1,i}(V)\), the current point by
Lemma~\ref{lem:remainingpole772} since \(D\ge B\).
Pair the kernel difference with the displacement of
\(\phi u_{i^-}\); the surviving factor is
\(u^p P^{D-s_{i^-}}\). Both moves give small values.

Next suppose that \(i^-\) is unchosen and \(i\) is chosen.
The preceding initial pair at \(i^-\) is the Frobenius
of the pair at \(i^{--}\), so the required terminal
points are \(\phi^2 u_{i^{--}},\phi^2 v_{i^{--}}\).
Use the side names for \(i\) in all three moves.
\begin{enumerate}
\item Move \(v_i\) to \(q_i\) at fixed \(u_i\). The
difference lies in the first kernel lattice. The factor
\(\delta E^D\) of the other displacement and the first
pairing bound give smallness as above.
\item Move \(u_i\) to \(\phi^2 u_{i^{--}}\) at fixed
\(q_i\). The difference belongs to
\(E^{-(b+W)}P^{-B}M_{2,i,A}(\ker\beta_U)\). Indeed the
upper bounds on \(U\) are \(b+W\) at \(i\) and \(B\) at
\(i^-\), because the sign alternates. Successive height
inclusions give the assertion in the extension lattice.
The quotient coordinates agree, so it holds in the
kernel. Here and in the other moves the kernel assertion
is Lemma~\ref{lem:adjacentlattice762}, transported by
Frobenius. Base change from \(\phi^l(R_{F^{-l}i})\)
to \(R_i\) is flat, with basis \(1,u,\ldots,u^{p^l-1}\),
and extension to \(A_i\) is flat.
The second Frobenius pairing at \(i\) costs only
\(\phi^2(\tau)\), since \(s_{i^{--}}=0\).
This element is a unit times \(EP\tau\), by applying
Frobenius to \(\mu/\tau\) as in
Lemma~\ref{lem:kummerlower760}. Pair with
\(d_{q_i}\in\delta E^A P^{B+1}M_{2,i,A}\); what remains
is in \(u E^{A-(b+W)-1}A_i(1)\).
\item Move \(q_i\) to \(\phi^2 v_{i^{--}}\). Both points
are in \(M_{2,i}(V)\). The fixed first point has
displacement divisible in that lattice for \(U\) by
\(\phi^2(\delta E^D)=u^{p^2}\phi^2(\tau)\phi^2(E)^D\).
The same pairing bound leaves a multiple of \(u^{p^2}\).
\end{enumerate}
All bounds apply also to the translates of kernel
differences in second-argument moves: \(H\) preserves
the \(A_i\)-lattices and the ideals of \(E,P\). Thus the
argument uses the given order of the graph and dual in
the cup, irrespective of their temporary names.

The terminal pair at every label is the single Frobenius
of the actual initial pair at \(i^-\), also when two
moves around the orbit return to the original label.
The cochain argument of Theorem~\ref{thm:labelledcup765}
therefore applies on \(H\) at precision \(\pi^m\).
Explicitly, the initial displacements give differential
\(p^{2e_*}\mathbf e_i\widetilde Q(c\cup c')\) modulo small
cochains, and the moves identify the initial and terminal
cochains there up to a coboundary. Summing over all
labels, correcting the resulting Frobenius difference
by a coboundary, and taking Frobenius constants gives a
coboundary over the original finite module, as in
Theorem~\ref{thm:twolatticecup768}. Thus restriction is
killed by \(p^{2e_*+\kappa}\). The extension
\(K(p^{1/p})/K\) has degree \(p\) by Eisenstein's
criterion. Corestriction shows that \(h=2e_*+\kappa+1\)
works over \(K\). An integral coefficient retraction fixing
\(1\) descends the coboundary to the original \(\pi^m\).
The bounds for \(e_*\) and \(\kappa\) give the estimate.
\end{proof}
\clearpage

\subsection{Taylor conditions at several Frobenius primes}
\label{sec:severalprimes779}
Keep \(K/\QQ_p\) unramified, \(p\ge5\), and use the
coefficient ring and Frobenius lattices of
Section~\ref{sec:horizontalradius774}. Put \(P_j=\phi^j(E)
=u^{p^j}-p\) for \(j\ge0\). Choose the compatible roots
\(\pi_j=p^{1/p^j}\) used to define \(u\).
For \(L\ge1\) let \(H_L=G_{K(\pi_L)}\). The rational
Taylor expansion at \(P_j\), for \(0\le j\le L\), is
computed with the map \(\theta_j=\theta\phi^{-j}\).
Its coefficient algebra for the variable \(u\) is
\(C[T]/(T^{p^j}-p)\), a product of fields.

\begin{proposition}[Simultaneous orders on a Frobenius lattice]
\leavevmode\label{prop:simultaneousorders779}
Let \(\beta:U\to\mathcal O\) be a crystalline lattice
surjection with trivial quotient. Fix a label \(i\) and
\(s\ge0\), and write \(M=M_{s,i}(U)\).
Let \(F_0\) be a product of nonnegative powers of the
\(P_v\), \(0\le v\le L\), excluding indices in a set
\(I\subseteq\{0,\ldots,L\}\). Suppose \(z\in M\) has
constant quotient and
\((g-1)z\in\delta F_0 M_A\) for \(g\in H_L\).
Prescribe an order \(T_j\ge1\) for every \(j\in I\).
For \(j<s\) put \(\mathcal D_j=\{0\}\); for \(j\ge s\)
take for \(\mathcal D_j\) the jumps of \(U\) at \(F^{-j}i\).
Set
\[
 I_j=\bigcup_{a\in\mathcal D_j}([-a,T_j-a)\cap\ZZ)
                           \ \setminus\{0\},\qquad
 c=\sum_{j\in I}\sum_{n\in I_j}(L+2+v_p(|n|)).
\]
There is \(z'\in M\) such that
\[
 z'-p^c z\in F_0 M\cap\ker\beta,\qquad
 (g-1)z'\in\delta F_0\prod_{j\in I}P_j^{T_j}M_A
                                                \quad(g\in H_L).
\]
\end{proposition}
\begin{proof}
For each \(j\), at finite rational precision in \(P_j\)
one has \(u=T(1+P_j/p)^{1/p^j}\), by the truncated
binomial formula, as in Lemma~\ref{lem:nextjet772}.
The inverse takes \(T\) to \(u(1+P_j/p)^{-1/p^j}\).
The two algebras agree by these inverse maps, or by
division by the distinguished monic polynomial \(P_j^n\)
at precision \(n\). In the internal period factor the
lift of \(\pi_j\) satisfying \(T^{p^j}=p\) is unique
after inverting \(p\), since the derivative at \(\pi_j\)
is nonzero. Thus the Taylor constants are fixed by \(H_L\)
for \(j\le L\), including after extension of coefficients.

At \(P_j\), the transition from \(M_{a,i}\) to \(M_{a+1,i}\)
has denominators only at \(P_a\). Thus if \(s>j\) the
rational lattice equals that of \(M_{j+1,i}\), and if
\(s\le j\) it equals that of \(M_{j,i}\). Transport
Lemma~\ref{lem:horizontalcoordinates767} from \(F^{-j}i\)
by \(\phi^j\). It gives a constant horizontal basis in
which the lattice at \(P_j\) is a direct sum of the
power series modules shifted by \(-a\), for precisely
the \(a\) stated in the proposition, with multiplicities.
The higher Frobenius lattice is regular in all directions.
The retraction of Lemma~\ref{lem:integraljet767} in a
basis of \(M\) preserves finite Taylor congruences here.
Indeed one clears poles, represents the truncated
coefficients by polynomials over \(C\) using the
preceding isomorphism, clears powers of \(p\), and applies
the \(R_i\)-linear retraction. The required injection
of \(A_i/P_j^n\) into the rational period quotient follows
from the \(E\)-adic injection by \(\phi^j\), also on
finite free coefficients. This likewise detects
divisibility of a vector of \(M_A\) by \(P_j^n\).

There exists \(g_0\in H_L\) acting by
\(g_0u=[\epsilon]^{p^L}u\) and fixing all \(p\)-power
roots of unity. The Kummer image over
\(K(\mu_{p^\infty})\) is all of \(\ZZ_p\), as proved
in Lemma~\ref{lem:nextjet772}; take the element with
parameter \(p^L\). Write \(\omega=g_0E/E\).
Since \(\mu/(E\tau)\) is a unit and
\(\theta(\sum_{h=0}^{p^L-1}[\epsilon]^h)=p^L\),
\[
 v_p(\theta(\omega-1))=L+1+\frac1{p-1},\qquad
 v_p(\theta_j(\delta))=\frac{1+1/(p-1)}{p^j}.
\]
The second formula uses
\(\theta_j(\delta)=\pi_j(\zeta_{p^{j+1}}-1)\).
Now \(g_0P_j/P_j=\phi^j(\omega)\), and
\(\phi^j(\delta)\in\delta A\).
Consequently, for \(n\ne0\) the integral element
\[
        Z_{j,n}=\big((g_0P_j/P_j)^n-1\big)/\delta
\]
has nonzero \(\theta_j\)-value of valuation
\[
 L+1+\frac{1}{p-1}+v_p(|n|)
                         -\frac{1+1/(p-1)}{p^j}.
\]
Here one uses the binomial formula for a unit differing
from one by valuation greater than \(1/(p-1)\), including
inversion for \(n<0\).
This is at most \(c_n=L+2+v_p(|n|)\).
Lift \(p^{c_n}/\theta_j(Z_{j,n})\) via \(\theta_j\) to
\(\alpha\in A\).

When the already imposed product is \(F\), process the
integers for the next \(j\) in increasing order using
\[
             v\ \longmapsto\ p^{c_n}v-
                   \rho\big(\alpha(g_0-1)v/\delta\big).
\]
The correction is in \(FM\) and has zero quotient.
Indeed every generator \(\phi^s m\) has displacement
in \(\phi^s(\delta)M_A\), the ring coefficients have
displacement in \(\delta A_i\), and for all \(g\in H_L\)
the ratios \(gP_v/P_v\) belong to \(1+\delta A_i\).
Thus \((g-1)(FM)\subseteq\delta FM_A\), preserving the
hypothesis for the next operation. In a horizontal
coordinate agreeing with a constant up to degree \(n\)
(terms below degree \(n\) agree), the degree-\(n\) term
of the divided displacement is \(Z_{j,n}\) times its
degree-\(n\) coefficient. Multiplication by \(\alpha\)
and retraction therefore cancel that term after scaling
by \(p^{c_n}\). Other coordinates retain their attained
orders, and at degree zero one takes the current constant
coefficient. This is the induction of
Theorem~\ref{thm:deepintegrallift767} for the shifted
intervals \([-a,T_j-a)\). It gives agreement with a
constant modulo order \(T_j-a\) in each direction.
Since \(\delta\) has nonzero specialization, the final
divided displacement belongs to \(P_j^{T_j}M_A\).

For \(v\ne j\), the specialization of \(P_v\) at
\(\theta_j\) is \(\pi_j^{p^v}-p\ne0\). Hence any product
of the other primes is a nonzero divisor modulo \(P_j\),
also after finite free extension and on the labelled
factor. Cancelling \(P_j\) repeatedly gives
\(FM_A\cap P_j^{T_j}M_A=FP_j^{T_j}M_A\).
Use this product for the subsequent indices. The
corrections at every stage belong to \(F_0M\cap\ker\beta\),
so their scaled sum gives the first assertion as well.
\end{proof}

\begin{theorem}[A common point across a string of bounded jumps]
\leavevmode\label{thm:longpoint779}
Let \(l\ge1,\ p^l\ge3\), and suppose \(E^a\) lifts at
\(i\) in \(M_{0,i}\) and \(p^\kappa\) lifts at
\(F^{-l}i\) in \(M_{0,F^{-l}i}\). Let \(h_j\ge0\)
bound the upper jumps at \(F^{-j}i\) for \(0\le j<l\).
Suppose the intermediate jumps (\(1\le j<l\)) are all
at least \(-\eta\), with \(\eta\ge0\). Put
\[
                     N=a+\sum_{j=0}^{l-1}p^j h_j .
\]
Let \(A>N\), and suppose every jump \(t<-\eta\) at the
current label satisfies \(h_0-t>N\) and \(-t\le A\).
There exists \(v\in M_{l,i}\) such that
\[
 \begin{gathered}
   \beta(v)=p^{\kappa+4\lceil A/p^l\rceil},\qquad
   (g-1)v\in\delta E^A M_{l,i,A}\quad(g\in G_K),\\
                \left(\prod_{j=0}^{l-1}P_j^\eta\right)v
                          \in M_{0,i}.
 \end{gathered}
\]
\end{theorem}
\begin{proof}
Divide a lift of \(E^a\) by \(E^a\) to get \(x\) over \(1\).
Take \(f=E^a\prod_{j<l}P_j^{h_j}\), a monic polynomial
of degree \(N\) in \(E=u-p\).
The successive upper-height inclusions give \(y=fx\in M_{l,i}\).
At \(E=0\) all factors \(P_j\), \(j>0\), are rational
units. The coordinates of \(y\) in a basis adapted to
the jumps \(t\) have orders at least \(h_0-t\).
The coefficient \([E^N]y\) consequently vanishes in all
directions \(t<-\eta\), and its quotient is \(1\).
Realize \(p^{4\lceil A/p^l\rceil}[E^N]y\) modulo \(E^A\)
by an integral \(\xi\in M_{l,i}\) using
Theorem~\ref{thm:frobeniusjet774}. As in
Theorem~\ref{thm:secondpoint775}, its quotient differs
from the prescribed constant by \(E^A w,\ w\in R_i\).
For \(y_0=\phi^l\) of the indicated lift of \(p^\kappa\),
set \(v=p^\kappa\xi-E^A w y_0\). The adapted coordinates
of \(v\) are all regular and have order at least \(A\)
where \(t<-\eta\). Its quotient is as asserted. Its
divided displacement is integral in \(M_{l,i,A}\), is
of order \(A\) in rational horizontal coordinates, and
is therefore divisible by \(E^A\) integrally, by the
entrywise divisibility argument of that theorem.

Multiplying by \(\prod_{1\le j<l}P_j^\eta\) sends
\(M_{l,i}\) into \(M_{1,i}\) by the intermediate
lower bounds. In the resulting vector, the directions
at \(E=0\) still have the orders just obtained. After
the further factor \(E^\eta\), the order in every
direction is at least \(-t\): for \(t\ge-\eta\) this
uses regularity, and otherwise \(A\ge-t\).
After clearing its possible \(E\)-denominator in
\(M_{0,i}\), the horizontal divisibility descends over
\(\mathcal O[[E]]\) by Lemma~\ref{lem:labelledheight764}.
\end{proof}
\clearpage

\subsection{Cup products when the larger weight is sparse}
\label{sec:sparsecups780}
At an unramified spherical place use \(b,G,W\) as in
Section~\ref{sec:mainorbit776}, increasing \(W\) by a fixed
factor for the scalar targets if necessary. At a label of
weight \(G\), every jump of the first two extensions has
absolute value at most \(\eta=2(G+W+1)\).
At a label of weight \(b\), each side has short jumps in
\([-\eta,\eta]\) and at most one further jump of absolute
value in \([b-W,b+W]\). The two possible far jumps have
opposite signs. In a scalar test the dual has no far jump.
We assign its side the sign opposite to the graph's far
window even if that window is empty on the dual. Thus at
every such label the side assigned positive sign has all
jumps at least \(-\eta\), and the other has all jumps at
most \(\eta\). The third extension has upper jumps at most
\(\beta+W\) at a label of weight \(\beta\), since its
kernel has differences \(j_T-j\) at that label.

\begin{theorem}[Sparse occurrences of the larger weight on an unramified orbit]
\leavevmode\label{thm:sparsecups780}
Let \(p\ge5,\ J\ge1\). In the periodically repeated
Frobenius orbit of labels, suppose no three occurrences
of weight \(b\) have both successive distances at most
\(J\). The distances must be positive and this condition
includes repetitions around the orbit. Assume
\[
               16(J+1)p^J\eta\le b .
\]
At sufficiently large nodes, for either the main or a
scalar target, the original local cup is killed by \(p^h\)
with
\[
 h\ll_{q_0,d}
  1+\frac{b}{p^{J+1}}+(J+1)^2\eta
                  +(J+1)\frac{\log(2b)}{\log p}.
\]
There is no restriction on changes of sign among the
larger weights. No change in the finite precision of the
cup is required.
\end{theorem}
\begin{proof}
Work first with coefficients containing the local field
\(K\). Call a label a cut if the next \(J\) labels all
have weight \(G\). If there is any occurrence of weight
\(b\), either it is a cut or its next such occurrence
within \(J\) steps is a cut, by the hypothesis on triples.
If there are no occurrences, every label is a cut.
Let \(c,i\) be consecutive cuts in positive cyclic order,
and write \(i=F^l c\), \(l>0\).
If \(l>1\), the label immediately after \(c\) is not a
cut. Hence the first larger weight following \(c\) is
exactly \(J+1\) steps away. No intervening label can
be a cut, since it sees this larger weight within its
next \(J\) steps. If the first larger weight is not
itself a cut, there is a second one within \(J\) steps,
and it is a cut for the same reason as above. Thus
\(l\le L=2J+1\). The indices after \(c\) through \(i\)
contain at most two larger weights. When there are two,
the first one, denoted by \(v\), is \(F^{J+1}c=F^{-r}i\)
with \(1\le r\le J\), and the second is \(i\).

At every cut a common \(p^\kappa\) lifts integrally on
all three extensions. In fact the weighted upper-jump
sum of Theorem~\ref{thm:orbitlift764} is at most
\[
       (G+W)/(p-1)+(b+W)/(p^J(p-1)).
\]
Consequently one may take \(\kappa\ll1+\eta+b/p^{J+1}\).
Choose such points on all three extensions. The third
one shows that \(\widetilde Q=p^\kappa Q\) on the level
zero kernel lattices takes values in the image of
\(\tau^{-1}A_i(1)\). Its \(l\)-fold Frobenius from \(c\)
gives the bound \(\tau_l^{-1}A_i(1)\) on the level
\(l\) kernels at \(i\), with \(\tau_l=\phi^l(\tau)\).
One has
\[
       \tau_l\ \in\ A_i^\times\,\tau\prod_{j=0}^{l-1}P_j ,
\]
by iterating \(\mu/\tau\in A^\times E\).

Put \(D=\eta+2\) and restrict temporarily to
\(H_L=G_{K(\pi_L)}\).
At each cut use Proposition~\ref{prop:simultaneousorders779}
on both points in \(M_0\) with \(T_j=D\) for
\(0\le j<L\). Their displacements now have factor
\(\delta\prod_{j<L}P_j^D\) in the corresponding kernel.
The increase in the quotient exponent is bounded by
\[
       C\left((L+1)^2\eta+
                       (L+1)\log(2b)/\log p\right)=:C_1 .
\]
Indeed at each prime the integer indices come from at
most two intervals, the short cluster and the far window,
of total length \(O(\eta)\), with absolute values
\(\le 8b\). For a light label the whole cluster already
lies in \([-\eta,\eta]\). Summing \(v_p(|n|)\) on each
interval costs its length divided by \(p-1\) plus at
most \(\log(8b)/\log p\), as in
Theorem~\ref{thm:mainorbit776}. The same bound \(C_1\),
with an absolute enlargement of \(C\), will suffice
for all subsequent applications at \(j<l\) on \(M_l\),
where the shifted index is simply zero.

We construct one auxiliary point for each pair of cuts.
All choices below are made before scaling the points to
a common quotient. Write
\(S=\prod_{0\le j<l}P_j^\eta\) and omit the endpoint
index \(i\) on its lattices.

Suppose first that all larger weights after \(c\) through
\(i\) assign the same positive side \(U\) (choose either
if there are none). Denote the other side by \(V\). The
lower bound on \(U\) and upper bound on \(V\), respectively,
give by successive height inclusions
\[
       S M_l(U)\subseteq M_0(U),\qquad
       S M_0(V)\subseteq M_l(V).
\]
Take the \(l\)-fold Frobenius of an original integral
lift of \(p^\kappa\) on \(U\) at \(c\). In \(M_l(U)\)
impose orders \(D\) at \(0\le j<l\) by the proposition.
This gives \(y\in M_l(U)\), over cost at most \(\kappa+C_1\),
whose displacement belongs to
\(\delta\prod_{j<l}P_j^D M_{l,A}(U)\).

In the other case the two larger weights have opposite
signs. Name \(V\) the side assigned negative at their
first occurrence \(v\), hence positive at \(i\). Let
\(U\) be the other side. At \(v\), \(E^a\) lifts on
\(V\) for \(a=\lceil(b+W)/(p-3)\rceil\), by
Theorem~\ref{thm:uniformlift760} and the coefficient
pullback of Corollary~\ref{cor:actualperiodmultiples768}.
At \(c\) we have the pure lift of \(p^\kappa\). Apply
Theorem~\ref{thm:longpoint779} at \(v\), with
\(l'=J+1,\ h_0=\cdots=h_J=\eta\) and
\(A=2(b+W)+2\eta+4\). The preceding \(J\) labels are
light and all their jumps lie in \([-\eta,\eta]\).
At \(v\), each jump \(t<-\eta\) satisfies
\(-t\ge b-W\). The required strict inequality follows
from
\[
      a+\sum_{j=1}^{J}p^j\eta
      \le (b+W)/2+1+Jp^J\eta < b-W .
\]
The last inequality uses the stated size assumption and
\(W=o(b)\). In particular the degree is less than \(A\).
We get a point of \(M_{l',v}(V)\) over
\(p^{\kappa+4\lceil A/p^{J+1}\rceil}\), with
displacement divisible there by \(\delta E^A\); after
multiplication by \(\prod_{0\le j<l'}P_j^\eta\) it lies
in \(M_{0,v}(V)\).

Transport this point by \(\phi^r\) to a point \(q_0\)
of \(M_l(V)\) at \(i\). The product
\(\prod_{r\le j<l}P_j^\eta\) puts it in \(M_r(V)\).
All the lower jumps of \(V\) between \(v\) and \(i\),
including \(i\) but excluding \(v\), are at least
\(-\eta\). Thus \(S q_0\in M_0(V)\), and its displacement
in \(M_l(V)\) before multiplying by \(S\) belongs to
\(\delta P_r^A M_{l,A}(V)\).
Apply Proposition~\ref{prop:simultaneousorders779}
with \(F_0=P_r^A\) and impose \(D\) at all \(j<l\)
except \(r\). Its correction in \(P_r^A M_l(V)\) still
belongs to \(S^{-1}M_0(V)\). Indeed at \(r\) the lower
jump is at least \(-(b+W)\); at all other indices it
is at least \(-\eta\). The lower-height inclusions prove
the membership since \(A\ge b+W\). The resulting point
\(q\in M_l(V)\cap S^{-1}M_0(V)\) has displacement in
\[
            \delta P_r^A\prod_{j<l,\ j\ne r}P_j^D M_{l,A}(V).
\]
Its quotient exponent is at most
\(\kappa+4\lceil A/p^{J+1}\rceil+C_1\).

Scale the initial points at every cut and all auxiliary
points to quotient \(p^{e_*}\), taking \(e_*\) to be the
largest exponent obtained. At a noncut define both
initial points by applying Frobenius the appropriate
number of times to the pair at the preceding cut. On
the following paths a value is small if it belongs to
the image of \(uA_i(1)\) at the given precision. Initial
cups at cuts are small by the level zero pairing and
the factor \(\delta=u\tau\); this persists under Frobenius
at all other labels. Write \(u_i,v_i\) for the cut points
on \(U,V\).

In the case of common sign move \(u_i\) to \(y\), then
\(v_i\) to \(\phi^l v_c\), then \(y\) to \(\phi^l u_c\).
For the first move the difference belongs to
\(S^{-1}M_0(\ker\beta_U)\). Its poles are cleared by
the displacement of \(v_i\) in the level zero pairing.
For the second the difference is in
\(S^{-1}M_l(\ker\beta_V)\), while \(d_y\) supplies
\(\delta\prod_{j<l}P_j^D\). The formula for \(\tau_l\)
and \(D\ge\eta+1\) leave at least a factor \(u\).
The third difference is in the level \(l\) kernel, and
\(d_{\phi^l v_c}\) has factor \(u^{p^l}\tau_l\).

In the case of opposite signs move \(v_i\) to \(q\),
then \(u_i\) to \(\phi^l u_c\), then \(q\) to \(\phi^l v_c\).
The first difference lies in \(S^{-1}M_0(\ker\beta_V)\)
and is handled by \(d_{u_i}\). For the middle move the
upper jumps of \(U\) are at most \(\eta\) except at the
index \(r\), where they are at most \(b+W\). Therefore
multiplication by \(P_r^{b+W}\prod_{j<l,\ j\ne r}P_j^\eta\)
puts its difference in \(M_l(\ker\beta_U)\). The
displacement of \(q\) supplies each of these factors and,
in addition, the full product in \(\tau_l/\tau\), since
\(A\ge b+W+1\). This gives smallness. Both points in the
final move are in \(M_l(V)\), and \(d_{\phi^l u_c}\)
supplies \(u^{p^l}\tau_l\). In both cases the differences
are in the kernels because the quotient constants agree.
All bounds hold for translates on \(H_L\).

The final pair at a cut is the single Frobenius of the
actual initial pair at its preceding label. At a noncut
this comparison is the identity. Summing the cochains
over the labels and taking Frobenius constants as in
Theorems~\ref{thm:labelledcup765} and~\ref{thm:mainorbit776}
shows that the restricted cup is killed by
\(p^{2e_*+\kappa}\). Corestriction from \(K(\pi_L)\),
of degree \(p^L\) by Eisenstein's criterion, adds \(L\)
to the exponent. The coefficient retraction descends
the original finite modulus as before. Since \(A\ll b\),
the bounds for \(\kappa,C_1,e_*,L\) prove the result.
\end{proof}
\clearpage

\subsection{Pure lifts and Taylor orders at ramified places}
\label{sec:ramifiedorders782}
Let \(K/\QQ_p\) be finite Galois, with maximal unramified
subfield \(K_0\), ramification index \(e\), and uniformizer
\(\varpi\). Coefficients contain all images of \(K\).
Use \(u=[\underline\varpi]\), its Eisenstein polynomial
\(E(u)\) over \(K_0\), and the normalized crystalline
lattices of Lemma~\ref{lem:adjacentlattice762}. Now the
idempotents \(i\) correspond to embeddings of \(K_0\).
At a fixed \(i\), its ring is \(R_i=\mathcal O[[u]]\);
different extensions of \(i\) to \(K\) can carry different
Hodge jumps. Let \(h_i,r_i\) be nonnegative integers
bounding these jumps above by \(h_i\) and below by \(-r_i\)
at \emph{all} those embeddings. With components understood,
write \(P_j=\phi^j(E)\), \(\delta=u\tau\), and set
\(d_K=\lceil v_p(E'(\varpi))\rceil\le1+v_p(e)\).
We use all primes \(p\), including two.

\begin{lemma}[Height and horizontal coordinates over the residue-field labels]
\leavevmode\label{lem:ramifiedcoordinates782}
At each label \(i\),
\[
        E^{h_i}M_{0,i}\subseteq M_{1,i}
                    \subseteq E^{-r_i}M_{0,i}.
\]
At \(P_j\) the rational horizontal coordinates of
\(M_{s,i}\) are regular if \(s>j\). If \(s\le j\), on
each rational factor the shifts are the negatives of
the jumps at the corresponding embedding of \(K\) above
\(F^{-j}i\). The Taylor constants in these coordinates
are fixed by \(G_{K(\varpi^{1/p^L})}\) for \(j\le L\).
\end{lemma}
\begin{proof}
First work at \(E\). On the indicated coefficient
component \(E\) is separable over \(C=\operatorname{Frac}
(\mathcal O)\); its roots are the \(e\) extensions of
the chosen embedding applied to \(\varpi\). The rational
\(E\)-adic completion is the product of the power series
rings for those roots, each with parameter \(E\).
In each factor the analytic comparison of
Lemma~\ref{lem:horizontalcoordinates767} identifies
\(M_1\) with the regular horizontal lattice. In fact the
zeros of \(\phi^t(E)\), \(t>0\), have valuation
\(1/(ep^t)\), different from \(1/e\); hence
\(\phi(\lambda)=\prod_{t\ge1}\phi^t(E/E(0))\) is a unit
at these roots. For convergence there, each normalized
factor tends to one geometrically in the exponent \(p^t\),
since \(E(0)\) has valuation one. The constant space is
the crystalline period space after extension from \(K_0\)
to \(K\). This space identifies with the filtered de Rham
periods since the representation is crystalline.
After choosing a basis adapted to the filtration at the
given embedding, the level zero \(B_{\rm dR}^+\)-lattice
has shifts \(-t\) for its Hodge jumps \(t\), by filtered
comparison. As in Lemma~\ref{lem:horizontalcoordinates767},
the level zero comparison uses that \(\tau\) specializes
to \(\zeta_p-1\ne0\). The equality of rational lattices
descends from the period ring factorwise by faithful
flatness over a discrete valuation ring, since \(E\)
is a uniformizer there. These comparisons, including a
common twist and untwist, are those of the rational
Breuil--Kisin comparison.~\cite[Section~2.1]{III-II-bartlettcrys767}

The two transition matrices are in \(R_i[1/E]\) by
Lemma~\ref{lem:adjacentlattice762}. If an element of
\(R_i\) vanishes modulo \(E^n\) at all rational factors,
it already vanishes in \(R_i/E^n\): Weierstrass division
by the distinguished monic polynomial \(E^n\) makes
this quotient finite free over \(\mathcal O\). The
rational shifts thus give the two integral inclusions.
Freeness of \(M_{s,i}\) is proved as in
Lemma~\ref{lem:labelledheight764}, decomposing the
residue-field coefficient ring and then applying
Frobenius. At \(P_j\) all transitions except the one
from level \(j\) to \(j+1\) are regular isomorphisms
rationally; the root valuations \(1/(ep^v)\) are
different for different \(v\). Transport the description
at level zero and one by \(\phi^j\).

More explicitly, \(P_j(u)\) is separable over \(C\).
At each finite order, division by \(P_j^n\) and Hensel
lifting the simple roots identify the rational Taylor
algebra with \((C[T]/P_j(T))[P_j]/P_j^n\): \(T\) lifts
the residue of \(u\) as an exact root of \(P_j(T)=0\),
and \(u=T+O(P_j)\) is recovered successively using
\(P'_j(T)\) invertible. In the internal period factor
the root reducing to \(\varpi^{1/p^j}\) is unique.
The indicated subgroup fixes that root by uniqueness,
and fixes \(K\) and the external coefficient field.
Thus it fixes the Taylor constants, including those used
to adapt the de Rham filtration separately in each
factor. It acts on the parameter by substitution of
\(gP_j\). This gives the last assertion.
\end{proof}

\begin{theorem}[A pure quotient lift after \(J\) bounded labels]
\leavevmode\label{thm:ramifiedcut782}
For a crystalline surjection \(U\to\mathcal O\) with
trivial quotient suppose \(h_i\le H\) everywhere.
Suppose \(h_{F^t i}\le\eta\) for \(1\le t\le J\), with
\(J\ge1\) and \(p^{J+1}\ge8e\). There is a lift of
\(p^\kappa\) in \(M_{0,i}\) with
\[
 \kappa\ll 1+J(d_K+1)\eta+\frac{eH}{p^{J+1}} .
\]
The implied constant is absolute.
\end{theorem}
\begin{proof}
Set \(I_j=\beta(M_{0,j})\subset R_j\). At precision
\(p^\nu\) let \(a_j\) be the least exponent of a pure
power of \(u\) in \(I_j+(p^\nu)\). It exists by the
residue-field and Witt comparison argument in
Theorem~\ref{thm:orbitlift764}. That argument uses only
freeness, nonzero images on each residue-field component
after reduction, and \(R_j/\pi\) a discrete valuation
ring, all of which remain valid here. By the height
lemma \(E^{h_{Fj}}I_{Fj}\subseteq R_{Fj}\phi(I_j)\).
Write \(u^e=E+pB(u)\) with integral \(B\). Modulo
\(p^\nu\), \(u^{e(h+\nu-1)}\) is a multiple of \(E^h\)
by binomial expansion. Decomposing a power series into
its \(p\) residue classes of exponents gives
\[
           p a_j\le a_{Fj}+e(h_{Fj}+\nu-1).
\]
Taking the maximum around the finite cycle yields
\(\max a_j\le e(H+\nu-1)/(p-1)\).

Starting at \(F^J i\), make \(J\) passages back to
\(i\) by Lemma~\ref{lem:transportpower761}. To verify
the labelled version, on a step from \(Fj\) to \(j\)
the image in the Frobenius lattice is \(R_{Fj}\phi(I_j)\).
Multiply the current power by \(E^{h_{Fj}}\) using the
height lemma. The divided displacements and descent
modulo \(p^\nu\) in that proof act on this labelled
ideal, because \(G_K\) fixes the residue-field constants
and the coefficient factors. For each successive exponent
\(t\) of \(E\), they remove one \(E\), add one \(u\),
and cost \(d_K+v_p(t)\). The residue-class decomposition
on extraction of \(\phi\) then replaces the parameter
exponent \(a\) by \(\lfloor(a+h_{Fj})/p\rfloor\).
Therefore at label \(i\) we have \(p^{c}u^{a_*}\)
in \(I_i+(p^\nu)\), with
\[
 c\le J(d_K+1)\eta,\qquad
 a_*\le\frac{e(H+\nu-1)}{p^J(p-1)}+\frac{\eta}{p-1}.
\]
The remaining parameter powers cost
\(a_*+v_p(a_*!)\le p a_* /(p-1)\) by
Lemma~\ref{lem:kummerlower760}, likewise on this component.
Let \(\nu=C+1\). The coefficient of \(H+C\) in the last
cost is \(pe/(p^J(p-1)^2)\le 4e/p^{J+1}\le1/2\).
Taking \(C=\lceil2(J(d_K+1)\eta+2\eta+4eH/p^{J+1}+1)\rceil\)
makes the total cost at most \(C\), including for \(p=2\).
Thus \(p^C\in I_i+(p^{C+1})\), and division by an element
of \(1+pR_i\) gives \(p^C\in I_i\).
\end{proof}

\begin{proposition}[Taylor orders with the different included]
\leavevmode\label{prop:ramifiedorders782}
The assertion of Proposition~\ref{prop:simultaneousorders779}
holds over the present \(K\) at every prime, with
\(H_L=G_{K(\varpi^{1/p^L})}\). For \(j\ge s\) use for
\(\mathcal D_j\) the union of the jumps above \(F^{-j}i\);
for \(j<s\) still use \(\{0\}\). In its cost replace
\(L+2+v_p(|n|)\) by \(L+3+d_K+v_p(|n|)\).
\end{proposition}
\begin{proof}
There exists \(g_0\in H_L\) with \(g_0u=[\epsilon]^{p^L}u\).
Indeed \(X^{p^a}-\varpi\) is Eisenstein over \(K\).
The embeddings on its root fields sending
\(\varpi^{1/p^a}\) to
\(\zeta_{p^a}^{p^L}\varpi^{1/p^a}\) are compatible under
the \(p\)-th power maps. Extend their union to the
algebraic closure; the resulting automorphism fixes the
chosen root at level \(L\). For \(\omega=g_0E/E\), the
polynomial divided-difference formula in
Lemma~\ref{lem:transportpower761} gives
\[
 v_p(\theta(\omega-1))=
  L+1/e+1/(p-1)+v_p(E'(\varpi)).
\]
To see the extra \(L\), write the displacement of \(u\)
as \(u\mu\sum_{h=0}^{p^L-1}[\epsilon]^h\). The sum has
specialization \(p^L\), and the divided difference of
\(E\) specializes to \(E'(\varpi)\). As
\(\mu/(E\tau)\) is a unit, \(gE/E\in1+\delta A\) for
every \(g\in G_K\). Also
\(v_p(\theta_j(\delta))=(1/e+1/(p-1))/p^j\).
It follows that
\(\big((g_0P_j/P_j)^n-1\big)/\delta\) belongs to \(A\)
and its nonzero \(\theta_j\)-value has valuation at
most \(L+3+d_K+v_p(|n|)\). The binomial valuation rule
applies at two as well: the displayed input valuation
strictly exceeds \(1/(p-1)\). Surjectivity of \(\theta_j\)
gives the integral multiplier of the divided difference
with the prescribed power of \(p\) as specialization.

Use the retraction in an \(R_i\)-basis, whose construction
in Lemma~\ref{lem:integraljet767} uses faithful flatness
and compactness, valid also for this Eisenstein parameter.
To check its Taylor property on several rational factors,
clear poles by a common \(P_j\)-power. The finitely many
component congruences are encoded in \(R_i[1/p]/P_j^n\)
by the finite Taylor isomorphism in the height lemma,
using the product idempotents there. Clear one power of
\(p\) to get integral linear congruences. The map
\(A_i/P_j^n\) into the corresponding rational period
quotient is injective as before: internally
\(A/E=\mathcal O_{\mathbb C_p}\), the successive quotients
for \(E^n\) are \(p\)-torsion free, and the rational map
is the de Rham completion. Transport by \(\phi^j\) and
pass to the finite free coefficients. Retraction thus
preserves these congruences, with no precision loss on
the original lattice.

Now use exactly the degree-raising operators in
Proposition~\ref{prop:simultaneousorders779}, with the
multiplier just constructed. The horizontal comparison
in the height lemma gives the intervals to process,
and the same operator at an integer \(n\) works on all
factors simultaneously. Every earlier product of
\(P_v\)'s is preserved, since \(gP_v/P_v\in1+\delta A\).
The image of another \(P_v\) modulo \(P_j\) is a
nonzero divisor: its internal specialization is nonzero
by the different root valuations, and the assertion
survives coefficient extension and taking an idempotent.
Therefore intersecting the divisibilities multiplies
their ideals. This proves both conclusions of the
proposition at the stated cost.
\end{proof}
\clearpage

\section{Higher-rank local tests at spherical and ramified primes}\label{nav:higher-rank-tests}

\subsection{The local conditions for the higher-rank summand}
\label{sec:highlocal788}
Let \(R_\xi\) be the rank \(s=r_x\) summand of a true
system supplied by Theorem~\ref{thm:longcharacter787}.
Use \(t=5s,\ k=s+1-m_x,\ h=s^8,\ L=(t-1)/2\).
At an infinite place \(\tau\), write
\(B_\tau\in\{b_x,a_x/x_{1,x}\}\) for the base weight
and \(w_j=B_\tau+(2j+1)h\).
The notation \(\sigma_p\), \(0\le p<s\), is the old row
at this place in Lemma~\ref{lem:archstring787}. All
Hodge--Tate weights here use weight \(-1\) for \(\kappa\).

\begin{proposition}[Labelled weights and exact widths]
\leavevmode\label{prop:highwidth788}
At the two conjugate labels over \(\tau\), the Hodge--Tate
multisets of \(R_\xi\), in one order, are
\[
  \{L+(\sigma_p+w_j):0\le p<s\},\qquad
  \{L-(\sigma_p+w_j):0\le p<s\}.
\]
For \(0\le a<t\), the multiset for \(u_\xi\otimes\kappa^a\)
contains zero. At either label its exact width
(largest minus smallest weight) is
\[
 W_{j,\tau}=
 \begin{cases}
 B_\tau+2m_xh+k/2-1,&j<m_x,\\
 B_\tau+(2m_x+1)h+k/2-1+L,&j=m_x .
 \end{cases}
\]
All normalized target systems at the same weight, including
those with additional finite auxiliary level, have the same
labelled multiset as \(u_\xi\) before applying the common
Tate twist.
\end{proposition}
\begin{proof}
We explain the local compatibility needed at the coefficient
prime. Apply the construction of
Theorem~\ref{thm:splitattachment722} block by block to the
cohomological base change of the definite target. If
\((\rho,b)\) has cuspidal degree \(d\), then for a constituent
with norm exponent \(z=(b-1)/2-i\) one uses
\(\rho|\det|^{z+(d-n)/2}\) in the rank \(d\) Galois
normalization. Its holomorphic infinitesimal entries, before
the shift by \((d-n)/2\), form a sublist of the regular
target entries of parity \((n-1)/2\bmod\ZZ\). At the
conjugate labels the same is true. Thus this twist is
regular algebraic of the rank \(d\) parity. Its cuspidal
factor \(\rho\) is conjugate self-dual by the cohomological
base change. After the norm twist it is polarized by
a norm power on \(E_x^+\): the required exponent is
\(2z+d-n\). Multiplying the polarizing character by
the quadratic character of \(E_x/E_x^+\), if needed,
sets its prescribed common real sign without changing
its base change. The Galois theorem for regular algebraic
cuspidal polarized pairs now gives de Rham representations
at every place over \(\ell\), with the labelled weights
of these algebraic parameters. It gives crystalline
representations at such places where the local cuspidal
representation is unramified. In degree one this is
the corresponding statement of local class field theory
for algebraic Hecke characters. The good split
Frobenius polynomials identify these semisimple
representations with the ones used in
Theorem~\ref{thm:splitattachment722}.
\cite[Theorem~2.1.1]{III-II-blggt788}

For clarity about signs and translations, the labelled
weight formula takes the infinitesimal entries at a fixed
label to their negatives or to themselves according to the
dual convention, with a common shift by the half sum
normalization. In fact for decreasing highest weight
\((a_i)\) the two lists in the usual dual convention are
formed from \(a_i+(d+1)/2-i\) and \(a_i+d-i\):
the difference is constant. Dualizing reverses and
negates the list. The shift \((d-n)/2\) above makes the
translation common to the whole rank \(n\) list, and
the normalizing character adds a single labelled weight.
The character block of length \(t\) has the known weights
\(0,\ldots,t-1\), so its center after all shifts is \(L\).
The old entries in Lemma~\ref{lem:archstring787} are
at displacements \(\sigma_p+w_j\) from that center
before the possible reversal of sign. Removing the
character block gives one of the two displayed multisets.
At the conjugate label the sign reverses. Indeed
\(u_\xi^c\simeq u_\xi^\vee\otimes\kappa^{1-t}\) by
Corollary~\ref{cor:normalizedtransfer722}; the
character block is stable under this operation, and
semisimplicity permits cancellation. The paired
weights are therefore \(d'\) and \(t-1-d'\).
The target calculation used only its infinitesimal
character and the common normalizer, so it applies
also to the other systems of its weight.

It remains to take the width of the union of
\([-L,L]\cap\ZZ\) and the displacements \(\sigma_p+w_j\).
Reflection and translation preserve this width.
The largest displacement is \(w_j+k/2-1>L\).
The smallest old displacement is \(h+w_j-w_{m_x}
=(1+2(j-m_x))h\). For \(j<m_x\) it lies below \(-L\);
for \(j=m_x\) the smallest entry of the union is \(-L\).
Subtraction gives the two formulas. Translating the
full multiset by \(-a\) retains zero because it
initially contains all of \(0,\ldots,t-1\).
\end{proof}

\begin{theorem}[The corrected finite local lengths at every coefficient prime]
\leavevmode\label{thm:highfinite788}
Fix \(0\le a<t\) and any stable lattice \(T_a\) in
\(R_\xi\otimes\kappa^a\), over a sufficiently large
finite coefficient field with valuation ring \(\mathcal O\)
and uniformizer \(\pi\). For \(v\mid\ell\) and \(N\ge1\)
let \(\mathcal F_{v,N}\) be the reduction in
\(H^1(E_{x,v},T_a/\pi^N)\) of the integral classes whose
rationalizations lie in the Bloch--Kato finite local
condition. Then
\[
 \sum_{v\mid\ell}
 \bigl(\operatorname{length}_{\mathcal O}\mathcal F_{v,N}
    -\operatorname{length}_{\mathcal O}
                 H^0(E_{x,v},T_a/\pi^N)\bigr)
                  =N f_x s .
\]
This holds at every node and every permitted twist on
the period interval.
\end{theorem}
\begin{proof}
All old displacements have absolute value at least
\(h>L\), as in the proof of Lemma~\ref{lem:archstring787}.
Consequently \(L-a+(\sigma_p+w_j)\) and
\(L-a-(\sigma_p+w_j)\) have opposite strict signs:
\(|L-a|\le L\). Each conjugate pair of embeddings
therefore contributes exactly \(s\) negative weights.
For a de Rham representation the sum of the dimensions
of \(D_{\rm dR}/\mathrm{Fil}^0\) over \(v\mid\ell\)
counts these negative weights over all labels, with
multiplicity. Each embedding induces one place over
\(\ell\), and there are \(f_x\) conjugate pairs.
By Lemma~\ref{lem:localevaluation728}, at each place
the difference of the two integral lengths on the
left is \(N\) times its tangent dimension. This proves
the identity, including at primes of ramification of
the field or of the tame level.
\end{proof}

\begin{theorem}[Vanishing of the integral lifting defect on short split orbits]
\leavevmode\label{thm:highlift788}
Take a place \(v\mid\ell>2\) of \(E_x\) unramified
over \(\QQ_\ell\), split over \(E_x^+\), and spherical
for the chosen level and splitting character. The
auxiliary levels in a trace relation are allowed to
vary away from \(v\) and its conjugate. Suppose
\[
 W_{j,\tau}\le\ell-1\quad
       \text{for all labels inducing }v,\qquad
 W_{j,\tau_0}\le\ell-2\quad\text{for at least one.}
\]
For the unshifted trace relations take as target either
a lattice in \(R_\xi\) or one in a prescribed cyclotomic
summand of \(u_\xi\). More generally one can twist all
the representations in the local relation by
\(\kappa^a,\ 0\le a<t\), with the target chosen among
the corresponding summands. At every finite precision
every cocycle of Lemma~\ref{lem:graphrealization724} for
these data lifts to an integral cocycle with crystalline
rational extension. Its class in the quotient by
\(\mathcal F_{v,N}\) for that target is zero.
\end{theorem}
\begin{proof}
At a split place the cohomological base change matches
the local \({\rm GL}_n\) representation. A spherical
parameter has trivial inertia and zero monodromy, and
the local parameter of the isobaric sum is the direct
sum of its twisted cuspidal parameters. Hence every
one of those cuspidal factors is unramified at \(v\).
The crystalline clause used in
Proposition~\ref{prop:highwidth788} applies, including
at the additional levels which avoid this place. The
avatar of the locally unramified algebraic normalizing
character is crystalline there, as are Tate twists.
Taking a direct summand preserves crystallinity.

At each label inducing \(v\), take the interval from
the smallest to the largest weight of the full
normalized and commonly twisted target representation.
It contains zero and all weights of every lattice in
the relation, since by Proposition~\ref{prop:highwidth788}
the target summand is a subrepresentation and the
other full target systems have the same weights.
Lemma~\ref{lem:graphrealization724} realizes the
cocycle extension as a subquotient at the given
precision of lattices with these rational constituents
and the trivial line. The hypotheses on the widths
are precisely those of the labelled integral
Fontaine--Laffaille lifting theorem,
Theorem~\ref{thm:labelledlift735}. Applying it to
this extension gives a lift with the stated end
identifications and the stated cohomology class.
\end{proof}
For example, away from the old level and field
discriminant, a prime split in \(\QQ(\sqrt{-d})\)
has all its places split over \(E_x^+\). A decomposition
group there fixes that quadratic subfield and so cannot
contain complex conjugation. The levels chosen away from
\(\ell\) introduce no further exception above it. Thus
the width bound with \(B_\tau\) replaced by the larger
of the two base weights suffices simultaneously at all
these places. On an orbit containing only the smaller
base weight one may use that smaller value. The exact
finite-condition length above and this integral lifting
result have different ranges. The congruence-class
loss for long or ramified orbits still requires an
integral estimate.
\clearpage

\subsection{The higher-rank trace defects at small unramified primes}
\label{sec:highsmall789}
Use the sparse choice of infinite places of
Lemma~\ref{lem:spacing781} on the original period field
\(E_x\). Put \(s=r_x,\ t=5s,\ Q_x=k_x t,\ b=b_x\)
and \(G=xb/x_{1,x}\). At a coefficient prime \(p\)
use the selected twist and any of the auxiliary levels
away from \(p\) in the trace relations. The targets
\(T_i\) in this section consist of a stable lattice
in \(R_\xi\) and the \(t\) scalar lattices
\(\mathcal O(-a)\), \(0\le a<t\), from
Theorem~\ref{thm:longcharacter787}.

\begin{lemma}[Crystallinity at the spherical coefficient places]
\leavevmode\label{lem:inertcrys789}
At a place \(v\mid p\) where the level and the smooth
part of the normalizing character are unramified, the
normalized definite target systems are crystalline.
This includes the spherical places inert over \(E_x^+\).
In particular Theorem~\ref{thm:highlift788} also holds
at these places when \(E_{x,v}/\QQ_p\) is unramified
and its stated width conditions hold.
\end{lemma}
\begin{proof}
The identity Hermitian space is quasi-split at the
nonsplit finite places on these fields, as checked in
Theorem~\ref{thm:splitattachment722}. The extension at
such a place is unramified quadratic and the level
here is hyperspecial. Labesse's Corollary~5.3 identifies
the base change locally at unramified as well as at
split places, for the definite cohomological space
used there. Thus a hyperspecial constituent has
unramified standard base change at \(v\).
\cite[Corollaire~5.3]{III-II-labessesph789}
Its local parameter is the direct sum of those of
the norm twists of the cuspidal factors. Trivial
inertia and zero monodromy on a direct sum hold on
each summand. The local cuspidal representations
are therefore unramified, and the crystalline clause
of Proposition~\ref{prop:highwidth788} applies at \(p\).
The locally unramified algebraic normalizer has
crystalline avatar. This supplies exactly the
crystallinity used in Theorem~\ref{thm:highlift788};
its remaining hypotheses and proof use only the local
field and the labelled widths.
\end{proof}

For a spherical \(v\mid p\), write \(K=E_{x,v}\).
For a precision \(\pi^N\) in a coefficient valuation
ring \(\mathcal O\), let \(C_i\subset H^1(K,T_i/\pi^N)\)
be the image of the cocycles satisfying the matrix
relations of Lemma~\ref{lem:graphrealization724} for
the given automorphic family. Let \(F_i\) be the
reduction of the integral Bloch--Kato finite condition
for \(T_i\). Define the nonnegative length
\[
       d_{v,i}=\frac{1}{v_\pi(p)}
             \operatorname{length}_{\mathcal O}\!
                   \left(C_i/(C_i\cap F_i)\right).
\]

\begin{proposition}[Cup products with several large jumps at one label]
\leavevmode\label{prop:highsparse789}
For \(5\le p\le Q_x\), \(p\nmid q_1d\), and spherical
\(v\mid p\), every cup of a class in \(C_i\) with the
reduction of an integral finite-condition class for
\(T_i^\vee(1)\) is annihilated by \(p^H\) with
\[
                    H\ll_{q_0,d} b x^{-120}.
\]
This bound is uniform in \(i,\xi,j\), in the coefficient
field and embedding, and in the auxiliary levels and
precision.
\end{proposition}
\begin{proof}
The graph construction of
Lemma~\ref{lem:graphrealization724} gives a lattice
extension of the trivial line with kernel \(Y\)
and an equivariant map \(Y\to T_i/\pi^N\) carrying
its cocycle to the specified one. The graph extension
splits rationally, and its rational kernel is a direct
sum of copies of the given systems. The entire extension
is therefore crystalline by Lemma~\ref{lem:inertcrys789}.
Use the integral
dual-test extension for the second side. For the
contraction, lift the map on \(Y\) to a linear map
\(\psi:Y\to T_i\). The cocycle
\((g\psi-\psi)/\pi^N\) defines a rationally split
extension by \(\operatorname{Hom}_{\mathcal O}(Y,T_i)\).
Evaluation after multiplication by \(\pi^N\) and
addition of the quotient coordinate times \(\psi\)
recovers the original pairing modulo \(\pi^N\).
This is the construction of
Corollary~\ref{cor:periodlabels764} with the full
target lattice in the Hom space.

We verify the estimates on jumps that enter the
Frobenius argument of Theorem~\ref{thm:sparsecups780}.
At a label of base weight \(\beta\in\{G,b\}\),
Proposition~\ref{prop:highwidth788} puts all the
graph jumps into two intervals. One is \([-W,W]\).
The other is \(\epsilon\beta+[-W,W]\) for a single
sign \(\epsilon\) at that label. Here
\(W=\lceil C x^{42}\rceil\) works on the whole
interval since the shifts \((2j+1)s^8\), \((2m_x+1)s^8\),
\(s\), and \(t\) are all \(O(x^{41})\).
The dual jumps are in the short interval and in the
window of the opposite sign, by duality followed by
the shift \(-1\). The far window is empty on the
dual for a scalar target. In the Hom space every
jump is a difference of two jumps at the same label.
Two far terms there have the same sign before taking
their difference. Thus increasing \(C\) once puts
all upper jumps of the third extension below
\(\beta+W\). These assertions allow every one of the
\(k_x-1\) old entries in the far cluster.

Here are the details that permit these intervals in
the proof of Theorem~\ref{thm:sparsecups780}, which
was stated for the torus data. Put
\(\eta=2(G+W+1)\) and
\(J=\lfloor120\log x/\log p\rfloor\).
Use the cuts followed by \(J\) light labels. By
Lemma~\ref{lem:spacing781} no three heavy labels
have two successive positive gaps at most \(J\),
since \(2J<l_0\). Thus consecutive cuts are at distance
at most \(2J+1\). At most two heavy labels occur after
the initial cut and through the terminal cut. The upper jumps
of each of the three extensions give a common
integral lift of \(p^\kappa\) at each cut with
\(\kappa\ll1+\eta+b/p^{J+1}\), by
Theorem~\ref{thm:orbitlift764}. On a light label all
jumps on the first two sides are in \([-\eta,\eta]\).
On a heavy label the side whose far window is positive
has lower jumps at least \(-\eta\), and the opposite
side has upper jumps at most \(\eta\).

The simultaneous order corrections of
Proposition~\ref{prop:simultaneousorders779} use the
\emph{union} of the shifted integer intervals of
length \(\eta+2\) over the jumps. Each of our two
clusters gives a single containing interval of length
\(O(\eta)\), regardless of the number of repeated or
distinct jumps in it. Their endpoints have absolute
value at most \(8b\). The sum of \(v_p(|u|)\) over
the nonzero integers in either interval of length
\(M\) is \(O(M/(p-1)+\log(8b)/\log p)\), by counting
multiples of \(p^a\). Thus the correction over a
path still costs
\[
        O\big((J+1)^2\eta+(J+1)\log(2b)/\log p\big).
\]
For paths with a common heavy sign the two one-sided
height inclusions in Theorem~\ref{thm:sparsecups780}
use only the lower and upper bounds just verified.
For opposite signs the negative jumps \(u<-\eta\)
at the first heavy label all satisfy \(-u\ge b-W\).
With \(a'=\lceil(b+W)/(p-3)\rceil\) one has
\[
       a'+\sum_{i=1}^Jp^i\eta<b-W
              \quad\text{if }16(J+1)p^J\eta\le b
\]
at large \(x\). The degree test in
Theorem~\ref{thm:longpoint779} therefore holds for
every negative direction, taking Taylor order
\(2(b+W)+2\eta+4\). That theorem incurs the cost
\(O(1+b/p^{J+1})\) once for the whole vector.
These are all uses of the jumps in constructing
the paths. The three changes of points and the
contraction estimates in Theorem~\ref{thm:sparsecups780}
now apply to the same lattices and precision. Their
kernel inclusions depend only on the two one-sided
bounds and the bound for the third extension. Taking
Frobenius constants and corestricting from an extension
of degree at most \(p^{2J+1}\) gives
\[
     H\ll 1+b/p^{J+1}+(J+1)^2\eta
                                  +(J+1)\log(2b)/\log p .
\]
All operations can take place after extending the
coefficient ring to contain \(K\), since an integral
coefficient retraction fixing one descends the same
annihilation at the original precision.

Finally \(p^J\le x^{120}<p^{J+1}\), \(J\ll\log(2x)\),
\(G/b\le x^{1-10^6}\) and \(b\ge x^{400}\).
These inequalities imply the degree condition and bound
the last display by \(O_{q_0,d}(b x^{-120})\), including
the logarithmic term. They hold for the same choices
that enter Theorem~\ref{thm:fixedconductor781}.
\end{proof}

\begin{theorem}[A small-prime sum for the actual higher-rank local quotients]
\leavevmode\label{thm:highsmall789}
On the same fields and interval, at each sufficiently
large node one has
\[
 \sum_{\substack{p\le Q_x\\p\nmid6q_0d}}(\log p)
   \sup\sum_{\substack{v\mid p\\v\ {\rm spherical}}}
                     \sum_i d_{v,i}
       \ll_{q_0,d} f_x(s+t)b_x x^{-110}.
\]
The supremum includes the higher-rank true system,
its simultaneous twist, the auxiliary systems at that
level, the embeddings and all precisions. After
division by \(sB_x^{\rm det}\) the right hand side
is \(O_{q_0,d}(x^{-105})\).
\end{theorem}
\begin{proof}
The integral finite conditions for \(T_i\) and
\(T_i^\vee(1)\), reduced at precision \(\pi^N\),
are exact annihilators under finite local Tate
duality. This is the order and pairing calculation
in Proposition~\ref{prop:freetests741}, including
the integral torsion classes. Thus \(C_i/(C_i\cap F_i)\)
injects by cup products into a product of copies of
\(\mathcal O/\pi^N\), using generators of the dual
finite condition. One can take the images of a basis
of its integral finite-condition raw cocycles. Their
rank is the target rank \(d_i\) plus the de Rham
tangent dimension of the dual, by
Lemma~\ref{lem:localevaluation728}, and is at most
\((1+[K:\QQ_p])d_i\). Applying
Proposition~\ref{prop:highsparse789} to these integral
representatives handles every generator, including
those defining torsion cohomology classes. The
\(p^H\)-torsion in one copy of \(\mathcal O/\pi^N\)
has normalized length at most \(H\). Summing the
coordinate bounds gives
\[
        \sum_{v,i} d_{v,i}
          \ll_{q_0,d} 4 f_x(s+t)b_x x^{-120},
\]
since \(\sum_i d_i=s+t\) and the local degrees over
\(p\) sum to \(2f_x\). All prime divisors of \(q_1\)
divide \(6q_0\). The hypotheses of the proposition
therefore hold at the spherical places retained here.
The elementary estimate \(\sum_{p\le Q_x}\log p\ll
Q_x\ll x^{10}\) proves the prime sum, also after the
indicated supremum. Finally \(s+t=6s\) and
\(B_x^{\rm det}=f_xxb_x(\log b_x)^3/(m_x k_x)\),
with \(m_x k_x\ll x^{11/2}\). Substitution proves
the normalized estimate.
\end{proof}
This controls the loss of the finite local condition
for all the listed summands and primes. The determinant
comparison still requires the other primes and the
passage from these local quotients to the full tested
Hecke lattice.
\clearpage

\subsection{The higher-rank local quotients at the fixed conductor}
\label{sec:highconductor790}
Put \(\mathcal R=\{p:p\mid 6q_0d\}\).
We use the fields, the sparse set of infinite places, and the
local conditions \(C_i,F_i,d_{v,i}\) of
Section~\ref{sec:highsmall789}. The ramification index of a
completion above a fixed \(p\) can grow along this family.

\begin{lemma}[The sign of the large weights above a conductor prime]
\leavevmode\label{lem:highsign790}
Fix \(p\in\mathcal R\), a place \(v\mid p\), a coefficient
embedding, and a lift index \(j\). Put \(K=E_{x,v}\) and let
\(K_0\) be its maximal unramified subfield. At every full
embedding of \(K\) having base weight \(b=b_x\), the large
weights of all the graph representations in the local
trace relation lie in \(\epsilon b+[-W,W]\) for a common
sign \(\epsilon\) and a common integer \(W\ll x^{42}\).
At these embeddings the large weights of the dual test for
\(R_\xi\) lie in the opposite window. All other weights
of either representation at these embeddings lie in
\([-W,W]\), and a scalar dual has only these latter weights.
The statement holds after a common twist by \(\kappa^a\),
\(0\le a<t\), and replacement of the dual by the dual of the
twisted target with its Tate twist.
\end{lemma}
\begin{proof}
Identify the complex labels using the specified coefficient
embedding. The field has CM type \(\Phi\) induced by a fixed
embedding of \(K_d=\QQ(\sqrt{-d})\).
The labelling in Proposition~\ref{prop:highwidth788} can be
tracked simultaneously at all infinite places. Indeed, at
each real place \(\tau\) the complexified unitary group is
identified with \({\rm GL}_n\) by the coordinate indexed by
\(\Phi(\tau)\). The two coordinates of the complexified
unitary group are related by inverse transpose, up to the
change of basis of the Hermitian form. Consequently passage
to the other coordinate negates the infinitesimal entries
and reverses their order. The prescribed positive columns
act on exterior powers of the \(\Phi(\tau)\) coordinate,
and the negative columns act on their duals. Taking the
successive differences in this order is exactly the calculation
of Lemma~\ref{lem:archstring787} at every \(\Phi(\tau)\).

When passing to the labelled Galois weights the common
translation places the center of the string at \(L=(t-1)/2\).
Taking a dual, if it is used in the Galois normalization,
negates the displacements from its center at every label.
It follows that for one fixed sign \(\delta\) the old
displacements at \(\sigma\) are
\[
    \delta\,\epsilon_\Phi(\sigma)(\sigma_h+w_j),
       \qquad 0\le h<s ,
\]
where \(\epsilon_\Phi\) is \(1\) on \(\Phi\) and \(-1\) on its
conjugate, and \(\sigma_h\) is the row entry at the underlying
infinite place. The normalized scalar entries are \(0,\ldots,t-1\).
For \(h<k_x-1\), \(\sigma_h+w_j\) equals the base weight plus
an error \(O(x^{41})\). The other displacements are
\((1+2(j-(h+2-k_x)))s^8=O(x^{41})\), since
\(s\ll x^5\) and \(m_x\ll\sqrt{x}\).
The same highest weights occur at each auxiliary level
in the local graph. A common Tate twist changes each
entry by at most \(t\). Duality followed by the Tate
twist changes a weight \(u\) into \(-1-u\).

The field calculation in Lemma~\ref{lem:conductorsign783}
shows that all the full embeddings of this completion
carrying \(b\) induce the same CM orientation. For primes
dividing \(q_0\) this follows from splitting in \(K_d\).
At \(p=d\) the ratios of images of a \(q_1\)-th root of
unity belong to \(\langle d\rangle\bmod q_1\), and its
image modulo \(q_0\) omits \(-1\). These statements also
hold for the translated set of large weights. Multiplying
the common orientation by \(\delta\) gives the asserted
sign. Enlarging \(W\) absorbs the center, the twists and
the bounded displacements.
\end{proof}

\begin{proposition}[A conductor bound uniform in rank and precision]
\leavevmode\label{prop:highram790}
At every spherical \(v\mid p\), \(p\in\mathcal R\), a cup of a
class in \(C_i\) with a reduced integral finite-condition
class for \(T_i^\vee(1)\) is annihilated by \(p^H\), with
\[
                        H\ll_{q_0,d} b_x x^{-108}.
\]
The constant is uniform in the level choices away from \(p\),
the coefficient embedding, the test, and the precision.
\end{proposition}
\begin{proof}
Use the graph, the integral finite dual extension, and the
extension by \(\operatorname{Hom}(Y,T_i)\) constructed in
Proposition~\ref{prop:highsparse789}. The first and third
extensions are rationally split, the rational systems at \(v\)
are crystalline by Lemma~\ref{lem:inertcrys789}, and the
dual extension is crystalline by its finite local condition.
Write \(G=xb_x/x_{1,x}\) and set \(\eta=2(G+W+1)\), with
\(W\) enlarged by an absolute factor. Mark a residue-field
embedding when at least one full embedding over it has weight
\(b_x\). At unmarked embeddings all upper and lower jumps
on the three extensions have absolute value at most \(\eta\).
At each full embedding, the jumps of the Hom kernel are
differences of target and graph jumps. In any such difference
two large weights have equal sign before subtraction.
All upper jumps on the three extensions are therefore at
most \(b_x+W\) everywhere. Lemma~\ref{lem:highsign790} gives
a choice of side \(U\), among graph and dual, with lower jumps
at least \(-\eta\) at every full embedding. The other side
\(V\) has upper jumps at most \(\eta\) at all of them. This
also holds for a scalar dual, whose weights are all short.

We give the dependence of the ramified comparison on these
data. Take \(J=\lfloor120\log x/\log p\rfloor\), \(A=2J+1\)
and mark cuts followed by \(J\) consecutive unmarked
residue-field embeddings. Lemma~\ref{lem:conductorsign783}
gives the spacing condition \(2J<l_0\). The gap argument of
Theorem~\ref{thm:sparsecups780} then bounds the number of
steps between consecutive cuts by \(A\).
For \(e=[K:K_0]\) and for \(d_K\) as in
Section~\ref{sec:ramifiedorders782}, one has
\(e\le 2 f_x\ll_{q_0,d} x^{12}\) and \(d_K\le 1+v_p(e)\).
In particular \(p^{J+1}\ge 8e\) at large nodes.
Theorem~\ref{thm:ramifiedcut782} gives integral points over
\(p^\kappa\) at each cut for the three extensions, with
\[
    \kappa\ll 1+J(d_K+1)\eta+e(b_x+W)/p^{J+1}.
\]
Impose Taylor orders \(D=\eta+2\) at each of the first
\(A\) primes \(P_j=\phi^j(E)\) on the two cut points.
For a single residue-field embedding collect the jumps
over all its full embeddings. The short jumps give one
interval of indices of length \(O(\eta)\). On each side
all far jumps, if present, give one further interval
of length \(O(\eta)\), by the sign lemma. This uses the
actual multiplicities only to specify the jumps. The order
correction in Proposition~\ref{prop:ramifiedorders782}
uses their union, taking each integer index once.
For an interval of length \(M\) inside \([-8b_x,8b_x]\)
the sum of the valuations of its nonzero integers is
\(O(M/(p-1)+\log(2b_x)/\log p)\), by counting multiples
of each \(p^a\). The cost over these primes is therefore
bounded by
\[
    C_1\ll_{q_0,d}(A+1)(A+d_K+4)\eta
                        +(A+1)\log(2b_x)/\log p .
\]

For completeness, let \(c,i\) be consecutive cuts at distance
\(a\), and \(S=\prod_{j<a}P_j^\eta\). In the notation of
Lemma~\ref{lem:ramifiedcoordinates782}, the one-sided bounds
give \(S M_{a,i}(U)\subseteq M_{0,i}(U)\) and
\(S M_{0,i}(V)\subseteq M_{a,i}(V)\). The \(a\)-fold
Frobenius of an initial uncorrected \(U\)-point at \(c\)
is integral in \(M_{a,i}(U)\). Imposing orders \(D\) there
costs at most \(C_1\), since that lattice is regular at
\(P_j\) for \(j<a\). Scale all points to their maximum
quotient exponent \(e_*\le\kappa+C_1\), up to rounding.
At \(i\), replace the \(U\)-point by this corrected
Frobenius point, replace the \(V\)-point by Frobenius of
the corrected one at \(c\), and then do the same on \(U\).
The three kernel differences belong respectively to
\(S^{-1}M_{0,i}(U)\), \(S^{-1}M_{a,i}(V)\) and \(M_{a,i}(U)\).
In the first two moves the unchanged point has displacement
divisible by \(u\tau\prod_{j<a}P_j^D\).
In the last move its displacement contains
\(u^{p^a}\phi^a(\tau)\).
The contraction from the third extension, multiplied by
\(p^\kappa\), has denominator \(\tau\) at level zero and
\(\phi^a(\tau)\) at level \(a\). Their ratio is a unit
times \(\prod_{j<a}P_j\). In each move a factor \(u\)
survives after clearing these denominators and \(S\).
The same divisibilities hold on translates in
\(G_{K(\varpi^{1/p^A})}\), as used in
Proposition~\ref{prop:conductorcups783}.
The Frobenius cochain descent of that proposition thus
kills the cup by \(p^{2e_*+\kappa}\) on this subgroup.
Eisenstein's criterion gives index \(p^A\), so
corestriction adds \(A\) to the exponent. Extending
coefficients for the residue-field components and then
applying an integral coefficient projection fixing \(1\)
gives the same annihilator in the original coefficient ring.

Finally \(A+d_K=O_{q_0,d}(\log(2x))\), \(p^{J+1}>x^{120}\),
\(G/b_x\le x^{1-10^6}\) and \(b_x\ge x^{400}\).
They bound \(e(b_x+W)/p^{J+1}\) by \(O_{q_0,d}(b_xx^{-108})\),
and bound all the other terms in \(2e_*+\kappa+A\) by the
same quantity.
\end{proof}

\begin{theorem}[The integral local comparison through the product of the ranks]
\leavevmode\label{thm:highallsmall790}
At every sufficiently large node in the original interval,
\[
 \sum_{p\le Q_x}(\log p)\sup
       \sum_{\substack{v\mid p\\v\ {\rm spherical}}}\sum_i d_{v,i}
              \ll_{q_0,d} f_x(s+t)b_x x^{-108},
 \qquad Q_x=k_xt .
\]
In particular the sum over primes of the supremum over
integer nodes \(x\ge V_0\) in the interval of these local
terms, divided at each node by \(s B_x^{\rm det}\), is
\(O_{q_0,d}(V_0^{-102})\). Only nodes with \(p\le Q_x\)
are counted in this last assertion.
\end{theorem}
\begin{proof}
For \(p\in\mathcal R\), the integral annihilator and the
generator count in Theorem~\ref{thm:highsmall789} give
\(d_{v,i}\ll_{q_0,d}(1+[E_{x,v}:\QQ_p])\operatorname{rank}(T_i)
b_x x^{-108}\) by the proposition. The sum over all \(v\)
of their local degrees is \(2f_x\), and the sum over \(i\)
of the ranks is \(s+t\). The finite sum of \(\log p\) over
\(\mathcal R\) depends only on \(q_0,d\). The complementary
sum for \(p\le Q_x\) is bounded in
Theorem~\ref{thm:highsmall789} by the same expression with
\(x^{-110}\). This proves the first assertion.
Since \(s+t=6s\) and \(m_x k_x\ll x^{11/2}\), division by
\(sB_x^{\rm det}=s f_x x b_x(\log b_x)^3/(m_xk_x)\)
gives \(O_{q_0,d}(x^{-103})\) at a node. For each prime
the supremum of nonnegative normalized terms is bounded
by their sum over nodes. Interchange the two nonnegative
sums and use \(\sum_{x\ge V_0}x^{-103}\ll V_0^{-102}\).
\end{proof}
\clearpage

\subsection{Repeated large weights and the middle range of coefficient primes}
\label{sec:highmiddle791}
Write \(Z=x_{1,x}\), \(q=q_{1,x}\), \(b=b_x\) and \(G=xb/Z\).
For each coefficient prime define
\[
 D_{x,p}=\sup\sum_{\substack{v\mid p\\v\ {\rm spherical}}}
                         \sum_i d_{v,i},
\]
with the data and the supremum of Theorem~\ref{thm:highsmall789}.
The modulus \(q\), which grows along the interval, controls the
primes at which the sparse argument can encounter three successive
large weights. This continues to hold above \(Q_x=k_xt\).

\begin{lemma}[Three successive large weights]
\leavevmode\label{lem:threesuccessive791}
On the sparse sets chosen in Lemma~\ref{lem:spacing781}, at any
prime \(p\nmid qd\), three successive large weights in the
periodic Frobenius orbit imply \(p\equiv 1\) or \(-1\bmod q\).
The implication is uniform under translation of the set.
\end{lemma}
\begin{proof}
Let \(M=2^s=m_x^{22}\) in the notation of the spacing
construction, using \(s\) here only for the exponent of two.
The selected place classes have the form \(-u g^e\bmod q\),
where \(u\) is a fixed unit, \(g=1+3l_0q_0\) has order
\(M\), and the selected exponents belong to
\([0,M/(8L))\), \(L=2l_0\). They contain no nonconstant
three-term arithmetic progression. Indeed their construction
in Lemma~\ref{lem:sphere48} gives this property with both
positive coefficients equal to one. Thinning the set
preserves the property.

Take the three embeddings in the order of arithmetic
Frobenius, reversing them if necessary. Its action on a
\(q\)-th root of unity raises the root to the power \(p\).
Reorienting a label to the CM type changes its exponent on
the root by a sign. The ratios of the second and third
place classes to the first are thus \(\pm p\) and
\(\pm p^2\bmod q\), respectively. If their three selected
exponents are \(e_0,e_1,e_2\), raising the first ratio to
the fourth power and the second to the second power gives
\[
                2(2e_1-e_0-e_2)\equiv0\pmod M .
\]
The integer on the left has absolute value less than \(M\)
by the range of the exponents. Thus \(2e_1=e_0+e_2\) and
the progression property gives \(e_0=e_1=e_2\). In particular
the first ratio is one modulo the full \(q\), proving the
claim. The argument allows repeated visits to the same label.
\end{proof}

\begin{proposition}[The total loss in the two exceptional classes]
\leavevmode\label{prop:highrare791}
At a sufficiently large node,
\[
 \sum_{p\equiv\pm1\bmod q}(\log p)D_{x,p}
       \ll_{q_0,d} f_x(s+t)b\,\frac{(\log(2b))^2}{q}.
\]
Here and below \(s=r_x\) denotes the rank. The sum includes
all primes in the two residue classes.
\end{proposition}
\begin{proof}
For \(x\) sufficiently large, a prime congruent to \(\pm1\bmod q\)
is at least \(q-1\), hence greater than any prime divisor of
\(6q_0d\). The completion is unramified and \(p\ge5\).
By Proposition~\ref{prop:highsparse789} and its calculation of
the labelled weights, all three crystalline extensions used in
the pairing have upper jumps \(h_i\le Cb\). Only this last
weight bound is used here. In Theorem~\ref{thm:orbitlift764}
the forward sum is at most \(Cb/(p-1)\). The single-lattice
exponent \(c_i\) is zero when that sum is less than one and
otherwise at most seven times the sum. The common-lattice
exponent is at most \(c_i+v_p(h_i!)\le c_i+h_i/(p-1)\).
The two-point cochain comparison in
Theorem~\ref{thm:labelledcup765} therefore gives a cup
annihilator \(p^a\) with \(a\ll b/p\). The three exponents
are all zero if \(p>Cb+2\), on increasing \(C\).
The contraction here can be formed from
\(\operatorname{Hom}(Y,T_i)\), as in
Corollary~\ref{cor:periodlabels764}. The generator count
of Theorem~\ref{thm:highsmall789} now gives
\(D_{x,p}\ll f_x(s+t)b/p\), with vanishing beyond \(Cb+2\).
These bounds allow torsion tests at the given finite precision.

Write the positive integers in the two classes, apart from
one, as \(kq\pm1,\ k\ge1\). For \(p\le Cb+2\) we have
\(\log p\ll\log(2b)\). Since \(kq-1\ge kq/2\),
\[
   \sum_{\substack{p\le Cb+2\\p\equiv\pm1\bmod q}}
       \frac{\log p}{p}
        \ll \frac{\log(2b)}q
               \sum_{1\le k\le (Cb+3)/q}\frac1k
        \ll \frac{(\log(2b))^2}q .
\]
This proves the estimate for the actual prime sum.
\end{proof}

\begin{theorem}[The higher-rank comparison up to \(x_{1,x}/x^{12}\)]
\leavevmode\label{thm:highmiddle791}
For a node \(x\) put \(P_x=x_{1,x}/x^{12}\) and
\[
  \mathcal P_x=\{p:p\le P_x\}\ \cup\
               \{p:p\equiv\pm1\bmod q_{1,x}\}.
\]
On the original fields and interval, one has
\[
 \sum_p(\log p)\sup_{\substack{x\ {\rm in\ the\ interval}\\
                               p\in\mathcal P_x}}
     \frac{D_{x,p}}{s B_x^{\rm det}}
          \ll_{q_0,d}V_0^{-9/2}.
\]
An empty supremum contributes zero. The level and the
coefficient choices in \(D_{x,p}\) are taken at that node.
\end{theorem}
\begin{proof}
Fix \(x\) large. Theorem~\ref{thm:highallsmall790} covers
\(p\le Q_x\). Consider \(Q_x<p\le P_x\) outside the two
classes of the proposition. These primes are unramified in
\(E_x\). Lemma~\ref{lem:threesuccessive791} gives the
no-triple condition for the Frobenius argument with \(J=1\).
The part of the proof of Proposition~\ref{prop:highsparse789}
through its bound for general \(J\) uses the prime cutoff
only to obtain that condition and the size inequality
\(16(J+1)p^J\eta\le b\), where \(\eta=2(G+W+1)\) and
\(W\ll x^{42}\). Here \(x^{10^6}\le Z\le b^{10^{-6}}\),
so \(\eta\ll G\) and \(p\eta/b\ll x^{-11}\). Thus the size
inequality holds. More explicitly, between cuts there are
at most three steps. For a pair of large weights of opposite
signs, the critical degree inequality is
\[
      \left\lceil(b+W)/(p-3)\right\rceil+p\eta < b-W .
\]
It follows from \(p>Q_x\), \(W=o(b)\), and the preceding
bound on \(p\eta/b\). The order corrections for the short and
far windows and the three changes of points consequently give
the uniform exponent
\[
     H_p\ll_{q_0,d} 1+b/p^2+\eta+\log(2b)/\log p .
\]
All pairings here are at a spherical place of the original
field and use the original precision. The same local generator
count as in the proposition bounds \(D_{x,p}\) by
\(O_{q_0,d}(f_x(s+t)H_p)\).

Summing over these primes after removing the factor \(f_x(s+t)\)
costs at most
\[
       C_{q_0,d}\big(b/Q_x+P_x(1+\eta+\log(2b))\big).
\]
Indeed \(\sum_{p\le y}\log p\ll y\) and dyadic summation gives
\(\sum_{p>Q_x}(\log p)/p^2\ll1/Q_x\). We have
\(G\ge b^{1-10^{-6}}\), so \(1+\eta+\log(2b)\ll G\).
Now \(P_xG=b x^{-11}\) and \(Q_x\asymp x^{10}\).
The displayed sum is \(O_{q_0,d}(b x^{-10})\).
For the exceptional classes at every prime the proposition
gives \(O_{q_0,d}(b(\log(2b))^2/q)\) after the same factor
is removed. Since \(q\ge m_x^{22}\ge x^{11}\), adding these
terms and the small-prime bound gives
\[
 \frac{1}{s B_x^{\rm det}}
      \sum_{p\in\mathcal P_x}(\log p)D_{x,p}
   \ll_{q_0,d}
       \frac{x^{9/2}}{(\log b)^3}
       \left(x^{-10}+x^{-11}(\log(2b))^2\right)
   \ll_{q_0,d} x^{-11/2}.
\]
Here \(s+t=6s\) and \(m_x k_x/x\ll x^{9/2}\).
For each prime bound the interval supremum of the nonnegative
terms by their sum over its eligible nodes. Interchanging
the sums and summing \(x^{-11/2}\) over integers
\(x\ge V_0\) proves the assertion.
\end{proof}
\clearpage

\subsection{Integral lifting at primes above the inverse weight gap}
\label{sec:hightail792}
Keep the trace relations and \(D_{x,p}\) of
Sections~\ref{sec:highsmall789} and~\ref{sec:highmiddle791}.
The normalized graph systems at the chosen automorphic twist
have scalar weights \(0,\ldots,t-1\). We use this normalization
throughout. Let \(h=s^8\). The old weights have absolute
value at least \(h-(t-1)/2\), which permits integral comparison
of both sides of a cup product near the upper end of the range
of coefficient primes.

\begin{lemma}[Choosing the order of the two moves at each label]
\leavevmode\label{lem:swappable792}
Use the crystalline extensions, integral contraction, unramified
field, lifts of \(E^a\) and integer \(B<p\) in
Theorem~\ref{thm:twolatticecup768}. At each embedding suppose
that each of the two initial extensions has a unit lift or has
no jump in \([-a,0)\). Suppose the two sides at this embedding
can be denoted by \(U,V\) so that all jumps of \(U\) are at
most \(B\) and all jumps of \(V\) are at least \(-(B+1)\).
The choice of \(U\) may depend on the embedding. Then the
exponent \(s_X+s_{X'}\) of that theorem still annihilates the
cup at its given precision.
\end{lemma}
\begin{proof}
Take \(D=a+B+1\) and \(P=\phi(E)\). On the two original
sides \(X,X'\), independently, construct the integral points
of Theorem~\ref{thm:deepintegrallift767} and scale them to
the respective maxima \(s_X,s_{X'}\) over all labels. Their
displacements have factors \(u\tau E^D\) in the first kernel
lattices. The initial cup is small, in the sense of
Theorem~\ref{thm:cuppath757}, since the pairing on the first
lattices has denominator \(E^a\tau\). On the adjacent
Frobenius lattices it has denominator \(P^a\mu\), with
\(\mu=\phi(\tau)\) in this notation.

Fix a label \(i\) and name the two points \(u_i,v_i\) using
the choice \(U,V\) at \(i\). A subscript \(i^-\) in the
following comparison denotes the point on that \emph{same}
original side at the preceding label, regardless of the
names chosen there. For \(\phi v_{i^-}\) the displacement
has factor \(u^p\mu P^D\) in its adjacent lattice.
Corollary~\ref{cor:secondlatticedepth767} supplies a point
\(q_i\) in that lattice of the same quotient, with displacement
in \(u\mu E^B P^a\) times its kernel. The cost \(v_p(B!)\)
is zero. First change \(v_i\) to \(q_i\). The difference
and its translates lie in \(E^{-(B+1)}M_{0,i}(V)\) on the
kernel, by the lower jump bound. The fixed \(U\)-displacement
clears this power, the pairing denominator, and leaves \(u\).
Next change \(u_i\) to \(\phi u_{i^-}\). The difference now
lies in \(E^{-B}M_{1,i}(U)\) by the upper jump bound, and
the displacement of \(q_i\) clears it and the adjacent
denominator. Finally change \(q_i\) to \(\phi v_{i^-}\).
Both are in the adjacent lattice, and the displacement of
\(\phi u_{i^-}\) leaves a factor \(u^p P^{D-a}\).

All these inclusions apply to the translates needed for
either order of the two arguments, since \(G_K\) preserves
the lattices and the displayed principal ideals. Thus a
first-argument change is small and a second-argument
change is a coboundary modulo the small subgroup, by the
two cochain identities of Theorem~\ref{thm:cuppath757}.
In the original \(X,X'\) coordinates the terminal pair
is exactly the Frobenius transform of the pair at \(i^-\)
on every label. Summing over the labels and taking Frobenius
constants as in Theorem~\ref{thm:twolatticecup768} gives
the original cup multiplied by \(p^{s_X+s_{X'}}\).
\end{proof}

\begin{theorem}[An explicit integer supporting the very large prime defects]
\leavevmode\label{thm:highproduct792}
There is a constant \(C=C(q_0,d)\) such that, at every sufficiently
large node, the positive integer
\[
     A_x=\prod_{\substack{n\in\ZZ\\|n-b_x|\le Cx^{42}}}n
\]
has the following property. At every prime \(p\ge 8b_x/h\)
and spherical place \(v\mid p\), the cohomology quotient
\(C_i/(C_i\cap F_i)\) vanishes when \(p\nmid A_x\).
At the remaining primes in this range it is killed by \(p^4\).
In particular every such prime contributing nonzero length
divides just one integer in the displayed product and divides
it to the first power.
\end{theorem}
\begin{proof}
Let \(b=b_x\), \(G=xb_x/x_{1,x}\). Choose an integer
\(W=O_{q_0,d}(x^{42})\) bounding the deviations from the
centers \(0,\pm G,\pm b\) in
Proposition~\ref{prop:highsparse789}, enlarged to cover
the upper jumps of the contraction, which are at most
\(b+W\). At a spherical place above \(p\ge 8b/h\), the
field is unramified and all three representations are
crystalline by Lemma~\ref{lem:inertcrys789}. Put
\[
             a=\lceil(b+W)/(p-3)\rceil,\qquad B=W .
\]
For large \(x\), \(a<h/4\) and \(B<p\), since \(h\ll x^{40}\),
\(W=o(b)\), and \(b\ge x^{400}\). By
Theorem~\ref{thm:uniformlift760} the quotient \(E^a\) lifts
on each of the three first lattices. For general coefficients
one applies that theorem to the pullback of the trivial
quotient along \(\ZZ_p\to\mathcal O\) and projects the lift
to each residue-field component, as in
Corollary~\ref{cor:actualperiodmultiples768}.

We check the two conditions on the initial points. The
graph weights consist of the nonnegative scalar entries and
the old entries \(L\pm(\sigma_u+w_j)\) at the indicated
labels, where \(L=(t-1)/2\). In the far window the
displacement has magnitude at least \(G-O(x^{41})\);
in the other old entries it is the absolute value of a
nonzero odd multiple of \(h\). Hence every negative old
graph weight has absolute value at least \(h-L\). For the
dual of the old target a nonzero weight is of the form
\(-1-L\mp(\sigma_u+w_j)\) and has absolute value at least
\(h-L-1\). In particular the graph and this dual extension
have no jump in \([-a,0)\). For a scalar dual the jumps of
the test extension are \(0,-1-u\) for \(0\le u<t\); all
are nonpositive. Its unit quotient thus lifts by
Theorem~\ref{thm:orbitlift764}, since all upper jump bounds
are zero. These facts include every integral dual test.

At each label the graph has a single sign for all its
far jumps and the old dual has the opposite sign. A scalar
dual has no far jump. Thus the side with negative far jumps
has upper jumps bounded by \(B\), and the opposite side
has lower jumps at least \(-(B+1)\), increasing the original
\(W\) once to absorb the shift of the dual. If the negative
window is empty, its side still has the required upper
bound. Lemma~\ref{lem:swappable792} applies. Put
\(D=a+B+1\). Each valuation index occurring in the exponent
on either side is a nonzero integer in some \([-a-u,D-u)\),
where \(u\) is a jump at the label in question. All these
indices lie within \(C x^{42}\) of a center in
\(\{0,b,-b,G,-G\}\), with \(C\) fixed independently of
the datum, since \(a<h/4\). We may take the same constant
in the definition of \(A_x\).

For large \(x\), \(G+Cx^{42}<8b/h\). Indeed
\(Gh/b=xh/x_{1,x}\to0\) and \(x^{42}h/b\to0\).
Thus the three intervals about \(0,\pm G\) contain no
nonzero multiple of \(p\). The remaining two intervals
are negatives of each other, and the positive one has
length \(2Cx^{42}<p\). It contains at most one multiple
of \(p\), which cannot be a multiple of \(p^2\), since
\(p^2\ge 64b^2/h^2>b+Cx^{42}\). The union therefore
contributes at most two to each valuation sum, and zero
if \(p\nmid A_x\). The lemma gives the claimed exponent
for every cup. The annihilator equality in
Proposition~\ref{prop:freetests741} gives the same exponent
for the quotient, and gives its vanishing when all cups
are zero.
\end{proof}

\begin{corollary}[The upper tail on the original period interval]
\leavevmode\label{cor:hightailsum792}
With the data and denominator of Theorem~\ref{thm:highmiddle791},
\[
   \sum_p(\log p)\sup_{\substack{x\ {\rm in\ the\ interval}\\
                                  p\ge 8b_x/s^8}}
       \frac{D_{x,p}}{sB_x^{\rm det}}
               \ll_{q_0,d} V_0^{-349}.
\]
\end{corollary}
\begin{proof}
At an eligible prime the generator count of
Theorem~\ref{thm:highsmall789} gives
\(D_{x,p}\le 16 f_x(s+t)\). Theorem~\ref{thm:highproduct792}
restricts its support to primes dividing \(A_x\). There are
\(O_{q_0,d}(x^{42})\) positive factors and each is at most
\(2b_x\). Unique factorization gives
\(\sum_{p\mid A_x}\log p\le\log A_x\ll_{q_0,d}x^{42}\log(2b_x)\).
The normalized sum at the node is therefore
\[
      \ll_{q_0,d}
       \frac{x^{42+9/2}}{b_x(\log b_x)^2}
       \ll_{q_0,d} x^{-350},
\]
using \(b_x\ge x^{400}\). Bounding the primewise suprema by
their sums over nodes proves the corollary.
\end{proof}
Together with Theorem~\ref{thm:highmiddle791}, this leaves
the spherical coefficient primes satisfying
\[
       \frac{x_{1,x}}{x^{12}}<p<\frac{8b_x}{s^8},
                 \qquad p\not\equiv\pm1\pmod{q_{1,x}},
\]
for the sum of the local quotients considered here.
The integral index of the full tested Hecke lattice in
Corollary~\ref{cor:remainingbudget784} still requires a
global comparison.
\clearpage

\subsection{All scalar tests in the higher-rank graph}
\label{sec:highscalars793}
Consider the full higher-rank graph, including its old weights
of both signs. For \(0\le c<t\), the scalar summand
\(\mathcal O(-c)\) of the true lift has dual test lattice
\(\mathcal O(c+1)\). Its test extension has only the weights
\(0,-1-c\) at every label. This permits an integral comparison
for the scalar targets throughout the interval of primes
left by Theorem~\ref{thm:highproduct792}.

\begin{lemma}[A scalar test against a graph with short weights of both signs]
\leavevmode\label{lem:shortgraph793}
Let \(K/\QQ_p\) be unramified, \(p\ge5\). Use the two crystalline
extensions \(X,X'\) and the integral contraction \(X''\) of
Theorem~\ref{thm:twolatticecup768}. Suppose \(E^a\) lifts in
the first lattice for \(X,X''\) at every label, and a unit
lifts there for \(X'\). Let \(B\ge0,\ \eta=B+1\) be integers.
Assume all the jumps of \(X'\) belong to \([-\eta,0]\).
At each label assume either that all jumps of \(X\) are at
least \(-B\), or that all its jumps are at most \(B\).
Put \(d=2a+2\eta+1\), \(a_X=a,\ a_{X'}=0\). For \(U=X,X'\)
define
\[
   s_U=\max_i
       \sum_{\substack{n\ne0\\ -a_U-j\le n<d-j\
                         {\rm\ for\ some\ jump}\ j\ {\rm at}\ i}}
                         v_p(|n|).
\]
Then \(p^{s_X+s_{X'}+v_p(d!)}\) kills the cup of the two
cocycles at its original modulus.
\end{lemma}
\begin{proof}
We give the changes of points, including the effect of
negative short graph weights. Write \(M_i,L_i\) for the
first and adjacent Frobenius lattices at \(i\) and
\(P=\phi(E)\). By Theorem~\ref{thm:deepintegrallift767}
there are graph points \(x_i\in E^{-a}M_i(X)\) over
\(p^{s_X}\) and test points \(x'_i\in M_i(X')\) over
\(p^{s_{X'}}\), with displacements in \(u\tau E^d\)
times the respective first kernel lattices. We suppress
the label idempotent on scalar quotients. A point over
\(E^a\) in \(X''\) gives the denominator \(E^a\tau\)
for the pairing of first lattices. Frobenius gives
\(P^a\mu\) for the pairing of adjacent lattices.
As in Theorem~\ref{thm:twolatticecup768}, these are
inclusions in the integral lattices inside the finite
Witt comparison, where the denominators are invertible.

Put \(e_0=v_p(d!)\), \(v_i=p^{e_0}x'_i\).
Apply Corollary~\ref{cor:secondlatticedepth767} to
\(\phi x'_{i^-}\in L_i(X')\). This point has displacement
in \(u^p\mu P^d\) times the adjacent kernel, which is
contained in \(u\tau P^d\) times that kernel. The corollary
with both exponents \(d\) gives \(q_i\in L_i(X')\) over
the same quotient as \(v_i\), with displacement in
\(u\mu E^d P^d\) times the adjacent kernel. On the graph,
\(\phi x_{i^-}\in P^{-a}L_i(X)\) has displacement in
\(u^p\mu P^d\) times its adjacent kernel.

The initial cup of displacements at \((x_i,v_i)\) is
small since one input supplies \(u\tau E^d\) and the first
denominator is \(E^a\tau\). First change \(v_i\) to \(q_i\).
The test lower bound gives \(L_i(X')\subset E^{-\eta}M_i(X')\),
so the kernel difference pairs with the graph displacement
into \(uE^{d-a-\eta}A_i(1)\), in the notation of
Theorem~\ref{thm:scalarmeromorphic770}. Next change \(x_i\)
to \(\phi x_{i^-}\) at fixed \(q_i\).
If the graph jumps are at least \(-B\), then
\(L_i(X)\subset E^{-B}M_i(X)\).
Its kernel difference belongs to
\(E^{-a-B}P^{-a} M_i(Y)\), where \(Y\) is its kernel.
Indeed multiplication by \(E^{a+B}P^a\) puts both points
in the first lattice with the same quotient. The displacement
of \(q_i\) belongs to \(u\mu E^{d-\eta}P^d\) times the first
test kernel. The first denominator leaves the factor
\[
        u\,(\mu/\tau)\,E^{d-\eta-2a-B}P^{d-a}.
\]
Both exponents are nonnegative. If instead the graph jumps
are at most \(B\), we have \(M_i(X)\subset E^{-B}L_i(X)\).
The same multiplication puts the graph difference in the
adjacent kernel. Using \(q_i\) in its adjacent lattice
and the adjacent denominator leaves
\(uE^{d-a-B}P^{d-2a}\), again small. Finally change \(q_i\)
to \(\phi v_{i^-}\). Both are in the adjacent test lattice
and the fixed graph displacement contributes
\(u^p\mu P^d\). Every translate satisfies the same inclusions.

The terminal pair is the Frobenius of the actual initial
pair at the preceding label. The first and second cochain
identities and descent to the Frobenius constants, as in
Theorems~\ref{thm:cuppath757} and~\ref{thm:labelledcup765},
therefore give a coboundary for the stated multiple. The
argument over coefficients containing all labels descends
by an integral coefficient projection fixing \(1\).
\end{proof}

\begin{theorem}[The full spherical sum for the scalar summands of a true lift]
\leavevmode\label{thm:highscalars793}
With the original field, twists and interval, write
\[
 D^{\rm sc}_{x,p}
   =\sup\sum_{\substack{v\mid p\\v\ {\rm spherical}}}
               \sum_{0\le c<t}d_{v,\mathcal O(-c)} .
\]
Then for sufficiently large \(V_0\)
\[
     \sum_p(\log p)\sup_{x\ {\rm in\ the\ interval}}
          \frac{D^{\rm sc}_{x,p}}{s B_x^{\rm det}}
                 \ll_{q_0,d} V_0^{-9/2}.
\]
Here the trace cocycles can come from the entire higher-rank
graph at the specified auxiliary levels. All finite precisions
and all integral dual classes are included.
\end{theorem}
\begin{proof}
At \(p\le P_x=x_{1,x}/x^{12}\), this is included in
Theorem~\ref{thm:highmiddle791} by nonnegativity. Fix a
node and write \(b=b_x,\ G=xb/x_{1,x}\). For \(p>P_x\)
the field is unramified. At each spherical place its graph
systems and their normalized summands are crystalline by
Lemma~\ref{lem:inertcrys789}. Form the two extensions plus
the contraction for a single test as in
Proposition~\ref{prop:highsparse789}.
Choose \(W=O_{q_0,d}(x^{42})\), enlarged to exceed \(t+2\),
so that all short weights have absolute value at most \(W\)
and at a label of base weight \(\beta\in\{b,G\}\) all the
graph's far weights lie in one signed window
\(\epsilon\beta+[-W,W]\). At such a label either the
lower bound \(-W\) or the upper bound \(W\) applies to
all graph weights. Both the graph and the contraction
have upper jumps at most \(b+W\), by increasing \(W\)
to cover weight differences. The scalar test has only
\(0,-1-c\) as jumps. Hence its unit lifts at every label,
by Theorem~\ref{thm:orbitlift764}. For \(X,X''\) use
Theorem~\ref{thm:uniformlift760} to lift \(E^a\) at all
labels, where \(a=\lceil(b+W)/(p-3)\rceil\). It applies
to these coefficient lattices by the pullback and projection
in Corollary~\ref{cor:actualperiodmultiples768}.

Apply the lemma with \(B=W\). Each integer in either of
its two Taylor sums, and each positive integer through
\(d\) in the factorial, lies within \(C(W+1+b/p)\)
of a member of \(\{0,b,-b,G,-G\}\). Indeed the index
interval about the negative of a jump has endpoints
\(-a_U-j,d-j\), and \(a,d=O(W+1+b/p)\). The exponent
for any scalar test is at most
\[
         3\sum_{n\in\mathcal I_x}v_p(|n|),
 \qquad
 \mathcal I_x=\{n\in\ZZ\setminus\{0\}:
     \min_{\beta\in\{0,b,-b,G,-G\}}|n-\beta|
             \le C(W+1+b/P_x)\}.
\]
This single set works for all \(p>P_x\) and all the data
at the node. Each integer in one Taylor union is counted
only once. A unit of normalized length in a dual coordinate
is measured using \(v_\pi(p)\), so the local generator count
and exact annihilator in Theorem~\ref{thm:highsmall789} give
\[
       D^{\rm sc}_{x,p}\ll_{q_0,d}
             f_x t\sum_{n\in\mathcal I_x}v_p(|n|).
\]
Here we summed the ranks of the scalar targets, whose sum
is \(t\).

There are \(O_{q_0,d}(W+1+b/P_x)\) nonzero integers in
\(\mathcal I_x\), all of absolute value at most \(2b\) for
large \(x\). By unique factorization,
\[
  \sum_{p>P_x}(\log p) D^{\rm sc}_{x,p}
   \ll_{q_0,d} f_x t\sum_{n\in\mathcal I_x}\log |n|
   \ll_{q_0,d} f_x t(W+1+b/P_x)\log(2b).
\]
Since \(P_x\ge x^{10^6-12}\), \(W\ll x^{42}\), and
\(b\ge x^{400}\), the factor \(W+1+b/P_x\) is
\(O_{q_0,d}(b x^{-120})\). Division by \(sB_x^{\rm det}\)
costs at most \(O(x^{9/2}/(f_x t b(\log b)^3))\) as a
multiplicative factor, by \(t=5s\) and \(m_xk_x/x\ll x^{9/2}\).
The normalized contribution is \(O_{q_0,d}(x^{-115})\).
Sum over the integer nodes to bound the primewise suprema
for \(p>P_x\), and add the bound already proved for
\(p\le P_x\).
\end{proof}
For the local sum over spherical places, only the target
\(R_\xi\) of rank \(s\) remains in the intermediate range
of Section~\ref{sec:hightail792}. Its dual contains large
weights, so the bounded test hypothesis in
Lemma~\ref{lem:shortgraph793} holds here specifically for
the scalar targets.
\clearpage

\subsection{Removing the poles of a family of integral dual tests}
\label{sec:rankpoles794}
Let \(K/\QQ_p\) be unramified of degree \(f\), with \(p\ge5\).
Let \(\mathcal O\) be the valuation ring of a finite coefficient
field \(L\) containing the images of \(K\), and \(\pi\) a uniformizer.
We measure finite length by
\(\ell_p(A)=\operatorname{length}_{\mathcal O}(A)/v_\pi(p)\).
The \(f\) labels are the embeddings of \(K\) into \(L\).
At a label \(i\) use \(R_i=\mathcal O[[u]],\ E=u-p\) and the
first Breuil--Kisin lattice \(M_i\) of
Lemma~\ref{lem:horizontalcoordinates767}. We first keep track
of all dual cocycles together.

\begin{proposition}[The cost of removing a pole at one labelled weight]
\leavevmode\label{prop:rankpole794}
Let \(X\to B\) be a surjection of crystalline \(\mathcal O\)-lattices
with trivial free quotient \(B\) and kernel \(Y\). Take integers
\(a,D\ge0\). Suppose that for each basis vector of \(B\) and each
label its multiples by \(E^a\) and by \(p^a\) have lifts in \(M_i(X)\).
Let \(S\) be an integer at least the maximum over labels of
\[
 \sum_{\substack{n\ne0\\-a-j\le n<D-j\
                       {\rm\ for\ some\ labelled\ weight}\ j\ {\rm of}\ X}}
          v_p(|n|).
\]
There is a sublattice \(B_*\subseteq B\) with
\[
 \ell_p(B/B_*)\le S\operatorname{rank} B+
        (a+S)\sum_i\sum_{-a\le t<0}
                  \operatorname{mult}_{i,Y}(t).
\]
Every \(c\in B_*\) has, at every label \(i\), an integral
point in \(M_i(X)\) over \(\mathbf e_i c\) with displacement in
\(u\tau E^D M_{i,A}(Y)\) for every \(g\in G_K\).
Here \(\mathbf e_i\) denotes the label idempotent.
\end{proposition}
\begin{proof}
Fix a horizontal basis \(d_j\) adapted to the filtration
at \(i\), allowing repeated indices for repeated weights.
An element of \(E^tM_i(X)\) has, in the direction of weight
\(j\), a coordinate in \(E^{t-j}L[[E]]\).
For \(-a\le t\le0\) let \(I_{i,t}\subseteq B\) be the constants
with a lift in \(E^tM_i(X)\) satisfying the following two
conditions. Its coordinate in every direction of weight \(j\)
agrees with an element of \(L\) modulo \(E^{D-j}L[[E]]\),
and its displacement belongs to \(u\tau E^DM_{i,A}(Y)\).
These are \(\mathcal O\)-linear conditions on the space of
points over constants. The kernel statement of
Lemma~\ref{lem:adjacentlattice762} applies to a trivial free
quotient of any finite rank, by preservation of kernels or
by taking the kernels of its coordinate maps successively.

The Taylor construction in Theorem~\ref{thm:deepintegrallift767}
applies to a point over any constant vector. At a nonzero
integer order the divided displacement has quotient zero,
so the correction multiplies the quotient vector by
\(p^{v_p(|n|)}\). At order zero one takes the actual Taylor
constant in each horizontal coordinate. Thus the proof gives
both conditions defining \(I_{i,t}\). Starting with an
\(E^a\)-lift divided by \(E^a\) yields \(p^S B\subset I_{i,-a}\).
Starting with an integral \(p^a\)-lift yields
\(p^{a+S}B\subset I_{i,0}\). We use the same bound \(S\)
at both starts, multiplying up to the indicated exponent.

Consider a step \(t\to t+1\), with \(t<0\). For \(j<t\)
the initial order \(t-j\) is positive, so the correcting
Taylor constant in that direction must be zero, since
\(D-j>t-j\). The order is therefore at least \(t+1-j\).
For \(j>t\), agreement with a constant gives order at least
\(\min(0,D-j)\ge t+1-j\). For \(j=t\), exactly one
coefficient in each such direction can obstruct the next
step, namely its constant term. Vanishing of all these
coefficients puts the point in \(E^{t+1}M_i(X)\). Indeed
after division by \(E^t\) the horizontal calculation shows
divisibility by \(E\) in the rational Taylor lattice. It
holds in \(M_i(X)\) since this module is free over \(R_i\)
and \(R_i/E=\mathcal O\) injects into \(L\).

The map taking these leading coefficients on the space
of allowed points takes values in \(L^h\), where
\(h=\operatorname{mult}_{i,Y}(t)\). The quotient \(B\)
adds only weight zero. The image is a finitely generated
\(\mathcal O\)-module of rank at most \(h\). In fact the
map on \(E^tM_i(X)\) in just the directions \(j=t\) factors
through its finite free \(\mathcal O\)-fiber modulo \(E\),
and \(\mathcal O\) is noetherian. The image on the allowed
points is free and surjects onto
\(I_{i,t}/I_{i,t+1}\) by projection to \(B\). Thus that
latter group needs at most \(h\) generators. It is killed
by \(p^{a+S}\) by the integral start above. Its normalized
length is at most \(h(a+S)\).

Set \(B_*=\bigcap_i I_{i,0}\). The quotient of \(p^SB\)
by its intersection with \(B_*\) embeds in
\(\bigoplus_i I_{i,-a}/I_{i,0}\).
Adding the length of \(B/p^SB\) and summing the successive
pole steps gives the stated bound. This counts the initial
Taylor cost \(S\operatorname{rank}B\) only once.
\end{proof}

\begin{theorem}[The main summand at the split places with negative far graph weights]
\leavevmode\label{thm:neggraph794}
Use the original fields and twists and \(P_x=x_{1,x}/x^{12}\).
For each datum let \(\mathcal V^-(x,p)\) consist of the spherical
places over \(p\) at which all labelled far weights of the
normalized graph are negative. For the main target, a lattice
in \(R_\xi\), one has
\[
 \sum_p(\log p)\sup_{\substack{x\ {\rm in\ the\ interval}\\p>P_x}}
 \ \sup_{\rm data}
 \frac{\displaystyle\sum_{v\in\mathcal V^-(x,p)}d_{v,R_\xi}}
      {s B_x^{\rm det}}
       \ll_{q_0,d} (\log(2V_0))^{-2}.
\]
At every prime splitting in \(K_d=\QQ(\sqrt{-d})\), the
negative orientation selects one place in each conjugate
pair before restriction to the spherical level.
\end{theorem}
\begin{proof}
Fix a node, prime and place in the assertion. Write \(b=b_x\),
\(G=xb/x_{1,x}\), and choose \(W=O_{q_0,d}(x^{42})\)
large enough for all the deviations in
Proposition~\ref{prop:highsparse789}. Enlarge it to bound the
short jumps on graph and dual and to make \(b+W\) an upper
jump bound on their integral contraction. Here \(p>P_x\)
is unramified in \(E_x\) and \(W+2<p\). All the representations
in question are crystalline at this spherical place. Put
\(a=\lceil(b+W)/(p-3)\rceil,\ D=a+W+1\).
Choose one integer \(S_p\), independent of the place and
the datum at the node, as the sum of \(v_p(|n|)\) over all
nonzero integers within \(C(W+1+b/P_x)\) of one of
\(0,b,-b,G,-G\). As in Theorem~\ref{thm:highscalars793},
this includes every index in the Taylor sums for these
choices of \(a,D\).

Let \(T\) be the main target lattice and
\(T'=T^\vee(1)\). For \(B\) take the lattice of all
integral finite-condition raw cocycles with values in \(T'\).
Its rank is at most \((1+f)s\) by Lemma~\ref{lem:localevaluation728}.
The lattice \(X'=T'\oplus B\) has action
\[
               g(y,c)=(gy+c(g),c).
\]
It is crystalline after inverting \(p\). Indeed for an
\(\mathcal O\)-basis of \(B\), take the direct sum of the
crystalline extensions defined by its basis cocycles and
push out the \(T'\) kernels by addition. The resulting
extension has the displayed action. For every vector
\(c\in B\), apply Theorem~\ref{thm:uniformlift760} to the
pullback along \(1\mapsto c\). It gives the two lifts
required in Proposition~\ref{prop:rankpole794}; mapping
to \(X'\) and projecting to the label gives the indicated
quotient vector. This uses only the common upper bound
\(b+W\), including the trivial quotient.

The far weights of \(T'\) at this place are positive.
At most \(m_x\) weights of \(T'\) per label can lie in
\([-a,0)\). Indeed the old row has \(k_x-1\) far entries
out of \(s\), leaving \(s+1-k_x=m_x\) entries, by
Proposition~\ref{prop:highwidth788}. Moreover every
negative old dual weight has absolute value at least
\(h-(t-1)/2-1\), \(h=s^8\). This exceeds \(a\) whenever
\(p\ge8b/h\), as checked in Theorem~\ref{thm:highproduct792}.
Proposition~\ref{prop:rankpole794} therefore produces integral
deep points for all \(c\) in a sublattice \(B_*\) whose
normalized colength is at most
\[
       S_p(1+f)s+
        f m_x(a+S_p)\mathbf 1_{\{p<8b/h\}} .
\]

On the graph all upper jumps are at most \(W<p-1\), so
Theorem~\ref{thm:orbitlift764} gives an integral unit point
at every label. Correct it to depth \(D\) with
Theorem~\ref{thm:deepintegrallift767} and scale to quotient
\(p^{S_p}\) simultaneously. The contraction lattice has its
\(E^a\)-lifts by Theorem~\ref{thm:uniformlift760}. For each
\(c\in B_*\) both initial points are now integral at all
labels, with the displacements of Lemma~\ref{lem:swappable792}.
Its three moves apply with the graph as upper-bounded side
and the dual as lower-bounded side. The dual point has
constant vector \(c\), and the graph point has constant
\(p^{S_p}\). In the cochain identity the displacement on
the dual is \(c(g)+gy-y\) for a kernel coordinate \(y\).
Thus the identity gives
\(p^{S_p}[z\cup\operatorname{red}c]=0\) for every local
graph class \(z\), also at the original modulus.
The cochain argument uses kernel differences at equal
quotient vectors, exactly as for a trivial line.
In particular \(p^{S_p}B_*\) annihilates all the graph
classes. This sublattice is the same for all \(z\) at
the given target and local datum, because its construction
used only \(T'\).

By the annihilator equality of Proposition~\ref{prop:freetests741},
pairing injects the quotient measured by \(d_{v,R_\xi}\)
into \(\operatorname{Hom}_{\mathcal O}(B/p^{S_p}B_*,
\mathcal O/\pi^N)\) at precision \(\pi^N\).
On a cyclic summand the length of this Hom is bounded
by that of its domain. Since \(\sum_{v\mid p} [E_{x,v}:\QQ_p]=2f_x\),
the preceding estimates give, uniformly at the node,
\[
  \sum_{v\in\mathcal V^-(x,p)}d_{v,R_\xi}
       \ll_{q_0,d} f_x\bigl(sS_p+
                   m_x (b/p)\mathbf 1_{\{p<8b/h\}}\bigr).
\]
We used \(a\ll b/p\) on the indicated range and \(m_x\le s\).
The proof may use coefficients containing the images of
the local field. Finite-condition raw cocycles commute
with finite free scalar extension, by expansion in a
coefficient basis and the rational finite condition.
Local cohomology with the extended finite module is
likewise the corresponding scalar extension. The original
graph classes after extension satisfy the same relations,
by flatness. Normalized length of a finite module is
unchanged under this extension, so the inequality applies
over the original coefficient ring.

The total of \((\log p)S_p\) for \(p>P_x\) is at most
\(C(W+1+b/P_x)\log(2b)\), by factoring the integers in
the fixed windows. As in Theorem~\ref{thm:highscalars793},
summing the corresponding normalized costs over nodes
is \(O_{q_0,d}(V_0^{-110})\). For the second term divide
by \(sB_x^{\rm det}=s f_x x b(\log b)^3/(m_x k_x)\).
Since \(m_x^2 k_x/(sx)\ll1\) and \(p<8b/h\le b\),
the result at this prime is \(O(1/(p(\log p)^3))\),
uniformly over eligible nodes. Each such prime exceeds
\(V_0^{10^6-12}\). On a dyadic interval \([y,2y)\),
\(\sum_{p<2y}\log p\ll y\) bounds the weighted sum by
\(O((\log y)^{-3})\). Summing those intervals gives the
claimed squared logarithm. Finally the orientation of the
far weights is fixed by the restriction of the label to
\(K_d\), as in the first part of Lemma~\ref{lem:highsign790}.
At a split prime all labels of one completion induce the
same embedding of \(K_d\); conjugation reverses it.
\end{proof}
\clearpage

\subsection{A two-prime comparison for the main summand on arbitrary orbits}
\label{sec:interior795}
The pole bound in Lemma~\ref{lem:remainingpole772} applies to
the main summand as soon as the far weights miss the possible
pole interval. The short weights may still include negative
entries on both sides. We give the comparison for these signs,
allowing the sign of the far weights to change at each step.

\begin{proposition}[Comparison with bounded residual poles on both sides]
\leavevmode\label{prop:residualpoles795}
Take the unramified local setting, three crystalline extensions,
and integral contraction of Theorem~\ref{thm:twolatticecup768}.
Assume \(p\ge5\) and \(E^a\) lifts on every label for all three
extensions. Choose nonnegative integers \(e,B\). For the first
two extensions at each label, suppose the negative jumps in
\([-a,0)\) have absolute value at most \(e\). At each label
suppose the two sides can be denoted \(U,V\) so that \(U\)
has lower jumps at least \(-B\) and \(V\) has upper jumps at
most \(B\). Put \(D=a+e+B+2\). Let \(S\) be the maximum over
both sides and all labels of
\(\sum v_p(|n|)\) over nonzero integers in the union of the
intervals \([-a-j,D-j)\) for the jumps \(j\). Let \(k\) bound
the maximum of \(\sum (2+v_p(|n|))\) over nonzero integers in
the union of \([-j,e-j)\) at any label on either side.
The cup is annihilated by \(p^{2S+2k+1}\).
\end{proposition}
\begin{proof}
Use \(E=u-p,\ P=\phi(E),\ \delta=u\tau,\ \mu=\phi(\tau)\).
At each label and on each side, Theorem~\ref{thm:deepintegrallift767}
constructs a point from the \(E^a\)-lift divided by \(E^a\).
Scale to constant \(p^S\) and call the points \(u_i,v_i\)
using the assignment at \(i\). Their displacements are in
\(\delta E^D\) times the first kernel lattices. By
Lemma~\ref{lem:remainingpole772} these points belong to
\(E^{-e}M_i\), and \(v_i\) also belongs to \(L_i\) because
its jumps are at most \(B\le D\). Here \(M_i,L_i\) are
the first and adjacent lattices. The contraction gives
denominators \(E^a\tau\) and \(P^a\mu\) in the first and
adjacent pairings, respectively.

Write \(H=G_{K(p^{1/p})}\). Apply Lemma~\ref{lem:nextjet772}
at \(P\) to \(v_i\) with order \(e\). Its cost uses the jumps
of the same side at the preceding label and is at most \(k\).
After scaling the result at \(i\) to quotient \(p^{S+k}\),
write it as \(q_i\). Then
\[
 q_i-p^k v_i\in E^D M_i(V)\cap\ker\beta,\qquad
 (g-1)q_i\in\delta E^D P^e M_{i,A}(\ker\beta_V)
                          \quad(g\in H).
\]
Here \(\beta\) denotes projection onto the trivial quotient.
We have \(q_i\in L_i(V)\) since \(E^D M_i(V)\subseteq L_i(V)\).
In the next comparison the subscripts \(i^-\) refer to the
same original side as at \(i\), before assigning names to
the two points at \(i^-\).

Start with the two initial points multiplied by \(p^k\).
Their displacements pair into \(uA_i(1)\) because \(D\ge a\).
First change \(p^k v_i\) to \(q_i\). Its difference is in
the first kernel lattice, where the unchanged displacement
and the first pairing denominator leave a factor \(u\).
Next change \(p^ku_i\) to \(p^k\phi u_{i^-}\) at fixed \(q_i\).
The first point has at most an \(E^e\) denominator and
the second lies in \(P^{-e}L_i(U)\). The lower jump bound
gives \(L_i(U)\subset E^{-B}M_i(U)\). Multiplying both
points by \(E^{e+B}P^e\) puts them in the first lattice,
where their quotient coordinates agree. The displacement
of \(q_i\) and the first pairing bound consequently give
\(u E^{D-e-B-a}A_i(1)\) for the move.
Finally change \(q_i\) to \(p^k\phi v_{i^-}\). Both are in
\(P^{-e}L_i(V)\). Pair the difference with the displacement
of the Frobenius point on \(U\) in \(u^p\mu P^D L_{i,A}\).
The second denominator leaves \(u^pP^{D-e-a}\).
All the remaining exponents are nonnegative. The same
inclusions hold for translates by \(H\).

The three comparisons leave elements of the small subgroup
at every label. The final ordered pair is Frobenius of the
preceding initial pair in the original two coordinates.
The cochain identities in Theorem~\ref{thm:cuppath757}
and descent on \(H\), as in Theorem~\ref{thm:mixedcup772},
kill the restricted cup by \(p^{2S+2k}\). Since
\(X^p-p\) is Eisenstein over the unramified field \(K\),
corestriction adds one to the exponent. Retraction of
any coefficient extension gives the result at the original
precision.
\end{proof}

\begin{theorem}[The interior prime range for all spherical places]
\leavevmode\label{thm:interior795}
With \(D_{x,p}\) as in Section~\ref{sec:highmiddle791}, put
\[
      R_x=8x_{1,x}/x,\qquad Y_x=b_x/x^{45}.
\]
On the original interval and for the original twists one has
\[
   \sum_p(\log p)\sup_{\substack{x\ {\rm in\ the\ interval}\\
                                R_x\le p\le Y_x}}
       \frac{D_{x,p}}{sB_x^{\rm det}}
                   \ll_{q_0,d}(\log(2V_0))^{-2}.
\]
This includes the main target and both orientations of the
far weights, including alternating orientations on a single
Frobenius orbit.
\end{theorem}
\begin{proof}
Write \(b=b_x,\ G=xb/x_{1,x},\ h=s^8\).
We record the actual size of the short window. In the row
of Proposition~\ref{prop:highwidth788} each non-far old
displacement is an odd multiple of \(h\) with multiplier
of absolute value at most \(2m_x+1\). The offsets from
either far center have size at most
\((2m_x+1)h+O(s+t)\). The scalar entries, translations
and dual shifts cost at most \(O(s+t+1)\). Thus a single
integer
\[
        W\ll_{q_0,d} (m_x+1)h+s+t+1
                    \ll_{q_0,d} x^{81/2}
\]
bounds the short jumps and the deviations from the centers
\(\pm b,\pm G\); take it large enough also to bound the
upper jumps of the contraction by \(b+W\), using
same-label differences. In particular \(W=o(G)\).
The coefficient prime is unramified in \(E_x\) and exceeds
\(W+2\) for large \(x\). Put \(a=\lceil(b+W)/(p-3)\rceil\).
Since \(p\ge8b/G\) and \(b/G\to\infty\), we have \(a<G/4\).
The far jumps of either side have absolute value at least
\(G-W>a\). All possible negative jumps in \([-a,0)\)
therefore have absolute value at most \(W\).

At every spherical place above \(p\), use the graph, dual
and contraction of Proposition~\ref{prop:highsparse789}.
Their upper jump bounds give the \(E^a\)-lifts, by
Theorem~\ref{thm:uniformlift760}. At any one label the two
far windows on the first two sides have opposite signs
when both occur. Hence one side has all jumps at least
\(-W\) and the other has all jumps at most \(W\), with
the same conclusion for a scalar dual. Apply the proposition
with \(B=e=W\). For the first Taylor cost take \(S_p\) to be
the sum of \(v_p(|n|)\) over nonzero integers within
\(C(W+1+b/P_x)\) of \(0,\pm b,\pm G\), where
\(P_x=x_{1,x}/x^{12}\). Indeed \(R_x>P_x\) and
\(D+a=O(W+1+b/p)\). For the second cost at any one label,
the union of \([-j,W-j)\) is covered by a short interval
and a single far interval of total length \(O(W+1)\),
with absolute endpoints at most \(3b\). Counting multiples
of powers of \(p\) as in Proposition~\ref{prop:highsparse789}
gives
\[
        k\ll_{q_0,d} W+1+\log(2b)/\log p .
\]
The actual sums and their multiplicities have been taken
before any supremum over levels or embeddings. Uniformly
over those choices, exact local duality and the generator
count in Theorem~\ref{thm:highsmall789} now give
\[
 D_{x,p}\ll_{q_0,d} f_x(s+t)
                 \bigl(S_p+W+1+\log(2b)/\log p\bigr).
\]

As in Theorem~\ref{thm:highscalars793}, unique factorization
bounds the weighted sum of the \(S_p\) at a node by
\(O_{q_0,d}((W+1+b/P_x)\log(2b))\). After division by
the determinant normalization the sum of the corresponding
node costs over \(x\ge V_0\) is \(O_{q_0,d}(V_0^{-110})\).
For the other terms use \(m_x k_x/x\ll x^{9/2}\),
\(s+t=6s\), and \(p\le b/x^{45}\). In particular
\[
       \frac{x^{9/2}(W+1)}b\ll_{q_0,d}\frac1p
       \qquad\text{since }9/2+81/2=45.
\]
The term involving \(\log(2b)/\log p\) after normalization
is also \(O_{q_0,d}(1/(p(\log p)^3))\), using
\(x^{9/2}/b\le1/p\) and \(\log p\le\log b\).
The resulting primewise majorant holds before taking the
interval supremum. Since \(p\ge R_x\ge x^{10^6-1}\),
summing with weight \(\log p\) by the dyadic calculation
of Theorem~\ref{thm:neggraph794} proves the bound.
\end{proof}
\clearpage

\subsection{The transition between the two base weights}
\label{sec:weighttransition796}
Let \(Z=x_{1,x}\), \(b=b_x\), \(G=xb/Z\) and \(P_x=Z/x^{12}\).
We use the integral base weights at the node in this notation,
as in Section~\ref{sec:mainorbit776}. When \(p\) is close to \(b/G\),
the possible pole order of a crystalline lift can reach the
smaller weight \(G\). Moreover some short jumps can be negative
on \emph{both} sides of the pairing. The next proposition
allows these poles on the actual higher-rank graph and its dual.

\begin{proposition}[A local cup bound across the smaller weight]
\leavevmode\label{prop:transitioncup796}
On the higher-rank family of Section~\ref{sec:highlocal788},
choose an integral width \(W\ll_{q_0,d}x^{81/2}\) as in
Theorem~\ref{thm:interior795}. For every sufficiently large
node, every prime \(p>P_x\), every spherical place \(v\mid p\)
and every target, the pairing of a graph class with an integral
finite-condition test at the original precision is killed by
\(p^h\) for an integer
\[
       h\ll_{q_0,d} W+1+\frac{b}{p^2}+\frac{G}{p}
                                  +\frac{\log(2b)}{\log p}.
\]
All implied constants are independent of the orbit, the trace
rank, the level, the coefficient field and the precision.
\end{proposition}
\begin{proof}
Fix a place and extend coefficients to contain the unramified
field \(K=E_{x,v}\). Write \(M_{\ell,i}\) for the \(\ell\)-th
Frobenius lattice at embedding \(i\), \(i^-=F^{-1}i\),
\(E=u-p\), \(P=\phi E\), \(\delta=u\tau\), and \(\tau_\ell=\phi^\ell\tau\).
We use the three integral extensions of
Proposition~\ref{prop:highsparse789}, associated to the graph,
the test and the contraction \(Q\). At every embedding with
base weight \(\beta_i=b\) or \(G\), all short jumps of the first
two extensions lie in \([-W,W]\). Their other jumps are within
\(W\) of \(+\beta_i\) on one side and \(-\beta_i\) on the
other. For a scalar test the latter assignment on the side
without far jumps is arbitrary once the graph sign is fixed.
All upper jumps of the third extension are at most \(\beta_i+W\).
These are the same-label bounds of Proposition~\ref{prop:highsparse789},
with the width calculated in Theorem~\ref{thm:interior795}.

Put
\[
 a=\left\lceil\frac{b+W}{p-3}\right\rceil,\quad
 J=W+1+\lceil b/p^2\rceil+\lceil G/p\rceil,\quad
 D=2a+3W+10 .
\]
Our scale assumptions \(x^{10^6}\le Z\le b^{10^{-6}}\)
give \(p>4(W+2)\), \(G>2(W+2)\), and \(b-W>a\) for large
\(x\). In particular \(a\le b/p+8J\) and \(a/p\le4J\).
Theorem~\ref{thm:uniformlift760}, applied to the trivial-line
pullback and then projected to a label, gives integral lifts
of \(p^k E^s\) on all three extensions whenever \(k+s\ge a\).

We select labels on the cyclic orbit and at each selected label
specify an integer \(s_i\), together with a common \(\kappa\),
so that \(p^\kappa E^{s_i}\) lifts.
\begin{enumerate}
\item If \(b/p-G\le40J\), select all labels and take
\(s_i=\min(a,\lfloor G-W-2\rfloor)\), \(\kappa=a-s_i\).
Then \(0\le\kappa\le51J\). A negative far jump has absolute
value at least \(G-W>s_i\).
\item If \(b/p-G>40J\), select each label of weight \(b\)
and each label followed by one of weight \(G\). At labels
of the second kind put \(s_i=0\). For each of the three
extensions the weighted sum of future upper bounds there
is at most
\[
            (G+W)/p+(b+W)/(p(p-1)).
\]
Theorem~\ref{thm:orbitlift764} thus supplies integral lifts
of a common \(p^\kappa\) with \(\kappa\ll J\); take the
largest required exponent, or zero for the empty set.
At every other selected label put \(s_i=a\). Its weight
is \(b\), so \(b-W>a\) excludes every negative far jump
from \([-s_i,0)\). The unselected labels have weight \(G\)
and are followed by weight \(b\). They are isolated. A gap
of two steps starts at a selected label with \(s_i=0\).
\end{enumerate}

At each selected label, divide a lift on each of the first
two sides by \(E^{s_i}\). Raise its displacement to order
\(\delta E^D\) in \(M_{0,i,A}\) by Theorem~\ref{thm:deepintegrallift767}.
Its Taylor operators apply to any constant quotient since their
corrections have zero quotient. The exponent paid is the sum
of \(v_p(|n|)\) over the nonzero integers in the union of
\([-s_i-j,D-j)\) for the jumps \(j\). This union lies in two
intervals of total length \(O(a+W+1)\) and absolute endpoints
at most \(5b\). In an interval of length \(N\), the nonzero
multiples of \(p^r\) number at most \(1+N/p^r\). Summing over
\(1\le r\le\log(5b)/\log p\) bounds the Taylor cost by
\(O(J+\log(2b)/\log p)\).

Lemma~\ref{lem:remainingpole772} puts the resulting point
in \(E^{-W}M_{0,i}\), since every possible negative jump of
absolute value at most \(s_i\) is short. At a label with
\(s_i=0\) the point is integral. On a side with upper
jumps at most \(W\) the point is also in \(M_{1,i}\).
On each side use Lemma~\ref{lem:nextjet772} to impose order
\(W\) at \(P\) over \(H=G_{K(p^{1/p})}\).
Its cost is the sum of \(2+v_p(|n|)\) over the nonzero
integers in the union of \([-j,W-j)\) at the preceding
label. The two containing intervals now have total length
\(O(W+1)\), so the additional exponent is
\(O(W+1+\log(2b)/\log p)\). The correction is in
\(E^D M_{0,i}\), hence preserves all three membership
properties just stated, using \(E^D M_{0,i}\subset M_{1,i}\)
when the upper jumps are at most \(W\).
Write \(d_z(g)=(g-1)z\). The selected points now have
\[
          d_z(g)\in\delta E^D P^W M_{0,i,A}(\ker\beta)
                         \qquad (g\in H),
\]
where \(\beta\) is projection to the trivial line.

On the third extension divide a lift of \(p^\kappa E^{s_c}\)
at a selected label \(c\) by \(E^{s_c}\). By the integral
contraction of Corollary~\ref{cor:periodlabels764} and Frobenius
transport, \(p^\kappa Q\) on the two \(\ell\)-th kernel
lattices at \(F^\ell c\) takes values in
\[
       \phi^\ell(E)^{-s_c}\tau_\ell^{-1}A_{F^\ell c}(1)
\]
after reduction at the given precision. We shall use
\(\tau_1/\tau\) equal to \(E\) times a unit and
\(\tau_2/\tau\) equal to \(EP\) times a unit.
Call a value small when it belongs to the image of \(uA_i(1)\).

Consider any two consecutive selected labels \(c,i\), allowing
wraparound. Name the positive far side at \(i\) by \(U\) and
the other by \(V\), using these names for the same original
systems at the preceding labels. Thus at \(i\) the lower jumps
on \(U\) are at least \(-W\) and the upper jumps on \(V\)
are at most \(W\). Denote the selected points by \(u_i,v_i\).
We describe comparisons between points with the same quotient.
All auxiliary points below are constructed first. We then
scale each selected or auxiliary point to the largest quotient
exponent \(e_*\) that occurs on the orbit. This does not affect
membership or displacement divisibility.

For a gap of one step, first change \(u_i\) to \(\phi u_c\)
at fixed \(v_i\). Both points lie in
\(E^{-W}P^{-W}M_{0,i}(U)\): the second lies in
\(P^{-W}M_{1,i}(U)\), and \(M_{1,i}(U)\subseteq
E^{-W}M_{0,i}(U)\) by the lower bound. Their difference
belongs to the corresponding kernel lattice. Pairing with
\(d_{v_i}\) in the first lattice leaves a multiple of
\(uE^{D-W-s_i}\). Next change \(v_i\) to \(\phi v_c\).
Both are in \(P^{-W}M_{1,i}(V)\). Pairing their kernel
difference with \(\phi d_{u_c}\) in the adjacent lattice
leaves a multiple of \(u^p P^{D-W-s_c}\).

For a gap of two steps, the intervening label \(i^-\) has
weight \(G\), \(i\) has weight \(b\), and \(s_c=0\). Hence
\(\phi^2u_c,\phi^2v_c\) lie in their second lattices.
Put \(S=E^W P^W\). There are two possible sign patterns.
\begin{enumerate}
\item Suppose the signs at \(i^-,i\) agree. Successive
one-sided height inclusions give \(S M_{2,i}(U)\subset
M_{0,i}(U)\) and \(S M_{0,i}(V)\subset M_{2,i}(V)\).
For \(U\) start with the twice-transported pure lift of
\(p^\kappa\) from \(c\). Apply Proposition~\ref{prop:simultaneousorders779}
in \(M_{2,i}(U)\) with \(L=1\), raising the orders at both
\(E,P\) to \(T=2W+2\). The rational Taylor lattices at
both primes are regular on \(M_{2,i}\), and \(T<p\), so the
exponent added is at most \(6(T-1)\). The resulting
\(q\in M_{2,i}(U)\) has \(Sq\in M_{0,i}(U)\) and
\(d_q\in\delta E^T P^T M_{2,i,A}\) on \(H\).

First change \(u_i\) to \(q\) at fixed \(v_i\). The two
points are in \(S^{-1}M_{0,i}(U)\), and the first pairing
leaves a multiple of \(uE^{D-W-s_i}\). At fixed \(q\)
change \(v_i\) to \(\phi^2 v_c\). The difference is in
\(E^{-2W}P^{-W}M_{2,i}(V)\) with zero quotient, by
\(v_i\in E^{-W}M_{0,i}(V)\) and the second height inclusion.
The pairing with \(d_q\) costs \(\tau_2\) and leaves a
multiple of \(u E^{T-2W-1}P^{T-W-1}\). Finally change
\(q\) to \(\phi^2u_c\) at fixed \(\phi^2v_c\). Both
moving points are integral in \(M_{2,i}(U)\), and the
fixed displacement supplies \(u^{p^2}\tau_2\).

\item Suppose the signs at \(i^-,i\) differ. At \(i^-\)
the upper and lower bounds on \(V\) are \(G+W\) and
\(-W\), and \(U\) has upper jump at most \(W\).
At \(i\) the upper bounds on \(V,U\) are \(W,b+W\)
respectively. Use Theorem~\ref{thm:longpoint779} on \(V\)
with \(\ell=2\), \(\eta=W\), \(h_0=W,\ h_1=G+W\) and
\(A=2(b+W)+3W+6\). The required \(E^a\)-lift at \(i\)
and pure \(p^\kappa\)-lift at \(c\) were proved above.
In this two-step case \(pG<b-40pJ\). Since \(W\le J\),
\(b/p\le pJ\), and \(a\le b/p+8J\), it follows that
\[
          a+p(G+W)<b-W,\qquad N=a+W+p(G+W)<A .
\]
Every jump \(j<-W\) at \(i\) has \(b-W\le -j\le b+W\);
in particular \(W-j>N\) and \(-j\le A\). The theorem
gives \(q_0\in M_{2,i}(V)\) with \(Sq_0\in M_{0,i}(V)\),
quotient \(p^{\kappa+4\lceil A/p^2\rceil}\), and
\(d_{q_0}\in\delta E^A M_{2,i,A}\). Apply
Lemma~\ref{lem:secondorders775} with order \(W+1\).
It gives \(q\in M_{2,i}(V)\) with displacement in
\(\delta E^A P^{W+1}M_{2,i,A}\) on \(H\), at cost
\(2W+v_p(W!)=2W\). The correction is in \(E^A M_{2,i}(V)\).
Here \(M_{2,i}(V)\subseteq P^{-W}M_{1,i}(V)\subseteq
P^{-W}E^{-(b+W)}M_{0,i}(V)\), by the two lower bounds.
Consequently \(Sq\in M_{0,i}(V)\) as well.

Change \(v_i\) to \(q\) at fixed \(u_i\). Their difference
is in the first kernel lattice after multiplication by
\(S\), and \(d_{u_i}\) supplies that factor as well as
\(u E^{s_i}\tau\). Change \(u_i\) to \(\phi^2u_c\) at fixed
\(q\). The successive upper bounds on \(U\) and the pole of
\(u_i\) put the difference in \(E^{-(b+2W)}P^{-W}M_{2,i}(U)\)
with zero quotient. The displacement of \(q\) and the
second Frobenius pairing leave a multiple of
\(uE^{A-b-2W-1}\). Finally change \(q\) to \(\phi^2v_c\).
Both are integral in \(M_{2,i}(V)\) and the fixed
displacement supplies \(u^{p^2}\tau_2\).
\end{enumerate}

All the displayed exponents after cancellation are
nonnegative. The kernel inclusions follow from
Lemma~\ref{lem:adjacentlattice762} after flat Frobenius
base change; \(H\) preserves each indicated \(A_i\)-lattice
and principal ideal, so translates of differences have the
same bounds. The estimates for the selected and auxiliary
points give \(e_*\ll J+\log(2b)/\log p\). At an unselected
label take both initial points to be the Frobenius of the
points at its predecessor. At a selected label the initial
cup of displacements is small since \(D\ge s_i\). Its
Frobenius at an unselected label is again small.
Thus the initial cup and every change of a point lie in
the small subgroup in the cochain comparison of
Theorem~\ref{thm:cuppath757}. The terminal pair at each
label is exactly the single Frobenius of the preceding
initial pair, including across a gap of two steps.
The descent of Theorem~\ref{thm:labelledcup765} on \(H\)
therefore kills the restricted cup by \(p^{2e_*+\kappa}\):
the two affine points contribute \(p^{e_*}\) each, and
the integral pairing used throughout was \(p^\kappa Q\).
Restriction followed by corestriction costs one further
power of \(p\) since \(X^p-p\) is Eisenstein over \(K\).
An integral coefficient retraction fixing one descends the
coboundary at the original modulus. This proves the stated
bound for \(h=2e_*+\kappa+1\).
\end{proof}

\begin{theorem}[All spherical places through \(b_x/x^{45}\)]
\leavevmode\label{thm:transitionrange796}
With \(D_{x,p}\) of Section~\ref{sec:highmiddle791}, on the
original fields and interval one has
\[
   \sum_p(\log p)\!
        \sup_{\substack{x\ {\rm in\ the\ interval}\\p\le b_x/x^{45}}}
          \frac{D_{x,p}}{sB_x^{\rm det}}
                 \ll_{q_0,d}(\log(2V_0))^{-2}.
\]
Both scalar and main targets are included. On the transition
range \(P_x<p<8Z/x\) alone, the left side is
\(O_{q_0,d}(V_0^{-100})\).
\end{theorem}
\begin{proof}
Theorem~\ref{thm:highmiddle791} treats \(p\le P_x\);
Theorem~\ref{thm:interior795} treats \(8Z/x\le p\le b_x/x^{45}\).
It suffices to check the transition at the same nodes.
By exact local duality and the count of generators in
Theorem~\ref{thm:highsmall789}, the proposition bounds
\(D_{x,p}\) by \(C f_x(s+t)h\).
Set \(R=8Z/x\). Chebyshev's inequality for the sum of
\(\log p\) and a dyadic decomposition imply
\(\sum_{p>P_x}(\log p)/p^2\ll1/P_x\) and
\(\sum_{p<R}(\log p)/p\ll 1+\log R\).
Also each prime contributes \(\log(2b)\) to the weighted
sum of \(\log(2b)/\log p\). Thus at a node the total, after
removing \(Cf_x(s+t)\), is bounded by
\[
        R(W+1+\log(2b))+b/P_x+G(1+\log R).
\]
Here \(R\le8b^{10^{-6}}\), \(W\ll x^{81/2}\),
\(b\ge x^{400}\), \(P_x\ge x^{10^6-12}\), and
\(G/b=x/Z\le x^{1-10^6}\). Since
\(f_x(s+t)/(sB_x^{\rm det})\ll x^{9/2}/(b(\log b)^3)\),
the normalized total is \(O_{q_0,d}(x^{-110})\);
for example \(x^{9/2}R(W+1)/b\ll x^{45}b^{-1+10^{-6}}\).
Bound the supremum at each prime by the sum over eligible
integer nodes and sum \(x^{-110}\) for \(x\ge V_0\).
Combining the three ranges gives the assertion.
\end{proof}
\clearpage

\subsection{The length of one horizontal Taylor condition}
\label{sec:jetlength797}
Fix an unramified \(K/\QQ_p\) of degree \(f\), \(p\ge5\).
Write \(C\) for a finite coefficient field containing the images
of \(K\), and \(\mathcal O\) for its valuation ring. As in
Proposition~\ref{prop:rankpole794}, let \(X\to B\) be a surjection
of crystalline lattices with kernel \(Y\) and trivial free quotient.
Fix a label \(i\), put \(R=\mathcal O[[u]]\), \(E=u-p\),
\(P=u^p-p\), and write \(M=M_{0,i}\).
At \(E\) use the horizontal coordinates of
Lemma~\ref{lem:horizontalcoordinates767}; at \(P\) use their
Frobenius transports from \(i^-\) as in Lemma~\ref{lem:nextjet772}.
At each prime choose a basis adapted to the filtration on \(Y\)
and extend it by horizontal lifts of a basis of \(B\) of weight
zero. Such lifts exist by strict exactness of the filtered
de Rham comparison on crystalline representations. The two
choices may be made separately. In these coordinates the
quotient coordinates of any point over a constant vector are
constant. Only kernel coordinates will impose conditions.

At \(E\) the rational Taylor coefficients lie in \(C\).
At \(P\) they lie in \(C_1=C[T]/(T^p-p)\), with
\(u=T(1+P/p)^{1/p}\). Call a coordinate of depth \(n\) if
it agrees with some member of \(C\) modulo order \(n\)
at the prime in question. In particular the approximating
constant at \(P\) belongs to \(C\) even when \(C_1\) has more
than one factor. This condition implies order at least \(n\)
for its Galois displacement under the full group \(G_K\).
The horizontal vectors and the external coefficients are fixed
and the prime is taken to a unit multiple of itself.

Prescribe integers for the individual coordinate depths at
both primes. Write \(X_c\) for the integral points of \(M(X)\)
over constant vectors that meet the prescriptions \(c\).
Write \(N_c\subset M(Y)\) for integral kernel points whose
coordinates themselves vanish to the prescribed orders. Thus
\(N_c\) is an \(R\)-module contained in \(X_c\), whereas
\(X_c\) is an \(\mathcal O\)-module. Denote by \(I_c\subset B\)
the image of \(X_c\). For a finite \(\mathcal O\)-module
let \(\ell_p\) be its \(\mathcal O\)-length divided by
\(v_\pi(p)\). Equivalently, it is its \(\ZZ_p\)-length
divided by \([C:\QQ_p]\).

\begin{lemma}[Raising one coordinate order over \(G_K\)]
\leavevmode\label{lem:operators797}
Raise a set of coordinates of common prescribed depth \(n\)
by one at one of the two primes, obtaining \(d\) from \(c\).
There is an additive map \(A:X_c\to X_d\) with
\(A z-p^k z\in N_c\). At \(E\) take \(k=v_p(|n|)\) for
\(n\ne0\) and \(k=0\) for \(n=0\). At \(P\) one may take
\(k=2+v_p(|n|)\) for \(n\ne0\), and \(k=1\) for \(n=0\).
For a nonzero step the map also eliminates the leading term
of any kernel coordinate vanishing through all orders below
\(n\). On kernel points, every other previously prescribed
vanishing order persists.
\end{lemma}
\begin{proof}
Use the \(R\)-linear retraction \(\rho\) of
Lemma~\ref{lem:integraljet767} on the lattice. It respects
finite rational Taylor congruences at \(E\) and \(P\);
the latter assertion, including right sides in \(C_1[P,P^{-1}]\),
is proved in Lemma~\ref{lem:nextjet772} by Weierstrass division
by \(P^N\), clearing coefficient denominators, then retracting.
In particular homogeneous orders in several coordinates at
both primes are preserved simultaneously.

For \(n\ne0\) the operator is
\[
         z\longmapsto p^k z-
             \rho\bigl(\alpha(gz-z)/\delta\bigr),\qquad \delta=u\tau.
\]
At \(E\) use \(g,\alpha\) from Theorem~\ref{thm:deepintegrallift767},
and at \(P\) those of Lemma~\ref{lem:nextjet772}. In the latter
case \(g\) fixes the Taylor constants in \(C_1\) because it
fixes \(p^{1/p}\). At degree \(n\) the factor
\(\alpha((gE/E)^n-1)/\delta\), or its companion with \(P\)
in place of \(E\), specializes to \(p^k\). Thus the indicated
coefficient is cancelled, also when \(n<0\). Every divided
displacement used here is integral in the kernel lattice
over \(A_i\), by crystalline descent. It has the prescribed
homogeneous orders at both primes: \(\delta\) is a rational
unit at both, and all approximating constants belong to
\(C\). Retraction therefore puts the correction in \(N_c\).
These observations also prove the assertions on kernel points.
At \(E\), degree zero requires only taking the actual constant
Taylor coefficient; the identity acts on \(X_c\).

At \(P\), degree zero requires a separate argument. Choose
\(g_*\in G_K\) fixing the \(p\)-power roots of unity and
sending the chosen system of roots of \(p\) to its product
with the compatible roots of unity of Kummer parameter one.
Existence was proved in Lemma~\ref{lem:nextjet772}. Use
\[
           A z=pz+\rho\left(\sum_{j=0}^{p-1}(g_*^jz-z)\right).
\]
The correction is integral, has zero quotient and has all
prescribed homogeneous orders. At a coordinate now regular at
\(P\), write its constant term as \(\sum_{m=0}^{p-1}a_m T^m\)
with \(a_m\in C\). In the internal period factor the lifted
root \(T\) is taken to \(\zeta_p^jT\), and \(g_*^jP/P\)
is a unit. Hence the sum of the \(p\) translates has constant
term \(pa_0\), by the geometric sum of roots of unity.
This is the same rational prescription on all external
coefficient factors, so it persists under \(\rho\). This
gives agreement to order one with an element of \(C\).
\end{proof}

\begin{theorem}[The length of a single Taylor step]
\leavevmode\label{thm:jetlength797}
Suppose a step raises \(r\) kernel coordinates of common
prescribed depth \(n\) by one, keeping all other prescriptions.
At \(E\) one has \(\ell_p(I_c/I_d)\le r k\), with \(k\) as
in Lemma~\ref{lem:operators797}. The same bound holds at
\(P\) provided all the kernel coordinates at \(E\) have a
common prescribed depth \(M_0\), where
\[
    M_0>0,\qquad p\mid M_0,\qquad M_0\ge -j
        \quad\text{for each weight \(j\) of \(Y\) at \(i\)} .
\]
\end{theorem}
\begin{proof}
At \(E\) a nonzero step imposes the vanishing of one Taylor
coefficient in each of \(r\) directions. The resulting map on
\(X_c\) has image in \(C^r\) with bounded denominators. Indeed
the points belong to a fixed integral free lattice and each
finite horizontal coefficient uses finitely many Taylor
coefficients of its basis and of the integral coordinates.
The image is consequently a finitely generated
\(\mathcal O\)-module of rank at most \(r\).
Its quotient \(I_c/I_d\) is killed by \(p^k\), by the lemma.
This proves the bound. At degree zero there is no loss at \(E\).

At \(P\), put \(Q=X_c/(X_d+N_c)\). This group surjects onto
\(I_c/I_d\), and the lemma gives \(p^kQ=0\). We first
compute a different quotient, using only kernel points.
For \(w\in N_c\) record the coefficient of \(P^n\) in the
\(r\) directions in question. Its kernel is \(N_d\).
The image \(L\subset C_1^r\) is stable under
\(\mathcal O_K[[u]]\), acting on it through specialization
at \(P\). Write \(K_1=K(\pi_1)\), \(\pi_1^p=p\). The
specialized ring is \(\mathcal O_{K_1}=\mathcal O_K[\pi_1]\).
Indeed \(T^p-p\) is Eisenstein, and the valuations of
\(a_j\pi_1^j\), \(0\le j<p\), for \(a_j\in K^\times\),
are distinct modulo integral valuations of \(K^\times\);
hence an integral sum has integral coefficients.
View \(C_1=C\otimes_K K_1\) as a \(K_1\)-space of
dimension \([C:K]\). The image \(L\) has bounded denominators:
the integral coefficients of a vector reduced modulo \(P^m\)
have representatives of degree \(<pm\) with coefficients in
\(\mathcal O\), and the fixed rational coordinate matrices
and the truncated substitution for \(u\) introduce only a
fixed denominator for any specified coefficient. Thus \(L\)
is a submodule of a finite free fractional lattice over the
discrete valuation ring \(\mathcal O_{K_1}\), and is free
of rank at most \(r[C:K]\). The image of \(E\) in this ring
is \(\pi_1-p\), a uniformizer. We obtain
\[
    \dim_{\mathbb F_p} N_c/(N_d+E N_c)
       =\dim_{\mathbb F_p} L/(\pi_1-p)L
       \le f r[C:K]=r[C:\QQ_p].
\]
This computation uses the unramified base \(K\), so it applies
also to ramified external coefficients.

Set \(d_0=r[C:\QQ_p]\), and take any \(d_0+1\) elements
of \(Q[p]\), represented by \(x_j\in X_c\).
Write \(px_j=x'_j+w_j\) with \(x'_j\in X_d,\ w_j\in N_c\).
The preceding dimension bound supplies a relation among
the \(w_j\) modulo \(N_d+E N_c\), with coefficients in
\(\ZZ_p\) not all divisible by \(p\). For the same
combination \(x\) of the \(x_j\) we therefore have
\[
                      px=z+Ew,\qquad z\in X_d,\quad w\in N_c .
\]
Apply in succession the \(E\)-operators of the lemma for
the degrees \(M_0+1,\ldots,M_0+p-1\), using the displayed
nonzero-step formula on all coordinates. Denote the composite
by \(A\). All these operators cost zero since \(p\mid M_0\).
They preserve \(X_d\), and \(Ax-x\in N_c\): whenever a
coordinate agrees with a constant to its prescribed order,
the correction has that homogeneous order, whether or not
the degree being cancelled is its next degree. Starting
with \(Ew\), whose kernel coordinates vanish to order
\(M_0+1\) at \(E\), they give \(A(Ew)=E^p w'\) with
\(w'\in N_c\). In fact division by \(E^p\) leaves order
\(M_0\) at \(E\), at least the required integral order
\(-j\) in each direction. Integrality descends from the
horizontal comparison after clearing the \(E\)-denominator,
as in Lemma~\ref{lem:labelledheight764}. At \(P\) this
division is by a rational unit and preserves the homogeneous
orders of the kernel point.

Now \(E^p=P+p r_0\) for some \(r_0\in\mathcal O_K[u]\)
by the binomial formula. Additivity gives
\[
      p(Ax-r_0w')=Az+Pw'\in X_d .
\]
Indeed \(Pw'\in N_d\). The element \(Ax-r_0w'\) belongs to
\(x+N_c\); since \(X_d\) is given there by rational Taylor
conditions, the factor \(p\) cancels. This proves a nonzero
\(\mathbb F_p\)-relation among the chosen elements of \(Q[p]\).
Thus \(\dim_{\mathbb F_p}Q[p]\le d_0\). Multiplication by
\(p^{j-1}\) embeds \(Q[p^j]/Q[p^{j-1}]\) into \(Q[p]\).
The exponent bound now gives \(\operatorname{length}_{\ZZ_p}Q
\le k d_0\). Pass to its quotient \(I_c/I_d\) and divide
by \([C:\QQ_p]\).
\end{proof}

\begin{corollary}[Integral tests with prescribed orders at two primes]
\leavevmode\label{cor:negativejets797}
Suppose the weights of \(Y\) at label \(i\) lie in
\([-\eta,\eta]\) and a single signed interval
\(\pm[\gamma_i,\gamma_i+w]\). Let \(m_i\) count those in
the first interval and let \(s=\operatorname{rank}Y\).
Take positive integers \(a,\eta,w\), put
\(J=a+\eta+w+1,\ D=3a+2\eta+2w+4\), and suppose
\[
   32J<\min_i(\gamma_i,p),\qquad
   \max_i(\gamma_i+w+\eta+1)\le(p-3)a .
\]
Assume that \(p\) divides no nonzero integer within \(32J\)
of \(0\) or any \(\pm\gamma_i\). For every basis vector of
\(B\) assume its \(E^a\)- and \(p^a\)-multiples lift at
every label of \(M(X)\). There is a sublattice \(B'\subset B\)
with
\[
       \ell_p(B/B')\ll (a+1)\sum_i m_i + f s(w+1)
\]
whose vectors have integral points at all labels, of
displacement in \(\delta E^D M_{i,A}(Y)\) over \(G_K\).
At a label with negative far interval one may also require
that the divided displacement have horizontal orders
\[
 \begin{array}{ll}
 E: & M_i^*=p\left\lceil(\gamma_i+w+D+1)/p\right\rceil
                 \quad\text{in every direction},\\
 P: & w+1-v+2a\,\mathbf 1_{\{|v|\le\eta\}}
                 \quad\text{in each direction of weight \(v\) at \(i^-\)} .
 \end{array}
\]
The point itself is regular in horizontal coordinates at
\(E\) on those labels.
\end{corollary}
\begin{proof}
Proposition~\ref{prop:rankpole794} applies with its \(S=0\)
by the excluded multiples, since the indices in the
intervals \([-a-j,D-j)\) are within \(32J\) of the indicated
centers. The only weights in \([-a,0)\) are short. Its proof
thus gives, at total normalized cost at most \(a\sum_i m_i\),
integral lifts whose kernel coordinates agree with constants
through depth \(D-j\) at \(E\). At \(P\), integrality already
implies depth \(-v\).

At a negative label next impose the uniform \(E\)-depth
\(M_i^*\), coordinate by coordinate. Directions which
already meet that requirement incur no cost. On a short line
the sum of \(v_p(|n|)\) over the steps is at most
\(v_p(M_i^*!)\le M_i^*/(p-1)\ll a+1\). Here \(D-j>0\)
and \(\gamma_i+w+\eta+1\le(p-3)a\). On a far line
\(-j\in[\gamma_i,\gamma_i+w]\). A multiple of \(p\) in
\([D-j,M_i^*)\) is at most \(M_i^*-p<\gamma_i+w+D+1\)
and at least \(\gamma_i+D\), contrary to the exclusion.
The first part of Theorem~\ref{thm:jetlength797} shows
that imposing the uniform \(E\)-depth costs \(O((a+1)m_i)\).
Now \(M_i^*\ge D-j\) on every line. In particular the
new depth retains the original divisibility and satisfies
all hypotheses on \(M_0\) in the second part of that theorem.

Raise the \(P\)-orders from \(-v\) to the stated values,
using that second part at each step. The nonzero integers
visited are within \(32J\) of \(0\) or \(\pm\gamma_{i^-}\),
so the cost per direction and step is at most two,
including degree zero. This costs at most
\(2s(w+1)+4a m_{i^-}\) at this label. The quotient
coordinates throughout are the given constants of \(B\).
Intersect the resulting sublattices over the labels; the
quotient by an intersection embeds in the direct sum of
the quotients, giving the claimed total. Agreement with
external constants gives the asserted displacement orders
even after division by \(\delta\), which is a unit in
both rational comparisons. The original integral
\(\delta E^D\)-divisibility persists by the integral
test of Lemma~\ref{lem:integraljet767}. Since \(M_i^*>0\),
the points themselves are regular at \(E\) on the specified
labels.
\end{proof}
\clearpage

\subsection{The spherical local sum for all tested summands}
\label{sec:allspherical798}
At a large coefficient prime a graph may have many coordinates
of nearly equal large weight. The graph itself can have
arbitrarily large rank in the trace construction. We impose
the extra Taylor conditions of Corollary~\ref{cor:negativejets797}
on the dual tests, whose kernel is the specified target dual.
The proof below uses the width of the large weights at a
\emph{single} label.

\begin{proposition}[Pairing the graph with the restricted dual tests]
\leavevmode\label{prop:windowpair798}
Fix an unramified spherical place \(K\) of degree \(f\) over
\(\QQ_p\), a main target \(T\) of rank \(s\), and the graph,
dual and integral contraction of Proposition~\ref{prop:highsparse789}.
Use \(\eta,w,\gamma_i,a,J,D\) satisfying the numerical and
excluded-multiple hypotheses of Corollary~\ref{cor:negativejets797}.
At label \(i\) suppose that both kernels have short jumps in
\([-\eta,\eta]\) and that their other jumps belong to opposite
signed intervals \(\pm[\gamma_i,\gamma_i+w]\). Suppose
\(w\le\eta\), and that the third kernel is
\(\operatorname{Hom}(Y,T)\), where \(Y\) is the graph kernel
and the second kernel is \(T'=T^\vee(1)\). If \(m_i\) counts
the short jumps of \(T'\) at \(i\), then
\[
              d_{v,T}\ll (a+1)\sum_i m_i+fs(w+1).
\]
The same bound works for every trace graph of the indicated
weights, at the original test precision.
\end{proposition}
\begin{proof}
Work first over coefficients containing \(K\). Take \(B\) to
be the integral finite-condition raw cocycles of \(T'\), and
use the crystalline extension \(T'\oplus B\) with
\(g(y,c)=(gy+c(g),c)\), constructed and justified in
Theorem~\ref{thm:neggraph794}. The third extension is the
integral lift of the pairing modulo \(\pi^N\): if
\(\psi:Y\to T\) is a linear lift, its map to
\(\operatorname{Hom}(Y,T)\) takes a kernel coordinate \(z\)
and a quotient coordinate \(e\) to \(\pi^Nz+e\psi\).
This map is equivariant by the defining cocycle. Evaluation
and pairing with \(T'\) give an integral equivariant trilinear
map to the Tate line, equal to the required pairing \(Q\)
modulo \(\pi^N\) when the third point has quotient one.

All three extensions have upper jumps at most \((p-3)a\).
For the third, a jump in the kernel is
\(-u-v-1\) for a graph jump \(u\) and a test jump \(v\)
at the same label. Two far jumps have opposite signs and
\(|u+v|\le w\). A short and a far jump give magnitude at
most \(\gamma_i+w+\eta+1\), and two short ones give
magnitude at most \(2\eta+1\). This also shows that the
third jumps are within \(2\eta+w+1\) of \(0,\pm\gamma_i\).
All quotients add only the zero weight. By
Theorem~\ref{thm:uniformlift760} every basis vector of a
trivial quotient admits lifts of its multiples by \(E^a\)
and \(p^a\), first pulling back along the vector and then
projecting to the label.

Choose \(c\in B'\) from Corollary~\ref{cor:negativejets797}
and use the points \(y_i\) there over \(c\). On the graph
divide an \(E^a\)-lift by \(E^a\) and correct it to order
\(\delta E^D\) by Theorem~\ref{thm:deepintegrallift767}.
All the relevant nonzero Taylor indices are within \(32J\)
of the excluded centers, so the quotient remains one.
Call this point \(x_i\in E^{-a}M_{0,i}(X)\). In a horizontal
coordinate of weight \(u\) its order is at least \(d_D(u)\),
where
\[
     d_t(l)=
       \begin{cases}
         t-l,& l<-a,\\
         \min(0,t-l),& l\ge-a .
       \end{cases}
\]
Indeed the initial order is \(-a-l\); if this is positive
the Taylor constant must vanish, and in all cases there
is agreement with a constant to order \(t-l\), as in
Lemma~\ref{lem:remainingpole772}. On the third extension
likewise choose an \(E^{-a}\)-point \(z_i\) over one
corrected at \(E\) with \(t=a\). Its corrected orders
are \(d_a(l)\). Its nonzero Taylor indices also miss the
multiples of \(p\), by the bounds for the third jumps.
The function \(d_a(l)\) is nonincreasing in \(l\).
All constructions so far use the full group \(G_K\).

Use \(\tau_1=\phi(\tau)\), \(P=\phi E\), and say that a
pairing is small if its reduction lies in \(uA_i(1)\) at
the original precision. Integral contraction with \(z_i\)
gives the first-lattice denominator \(E^a\tau\), and with
\(\phi z_{i^-}\) the adjacent denominator \(P^a\tau_1\).
The initial displacements pair to a small value since
both have factor \(\delta E^D\), where \(D\ge a\).
We compare the pair \((x_i,y_i)\) to
\((\phi x_{i^-},\phi y_{i^-})\) separately at each label.

If the test has positive far sign, the graph has upper
jump at most \(\eta\) and the test has lower jump at least
\(-\eta\) at \(i\). The point \(\phi y_{i^-}\) is integral
in \(M_{1,i}\); its displacement contains
\(u^p\tau_1 P^D\) in that lattice. By
Corollary~\ref{cor:secondlatticedepth767} there is a point
\(q\in M_{1,i}\) of the same quotient whose displacement
has factor \(u\tau_1 E^{a+\eta+1}P^{2a}\). Here \(D\ge2a\)
and \(v_p((a+\eta+1)!)=0\). First change \(y_i\) to \(q\).
Their difference lies in \(E^{-\eta}M_{0,i}(T')\) since
\(y_i\) is integral and the lower bound gives this inclusion
for the adjacent lattice. The fixed graph displacement
supplies \(u\tau E^D\). Next change \(x_i\) to
\(\phi x_{i^-}\). The upper bound on the graph puts their
difference in \(E^{-(a+\eta)}P^{-a}M_{1,i}(Y)\), including
the possible pole of each point. The displacement of \(q\)
cancels these factors and the adjacent denominator, leaving
a multiple of \(u\). Finally change \(q\) to \(\phi y_{i^-}\);
both are in the adjacent lattice and the displacement
of \(\phi x_{i^-}\) supplies \(u^p\tau_1P^D\).

Suppose instead that the test has negative far sign at \(i\).
First change \(x_i\) to \(\phi x_{i^-}\) at fixed \(y_i\).
Write \(v_i(g)=(g-1)y_i/\delta\) and
\(\Delta=\phi x_{i^-}-x_i\). These are kernel vectors.
Before reduction, contract \(\phi z_{i^-},\Delta,v_i(g)\)
by the integral trilinear map. The result belongs to the
first Tate lattice with at most \(E,P\) denominators.
Indeed \(v_i(g)\) belongs to the first integral kernel
lattice, \(x_i\) has at most an \(E^a\) denominator and
the transported points have at most \(P^a\) denominators
in the adjacent lattices. Passing between the first and
adjacent lattices introduces only powers of \(E\).

We verify that the contraction actually belongs to the
integral first Tate lattice. At \(E\), the Frobenius of
the third point is regular in horizontal coordinates, since
the adjacent lattice is regular there and \(P\) is a unit
in this rational comparison. A graph direction of weight
\(u\) in \(\Delta\) has order at least \(\min(0,d_D(u))\),
which is at least \(-(\gamma_i+w)\). Every direction in
\(v_i(g)\) has absolute order
\(M_i^*\ge\gamma_i+w+D+1\). Thus the product in the
horizontal Tate coordinate has order at least one.

At \(P\) take the horizontal bases transported from \(i^-\).
In a nonzero component of the trilinear map let \(u,v,l\)
be the weights of the graph kernel, test kernel and third
extension, respectively, at \(i^-\). The filtered de Rham
map vanishes on that component if \(u+v+l>-1\), since
the Tate line has filtration jump \(-1\). Its coefficient
in horizontal bases is constant. Hence \(l\le -u-v-1\)
on every component that contributes. For \(\phi z_{i^-}\)
the order is at least \(d_a(-u-v-1)\), by monotonicity.
The graph difference has order at least \(\min(-u,d_D(u))\):
the current graph point has only a possible \(E\)-pole
and the transported point has its corrected orders from
\(i^-\). Corollary~\ref{cor:negativejets797} gives the
third term in the following lower bound for the product order,
with \(l_0=-u-v-1\),
\[
           d_a(l_0)+\min(-u,d_D(u))
                + w+1-v+2a\,\mathbf 1_{\{|v|\le\eta\}} .
\]
If \(v\) is short, use \(d_a(l_0)\ge-a-l_0\) and
\(\min(-u,d_D(u))\ge-a-u\). The sum is at least \(w+2\).
If \(v\) is far and \(u\) is short, then
\(|l_0|\ge\gamma_{i^-}-\eta-1>a\) and \(d_a(l_0)=a-l_0\),
giving the same lower bound. If both are far,
\(|u+v|\le w\) and \(|u|\ge\gamma_{i^-}>D+a\).
Then \(\min(-u,d_D(u))=-u\) and
\(d_a(l_0)\ge-(w+1)-l_0\). This gives order at least one.
These inequalities apply to every component and hence to
their sum. Horizontal kernel bases can be chosen as part
of bases of the extensions, as in Section~\ref{sec:jetlength797};
the equal quotient coordinates cancel in each difference.

In the first Tate lattice the required horizontal order is
one at \(E\) and one at \(P\), since its single weight is
\(-1\) at both labels. The first assertion is
Lemma~\ref{lem:horizontalcoordinates767}, and the second
follows by transport since the height transition between
the first and adjacent lattices is invertible at \(P\).
For clarity, these two rational tests remove the integral
denominators: after multiplying the contraction by
\(E^bP^{b'}\) to put it in the first Tate lattice, the
tests give divisibility there by \(E^b\) and \(P^{b'}\)
individually, using the integral injections of
Lemmas~\ref{lem:integraljet767} and~\ref{lem:nextjet772}.
The image of \(P\) modulo \(E\) is \(p^p-p\), a nonzero
divisor on \(A_i/EA_i\); cancellation gives divisibility
by their product. Thus the trilinear contraction is in
\(\tau^{-1}A_i(1)\), the first Tate lattice. Multiplying
back by \(\delta=u\tau\) and reducing gives the small
pairing for this change of the graph point.

Now change \(y_i\) to \(\phi y_{i^-}\). The latter is
integral in the adjacent lattice, and the former lies
there because it is integral in the first lattice and
regular at \(E\) in the horizontal basis, by the
corollary. The height inclusion and integral divisibility
test give this membership in \(M_{1,i}\). The fixed graph
displacement supplies \(u^p\tau_1P^D\), so the adjacent
pairing is small.

The estimates hold under translates in the cochain identities:
the group preserves the integral lattices, the two prime
ideals and the horizontal order bounds. At every label
the terminal pair is exactly Frobenius of the previous
initial pair. Summing the cochain identities of
Theorem~\ref{thm:cuppath757} and taking Frobenius constants
as in Theorem~\ref{thm:labelledcup765} gives
\([z\cup\operatorname{red}c]=0\) for the original graph
class \(z\) and every \(c\in B'\). For fixed \(T\) the dual
sublattice works for every graph and every precision. By
Proposition~\ref{prop:freetests741} the local quotient in
the assertion now injects into
\(\operatorname{Hom}_{\mathcal O}(B/B',\mathcal O/\pi^N)\).
Its length is at most that of \(B/B'\), by cyclic
decomposition. Corollary~\ref{cor:negativejets797} gives
the assertion. Descent to the original coefficients follows
by flat scalar extension and invariance of normalized
length as in Theorem~\ref{thm:neggraph794}.
\end{proof}

\begin{theorem}[The full spherical local sum on the period interval]
\leavevmode\label{thm:allspherical798}
Use the fields, simultaneous twists and full range of
levels and precisions of Section~\ref{sec:highmiddle791}.
For \(D_{x,p}\), the supremum of the sum of normalized
cohomology lengths \(d_{v,T}\) over the spherical places
above \(p\) and all tested summands, one has
\[
       \sum_p (\log p)\sup_{x\ {\rm in\ the\ interval}}
                  \frac{D_{x,p}}{sB_x^{\rm det}}
                      \ll_{q_0,d}(\log(2V_0))^{-2}.
\]
\end{theorem}
\begin{proof}
Theorem~\ref{thm:transitionrange796} covers the primes
up to \(Y=b/x^{45}\), where \(b=b_x\). Write \(h=s^8\).
Corollary~\ref{cor:hightailsum792} covers \(p\ge 8b/h\)
and Theorem~\ref{thm:highscalars793} covers the scalar
targets. We need the main target on \(Y<p<8b/h\).
Write \(G=xb/x_{1,x}\), and \(\beta_i\in\{b,G\}\) for the
base weight. At a given twist index \(j\) use the explicit
row of Lemma~\ref{lem:archstring787}. Its \(k_x-1\) far
entries have displacements \(\beta_i+(2j+1)h+k_x/2-1-r\),
\(0\le r<k_x-1\), before reflection and translation.
The remaining \(m_x=s+1-k_x\) old entries have displacements
which are odd multiples of \(h\) of magnitude at most
\((2m_x+1)h\). The scalar graph entries lie in
\([0,t-1]\); the translation is \((t-1)/2\) and passage
to the target dual negates and subtracts one. In particular
the far entries on either side have magnitudes within
\(s+t+2\) of \(\beta_i+(2j+1)h\), with opposite signs
on the graph and dual. This also holds on every constituent
of the rational graph kernel by
Proposition~\ref{prop:highsparse789}. Put
\[
 \begin{aligned}
 c_0&=2(s+t+3),&
 \gamma_i&=\lfloor\beta_i+(2j+1)h\rfloor-c_0,&
 w&=2c_0+3,\\
 \eta&=\lceil 4((m_x+2)h+s+t+w+2)\rceil,&
 a&=\lceil(b+3\eta+w+2)/(p-3)\rceil .
 \end{aligned}
\]
These intervals contain the weights just calculated; the
number of short weights on the dual is \(m_x\) per label.
Also \(\gamma_i+w+\eta+1\le(p-3)a\). Uniformly in the
node, \(\eta\ll x^{81/2}\), \(w\ll x^5\), \(a\ll x^{45}\),
\(p> b/x^{45}\ge x^{355}\) and
\(\gamma_i\gg G\ge b^{1-10^{-6}}\). Hence \(32J<\min(p,\gamma_i)\)
for sufficiently large \(x\).

Take a fixed \(C=C(q_0,d)\) large enough that the intervals
of radius \(32J\) about all \(0,\pm\gamma_i\), for every
twist at the node, are contained within \(C x^{45}\) of
\(\{0,\pm b,\pm G\}\). Call a prime in the remaining
range exceptional if it divides a nonzero integer in
these last five intervals. The sum of \(\log p\) over
them is at most \(O_{q_0,d}(x^{45}\log(2b))\), by factoring
the product of the absolute values of those integers.
At \emph{every} such prime the bound using only the largest
upper jump in the proof of Proposition~\ref{prop:highrare791}
gives \(D_{x,p}\ll f_x(s+t)b/p\). Indeed the future sums
in Theorem~\ref{thm:orbitlift764} are \(O(b/(p-1))\), as
are the factorial valuations, so the cup costs \(O(b/p)\)
per generator. Here \(b/p\le x^{45}\). The normalized
exceptional contribution is consequently
\[
    \ll_{q_0,d} \frac{x^{90+9/2}}{b(\log b)^2}
         \ll_{q_0,d}x^{-300}
\]
at the node. The normalizing identity is
\(B_x^{\rm det}=f_x x b(\log b)^3/(m_x k_x)\), with
\(m_x k_x/x\ll x^{9/2}\) and \(t=5s\).

At every remaining prime the proposition and
\(\sum_{v\mid p}[E_{x,v}:\QQ_p]=2f_x\) give, for
the main target,
\[
       \sum_{v\mid p,\ v\ {\rm spherical}} d_{v,T}
               \ll f_x\bigl(m_x b/p+s(w+1)\bigr),
\]
using \(a+1\ll b/p\) on this range. The first term after
normalization is at most
\(C/(p(\log p)^3)\), since \(m_x^2k_x/(sx)\ll1\)
and \(p\le b\). This bound precedes the interval supremum.
All eligible \(p\) exceed \(V_0^{355}\). On a dyadic
interval starting at \(y\), Chebyshev's inequality bounds
the sum of \(\log p/(p(\log p)^3)\) by
\(O((\log y)^{-3})\); summation gives
\(O((\log(2V_0))^{-2})\).
For the second term sum over primes \emph{at the node}
first: \(\sum_{p<8b/h}\log p\ll b/h\).
After normalization this contributes
\[
          \ll\frac{x^{9/2}(w+1)}{h(\log b)^3}
                 \ll_{q_0,d} x^{-30},
\]
since \(s\asymp x^5\). For this term and the exceptional
term, the supremum at each prime is bounded by the sum
over its eligible integer nodes. Summing for \(x\ge V_0\)
and adding the previously covered ranges proves the theorem.
\end{proof}
The sum in this theorem is the spherical cohomology term.
The tested projection in Corollary~\ref{cor:remainingbudget784}
uses an index in a global Hecke lattice. Bounding that index
also requires a comparison with the integral theta periods.
\clearpage

\subsection{Tame inertia and integral lifts at a level prime}
\label{sec:tametests800}
At either original level prime the local field is \(\QQ_p\).
Trivializing a tame type can produce a ramification index
much larger than \(p\). The number of different inertia
characters, however, controls the number of parameter powers
for which lowering the exponent costs a power of \(p\).

\begin{lemma}[Lifting a constant with tame descent data]
\leavevmode\label{lem:tamemonomials800}
Let \(U/\QQ_p\) be unramified with \(\mu_e\subset U\), where
\(p\nmid e\), and put \(F=U(p^{1/e})\). Write \(E(u)=u^e-p\).
Let \(X\to\mathcal O\) be a surjection of \(G_{\QQ_p}\)-lattices
with trivial quotient, crystalline over \(F\), with upper
Hodge jump at most \(h\ge0\). At each residue label of \(F\)
suppose the tame generator on the rational crystalline
periods is annihilated by a monic polynomial of degree at
most \(d\) with integral coefficients, having \(1\) as a
root. We allow a different polynomial at each label.
If \(e\le p^m,\ m\ge1,\ \kappa=md\), and
\(h+\kappa<p-1\), then \(p^\kappa\) has an integral lift
in each first Breuil--Kisin lattice of \(X\) over \(F\).
\end{lemma}
\begin{proof}
Enlarge the coefficient ring to contain the images of \(F\).
On each component let \(R_i=\mathcal O[[u]]\), and let
\(I_i\) be the image ideal in the module of the trivial
quotient for the normalized first lattice. A generator
\(\sigma\) of \(\operatorname{Gal}(F/U)\) acts on \(u\)
by \(\zeta_i u\), for a primitive \(e\)-th root \(\zeta_i\).
Indeed one may take compatible roots of \(p^{1/e}\)
and send their \(p^n\)-th roots to their multiples by
\(\zeta^{p^{-n}\bmod e}\). These assignments give the tame
action on the prism, fixing its residue-field coefficients.
The associated integral lattices carry this semilinear
descent. To see the descent on a lattice, apply functoriality
of the crystalline lattice functor to the automorphism
of the field and its uniformizer tower. The action on
the given representation identifies the transported
lattice with the original one, compatibly with projection.
On \(M_i/uM_i\) after inverting \(p\) this is the action
on the crystalline periods under specialization of the
Kisin module at zero. In an effective twist this is the
equivariant specialization for a tame descent datum;
untwisting by a Tate line preserves the inertia type.
Equivalently one may make this identification after
extending the tame group to all roots of order
\(|k_U^\times|\) and restrict to the given group, using
functoriality of the rational period comparison.
See the tame descent action and the equivariant
specialization in the proof of Proposition~5.17 of
Caraiani and Levin.~\cite{III-II-tamedescent800}

Set \(\nu=\kappa+1\). At each label let \(a_i\) be the
least \(a\) with \(u^a\in I_i+p^\nu R_i\). It exists by
the nonzero reduction of the quotient ideal: modulo a
coefficient uniformizer \(\pi\) it contains a power of
\(u\), as follows by base change to the Witt comparison
of the surjection. Raising this membership to
\(\nu v_\pi(p)\) in \(R_i/I_i\) gives one modulo \(p^\nu\).
The height inclusion yields
\(E^h I_i\subset R_i\phi(I_{i^-})\).
Since \(u^{e(h+\nu-1)}=(E+p)^{h+\nu-1}\) is divisible
by \(E^h\) modulo \(p^\nu\), extracting the component of
a pure power according to its exponent modulo \(p\) gives
\[
                 p a_{i^-}\le a_i+e(h+\nu-1).
\]
This is the finite free decomposition of \(R_i\) over the
Frobenius image with basis \(1,u,\ldots,u^{p-1}\), used
in Lemma~\ref{lem:frobeniusideal759}. Taking the maximum
over the orbit gives \(a_i\le e(h+\kappa)/(p-1)<e\).

Fix \(i\); write \(M=M_i(X)\), \(q\) for its indicated
annihilating polynomial. Then \(q(\sigma)M\subset uM\).
Indeed \(M/uM\) is free over \(\mathcal O\), and
\(\sigma\) is linear on this quotient. The rational
annihilation just identified therefore holds integrally
on the quotient. Suppose \(p^t u^k\in I_i+p^\nu R_i\),
\(1\le k\le a_i\). Choose \(x\in M\) whose projection is
\(p^t u^k+p^\nu f(u)\). The integral vector \(q(\sigma)x/u\)
projects to
\[
                 q(\zeta_i^k)p^t u^{k-1}+p^\nu f_1(u)
                    \quad(f_1\in R_i).
\]
For the error term, evaluating \(q(\sigma)f\) at zero
gives \(q(1)f(0)=0\), so division by \(u\) is integral.
Thus the exponent drops by one at no cost whenever
\(q(\zeta_i^k)\) is a unit. At another step
Lemma~\ref{lem:kummerlower760} gives
\(p^{t+1+v_p(k)}u^{k-1}\in I_i+p^\nu R_i\).
Its divided-displacement proof applies in a residue
component since \(G_F\) fixes its idempotent and the
external coefficients. Here \(1+v_p(k)\le m\) as
\(k<e\le p^m\).

The residues of \(\zeta_i^k\), \(1\le k\le a_i<e\),
are distinct because reduction is injective on roots
of unity of order prime to \(p\). The reduction of \(q\)
is monic, so at most \(d\) steps require the latter
operation. Starting with \(u^{a_i}\), we obtain
\(p^\kappa\in I_i+p^{\kappa+1}R_i\), multiplying by a
power of \(p\) if the actual cost is smaller. This gives
\(p^\kappa(1-pg)\in I_i\) with \(g\in R_i\); the second
factor is a unit. This proves the claim before reduction
at any test precision.
\end{proof}

\begin{theorem}[The local cup for a common tame spectrum]
\leavevmode\label{thm:tamecups800}
Over \(\QQ_p\), \(p\ge5\), take a trace graph, a target
\(T\), the dual test and their integral contraction as
in Proposition~\ref{prop:highsparse789}. Suppose the
rational systems in the graph and \(T\) are de Rham,
their Weil--Deligne monodromy is zero, and tame inertia
on them has a common eigenvalue set \(\Sigma\) at a
generator (an enclosing set suffices). Let \(p\nmid e\)
be an integer killing these finite tame characters.
Assume \(e\mid p^f-1\) with \(p\nmid f\), and
\(e\le p^m,\ m\ge1\). Put \(d=(|\Sigma|+2)^2\).
For two positive integers \(\eta,w\) and a positive
base \(\gamma>2\eta+w\), suppose all jumps on graph
and dual are in \([-\eta,\eta]\) and opposite signed
intervals of magnitudes \([\gamma,\gamma+w]\).
One of the two far intervals may be empty.
If \(p>\gamma+6(\eta+w+1)+md+2\), then every such
finite cup has an annihilator \(p^b\) with
\[
                     b\ll md+\eta+w+1 .
\]
\end{theorem}
\begin{proof}
Take \(U\) unramified of degree \(f\), and \(F=U(p^{1/e})\).
The Weil--Deligne restrictions have trivial inertia over
\(F\) and zero monodromy. They are crystalline over \(F\)
by the semistable comparison and descent of crystalline
periods, as in Theorem~\ref{thm:oldtorustraces743}.
The graph and contraction extensions split rationally
and hence are crystalline there. The test extension
also becomes crystalline: its class is finite over
\(\QQ_p\), and restriction preserves the kernel defined
by \(B_{\rm cris}\); the restricted kernel representation
is crystalline. All the embeddings of \(F\) restrict to
the original embedding of \(\QQ_p\). Thus each residue
label has the same enclosing short and far intervals,
including after a coefficient extension.

Enlarge coefficients to split the desired finite type.
A tame generator on the periods has finite order prime
to \(p\), so is semisimple. The kernel of the graph
has spectrum contained in \(\Sigma\cup\{1\}\); this
allows arbitrary repetitions of the rational trace
systems. The test kernel \(T^\vee(1)\) has inverse
spectrum, since the Tate line has trivial smooth inertia.
The contraction kernel \(\operatorname{Hom}(Y,T)\) has
spectrum in the set of all quotients of two members of
\(\Sigma\cup\{1\}\). The spectra on the three extensions
add at most the eigenvalue one. The products of \(X-\lambda\)
over the respective distinct eigenvalues, including one,
are integral monic annihilators of degree at most \(d\),
uniformly over all labels. The upper jumps on all three
extensions are at most \(\gamma+w+2\eta+1\) by the
same-label differences in the contraction. Hence the
lemma yields pure integral lifts of \(p^\kappa\) on all
three extensions with \(\kappa=md\).

Let \(D=2\eta+3\) and \(H=G_{F(\varpi^{1/p})}\) with
\(\varpi=p^{1/e}\). On the first two extensions take these
pure lifts and impose order \(D\) at \(E=u^e-p\) on their
divided displacements over \(H\), using
Proposition~\ref{prop:ramifiedorders782} with \(L=1\),
the first lattice, and the single prime \(E\).
The cost uses the union of the intervals \([-j,D-j)\)
for the jumps \(j\). They lie in two intervals of total
length \(O(\eta+w+1)\), independently of the ranks and of
the number of embeddings of \(F\). Every nonzero index
has absolute value less than \(p\), and the different
exponent in that proposition is
\(\lceil v_p(e\varpi^{e-1})\rceil\le1\). Thus the
exponent cost is \(O(\eta+w+1)\). The points are integral
in \(M_{0,i}\), their quotient coordinates can all be
scaled to \(p^{e_*}\) for
\(e_*\le\kappa+O(\eta+w+1)\), and their displacements
on \(H\) belong to \(\delta E^D M_{0,i,A}(\ker\beta)\).
On a side of upper weight at most \(\eta\), the points
also lie in \(M_{1,i}\). Indeed their horizontal coordinates
at every root of \(E\) agree with constants modulo
order \(D-j>0\); regularity at all roots gives integral
membership by Lemma~\ref{lem:ramifiedcoordinates782}.

Use the third pure point and its Frobenius transport.
On the first kernel lattices \(p^\kappa Q\) costs only
\(\tau\) in the denominator, on the adjacent ones only
\(\phi\tau\). At a label \(i\), call the side of lower
jumps at least \(-\eta\) the first moving side.
Change its point to the Frobenius of its preceding
point. Their difference is in \(E^{-\eta}M_{0,i}\)
on the kernel, because the transported point is in
\(M_{1,i}\) there. The opposite displacement supplies
\(u\tau E^D\), so the pairing is small. Now change
the upper-bounded point: both it and the transported
point are in \(M_{1,i}\). The fixed Frobenius displacement
supplies \(u^p(\phi\tau)(\phi E)^D\) and again the pairing
is small. These bounds apply to translates, and the
initial displacements give a small cup. The labelled
cochain descent of Theorem~\ref{thm:labelledcup765}
therefore kills the restriction with exponent \(2e_*+\kappa\).
The final restriction index is \(p e f\), whose \(p\)-part
is \(p\). Corestriction adds only one to the exponent,
and an integral coefficient retraction returns the
coboundary at the original precision.
\end{proof}

\begin{corollary}[The size of the tame term on the marked fields]
\leavevmode\label{cor:tamebudget800}
At each of the two original level primes and each place
above it, restrict to trace families and targets having
zero monodromy and a common tame spectrum of at most
\((s+t)s+1\) roots of unity of orders dividing
\(e=\operatorname{lcm}_{1\le j\le s}(p^j-1)\).
The supremum can be taken separately over such spectra,
embeddings, twists and precisions. For these families the
sum over primes of the interval suprema of the cohomology
loss, with weight \(\log p/(sB_x^{\rm det})\) as in
Theorem~\ref{thm:allspherical798}, is \(O_{q_0,d}(V_0^{-350})\)
at the two level primes. This statement concerns each
fixed spectrum before taking the supremum.
\end{corollary}
\begin{proof}
Here \(p_0,p_1\in[b_x^{10},b_x^{11}]\) and the places
split completely by Section~\ref{sec:levelmass678}.
Take \(f=\operatorname{lcm}(1,\ldots,s)\); then \(e\mid p^f-1\)
and \(p\nmid f\). Also \(e\le\prod_{j=1}^s(p^j-1)
<p^{s(s+1)/2}\), so \(m=s^2\) works.
The row computation in Theorem~\ref{thm:allspherical798}
gives a single far sign on each side at a fixed split
place, with \(w=O(s+t+1)\), \(\eta=O_{q_0,d}(x^{81/2})\),
and \(\gamma\) within \(O_{q_0,d}(x^{81/2})\) of
\(b_x\) or \(xb_x/x_{1,x}\). For a scalar target put its
dual weights in the short interval. In particular
\(md\ll x^{30}\), \(p>\gamma+6(\eta+w+1)+md+2\), and
every cup costs \(O_{q_0,d}(x^{41})\).
Exact local duality and the integral raw-cocycle generator
count of Theorem~\ref{thm:highsmall789} bound the sum of
normalized lengths at each of these primes by
\(O_{q_0,d}(f_x(s+t)x^{41})\). The count takes place over
the original split completions, after the preceding descent.
Multiplication by \(\log(p_0p_1)\ll\log b_x\) and division
by \(sB_x^{\rm det}\) give
\[
              O_{q_0,d}\!\left(
                  \frac{x^{41+9/2}}{b_x(\log b_x)^2}
                          \right).
\]
Finally \(b_x\ge x^{400}\). Replace each nonnegative
primewise supremum by its sum over eligible nodes and sum
the displayed bound for integers \(x\ge V_0\).
\end{proof}
The finite character mark splits the tested index as in
Theorem~\ref{thm:typedindex799}. Applying the last corollary
to an entire marked type requires identifying the
Weil--Deligne spectra of the target systems occurring in
that type, including its competitors at auxiliary levels.
\clearpage

\subsection{Inertia of a marked factor at the original level}
\label{sec:markedlevel801}
At an original level place the marking specifies a representation of
the group acting on the residue frame. We first determine what this
requires of the local parameter of a form containing that marking.
The action on the complementary block will enter the argument.

\begin{lemma}[A squarefree block with spherical complement]
\leavevmode\label{lem:finitemark801}
Let \(K\) be a nonarchimedean local field with residue field
\(k=\mathbb F_q\). For pairwise distinct Frobenius orbits of
nontrivial regular characters \(\theta_\nu\) of
\(\mathbb F_{q^{d_\nu}}^\times\), let \(\rho_\nu\) denote the
corresponding finite cuspidal representations of
\(\mathrm{GL}_{d_\nu}(k)\). Here regular means that the orbit
has size \(d_\nu\). Put \(s=\sum_\nu d_\nu\), and let \(t\ge1\).
Write \(\odot\) for finite parabolic induction, inflating the
tensor product on a Levi subgroup to its parabolic. The representation
\[
              \sigma=(\mathop{\odot}_{\nu}\rho_\nu)
                            \odot{\bf1}_{\mathrm{GL}_t(k)}
\]
is irreducible. If an irreducible smooth representation \(\pi\)
of \(\mathrm{GL}_{s+t}(K)\) contains its inflation on the maximal
compact subgroup, then the monodromy of its Weil--Deligne parameter
is zero. The smooth inertia characters consist of \(t\) trivial
characters and the \(d_\nu\) conjugates belonging to each
\(\theta_\nu\), with the characters read by local reciprocity over
the unramified degree-\(d_\nu\) extension of \(K\).
\end{lemma}
\begin{proof}
In the finite Green parametrization, induction of irreducibles
with disjoint cuspidal supports joins their partition data and
gives one irreducible. The support of the trivial representation
of \(\mathrm{GL}_t(k)\) consists of \(t\) copies of the trivial
character in degree one. All the other supports are distinct from
it and from one another, proving the first assertion.
~\cite[\S2.3]{III-II-finitegreen801}

For \(l\ge1\) let \(K_l^1\) denote the kernel of reduction on
\(\mathrm{GL}_l(\mathcal O_K)\). Taking \(K_l^1\)-invariants is
exact on smooth representations in characteristic zero by averaging.
For normalized parabolic induction its result is the finite parabolic
induction of the principal congruence invariants on the Levi blocks.
Indeed, by Iwasawa decomposition an inducing function is determined
by its restriction to the maximal compact subgroup. Right invariance
under the kernel of reduction forces the values there into the
invariants on the blocks. The modulus character on the integral
Levi is one, giving the finite function model in both directions.

Use the supercuspidal support of \(\pi\) and realize \(\pi\) as
a subquotient of induction from that support. Each supercuspidal
factor has nonzero principal congruence invariants by the preceding
calculation. Such a factor has depth zero. By the depth-zero
classification for the general linear group it is compactly induced
from the extension to the center times a maximal compact subgroup
of a finite cuspidal representation. Its principal congruence
invariants give just that finite cuspidal representation. To see
the last point use the Cartan representatives for the restriction
of the induction to the maximal compact. If the diagonal exponents
have a strict gap, the intersection with the conjugated principal
congruence subgroup reduces to a subgroup containing the unipotent
radical across that gap. This follows entry by entry from conjugation
by the diagonal powers of a uniformizer. The cuspidal inducing
representation has zero invariants under this radical. Equal
diagonal exponents give the single central double coset.
~\cite[Prop.~6.8]{III-II-depthzero801}

Uniqueness of finite cuspidal support now shows that the supercuspidal
factors of \(\pi\) have finite types \(\rho_\nu\), once each, and
trivial degree-one type, \(t\) times. A segment of length \(a\) in
the Langlands classification uses \(a\) unramified twists of one
supercuspidal. Hence every segment using a \(\rho_\nu\) has length
one. A longer segment could only occur among the unramified
characters. We exclude this using the standard module surjecting
onto \(\pi\). An unramified segment of length \(a\) gives an
unramified determinant twist of the Steinberg representation.
Its principal congruence invariants are the finite Steinberg
representation \(\mathrm{St}_a(k)\). This follows by taking
invariants in the presentation by functions on full flags modulo
the sum of the functions pulled back from proper partial flags.
~\cite{III-II-segments801}

Thus the invariant space of the standard module is the finite
induction of the \(\rho_\nu\) with the product
\(V_t=\mathop{\odot}_j\mathrm{St}_{a_j}(k)\), where the
\(a_j\)'s are the segment lengths on the trivial line. Each
constituent of \(V_t\) has support on that line. The disjoint-support
factorization distinguishes these constituents after induction
with the \(\rho_\nu\). Exactness and the occurrence of \(\sigma\)
in \(\pi\) force \({\bf1}_{\mathrm{GL}_t(k)}\) to occur in \(V_t\).
By finite Frobenius reciprocity this requires a trivial vector in
each \(\mathrm{St}_{a_j}(k)\). For \(a_j>1\) the finite Steinberg
representation is irreducible of dimension \(q^{a_j(a_j-1)/2}>1\).
Consequently all segments have length one. Local Langlands for the
Langlands quotient gives the direct sum of their parameters.
The depth-zero cuspidal formula gives the stated inertia characters
and zero monodromy on each nontrivial orbit. The other summands
are unramified characters by local class field theory.
~\cite[Prop.~2.4.1(i)]{III-II-egh801}
\end{proof}

\begin{theorem}[Common inertia at each marked level place]
\leavevmode\label{thm:markedtype801}
Use the higher-rank fields, weights and simultaneous twists of
Section~\ref{sec:highsmall789}. Fix a node, a lift index, an
embedding of the coefficient field and a tuple \(\omega\) in
Theorem~\ref{thm:typedindex799}. At \(p=p_{x,0}\) or \(p_{x,1}\)
fix a place \(v\mid p\) of \(E_x\). Every normalized target
system occurring in this tuple, at any of the auxiliary trace
levels away from \(p\), has zero Weil--Deligne monodromy at \(v\).
There is one set \(\Sigma_v\) of at most \(s+1\) eigenvalues
on a tame generator that contains the spectra of all these
systems, their true summands, the scalar targets and their
common Tate twists. Every eigenvalue has order dividing
\(\operatorname{lcm}_{1\le d\le s}(p^d-1)\).
\end{theorem}
\begin{proof}
At a chosen place write \(c\) for the degree-\(s\) member of
the tuple. The completion is \(\QQ_p\). The level fixes the
ordered frame formed by the first \(s\) residue columns. Its
stabilizer in the full finite group is the subgroup of the
parabolic preserving their span on which the first diagonal
block is the identity. In particular it contains both the
unipotent radical and the whole complementary \(\mathrm{GL}_t\).
They act trivially at this place. On the quotient of the
parabolic, acting through the first diagonal block \(g\), the
marked action is translation multiplied by \(\psi_v(\det g)\).
This is the determinant factor in the frame equivariance of
Proposition~\ref{prop:testedkernel785}. Thus a copy of the
marked finite representation \(\chi_c\) in the frame action
is \(\chi_c\otimes(\psi_v^{-1}\circ\det)\) in the ordinary
translation action. Frobenius reciprocity supplies a nonzero map
from the finite induction of this representation with
\({\bf1}_{\mathrm{GL}_t}\) into the local principal congruence
invariants. The same argument works after projection to any
automorphic constituent with a nonzero component in the tuple,
since translation and projection to its isotypic space commute.

The prescription of Section~\ref{sec:levelmass678} labels
\(\chi_c\) by distinct Frobenius orbits. Tensoring by the
determinant character multiplies the degree-\(d\) character
on each torus by \(\psi_v^{-1}\circ N_{\mathbb F_{p^d}/\mathbb F_p}\),
since the determinant there is the norm. In the character group
of tame inertia this is multiplication by one Frobenius-fixed
character, and hence preserves the orbit sizes and disjointness.
The residue character \(\psi_v\) has order dividing \(Q\).
Indeed its old part is a power of the prescribed order-dividing-\(Q\)
level character, and the auxiliary twist is split at the coefficient
prime as in Theorem~\ref{thm:oldtorustraces743}. An orbit of size
one that became trivial would have had order dividing \(Q\),
excluded by the condition on the roots of \(c\). Thus the lemma
applies. The resulting set of smooth eigenvalues, including
one, is fixed by the tuple, the local character, and the coefficient
embedding. A degree-\(d\) orbit contributes \(d\) elements of
order dividing \(p^d-1\). At the opposite place complex conjugation
replaces the type by its dual with an unramified norm twist.
There we take the inverses of the same eigenvalues.

It remains to pass to the coefficient-\(p\) Galois systems.
At this split place the cohomological base change in
Theorem~\ref{thm:splitattachment722} has the parameter of the
local automorphic representation. Its expression as an isobaric
sum is a direct sum of the local parameters of its cuspidal factors
with norm twists. Hence each factor has zero monodromy and inertia
contained in the set just found. The regular algebraic polarized
cuspidal twists on which the Galois construction acts were
identified in Proposition~\ref{prop:highwidth788}.
Theorem~1.1 of Barnet-Lamb, Gee, Geraghty and Taylor gives equality
of the semisimplified Weil representations there, and bounds the
Galois monodromy partitions above by the automorphic partitions.
Every block on the automorphic side has size one, forcing the
same on the Galois side. Semisimplification preserves the
representation of finite inertia in characteristic zero.
The normalizing character has unramified smooth part at the
original levels, and Tate and norm twists have trivial smooth
inertia. Direct summands restrict to subsets of the same
spectrum. This also covers the Tate lines used for the scalar
tests. Raising the auxiliary level at other places leaves the
frame condition at \(v\) in force on each new constituent.
~\cite[Thm.~1.1]{III-II-blggttame743}
\end{proof}

\begin{corollary}[Local cohomology on the marked components]
\leavevmode\label{cor:markedlocal801}
For the data of the theorem let \(D^{\mathrm{mark}}_{x,p}\) be
the supremum, with the tuple fixed inside each test, of the sum
of the defects \(d_{v,T}\) of Section~\ref{sec:highsmall789}
over the target summands and places above \(p\) which are spherical
or lie at the two original levels. The supremum includes the tuple
itself. On the original period fields and interval,
\[
       \sum_p(\log p)\sup_x
            \frac{D^{\mathrm{mark}}_{x,p}}{sB_x^{\rm det}}
                  \ll_{q_0,d}(\log(2V_0))^{-2}.
\]
At the two original level primes alone the sum of interval
suprema is \(O_{q_0,d}(V_0^{-350})\).
\end{corollary}
\begin{proof}
At each original level place the theorem supplies a common
spectrum for the entire individual trace family and its target.
The graph of Lemma~\ref{lem:graphrealization724} takes direct
sums of the corresponding rational systems. This adds no new
eigenvalues. The bound \(s+1\le(s+t)s+1\) permits the use of
Corollary~\ref{cor:tamebudget800}, at the chosen place and
independently at its conjugate. Its bound for the sum over
places and targets is uniform in the tuple before the
primewise supremum. This proves the assertion about the two
levels. For the spherical part use
Theorem~\ref{thm:allspherical798} on the same data.
The supremum of a sum of nonnegative terms is at most the
sum of their suprema. Adding the two estimates gives the
displayed inequality. The global index \(\lambda_\omega\)
of the tested theta projection within a tuple, defined in
Theorem~\ref{thm:typedindex799}, still requires a comparison
with these local quantities.
\end{proof}
\clearpage

\clearpage
\section{References}\label{nav:bibliography-III}
\begingroup\def\section*#1{}\begin{thebibliography}{99}\small
\bibitem{III-II-II-duality13} J. S. Milne, \emph{Arithmetic Duality Theorems}, 2nd ed., 2006, Chapters I and II. \href{https://www.jmilne.org/math/Books/ADTnot.pdf}{Book}.
\bibitem{III-II-V-cft101} J. S. Milne, \emph{Class Field Theory}. \href{https://www.jmilne.org/math/CourseNotes/CFT.pdf}{Local and global Artin reciprocity and the correspondence with idele class quotients}.
\bibitem{III-II-II-bktate717} S. Bloch and K. Kato, \emph{\(L\)-functions and Tamagawa numbers of motives}, The Grothendieck Festschrift, Vol.~I, Progress in Mathematics \textbf{86}, Birkh\"auser, 1990, 333--400, Corollary 3.8.4 and Example 3.9. \href{https://virtualmath1.stanford.edu/~conrad/BSDseminar/refs/BKTamagawa.pdf}{Published chapter}.
\bibitem{III-II-cltame743} X. Caruso and T. Liu, \emph{Quasi-semi-stable representations}, Bull. Soc. Math. France \textbf{137} (2009), 185--223. \href{https://www.numdam.org/item/BSMF_2009__137_2_185_0/}{Article}.
\bibitem{III-II-ozekitame743} Y. Ozeki, \emph{Full faithfulness theorem for torsion crystalline representations}, New York J. Math. \textbf{20} (2014), 1043--1061. \href{https://www.sci.kanagawa-u.ac.jp/math-phys/yozeki/pdf/torcris.pdf}{Author's text}.
\bibitem{III-II-blggttame743} T. Barnet-Lamb, T. Gee, D. Geraghty and R. Taylor, \emph{Local-global compatibility for \(l=p\), II}, Ann. Sci. \'Ec. Norm. Sup\'er. (4) \textbf{47} (2014), 165--179. \href{https://virtualmath1.stanford.edu/~rltaylor/lg2.pdf}{Theorem~1.1}.
\bibitem{III-II-liutorsion748} T. Liu, \emph{Torsion \(p\)-adic Galois representations and a conjecture of Fontaine}, Ann. Sci. \'Ecole Norm. Sup. (4) \textbf{40} (2007), 633--674, \S\S2--3. \href{https://www.numdam.org/articles/10.1016/j.ansens.2007.05.002/}{Original paper}.
\bibitem{III-II-wintenberger749} J.-P. Wintenberger, \emph{Le corps des normes de certaines extensions infinies de corps locaux; applications}, Ann. Sci. \'Ecole Norm. Sup. (4) \textbf{16} (1983), 59--89. \href{https://www.numdam.org/item/ASENS_1983_4_16_1_59_0/}{Original paper}.
\bibitem{III-II-gaofinite750} H. Gao, \emph{Breuil--Kisin modules and integral \(p\)-adic Hodge theory}, Theorem~1.1.16, with appendices by Y.~Ozeki and with T.~Liu. \href{https://arxiv.org/pdf/1905.08555}{Author's text}.
\bibitem{III-II-bcgntmonodromy752} G. Boxer, F. Calegari, T. Gee, J. Newton and J. A. Thorne, \emph{The Ramanujan and Sato--Tate conjectures for Bianchi modular forms}, Forum Math. Pi \textbf{13} (2025), e10, \S\S2.1--2.2. \href{https://doi.org/10.1017/fmp.2024.29}{Journal article}.
\bibitem{III-II-caraianilevin753} A. Caraiani and B. Levin, \emph{Kisin modules with descent data and parahoric local models}, Definition~3.6 and Propositions~5.16--5.17. \href{https://www.ma.imperial.ac.uk/~acaraian/papers/kisinmodules.pdf}{Authors' text}.
\bibitem{III-II-scholzetilt757} P. Scholze, \emph{Perfectoid spaces}, Publ. Math. IH\'ES 116 (2012), 245--313. \href{https://people.mpim-bonn.mpg.de/scholze/PerfectoidSpaces.pdf}{The author's text}, Section~3 and Theorem~3.7.
\bibitem{III-II-cgmexact759} P. \v{C}oupek, E. Gazaki and A. Marmora, \emph{Exactness Property of Breuil--Kisin Functors and Bloch--Kato Selmer Groups}, Sections~2 and~3, especially Proposition~3.4, Theorem~3.6 and Remark~3.7. \href{https://irma.math.unistra.fr/~marmora/2603.26035v1.pdf}{Authors' manuscript}.
\bibitem{III-II-geeliudesc760} M. Emerton and T. Gee, \emph{Moduli Stacks of \'Etale $(\varphi,\Gamma)$-Modules and the Existence of Crystalline Lifts}, Appendix F by T. Gee and T. Liu, Definition F.9, (F.10), and Theorem F.11. \href{https://www.math.uchicago.edu/~emerton/pdffiles/moduli.pdf}{Authors' text}.
\bibitem{III-II-prismatic762} B. Bhatt and P. Scholze, \emph{Prismatic $F$-crystals and crystalline Galois representations}, Theorems~5.6 and~7.9. \href{https://arxiv.org/pdf/2106.14735}{Authors' text}.
\bibitem{III-II-bartletthodge764} R. Bartlett, \emph{Degenerating products of flag varieties and applications to the Breuil--M\'ezard conjecture}, Selecta Mathematica 30 (2024), article 17, Section~10.3 and Theorem~10.5. \href{https://doi.org/10.1007/s00029-023-00905-3}{Journal article}.
\bibitem{III-II-bartlettcrys767} R. Bartlett, \emph{On the irreducible components of some crystalline deformation rings}, Forum of Mathematics, Sigma 8 (2020), e22, Section~2.1, especially Lemma~2.1.8, (2.1.9), and Proposition~2.1.10. \href{https://doi.org/10.1017/fms.2020.12}{Journal article}.
\bibitem{III-II-blggt788} T. Barnet-Lamb, T. Gee, D. Geraghty and R. Taylor, \emph{Potential automorphy and change of weight}, Ann. of Math. (2) \textbf{179} (2014), 501--609, Theorem~2.1.1(3)--(4). \href{https://annals.math.princeton.edu/2014/179-2/p03}{Published paper}.
\bibitem{III-II-labessesph789} J.-P. Labesse, \emph{Changement de base CM et s\'eries discr\`etes}, in \emph{On the Stabilization of the Trace Formula}, International Press, 2011, 429--470, Corollaire~5.3. \href{https://webusers.imj-prg.fr/~michael.harris/CBCM.pdf}{Author's text}.
\bibitem{III-II-tamedescent800} A. Caraiani and B. Levin, \emph{Kisin modules with descent data and parahoric local models}, Def.~3.1, Prop.~3.8 and Prop.~5.17. \href{https://arxiv.org/abs/1510.07503}{Authors' paper}.
\bibitem{III-II-finitegreen801} A.-M. Aubert, \emph{On the Macdonald correspondence}, \S2.3, especially the factorization over distinct cuspidal supports. \href{https://arxiv.org/abs/2501.02332}{Author's paper}.
\bibitem{III-II-depthzero801} A. Moy and G. Prasad, \emph{Jacquet functors and unrefined minimal K-types}, Comment. Math. Helv. \textbf{71} (1996), 98--121, Prop.~6.8. \href{https://researchportal.hkust.edu.hk/en/publications/jacquet-functors-and-unrefined-minimal-k-types/}{Publication record}.
\bibitem{III-II-egh801} M. Emerton, T. Gee and F. Herzig, \emph{Weight cycling and Serre-type conjectures for unitary groups}, Duke Math. J. \textbf{162} (2013), Prop.~2.4.1(i). \href{https://www.math.uchicago.edu/~emerton/pdffiles/egh.pdf}{Authors' paper}.
\bibitem{III-II-segments801} A. V. Zelevinsky, \emph{Induced representations of reductive $p$-adic groups II. On irreducible representations of $\mathrm{GL}(n)$}, Ann. Sci. \'Ec. Norm. Sup\'er. \textbf{13} (1980), 165--210. \href{https://www.numdam.org/item/ASENS_1980_4_13_2_165_0/}{Published paper}.
\end{thebibliography}\endgroup
\end{document}
